\documentclass[11pt]{article}
\usepackage[a4paper,margin=26mm,headheight=14pt]{geometry}
\usepackage[T1]{fontenc}
\usepackage{lmodern}
\usepackage{amsmath,amssymb,amsthm,mathtools}
\usepackage{microtype}
\usepackage{booktabs,longtable,array}
\usepackage{enumitem}
\usepackage[round,authoryear]{natbib}
\usepackage{xcolor}
\usepackage[colorlinks=true,linkcolor=blue!45!black,citecolor=blue!45!black,urlcolor=blue!45!black]{hyperref}
\hypersetup{pdftitle={Density Bounds for Quadratic Forms under Product Measures},pdfauthor={},pdfsubject={Quadratic forms of independent random variables},pdfkeywords={quadratic forms, anti-concentration, bounded densities, product measures, Brascamp-Lieb}}
\usepackage{fancyhdr}
\pagestyle{fancy}
\fancyhf{}
\fancyhead[L]{\small Density bounds for quadratic forms}
\fancyfoot[C]{\thepage}
\setlength{\parskip}{3pt}
\setlength{\parindent}{1.2em}
\setlength{\emergencystretch}{2em}
\setlist{itemsep=3pt,topsep=5pt}
\makeatletter
\renewcommand*{\l@subsection}{\@dottedtocline{2}{1.5em}{3.2em}}
\makeatother
\newtheorem{theorem}{Theorem}[section]
\newtheorem{lemma}[theorem]{Lemma}
\newtheorem{proposition}[theorem]{Proposition}
\newtheorem{corollary}[theorem]{Corollary}
\newtheorem{conjecture}[theorem]{Conjecture}
\theoremstyle{definition}
\newtheorem{definition}[theorem]{Definition}
\newtheorem{example}[theorem]{Example}
\theoremstyle{remark}
\newtheorem{remark}[theorem]{Remark}
\DeclareMathOperator{\rank}{rank}
\DeclareMathOperator{\tr}{tr}
\DeclareMathOperator{\diag}{diag}
\DeclareMathOperator{\supp}{supp}
\DeclareMathOperator{\Cov}{Cov}
\DeclareMathOperator{\Var}{Var}
\DeclareMathOperator{\vol}{vol}
\newcommand{\R}{\mathbb R}
\newcommand{\E}{\mathbb E}
\newcommand{\Prob}{\mathbb P}
\newcommand{\eps}{\varepsilon}
\newcommand{\one}{\mathbf1}
\newcommand{\norm}[1]{\left\lVert#1\right\rVert}
\newcommand{\abs}[1]{\left|#1\right|}
\newcommand{\ip}[2]{\left\langle#1,#2\right\rangle}

\title{Density Bounds for Quadratic Forms under Product Measures}
\author{}
\date{}

\begin{document}
\maketitle
\begin{abstract}
Let $X\in\mathbb R^n$ have independent coordinates with Lebesgue densities bounded by one, and let $B$ be an arbitrary real matrix. The quadratic form $X^TBX+b^TX+c$ depends on $B$ only through its symmetric part $A=(B+B^T)/2$. We investigate density bounds under the hypothesis $s_3(A)\ge1$. For every fixed dimension we obtain a bound of order $\sqrt n$ uniformly over $A$, $b$, and the coordinate laws. We prove dimension-free bounds for semidefinite forms, as well as bounds for diagonal forms, forms with suitable coordinate pivots, and several classes of indefinite forms satisfying quantitative geometric conditions. For product Gaussian scale mixtures satisfying an inverse-variance moment condition, the density bound is less than three for matrices of arbitrary rank and signature. Further estimates cover bounded-rank forms with unit-interval inputs and one side of a rank-three indefinite level set. We give a no-pivot reduction and integrability estimates for singular conditional densities. A dimension-free bound for general indefinite forms with arbitrary bounded coordinate densities remains open.
\end{abstract}
\clearpage
\tableofcontents
\newpage

\section{Introduction}\label{sec:introduction}

Let $X=(X_1,\ldots,X_n)$ have independent coordinates with probability densities $f_i$ satisfying $0\le f_i\le1$ almost everywhere. For a real $n\times n$ matrix $B$, put $A=(B+B^T)/2$ and consider
\begin{equation}\label{eq:main-form}
 Q(X)=X^TBX+b^TX+c=X^TAX+b^TX+c,
\end{equation}
where $b\in\mathbb R^n$ and $c\in\mathbb R$. Write $s_j(A)$ for the singular values of the symmetric part in decreasing order; these are its eigenvalue moduli in decreasing order. The density norm below is the essential supremum. No moment, centering, identical-distribution, smoothness, bounded-support, connected-support, or log-concavity assumption is imposed unless stated in a particular theorem.

\begin{conjecture}[Uniform three-eigenvalue density bound]\label{conj:main}
There is a finite absolute constant $C$ such that, for every dimension and every form \eqref{eq:main-form} with $s_3(A)\ge1$,
\begin{equation}\label{eq:conjecture-density}
 \|p_Q\|_\infty\le C.
\end{equation}
Equivalently, $\mathbb P(Q\in J)\le C|J|$ for every interval $J\subset\mathbb R$.
\end{conjecture}

The spectral condition concerns $A$, not $B$. Indeed, a full-rank skew-symmetric matrix $B$ in dimension four has $s_3(B)=1$ after normalization, while $x^TBx=0$ for every $x$. All subsequent theorems about symmetric matrices therefore apply to arbitrary real matrices through their symmetric parts.

The homogeneous case $b=c=0$ is the motivating question. Arbitrary affine terms are included because conditioning on original coordinates creates them, often outside the range of the remaining quadratic matrix. Passing from the homogeneous question to an affine-uniform intermediate theorem is consequently a substantive requirement of many proofs. Absolute continuity itself is elementary when $A\ne0$ (Lemma~\ref{geom:ac}); the issue is a uniform density bound. A separate transversality argument gives a continuous density vanishing at infinity for each fixed matrix of rank at least three (Theorem~\ref{fg:fixed-matrix-continuity}); its constants may depend on the matrix and dimension.

The normalization is natural. If $s_3(A)\ge a>0$ and the coordinate densities are bounded by a common $K>0$, then scaling $Y=KX$ reduces the proposed bound to $CK^2/a$. The remaining eigenvalues may be arbitrarily small or arbitrarily large and may have either sign. Replacing $A$ by a spectral truncation is not innocuous: the discarded quadratic part depends on the same coordinates.

\subsection{Principal conclusions}

The principal proved classes and their limitations are as follows. Bounds are stated at spectral scale one unless another scale is displayed.

\begin{longtable}{@{}p{0.27\textwidth}p{0.25\textwidth}p{0.41\textwidth}@{}}
\toprule
\textbf{Setting} & \textbf{Conclusion} & \textbf{Scope}\tabularnewline
\midrule
\endfirsthead
\toprule
\textbf{Setting} & \textbf{Conclusion} & \textbf{Scope}\tabularnewline
\midrule
\endhead
Diagonal $A$ & $\|p_Q\|_\infty\le\pi$ & All affine terms; $\pi$ is sharp in this class (Theorem~\ref{geom:diagonal}).\tabularnewline
Arbitrary real $B$ & $\|p_Q\|_\infty\le\frac{2\sqrt{6\pi n}}{s_3(A)}$ & $A=(B+B^T)/2$; every fixed dimension, all remaining spectrum and affine terms (Corollary~\ref{thm:finite-dimension}).\tabularnewline
Positive or negative semidefinite $A$ & $\|p_Q\|_\infty\le17100/s_3(A)$ & Arbitrary rank, including tiny positive spectral tails and affine terms outside the range (Theorem~\ref{sd:main}).\tabularnewline
Semidefinite, exact rank three & $\|p_Q\|_\infty\le58/s_3(A)$ & Arbitrary affine terms; $3e/s_3(A)$ suffices when the affine term is in the range.\tabularnewline
One sign, separated nonzero spectrum & $\|p_Q\|_\infty\le130/\theta$ & Every nonzero eigenvalue has modulus at least $\theta$; arbitrary rank at least three.\tabularnewline
Indefinite with small absolute spectral diagonal & $\|p_Q\|_\infty<35/s_3(A)$ & Requires $\max_i(|A|)_{ii}\le s_3(A)/680$, where $|A|=(A^2)^{1/2}$. Small $|a_{ii}|$ alone is insufficient for this argument.\tabularnewline
Physical indefinite $2\times2$ pair & Bound \eqref{eq:full-pair} & Only the determinant of the pair and $s_3(A)$ enter; no rank restriction.\tabularnewline
Diffuse finite-rank part plus coordinate phases & Integrable Fourier decay & Geometric frame conditions are required; all coordinate phases are allowed (Theorem~\ref{four:diffuse}).\tabularnewline
Cross-coordinate block with distributed volume & Dimension-free density & Quantitative volume and row-marginal conditions; these conditions are not automatic (Theorem~\ref{four:volume-density}).\tabularnewline
Hub configurations & Several dimension-free bounds & The independent hub, transverse square, projected crossing, and law-regularity assumptions are stated separately in Section~\ref{sec:hubs}.\tabularnewline
Gaussian scale mixtures & $\|p_Q\|_\infty<3$ for the full form & Every rank and signature under the average component-peak condition (Theorem~\ref{gm:main}). A separate theorem permits adapted hubs.\tabularnewline
Exact rank three, unit-interval laws & $\|p_Q\|_\infty<372$ & Every signature and affine term; bounded extra rank is allowed with a rank-dependent factor (Theorem~\ref{geom:interval}).\tabularnewline
Exact rank three, signature $(2,1)$ & $p_Q\le\pi\sqrt2$ on the ruled side & Affine term in the range only; the opposite side is not included (Theorem~\ref{geom:ruled}).\tabularnewline
\bottomrule
\end{longtable}

The dimension-free semidefinite theorem settles a broad portion of the question. For indefinite forms, the fixed-dimensional estimate above does not furnish a constant independent of $n$. The Gaussian-limit argument in Proposition~\ref{geom:threshold} forces any universal constant to be at least $3\sqrt2$.

\subsection{Proof methods}

The proofs use three complementary kinds of information. Coordinate conditioning exploits scalar and two-variable changes of variables while preserving the input product law. Fourier arguments use geometric integral inequalities, decoupling, and the distribution of squared minors of a crossing block. Hub arguments retain the law of a distinguished coordinate when conditioning on that coordinate alone would create a nonuniform bound. Each method controls specific quantities; none permits arbitrary replacement of the joint law of spectral coordinates by a bounded joint density.

The no-pivot reduction in Section~\ref{sec:structure} isolates a class for which a dimension-free density bound is still needed. The integrability estimates in Section~\ref{sec:lower-levels} control logarithmic and half-power singularities. Section~\ref{sec:limitations} gives counterexamples to stronger estimates for general joint densities and describes the open uniform case.

Classical inputs are cited at their points of use and discussed in Section~\ref{sec:literature}.

\subsection{Notation and measure-theoretic conventions}

For a rectangular matrix $M$, $\|M\|$ and $\|M\|_F$ denote its operator and Frobenius norms. The elementary symmetric function of the eigenvalues of $MM^T$ is denoted by $e_k(MM^T)$. By Cauchy--Binet, it is the sum of squared $k\times k$ minors. A principal coordinate block is always formed from original coordinates of $X$. The term \emph{leverage} refers to a diagonal entry of an orthogonal projection, or to the analogous quadratic row norm relative to a positive Gram matrix. We set $\log_+x=\max\{0,\log x\}$ for $x>0$.

Our Fourier convention is $\widehat f(t)=\int e^{itx}f(x)\,dx$, so inverse Fourier transforms carry $(2\pi)^{-1}$ in one dimension. Conditional density formulas are interpreted through Tonelli against nonnegative output test functions. Integrable characteristic functions give bounded continuous density versions; an interval estimate alone gives an essentially bounded density without an additional continuity claim. Smooth approximation of the physical coordinate laws is used when a pointwise level-set calculation could otherwise encounter exceptional slices.

\section{Relation to earlier work}\label{sec:literature}

The question considered here is an endpoint question with a uniformity requirement. If independent real random variables have densities bounded by $K$, and $A$ has three eigenvalues of absolute value at least $a>0$, the proposed conclusion is a bound of the form $C K^2/a$ for the density of $X^{\mathsf T}AX+b^{\mathsf T}X+c$. The constant must be independent of the ambient dimension, the remaining eigenvalues, the eigenvectors, the linear term, and the supports and regularity of the coordinate densities. Several established theories address parts of this question. Their precise hypotheses matter when they are used in the proofs below.

\subsection{Linear images, convolution, and geometric inequalities}

The basic linear input is the theory of marginal densities of product measures. \citet[Theorem~1.1]{RV2015} prove that an orthogonal projection onto a $k$-dimensional subspace has density bounded by $(C K)^k$. \citet[Theorem~1.1]{LPP2016} obtain the sharper estimate
\[
 \|f_{PX}\|_{\infty}\leq 2^{k/2}K^k,
\]
as well as a refinement depending on the ambient dimension and unequal coordinate density bounds. In the one-dimensional case, the sharp constant follows from the convolution comparison of \citet{Rogozin1987} and the cube slicing theorem of \citet{Ball1986}; the higher-codimension geometric background includes \citet{Ball1989}.

The integral arguments use the Brascamp--Lieb inequality \citep{BL1976}, its Gaussian extremization theory \citep{Lieb1990}, and its geometric formulation and transport proof \citep{Barthe1998}. For unit vectors $u_j$ and nonnegative weights $c_j$ satisfying $\sum_j c_j u_j u_j^{\mathsf T}=I_k$, the geometric inequality reads
\[
 \int_{\mathbb R^k}\prod_j h_j(\langle u_j,x\rangle)^{c_j}\,dx
 \leq\prod_j\left(\int_{\mathbb R}h_j\right)^{c_j}.
\]
These are established analytic tools. Gaussian extremization for this integral does not assert that Gaussian coordinate densities maximize every quadratic pushforward under an $L^\infty$ constraint. Likewise, a bound on a projected density does not give a uniform bound on the conditional density on each affine fiber: conditioning includes division by the marginal density, which can be arbitrarily small.

Where the proof partitions a frame with small individual vectors into nondegenerate groups, the relevant structural input is the partition consequence of the solution of the Kadison--Singer problem by \citet{MSS2015}. Applications must retain its upper bound on each vector norm and the resulting quantitative spectral bounds for the groups.

\subsection{Fourier estimates for rough densities}

An $L^\infty$ bound on a probability density gives useful integral control of its characteristic function, but does not furnish a universal pointwise rate of decay at every frequency. The estimate in \citet[Lemma~3.1]{RV2015} concerns the nonincreasing rearrangement of the modulus of the characteristic function. It is therefore suited to integral estimates and geometric Brascamp--Lieb arguments. Sharp Hausdorff--Young estimates provide another input \citep{Beckner1975}.

For linear combinations of independent variables, \citet{FS2007} use an analytic argument to bound L\'evy concentration functions in terms of Diophantine properties of the coefficients, including a multidimensional extension. The arithmetic hypotheses and the linear map in that result differ from the quadratic pushforward considered here.

\citet[Theorems~1--2 and~8]{BCG2012} give quantitative separation from modulus one using density bounds together with variance or quantile information, and give entropy versions of these estimates. The variance or quantile parameters are essential in those bounds. Related convolution concentration estimates are developed in \citet{BC2015}.

\subsection{Bilinear forms and the role of independence}

A particularly close result is the density estimate of \citet[Section~2]{BGKS2024}. Although their Theorem~2.1 is titled ``Density of Quadratic Forms,'' the expression there is bilinear in two independent random matrices. In the scalar case their Corollary~2.4 gives, for independent vectors $X,Y$ with independent coordinates of density at most $K$,
\[
 \left\|f_{X^{\mathsf T}ZY+u^{\mathsf T}X+v^{\mathsf T}Y+w}\right\|_\infty
 \leq C K^2\min_{j\geq2}
 \frac{1}{\sqrt j\left(\prod_{i=1}^j\sigma_i(Z)\right)^{1/j}}.
\]
The minimum is over nonzero singular values. Conditioning reduces the proof to a linear density estimate and an integrable negative first moment in dimension at least two. This result explains why a well-conditioned crossing block with at least two directions is useful. It does not directly cover a rank-one interaction accompanied by a quadratic expression in the same light coordinates; the independence and rank needed by the displayed estimate are then absent.

\subsection{The parabola extension estimate}

For the Fourier convention
\[
 (E h)(s,\tau)=\int_{\mathbb R}h(x)e^{i(sx^2+\tau x)}\,dx,
\]
the bound
\[
 \int_{\mathbb R^2}|E h(s,\tau)|^6\,ds\,d\tau
 \leq\frac{4\pi^3}{\sqrt3}\|h\|_2^6
\]
is the sharp one-dimensional Schr\"odinger Strichartz inequality. The underlying restriction theory is due to \citet{Strichartz1977}; sharp constants and Gaussian extremizers are established in \citet{Foschi2007,HZ2006}. The normalization follows from the constant $12^{-1/12}$ in \citet[Theorem~1.1]{Foschi2007} and Plancherel. If $h$ is a probability density bounded by one, then $\|h\|_2\leq1$, giving the constant $4\pi^3/\sqrt3$ used below. An elementary coarea proof of this special case is included below.

\subsection{General polynomial and Gaussian results}

The polynomial anti-concentration inequality of \citet{CW2001} gives, for a degree-$d$ polynomial $P$ under a log-concave probability measure, a bound of the form
\[
 \mu\{|P|\leq\varepsilon\}
 \leq C d\left(\frac{\varepsilon}{\int |P|\,d\mu}\right)^{1/d}.
\]
The degree-two exponent is $1/2$. Log-concavity and the normalization of the polynomial are part of that statement. \citet{GM2022} establish dimension-free covariance bounds for polynomials of product measures, and obtain sublevel and Fourier estimates under additional log-concavity hypotheses. Recent work treats further polynomial anti-concentration questions on $L^p$ balls and symmetric measures \citep{GM2026}, and regularity of nondegenerate polynomial pushforwards of log-concave measures \citep{KZ2026}. These conclusions have different hypotheses or exponents from the linear-in-interval-length endpoint sought here.

For analytic maps, \citet[Theorems~1.1--1.2]{GHS2026} relate local $L^p$ integrability of pushforward densities to singularity invariants. The value $\epsilon_\star=\infty$ means membership in every finite $L^p$; boundedness is a stronger property, treated for algebraic maps in their Corollary~6.2. In the real setting, \citet{Reiser2018} proves continuity of pushforward densities of compactly supported smooth measures for flat maps with rational-singularity fibers. These statements concern specified map and input regularity and do not give a dimension-uniform constant for arbitrary bounded product densities.

For Gaussian input, rotation preserves independence and a quadratic polynomial can be diagonalized without leaving the input class. Density bounds for squared Gaussian norms and weighted chi-square variables are studied by \citet{GNSU2019,BNU2020}; signed, noncentral Gaussian quadratic forms are also treated in \citet{CK2026}. The elementary characteristic-function identity already gives
\[
 |\mathbb E e^{it(g^{\mathsf T}Ag+b^{\mathsf T}g)}|
 \leq\prod_i(1+4t^2\lambda_i(A)^2)^{-1/4}.
\]
Three nonzero eigenvalues make the right-hand side integrable at infinity. General product laws lack this rotational invariance, so the Gaussian computation does not settle their uniform density bound.

A related Gaussian quadratic estimate appears in \citet[Lemma~2.1]{APSSS2017}: for any nonzero real symmetric matrix $B$, standard Gaussian vector $g$, fixed vector $d$, and scalar $s$,
\[
 \mathbb P\!\left\{
 \frac{\|B(g+d)\|}{|(g+d)^{\mathsf T}B(g+d)-s|}\geq t
 \right\}\leq\frac{C}{t},\qquad t\geq1,
\]
with an absolute constant. Its proof uses Gaussian rotational invariance and a joint characteristic function of quadratic quantities, assumptions that do not hold for general bounded product densities.

A related regularity method of \citet{HMP2025} uses a Gaussian
representation of independent identically distributed input laws and
spectral quantities of a Hessian multiplied by a random diagonal
matrix.  Their results combine inverse moments of the associated
Dirichlet energy with spectral remainders and small influences.
Its smooth representation and quantitative regularity hypotheses
differ from the sole coordinate density bound here.  In particular,
these results do not yield a dimension-uniform estimate for arbitrary
bounded coordinate laws and noncentral forms.

\subsection{Scope of the present results}

The semidefinite theorem below concerns bounded product densities, with an explicit constant controlled by the third spectral scale. The indefinite results require the additional geometric or analytic hypotheses in their statements. The unrestricted dimension-uniform bound remains a conjecture; the cited results do not establish it for arbitrary product laws.

\section{Coordinate geometry and convex input laws}\label{sec:geometry}

The estimates in this section use the original independent coordinates. An orthogonal change of variables in a geometric integral does not imply independence of the rotated random coordinates.

\begin{lemma}[Absolute continuity]\label{geom:ac}
If $A\ne0$ and the coordinates of $X$ have densities, then $X^TAX+b^TX+c$ has an absolutely continuous distribution.
\end{lemma}
\begin{proof}
If $a_{ii}\ne0$, condition on all coordinates except $X_i$. The conditional polynomial is a nonconstant quadratic, and the inverse image of a Lebesgue null set under it is null: split the real line at its vertex and apply change of variables. If every diagonal entry vanishes, choose $a_{ij}\ne0$. Conditional on $X_{-i}$, the coefficient of $X_i$ is $2\sum_{j\ne i}a_{ij}X_j+b_i$. This is nonzero almost surely, since a nontrivial linear combination of independent absolutely continuous random variables is absolutely continuous. The conditional law is therefore absolutely continuous in this case as well. Fubini completes the argument.
\end{proof}

\begin{lemma}[A scalar quadratic]\label{geom:square}
Let $X$ have density $f\le1$, and let $g(x)=a(x-m)^2+\kappa$ with $a\ne0$. On the side $(y-\kappa)/a>0$,
\[
 p_{g(X)}(y)=\frac{f(m+r)+f(m-r)}{2\sqrt{|a|\,|y-\kappa|}}
 \le\frac1{\sqrt{|a|\,|y-\kappa|}},\qquad
 r=\sqrt{|y-\kappa|/|a|}.
\]
The density vanishes on the other side. Moreover, for every $u>0$,
\begin{equation}\label{eq:square-small}
 \sup_s\mathbb P(|g(X)-s|\le u)\le2\sqrt{2u/|a|},
 \qquad
 \mathbb P(|g(X)-s|\le u)
 \le\frac{4u}{\sqrt{|a|\max\{|s-\kappa|,u\}}}.
\end{equation}
\end{lemma}
\begin{proof}
The density formula follows by change of variables on the two monotone half-lines. After translation and scaling, the admissible values of $|x-m|$ form an interval whose squared endpoints differ by at most $2u/|a|$. Its length is at most $\sqrt{2u/|a|}$, which proves the first bound after accounting for the two signs. If $s-\kappa$ lies on the occupied side and $|s-\kappa|>u$, rationalizing the difference of the two square roots gives the second estimate with constant $4$. If $|s-\kappa|\le u$, the first estimate implies it. If $s-\kappa$ is on the unoccupied side at distance greater than $u$, the probability is zero.
\end{proof}

\begin{lemma}[A definite pair]\label{geom:definite}
Suppose that $W\in\mathbb R^2$ has a joint density bounded by $K$. If $B$ is positive or negative definite and $D=|\det B|$, then every translated quadratic with quadratic part $B$ has density bounded by $\pi K/\sqrt D$.
\end{lemma}
\begin{proof}
Changing sign if necessary, write the eigenvalues as $\mu_1,\mu_2>0$. Complete the square. The parametrization
\[
 w=w_0+\sqrt{\delta/\mu_1}\cos\theta\,e_1
       +\sqrt{\delta/\mu_2}\sin\theta\,e_2,
 \qquad \delta>0,\quad0\le\theta<2\pi,
\]
has Jacobian $1/(2\sqrt D)$ in $(\delta,\theta)$. Integrating the density on this ellipse proves the assertion. This lemma requires a joint density bound, but not independence.
\end{proof}

\begin{theorem}[Diagonal forms and a sharp diagonal constant]\label{geom:diagonal}
If $A$ is diagonal and $s_3(A)\ge s>0$, then, for arbitrary affine terms,
\[
 \|p_Q\|_\infty\le\pi/s.
\]
For the normalization $s=1$, the constant $\pi$ is attained by
$Q=X_1^2+X_2^2-X_3^2$ with independent uniform coordinates on $[-1/2,1/2]$.
\end{theorem}
\begin{proof}
Among three diagonal coefficients of modulus at least $s$, two have the same sign. Condition on every other coordinate. The remaining two-variable form is definite, its determinant has modulus at least $s^2$, and its joint density is at most one. Apply Lemma~\ref{geom:definite} and average. For sharpness, the density of $X_1^2+X_2^2$ equals $\pi$ on $(0,1/4)$, since the corresponding circles lie in the square. Averaging at $X_3^2$ gives density $\pi$ at zero, in the continuous limiting sense. The uniform upper bound and dominated convergence justify this evaluation.
\end{proof}

\begin{proposition}[Spectral threshold and a constant lower bound]\label{geom:threshold}
One nonzero eigenvalue does not ensure bounded density, and two eigenvalues of opposite signs do not ensure it. Any universal constant in Conjecture~\ref{conj:main} must satisfy $C\ge3\sqrt2$.
\end{proposition}
\begin{proof}
The square of a centered uniform variable has the inverse-square-root singularity of Lemma~\ref{geom:square}. For independent centered unit uniforms $U,V$, the density of $U^2-V^2$ diverges logarithmically at zero; for $t>0$ its defining integral contains
\[
 \int_t^{1/4}\frac{du}{\sqrt{u(u-t)}}\longrightarrow\infty.
\]
For the lower bound on $C$, let $Z_{j,m}$, $j=1,2,3$, be independent normalized sums of $m$ centered unit uniforms. The three corresponding coefficient vectors in $\mathbb R^{3m}$ are orthonormal. Thus
\[
 Q_m=Z_{1,m}^2+Z_{2,m}^2-Z_{3,m}^2
\]
is admissible with nonzero eigenvalues $1,1,-1$. The central limit theorem and the continuous mapping theorem give weak convergence to $Q_\infty=G_1^2+G_2^2-G_3^2$, where the $G_j$ are independent $N(0,1/12)$ variables. If every $Q_m$ has density at most $C$, the same interval bound passes to the weak limit, hence its density is also at most $C$. Here $G_1^2+G_2^2$ is exponential of rate $6$, and $G_3^2$ has density $\sqrt{6/\pi}\,t^{-1/2}e^{-6t}$ on $t>0$. Their convolution at zero is
\[
 p_{Q_\infty}(0)=6\sqrt{6/\pi}\int_0^\infty t^{-1/2}e^{-12t}\,dt=3\sqrt2.
\]
The density is continuous at zero. This argument uses weak convergence and interval domination, and does not require a local central limit theorem. It does not prove that $3\sqrt2$ is the sharp universal constant.
\end{proof}

\subsection{The logarithmic envelope of an indefinite physical pair}

\begin{theorem}[Uniform hyperbolic envelope]\label{geom:pair}
Let $W_1,W_2$ be independent with densities bounded by one. Let
$q(w)=(w-w_0)^TB(w-w_0)+\kappa$, where $B$ has eigenvalues $\mu_+>0>\mu_-$. Put
\[
 D=|\det B|,\qquad \Lambda=\mu_+^{-1}+|\mu_-|^{-1}.
\]
For almost every $s\ne\kappa$,
\begin{equation}\label{eq:hyperbolic-envelope}
 p_q(s)\le\frac1{\sqrt D}\left(5.2+\log_+\frac{32}{\Lambda|s-\kappa|}\right)
 \le\frac1{\sqrt D}\left(5.2+\log_+\frac{16\sqrt D}{|s-\kappa|}\right).
\end{equation}
Consequently
\begin{equation}\label{eq:hyperbolic-window}
 \mathbb P(|q(W)-s|\le u)
 \le\frac{2u}{\sqrt D}\left(6.2+\log_+\frac{16\sqrt D}{u}\right).
\end{equation}
\end{theorem}
\begin{proof}
Translation preserves the coordinate density caps. For $\delta=s-\kappa>0$, the two branches of the level set are
\[
 w_\epsilon(t)=\epsilon\sqrt{\delta/\mu_+}\cosh t\,e_+
       +\sqrt{\delta/|\mu_-|}\sinh t\,e_-,
 \quad \epsilon\in\{-1,1\},\quad t\in\mathbb R.
\]
The Jacobian in $(\delta,t)$ equals $1/(2\sqrt D)$. Thus
\begin{equation}\label{eq:hyperbolic-coarea}
 p_q(\kappa+\delta)=\frac1{2\sqrt D}\sum_{\epsilon=\pm1}
 \int_{\mathbb R}f_1(w_{\epsilon,1}(t))f_2(w_{\epsilon,2}(t))\,dt.
\end{equation}
Write one branch as $u e^t+v e^{-t}$. Orthogonality of $e_+,e_-$ gives
$|u|=|v|=\rho_0=\sqrt{\delta\Lambda}/2$. Choose a physical coordinate $i$ with $|u_i|\ge\rho_0/\sqrt2$. For $t\ge T_1=(3/4)\log2$, the term $v_i e^{-t}$ has modulus at most $|u_i|e^t/2$. The coordinate $w_i(t)$ is monotone there and satisfies
\[
 |w_i'(t)|\ge\tfrac12|u_i|e^t,\qquad
 \tfrac12|u_i|e^t\le|w_i(t)|\le\tfrac32|u_i|e^t,
 \qquad |w_i'(t)|\ge |w_i(t)|/3.
\]
For $T=\max\{T_1,\log(2/|u_i|)\}$, put $m=|u_i|e^T/2\ge1$. Since $f_1f_2\le f_i$ and $\int f_i=1$, change of variables and the decreasing weight $1/|y|$ give
\[
 \int_T^\infty f_1(w_1(t))f_2(w_2(t))\,dt
 \le3\int_{|y|\ge m,\,\text{one side}}\frac{f_i(y)}{|y|}\,dy
 \le3\log(1+1/m)\le3\log2.
\]
The penultimate inequality follows by filling the interval of length one closest to $m$; it also follows directly from layer cake. The negative tail is identical with $v$ in place of $u$. Its cutoff $T'$ and the first cutoff satisfy
\[
 T,T'\le\tfrac34\log2+\tfrac12\log_+\frac{32}{\delta\Lambda}.
\]
On the middle interval, the integrand is at most one. Each branch integral is therefore at most
\[
 \tfrac32\log2+6\log2+\log_+\frac{32}{\delta\Lambda}
 <5.2+\log_+\frac{32}{\delta\Lambda}.
\]
Substitute into \eqref{eq:hyperbolic-coarea}. Negative $\delta$ exchanges the positive and negative eigenvalues. Since $\Lambda\ge2/\sqrt D$, the second estimate in \eqref{eq:hyperbolic-envelope} follows. Finally, among intervals of length $2u$, the integral of $\log_+(K/|x|)$ is largest for the centered interval, and is at most $2u(1+\log_+(K/u))$. Integration proves \eqref{eq:hyperbolic-window}. All calculations may first be made against a nonnegative output test function, so the almost-everywhere formulas are legitimate under later conditioning.
\end{proof}

\begin{theorem}[A determinant-only bound in arbitrary dimension]\label{geom:fullpair}
Suppose $s_3(A)\ge s>0$ and $A$ has an indefinite principal block $B$ on two original coordinates. For $D=|\det B|$ and arbitrary affine terms,
\begin{equation}\label{eq:full-pair}
 \|p_Q\|_\infty\le D^{-1/2}
 \left(9.2+\log_+\frac{10000\sqrt D}{s}\right).
\end{equation}
There is no rank bound and no restriction on the other entries of $A$.
\end{theorem}
\begin{proof}
Condition on the complementary physical coordinates, and complete the square in the pair. The conditional critical value is a quadratic $R$ whose matrix is the Schur complement $S=A/B$. Lemma~\ref{lem:schur-count} gives $\|S\|\ge s$. Theorem~\ref{four:quarter} and layer cake imply, for every $K>0$,
\[
 \sup_t\mathbb E\log_+\frac K{|t-R|}
 \le\int_0^\infty\min\{1,5(K e^{-u}/\|S\|)^{1/4}\}\,du
 \le4+\log_+\frac{625K}{\|S\|}.
\]
Average \eqref{eq:hyperbolic-envelope} with $K=16\sqrt D$. This proves \eqref{eq:full-pair} by Fubini. Notice that a small eigenvalue of the pair itself need not be controlled separately: its determinant is the only pair parameter in the conclusion.
\end{proof}

\subsection{Convex superlevel sets and interval laws}

\begin{lemma}[The unit ball in three dimensions]\label{geom:ball}
Let $Z$ be uniform on the unit ball of $\mathbb R^3$, and let $q$ have invertible quadratic part $M$. Then
\[
 \|p_{q(Z)}\|_\infty\le C_B|\det M|^{-1/3},\qquad
 C_B=\frac3{4\pi}\,38.2.
\]
\end{lemma}
\begin{proof}
Rotate the ball to diagonalize $M$, and complete squares in the polynomial. Let $|m_1|$ be the largest eigenvalue modulus and $D=|m_2m_3|\le m_1^2$. At a point $(x_2,x_3)$ in the unit disc $\Delta$, the first coordinate ranges along a chord of half-length $h$. The marginal density of $(x_2,x_3)$ is $3h/(2\pi)$, while the conditional first coordinate is uniform on a segment of length $2h$. The chord factors cancel in Lemma~\ref{geom:square}, yielding
\[
 p_q(s)\le\frac3{4\pi}|m_1|^{-1/2}
 \int_\Delta|s-R(x_2,x_3)|^{-1/2}\,dx_2dx_3.
\]
Layer cake bounds the integral by the integral over $v>0$ of the area where $|s-R|<v^{-2}$. For a definite $R$, the area is bounded by $\min\{\pi,2\pi v^{-2}/\sqrt D\}$, whose integral is $2\pi\sqrt2\,D^{-1/4}$. For an indefinite $R$, enclose $\Delta$ in $[-1,1]^2$. The latter Lebesgue measure is four times a product probability measure with coordinate densities $1/2\le1$. Theorem~\ref{geom:pair} gives the bound
\[
 D^{-1/4}\int_0^\infty\min\{\pi,8v^{-2}(6.2+\log_+(16v^2))\}\,dv.
\]
Split this integral at $a=28/5$. It is at most
\[
 \pi a+\frac8a\bigl(8.2+\log16+2\log a\bigr)<38.2.
\]
Thus $p_q\le C_B|m_1|^{-1/2}D^{-1/4}$. Since
\[
 |m_1|^{-1/2}D^{-1/4}
 =|\det M|^{-1/3}(D/m_1^2)^{1/12}\le|\det M|^{-1/3},
\]
the assertion follows.
\end{proof}

\begin{theorem}[Three-dimensional convex superlevel estimate]\label{geom:convex}
Let $Z\in\mathbb R^3$ have density $g\le K$ whose positive superlevel sets are convex. For every quadratic polynomial with invertible quadratic part $M$,
\begin{equation}\label{eq:convex-density}
 \|p_{q(Z)}\|_\infty\le71.1\,|\det M|^{-1/3}K^{2/3}.
\end{equation}
\end{theorem}
\begin{proof}
First consider Lebesgue measure on a convex body $C$. For an ellipsoid $E=c+TB_3$, Lemma~\ref{geom:ball} and $\det(T^TMT)=(\det T)^2\det M$ imply that its pushforward density is at most
\[
 C_B\,\operatorname{vol}(B_3)^{2/3}|\det M|^{-1/3}\operatorname{vol}(E)^{1/3}.
\]
John's ellipsoid theorem gives $E\subset C\subset3E$ (with dilations about the ellipsoid center). Monotonicity of pushforward measures and $\operatorname{vol}(E)\le\operatorname{vol}(C)$ therefore bound the pushforward of $C$ by
\[
 C_K|\det M|^{-1/3}\operatorname{vol}(C)^{1/3},\quad
 C_K=3(4\pi/3)^{2/3}C_B<71.1.
\]
The geometric input is John's theorem \citep{John1948}. Apply this to $C_t=\{g>t\}$ and use layer cake. Since $\int_0^K\operatorname{vol}(C_t)\,dt=1$, H\"older gives
\[
 \int_0^K\operatorname{vol}(C_t)^{1/3}\,dt\le K^{2/3}.
\]
This proves the result. Approximation by full-dimensional level sets handles degenerate or null superlevel sets.
\end{proof}

\begin{theorem}[Bounded rank for unit-interval inputs]\label{geom:interval}
Let $X$ be uniform on a product of unit intervals. Suppose
\[
 A=R\Lambda R^T+R'\Lambda'R'^T,
\]
where $[R\ R']$ has orthonormal columns, $R$ has three columns, $R'$ has $k$ columns, and $\Lambda=\operatorname{diag}(\lambda_1,\lambda_2,\lambda_3)$ is invertible. Decompose $b=R\beta+R'\beta'+\gamma u$, with $u$ a unit vector orthogonal to both ranges if $\gamma\ne0$, and put $k'=k+\mathbf1_{\{\gamma\ne0\}}$. Then
\begin{equation}\label{eq:interval-density}
 \|p_Q\|_\infty\le142.2\,[\sqrt2(2\pi)^{1/3}]^{k'}
 |\lambda_1\lambda_2\lambda_3|^{-1/3}.
\end{equation}
In particular, every exact rank-three form with $s_3(A)\ge1$ and arbitrary affine term has density bounded by $372$ for these inputs.
\end{theorem}
\begin{proof}
Write $Y=R^TX$ and $Y'=(R'^TX,\langle u,X\rangle)$, omitting the last coordinate if $\gamma=0$. The joint density $\rho$ of this orthogonal projection is log-concave by the Pr\'ekopa marginal theorem \citep{Prekopa1973}, and is bounded by $2^{(3+k')/2}$ by the product-projection theorem of \citet{LPP2016}. The polynomial is $q(Y)+\varphi(Y')$, where $q$ has quadratic part $\Lambda$. For almost every $y'$ with marginal density $\rho_{Y'}(y')>0$, the conditional density of $Y$ is log-concave and bounded by $2^{(3+k')/2}/\rho_{Y'}(y')$. Theorem~\ref{geom:convex}, followed by integration in $y'$, gives
\[
 \|p_Q\|_\infty\le71.1\,|\det\Lambda|^{-1/3}
 2^{(3+k')/3}\int\rho_{Y'}(y')^{1/3}\,dy'.
\]
Center $Y'$ at its mean $m$. Hoeffding's bounded-variable exponential estimate \citep{Hoeffding1963} and orthonormality give
$\mathbb E e^{\langle\theta,Y'-m\rangle}\le e^{|\theta|^2/8}$.
Introducing an independent standard Gaussian in $\mathbb R^{k'}$ yields
$\mathbb E e^{|Y'-m|^2}\le2^{k'/2}$.
H\"older, applied after multiplication and division by $e^{|y'-m|^2/3}$, now gives
\[
 \int\rho_{Y'}^{1/3}\le
 \left(\int\rho_{Y'}e^{|y'-m|^2}\right)^{1/3}
 \left(\int e^{-|y'-m|^2/2}\,dy'\right)^{2/3}
 \le2^{k'/6}(2\pi)^{k'/3}.
\]
Combining the powers of two proves \eqref{eq:interval-density}. The factor $\sqrt2(2\pi)^{1/3}$ is less than $2.61$; when $k=0$, one has $k'\le1$ and the stated constant follows.
\end{proof}

\begin{proposition}[Reductions in the class of coordinate laws]\label{geom:laws}
For fixed $A$, let $\mathcal D_A(f_1,\ldots,f_n)$ denote the supremum of the density norm over all affine terms. This functional is translation invariant, convex in each coordinate law, and nonincreasing when one coordinate law is convolved with an arbitrary probability measure. The supremum over all densities $f_i\le1$ equals the supremum over $f_i=\mathbf1_{E_i}$ with $|E_i|=1$. A theorem uniform over a matrix class stable under diagonal congruences $A\mapsto DAD$, $D\ge I$, transfers from unit-interval laws to independent mixtures of widened unit-interval laws and their convolutions, with the same constant.
\end{proposition}
\begin{proof}
Translation changes only the affine coefficients. For fixed output interval and fixed affine term, its probability is linear in each $f_i$ separately. Taking suprema gives convexity. A convolution is a mixture of translations, giving the contraction assertion. To obtain the extremal-set statement, fix an output interval and all laws but one; its probability is $\int f_i(x)g_i(x)\,dx$ with $0\le g_i\le1$. The bathtub principle says its supremum over $0\le f_i\le1$, $\int f_i=1$, is achieved or approximated arbitrarily well by the indicator of a measurable set of measure one. Replace the coordinates successively and then take the supremum over intervals and affine terms. No interval shape for $E_i$ follows from this argument.

Finally, widening the $i$th interval by $r_i\ge1$ replaces $A$ by $DAD$, $D=\operatorname{diag}(r_i)$. The number of positive eigenvalues at least one cannot decrease: on a subspace where $x^TAx\ge|x|^2$, the transformed subspace $D^{-1}E$ satisfies $y^TDADy\ge|Dy|^2\ge|y|^2$. Apply the same argument to $-A$ for negative eigenvalues. Rank is unchanged. The unit-interval theorem thus applies after widening; convexity and the convolution contraction finish the proof.
\end{proof}

The extremal-set argument above is an elementary special case of the more general extremal-law framework in \citet{MMX2017}. The convexity hypothesis in Theorem~\ref{geom:convex} and the restrictions on the input laws in Theorem~\ref{geom:interval} are substantive. They do not extend those results to arbitrary unions of separated intervals.

\subsection{The ruled side of an exact rank-three form}

\begin{theorem}[A one-sheet bound]\label{geom:ruled}
Let $u_1,u_2,u_3$ be orthonormal and $\lambda_1,\lambda_2,\mu\ge1$. For
\[
 Q=\lambda_1\langle u_1,X\rangle^2+
   \lambda_2\langle u_2,X\rangle^2-\mu\langle u_3,X\rangle^2,
\]
the restriction of the law to $(0,\infty)$ has density bounded by $\pi\sqrt2$. Equivalently, for every interval $J\subset(0,\infty)$,
$\mathbb P(Q\in J)\le\pi\sqrt2\,|J|$.
An affine term in the range of the quadratic matrix is allowed after shifting the critical value. Sign reflection gives the corresponding half-line for signature $(1,2)$.
\end{theorem}
\begin{proof}
First assume smooth compactly supported coordinate densities, bounded by one. Write $Y=(\langle u_1,X\rangle,\langle u_2,X\rangle,\langle u_3,X\rangle)$ and let $\rho$ be its density. Put $D=\operatorname{diag}(\sqrt{\lambda_1},\sqrt{\lambda_2},\sqrt\mu)$ and $\Delta=\lambda_1\lambda_2\mu$. In the coordinates $z=Dy$, the quadratic is $z_1^2+z_2^2-z_3^2$. For $s>0$, a ruling of its level set is
\[
 z(\theta,t)=\sqrt s\,e_\theta+t(e_\theta^\perp+e_3),\qquad
 e_\theta=(\cos\theta,\sin\theta,0),\quad
 e_\theta^\perp=(-\sin\theta,\cos\theta,0).
\]
At fixed $t=z_3$, the planar radius is $\sqrt{s+t^2}$ and the planar polar angle differs from $\theta$ by $\arctan(t/\sqrt s)$. Its angular Jacobian is one. Polar coarea, equivalently change of variables against an output test function, gives
\[
 p_Q(s)=\frac1{2\sqrt\Delta}\int_0^{2\pi}\int_{\mathbb R}
 \rho(D^{-1}z(\theta,t))\,dt\,d\theta.
\]
The line velocity is $v_\theta=D^{-1}(e_\theta^\perp+e_3)$, with
\[
 |v_\theta|^2=\frac{\sin^2\theta}{\lambda_1}+
 \frac{\cos^2\theta}{\lambda_2}+\frac1\mu.
\]
The inner integral is $1/|v_\theta|$ times a two-dimensional orthogonal marginal density of $Y$. It is also a two-dimensional orthogonal projection density of the original product vector $X$, and is at most $2/|v_\theta|$ by \citet{LPP2016}. In this smooth case, the marginal is continuous, so its essential bound is pointwise. Since
\[
 \Delta |v_\theta|^2=\lambda_2\mu\sin^2\theta+
 \lambda_1\mu\cos^2\theta+\lambda_1\lambda_2\ge2,
\]
the claimed bound follows by integration in $\theta$.

To remove regularity, approximate each cap-one density in $L^1$ by smooth compactly supported probability densities with the same cap. Explicitly, truncate it, replace the missing mass by an indicator on a disjoint interval of the required length, and convolve with a nonnegative compactly supported mollifier. A diagonal choice of truncation and smoothing gives the approximation. The product densities converge in $L^1$ by the telescoping product estimate, and their pushforwards under $Q$ converge in total variation. Pass the interval inequality to the limit for each fixed $J\subset(0,\infty)$. This step avoids evaluating an almost-everywhere bounded marginal on a potentially exceptional curve. If $b$ belongs to the range of $A$, a translation in the original coordinates removes $b$ and changes only the critical value; the coordinate caps are preserved.
\end{proof}

This theorem controls one side of the critical value. It supplies no estimate on the two-sheet side and does not accommodate arbitrary affine leakage into the kernel by itself.

\section{Fourier estimates, product projections, and cross blocks}
\label{sec:fourier}

Throughout this section the coordinates of \(X=(X_1,\ldots,X_n)\) are
independent, and their probability densities satisfy \(0\leq f_i\leq1\).
We use the Fourier transform
\[
 \widehat h(\xi)=\int_{\mathbb R^d}e^{i\langle \xi,x\rangle}h(x)\,dx,
 \qquad
 h(x)=(2\pi)^{-d}\int_{\mathbb R^d}
                 e^{-i\langle \xi,x\rangle}\widehat h(\xi)\,d\xi.
\]
Statements about essential suprema do not depend on the choice of a
density representative.  When the characteristic function is integrable,
Fourier inversion supplies a continuous density vanishing at infinity.

\subsection{Classical analytic inputs and an inversion formula}

We use the geometric rank-one Brascamp--Lieb inequality
\citep{BL1976,LPP2016}: if \(v_j\) are unit vectors in \(\mathbb R^d\),
\(c_j>0\), and \(\sum_jc_jv_jv_j^{\mathsf T}=I_d\), then
\begin{equation}\label{four:geometric-bl}
 \int_{\mathbb R^d}\prod_jH_j(\langle v_j,z\rangle)^{c_j}\,dz
 \leq\prod_j\left(\int_{\mathbb R}H_j\right)^{c_j}.
\end{equation}
The functions \(H_j\) are arbitrary nonnegative measurable functions.
We also use Lieb's Gaussian exhaustion theorem \citep{Lieb1990}.
For nonzero row vectors \(b_j\), positive exponents \(c_j\) with
\(\sum_jc_j=d\), the best constant in the corresponding rank-one
Brascamp--Lieb inequality is
\begin{equation}\label{four:gaussian-bl}
 \sup_{\lambda_j>0}
 \left[
 \frac{\prod_j\lambda_j^{c_j}}
      {\det\left(\sum_jc_j\lambda_jb_jb_j^{\mathsf T}\right)}
 \right]^{1/2}.
\end{equation}
The supremum may be infinite in general.  The determinant estimates
below show finiteness for the data considered here.

The sharp Hausdorff--Young and convolution inequalities
\citep{Beckner1975} are used with the following normalization:
\begin{align}
 \|\widehat h\|_{p'}&\leq
 (2\pi)^{1/p'}C_p\|h\|_p,
 &C_p&=p^{1/(2p)}(p')^{-1/(2p')},\quad 1\leq p\leq2,
 \label{four:hy}\\
 \|h_1*h_2\|_r&\leq\frac{C_pC_q}{C_r}
                     \|h_1\|_p\|h_2\|_q,
 &&\frac1p+\frac1q=1+\frac1r .
 \label{four:young}
\end{align}
In the convolution inequality \(1\leq p,q,r\leq\infty\), and
each \(C_t\) is given by the same conjugate-exponent formula.
Endpoint constants have their limiting values.  In particular
\eqref{four:hy} applies to complex functions.

\begin{lemma}[Integrable power envelopes]\label{four:inversion}
Let \(Y\) be real and suppose, for \(a>0\) and \(b>1\), that
\[
 |\varphi_Y(t)|\leq\min\{1,(a/|t|)^b\}.
\]
Then \(Y\) has a continuous density vanishing at infinity and
\begin{equation}\label{four:inversion-bound}
 \|p_Y\|_\infty\leq\frac{ab}{\pi(b-1)}.
\end{equation}
More generally, a bound \(K|t|^{-b}\) can be used with \(a=K^{1/b}\).
\end{lemma}
\begin{proof}
Splitting the integral at \(a\) gives
\[
 \int_{\mathbb R}|\varphi_Y(t)|\,dt
 \leq2a+2a^b\int_a^\infty t^{-b}\,dt
 =\frac{2ab}{b-1}.
\]
The inverse Fourier transform of \(\varphi_Y\) is continuous and vanishes
at infinity, and is the density of its probability law.  Its supremum
is at most \((2\pi)^{-1}\|\varphi_Y\|_1\).
\end{proof}

\subsection{Densities of projections and modulated scalar transforms}

\begin{lemma}[Product projection densities]\label{four:projection}
Let \(U\) be an \(n\times k\) real matrix with
\(U^{\mathsf T}U=I_k\), and put
\(c_i=\|U^{\mathsf T}e_i\|^2\).  Then \(U^{\mathsf T}X\) has a
density satisfying the explicit estimate
\begin{equation}\label{four:projection-elementary}
 \|p_{U^{\mathsf T}X}\|_\infty
 \leq\prod_{i:c_i<1}(1-c_i)^{-(1-c_i)/2}
 \leq e^{k/2}.
\end{equation}
The sharper classical product-marginal theorem \citep{LPP2016} gives
\begin{equation}\label{four:projection-sharp}
 \|p_{U^{\mathsf T}X}\|_\infty\leq2^{k/2}.
\end{equation}
For any rank-\(k\) map \(T:\mathbb R^n\to\mathbb R^k\), consequently,
\begin{equation}\label{four:linear-density}
 \|p_{TX}\|_\infty
 \leq\frac{2^{k/2}}{\sqrt{\det(TT^{\mathsf T})}}
 \leq\frac{e^{k/2}}{\sqrt{\det(TT^{\mathsf T})}}.
\end{equation}
\end{lemma}
\begin{proof}
To prove \eqref{four:projection-elementary}, complete \(U\)
to an orthogonal matrix \((U,W)\), with \(W\) having \(n-k\) columns.
Writing \(w_i^{\mathsf T}\) for its rows and \(b_i=\|w_i\|^2=1-c_i\),
Fubini's theorem gives a marginal density
\[
 p(y)=\int_{\mathbb R^{n-k}}
             \prod_i f_i((Uy)_i+\langle w_i,z\rangle)\,dz.
\]
Terms with \(b_i=0\) are at most one and can be omitted.  The other
vectors \(v_i=w_i/\sqrt{b_i}\) satisfy
\(\sum_i b_iv_iv_i^{\mathsf T}=I_{n-k}\).  Apply
\eqref{four:geometric-bl} to
\[
 H_i(s)=f_i((Uy)_i+\sqrt{b_i}s)^{1/b_i}.
\]
Because \(b_i\leq1\), \(f_i^{1/b_i}\leq f_i\), so
\[
 \left(\int H_i\right)^{b_i}
 \leq b_i^{-b_i/2}.
\]
This proves the product bound uniformly in \(y\).  The entropy
inequality
\begin{equation}\label{four:entropy}
 -(1-x)\log(1-x)\leq x,\qquad 0\leq x\leq1,
\end{equation}
with the endpoint interpreted continuously, shows that this product
is at most \(e^{\sum_i c_i/2}=e^{k/2}\).
If \(k=n\), the density is simply the orthogonal image of the product
density and is at most one.

The product-marginal theorem of \citet[Theorem 1.1]{LPP2016}
gives, for \(k<n\), the bound
\[
 \min\left\{
       \left(\frac{n}{n-k}\right)^{(n-k)/2},\,2^{k/2}
      \right\}.
\]
Its hypotheses hold for the independent product densities with cap
one considered here; see also \citet{RV2015} for a dimension-free
projection theorem.
Finally, write \(T=(TT^{\mathsf T})^{1/2}U^{\mathsf T}\) and apply
the change-of-variables formula to the invertible first factor.
\end{proof}

\begin{lemma}[Modulated scalar Fourier estimate]\label{four:lp}
Let \(f\) be a probability density with \(0\leq f\leq1\), and let
\(\omega\) be measurable with \(|\omega|\leq1\).
For \(h=f\omega\) and every real \(p\geq2\),
\begin{equation}\label{four:lp-sharp}
 \int_{\mathbb R}|\widehat h(\xi)|^p\,d\xi
 \leq \frac{2\pi}{\sqrt p}
       \left(\frac{p}{p-1}\right)^{(p-1)/2}
 \leq\frac{2\pi\sqrt e}{\sqrt p}
 \leq\frac{4\pi}{\sqrt p}.
\end{equation}
In particular the estimate is uniform in every measurable phase
\(\omega(x)=e^{iV(x)}\), and in translations of the Fourier variable.
\end{lemma}
\begin{proof}
Set \(p'=p/(p-1)\).  Since \(|h|\leq f\leq1\) and
\(\int f=1\), we have \(\|h\|_{p'}\leq1\).  Apply
\eqref{four:hy} with input exponent \(p'\), then raise to the
power \(p\):
\[
 \|\widehat h\|_p^p
 \leq2\pi(p')^{p/(2p')}p^{-1/2}
 =2\pi p^{-1/2}
       \left(1+\frac1{p-1}\right)^{(p-1)/2}.
\]
The inequality \(\log(1+x)\leq x\) proves the second bound, and
\(\sqrt e<2\) proves the last one.  Translation invariance of Lebesgue
measure proves the assertion about Fourier shifts.
\end{proof}

\begin{lemma}[Difference densities and Fourier powers]\label{four:scalar-norms}
Let \(f\) be as in Lemma~\ref{four:lp}, put
\(\widetilde f(x)=f(-x)\) and \(g=f*\widetilde f\).  Then
\(g\) is a probability density bounded by one, and, for \(0<c\leq1\),
\begin{equation}\label{four:difference-norm}
 \|g\|_{1/c}\leq
 2^c c^{c/2}(1+c)^{-(1+c)/2}.
\end{equation}
If \(q\geq1\) and \(0<d\leq q/2\), then
\begin{equation}\label{four:fourier-power-norm}
 \big\||\widehat f|^q\big\|_{1/d}
 \leq (2\pi)^d q^{q/2-d}(q-d)^{-(q-d)/2}d^{d/2}.
\end{equation}
\end{lemma}
\begin{proof}
The identity defining \(g\) gives \(\int g=1\) and
\(g(x)\leq\|f\|_\infty\int\widetilde f\leq1\).
Use \eqref{four:young} with both input exponents \(2/(1+c)\)
and output \(1/c\).  Their norms are at most one.  Substitution in
\(C_{2/(1+c)}^2/C_{1/c}\) gives
\(2^cc^{c/2}(1+c)^{-(1+c)/2}\), proving
\eqref{four:difference-norm}.  This remains valid at \(c=1\) by
continuity.
For the second estimate, the Fourier exponent \(q/d\) is at least
two and its conjugate input exponent is \(q/(q-d)\).  Apply
\eqref{four:hy} and raise to the power \(q\).  Since the input
norm is at most one, the resulting constant is
\[
 (2\pi)^d
 \left(\frac q{q-d}\right)^{(q-d)/2}
 \left(\frac qd\right)^{-d/2},
\]
which is \eqref{four:fourier-power-norm}.
\end{proof}

\subsection{Fresnel linearization with separable phases}

\begin{lemma}[Regularized Fresnel formula]\label{four:fresnel}
Let \(\Lambda\) be an invertible symmetric \(r\times r\) real matrix,
\(t\ne0\), and \(\varepsilon>0\).  Then
\begin{align}
 K_{\varepsilon,t}(\xi)
 &:=\int_{\mathbb R^r}
  e^{-\varepsilon|y|^2+it y^{\mathsf T}\Lambda y-i\langle\xi,y\rangle}\,dy
 \nonumber\\
 &=\pi^{r/2}\det(\varepsilon I-it\Lambda)^{-1/2}
       \exp\left[-\frac14
             \xi^{\mathsf T}(\varepsilon I-it\Lambda)^{-1}\xi\right],
 \label{four:fresnel-formula}
\end{align}
where the square root is the product of the branches having positive
real part in the scalar Gaussian factors.  Moreover,
\begin{equation}\label{four:fresnel-kernel}
 |K_{\varepsilon,t}(\xi)|
 \leq\pi^{r/2}|t|^{-r/2}|\det\Lambda|^{-1/2}.
\end{equation}
If \(g,\widehat g\in L^1(\mathbb R^r)\), it follows that
\begin{equation}\label{four:fresnel-integral}
 \left|\int_{\mathbb R^r}e^{it y^{\mathsf T}\Lambda y}g(y)\,dy\right|
 \leq(2\pi)^{-r}\|\widehat g\|_1
       \pi^{r/2}|t|^{-r/2}|\det\Lambda|^{-1/2}.
\end{equation}
\end{lemma}
\begin{proof}
Diagonalize \(\Lambda\) orthogonally.  Each one-dimensional integral
is
\[
 \int_{\mathbb R}e^{-(\varepsilon-it\lambda)y^2-i\eta y}\,dy
 =\sqrt\pi(\varepsilon-it\lambda)^{-1/2}
          e^{-\eta^2/[4(\varepsilon-it\lambda)]}.
\]
The Gaussian identity for a positive real coefficient extends to the
right half-plane by analytic continuation, with domination on compact
subsets; this proves the displayed identity for every \(\varepsilon>0\).
Multiplying the scalar factors proves \eqref{four:fresnel-formula}.
The real part of \((\varepsilon-it\lambda)^{-1}\) is
\(\varepsilon/(\varepsilon^2+t^2\lambda^2)>0\), so the exponential
factors have modulus at most one.  Also
\(|\varepsilon-it\lambda|^{-1/2}\leq|t\lambda|^{-1/2}\).
These observations prove \eqref{four:fresnel-kernel}.

Insert the inverse Fourier transform of \(g\) in the integral with
the Gaussian cutoff.  Fubini is justified by
\(\widehat g\in L^1\) and the integrable Gaussian.  Bound its inner
integral by \eqref{four:fresnel-kernel}.  Finally let
\(\varepsilon\downarrow0\); the left side converges by domination
with \(|g|\).  This proves \eqref{four:fresnel-integral}.
\end{proof}

\begin{theorem}[Diffuse quadratic components and separable phases]
\label{four:diffuse}
Let \(U^{\mathsf T}U=I_r\), let
\(c_i=\|U^{\mathsf T}e_i\|^2\leq1/2\), and let
\(\Lambda\) be invertible and symmetric.  For arbitrary real
measurable functions \(V_i\), finite almost surely, define
\[
 Y=(U^{\mathsf T}X)^{\mathsf T}\Lambda(U^{\mathsf T}X)
          +\sum_iV_i(X_i),\qquad
 B(U)=\prod_{i:c_i>0}(1-c_i)^{-(1-c_i)/2}.
\]
Writing \(\Delta=|\det\Lambda|\), for every \(r\geq1\),
\begin{equation}\label{four:diffuse-decay}
 |\varphi_Y(t)|
 \leq\min\{1,B(U)\pi^{r/2}\Delta^{-1/2}|t|^{-r/2}\},
 \qquad B(U)\leq e^{r/2}.
\end{equation}
If \(r>2\), then \(Y\) has a continuous density vanishing at infinity,
and
\begin{equation}\label{four:diffuse-density}
 \|p_Y\|_\infty
 \leq\frac r{r-2}B(U)^{2/r}\Delta^{-1/r}
 \leq\frac{er}{r-2}\Delta^{-1/r}.
\end{equation}
The same conclusion holds if the indicated representation is available
only for a principal block in the original coordinates of a quadratic
form; all other blocks and couplings, as well as affine terms, may be
arbitrary.
\end{theorem}
\begin{proof}
For fixed \(t\), let \(h_{i,t}(x)=f_i(x)e^{itV_i(x)}\).  Push the
complex product measure \(\prod_i h_{i,t}(x_i)\,dx_i\) forward by
\(x\mapsto U^{\mathsf T}x\).  Integrating over the orthogonal
complement, as in Lemma~\ref{four:projection}, gives an \(L^1\)
density \(g_t\) with \(\|g_t\|_1\leq1\).  Fubini gives
\[
 \widehat g_t(\xi)
 =\prod_i\widehat h_{i,t}(\langle U^{\mathsf T}e_i,\xi\rangle).
\]
Rows with \(c_i=0\) contribute factors of modulus at most one.
For the other rows set \(v_i=U^{\mathsf T}e_i/\sqrt{c_i}\).
Then \(\sum_i c_iv_iv_i^{\mathsf T}=I_r\).
The geometric Brascamp--Lieb inequality applied to
\(H_i(s)=|\widehat h_{i,t}(\sqrt{c_i}s)|^{1/c_i}\) yields
\[
 \|\widehat g_t\|_1
 \leq\prod_{i:c_i>0}
          c_i^{-c_i/2}\|\widehat h_{i,t}\|_{1/c_i}.
\]
The output exponent \(1/c_i\) is at least two.  Sharp
Hausdorff--Young with conjugate exponent \(1/(1-c_i)\) gives
\[
 \|\widehat h_{i,t}\|_{1/c_i}
 \leq(2\pi)^{c_i}
       (1-c_i)^{-(1-c_i)/2}c_i^{c_i/2}.
\]
The \(c_i\) factors cancel, and \(\sum_i c_i=r\), so
\begin{equation}\label{four:complex-projection}
 \|\widehat g_t\|_1\leq(2\pi)^rB(U).
\end{equation}
The entropy inequality \eqref{four:entropy} gives \(B(U)\leq e^{r/2}\).

The characteristic function is exactly
\(\int e^{it y^{\mathsf T}\Lambda y}g_t(y)\,dy\).
Lemma~\ref{four:fresnel} and \eqref{four:complex-projection} prove
\eqref{four:diffuse-decay}; the bound one follows from the probability
law of \(Y\).  Lemma~\ref{four:inversion}, with
\(b=r/2\) and \(a=\pi B(U)^{2/r}\Delta^{-1/r}\), proves
\eqref{four:diffuse-density}.

For the principal-block statement, condition on the complementary
original coordinates.  The remaining coordinates keep their original
independent distributions.  The outside couplings induce only linear
terms in them, and these are absorbed in the functions \(V_i\).
The bound \eqref{four:diffuse-decay} is uniform in all conditioning
values, so averaging preserves it and its integrability.
\end{proof}

\subsection{Decoupling and volume-sampling Brascamp--Lieb inequalities}

\begin{lemma}[A cross-block decoupling inequality]\label{four:decoupling}
Let \(I,J\) be disjoint original coordinate sets, let \(M\) be a real
\(|I|\times|J|\) matrix, and let
\[
 Y=F(X_I)+2X_I^{\mathsf T}MX_J+G(X_J),
\]
where \(F,G\) are arbitrary real measurable functions finite almost
surely.  If \(X'_J\) is an independent copy and \(D=X_J-X'_J\), then
\begin{equation}\label{four:decoupling-bound}
 |\varphi_Y(t)|^2
 \leq\mathbb E\prod_{i\in I}
      |\widehat f_i(2t(MD)_i)|.
\end{equation}
For every \(q\geq1\), consequently,
\begin{equation}\label{four:decoupling-q}
 |\varphi_Y(t)|^{2q}
 \leq\mathbb E\prod_{i\in I}
      |\widehat f_i(2t(MD)_i)|^q .
\end{equation}
Each coordinate of \(D\) has its own difference density and these
coordinates are independent.
\end{lemma}
\begin{proof}
Write the characteristic function by first integrating \(X_J\).
Cauchy--Schwarz in the probability law of \(X_I\), after taking
absolute values of the factor \(e^{itF(X_I)}\), gives
\[
 |\varphi_Y(t)|^2
 \leq\mathbb E_{X_I}
 \left|\mathbb E_{X_J}
   e^{itG(X_J)+2itX_I^{\mathsf T}MX_J}\right|^2.
\]
Expand the square using \(X_J'\).  Integration in \(X_I\) gives the
product of the scalar transforms in \eqref{four:decoupling-bound},
multiplied by \(e^{it(G(X_J)-G(X'_J))}\).
The triangle inequality proves \eqref{four:decoupling-bound}.
Jensen's inequality for the probability law of \(D\) proves
\eqref{four:decoupling-q}.  Independence and the difference-density
description follow directly from the independent coordinate pairs.
\end{proof}

For a real \(m\times d\) matrix \(M\) of rank at least \(k\), write
\begin{equation}\label{four:volume-definition}
 Z_k=e_k(M^{\mathsf T}M)
 =\sum_{\substack{S\subseteq[m],\,T\subseteq[d]\\|S|=|T|=k}}
           \det(M_{S,T})^2>0.
\end{equation}
Here \(e_k\) is the elementary symmetric polynomial of the squared
singular values.  The equality follows from Cauchy--Binet.
The squared-minor probability law and its marginals are
\begin{equation}\label{four:volume-marginals}
 w_{S,T}=\frac{\det(M_{S,T})^2}{Z_k},\qquad
 \alpha_i=\sum_{S\ni i,T}w_{S,T},\qquad
 \beta_j=\sum_{S,T\ni j}w_{S,T}.
\end{equation}
Thus \(0\leq\alpha_i,\beta_j\leq1\) and
\(\sum_i\alpha_i=\sum_j\beta_j=k\).

\begin{lemma}[Volume-sampling Fourier estimate]\label{four:volume}
In the setting of Lemma~\ref{four:decoupling}, suppose \(Z_k>0\).
Let \(q\geq1\) and assume
\begin{equation}\label{four:volume-gate}
 \alpha_i\leq q/2\quad\hbox{for every row }i.
\end{equation}
Then, for \(t\ne0\),
\begin{align}
 |\varphi_Y(t)|^{2q}
 &\leq e^k q^{-k/2}
          \left(\frac{\pi}{|t|}\right)^kZ_k^{-1/2},
 \label{four:volume-sharp-decay}\\
 |\varphi_Y(t)|^{2q}
 &\leq e^{k/2}q^{-k/2}
          \left(\frac{2\pi}{|t|}\right)^kZ_k^{-1/2}.
 \label{four:volume-classical-decay}
\end{align}
The first bound is stronger.  Neither bound requires a condition on
the column marginals.
\end{lemma}
\begin{proof}
Put \(h_i(u)=|\widehat f_i(2tu)|^q\), and let \(g_j\) be the
difference density of \(D_j\).  By \eqref{four:decoupling-q} it
suffices to bound
\[
 \int_{\mathbb R^d}
        \prod_i h_i((Mx)_i)\prod_jg_j(x_j)\,dx.
\]
Use row maps \(x\mapsto(Mx)_i\) with exponents \(\alpha_i\),
and coordinate maps \(x\mapsto x_j\) with exponents
\(1-\beta_j\).  Their sum is \(d\).  Rows with \(\alpha_i=0\)
can be dropped because \(h_i\leq1\).  Coordinates with
\(\beta_j=1\) can be dropped because \(g_j\leq1\).
Every retained row map is nonzero: a zero row occurs in no
nonzero minor.

We first bound the Brascamp--Lieb constant of this datum.  Let
\(a_i>0\) be Gaussian row parameters after absorbing the
\(\alpha_i\) factors, and let \(b_j>0\) be the similarly absorbed
parameters of the active coordinate maps.  Set \(a_i=0\) for
omitted rows and \(b_j=0\) when \(\beta_j=1\).
Expansion of a diagonal-plus-positive-semidefinite determinant,
followed by Cauchy--Binet, gives
\begin{align}
 \det(D_b+M^{\mathsf T}D_aM)
 &=\sum_{T\subseteq[d]}
       \left(\prod_{j\notin T}b_j\right)
       \det\big((M^{\mathsf T}D_aM)_{T,T}\big)
 \nonumber\\
 &=\sum_{\substack{S\subseteq[m],\,T\subseteq[d]\\|S|=|T|}}
       \det(M_{S,T})^2
       \prod_{i\in S}a_i\prod_{j\notin T}b_j .
 \label{four:determinant-expansion}
\end{align}
This is a polynomial identity and is valid when some diagonal
parameters vanish.  Retain the terms with \(|S|=|T|=k\).
Weighted arithmetic--geometric mean with probabilities
\(w_{S,T}\) gives
\begin{equation}\label{four:determinant-amgm}
 \det(D_b+M^{\mathsf T}D_aM)
 \geq Z_k\prod_i a_i^{\alpha_i}
                    \prod_j b_j^{1-\beta_j}.
\end{equation}
Every column with \(\beta_j=1\) belongs to every \(T\) of positive
weight, so its zero parameter never occurs in the product of an
active summand.  The same observation applies to omitted rows.
Thus all active products in the arithmetic--geometric mean step
are positive, and no inverse of a zero parameter is used.

Substitute \eqref{four:determinant-amgm} in
\eqref{four:gaussian-bl}, restoring the absorbed exponents.
The Brascamp--Lieb constant is at most
\begin{equation}\label{four:volume-bl-constant}
 Z_k^{-1/2}\prod_i\alpha_i^{-\alpha_i/2}
                 \prod_j(1-\beta_j)^{-(1-\beta_j)/2}.
\end{equation}
Zero exponents in this formula have value one by continuity.
Consequently the integral under consideration is at most this
constant times
\(\prod_i\|h_i\|_{1/\alpha_i}
  \prod_j\|g_j\|_{1/(1-\beta_j)}\).
The latter column norms are at most one.

The scalar estimate \eqref{four:fourier-power-norm}, including the
change of variable \(2tu\), gives
\begin{equation}\label{four:volume-hy}
 \|h_i\|_{1/\alpha_i}
 \leq\left(\frac{\pi}{|t|}\right)^{\alpha_i}
       \left(\frac{\alpha_i}{q}\right)^{\alpha_i/2}
       \left(1-\frac{\alpha_i}{q}\right)^{-(q-\alpha_i)/2}.
\end{equation}
Its use is legitimate precisely because \(q/\alpha_i\geq2\).
After cancellation of the powers of \(\alpha_i\), the row
entropy factor is at most \(e^{k/2}\), by
\eqref{four:entropy} applied to \(x=\alpha_i/q\).
The column entropy factor in \eqref{four:volume-bl-constant}
is also at most \(e^{k/2}\).  Since \(\sum_i\alpha_i=k\),
this proves \eqref{four:volume-sharp-decay}.

Alternatively, use only the last bound in
\eqref{four:lp-sharp} to obtain
\[
 \|h_i\|_{1/\alpha_i}
 \leq\left(\frac{2\pi}{|t|}\right)^{\alpha_i}
       \left(\frac{\alpha_i}{q}\right)^{\alpha_i/2}.
\]
The same cancellation, followed by just the column entropy
bound, proves \eqref{four:volume-classical-decay}.
The stronger estimate follows also numerically from
\(e<2\sqrt e\).
\end{proof}

\begin{theorem}[General volume density criterion]\label{four:volume-density}
Under the hypotheses of Lemma~\ref{four:volume}, suppose in addition
that \(k>2q\).  Then \(Y\) has a continuous density vanishing at
infinity and
\begin{equation}\label{four:volume-density-bound}
 \|p_Y\|_\infty
 \leq\frac{ke}{(k-2q)\sqrt q}\,Z_k^{-1/(2k)}.
\end{equation}
Using the second estimate of that lemma gives the also valid bound
\begin{equation}\label{four:volume-density-classical}
 \|p_Y\|_\infty
 \leq\frac{2k\sqrt e}{(k-2q)\sqrt q}\,Z_k^{-1/(2k)}.
\end{equation}
Both conclusions apply to a cross block \(M=A_{I,J}\) in a larger
quadratic form with arbitrary self blocks, outside couplings, and
affine terms.
\end{theorem}
\begin{proof}
The first Fourier envelope has exponent \(b=k/(2q)>1\).
Its crossover scale is
\(a=e\pi q^{-1/2}Z_k^{-1/(2k)}\).
Lemma~\ref{four:inversion} proves
\eqref{four:volume-density-bound}.  The second crossover scale is
\(2\pi\sqrt e\,q^{-1/2}Z_k^{-1/(2k)}\), giving
\eqref{four:volume-density-classical}.
For a larger form, condition on the coordinates outside \(I\cup J\).
They change only the functions \(F,G\) and an additive constant.
The Fourier envelopes remain uniform and can therefore be averaged.
\end{proof}

\begin{corollary}[Three-volume and five-volume cases]\label{four:d-cases}
The following explicit cases of Theorem~\ref{four:volume-density}
hold:
\begin{enumerate}
\item If \(Z_3>0\) and \(\max_i\alpha_i^{(3)}\leq1/2\), then
\[
 \|p_Y\|_\infty\leq3e\,Z_3^{-1/6}
                         \leq6\sqrt e\,Z_3^{-1/6}.
\]
\item If \(Z_3>0\) and
\(\max_i\alpha_i^{(3)}\leq h\) with \(1/2\leq h<3/4\), then
\[
 \|p_Y\|_\infty
 \leq\frac{3e}{(3-4h)\sqrt{2h}}\,Z_3^{-1/6}
 \leq\frac{6\sqrt e}{(3-4h)\sqrt{2h}}\,Z_3^{-1/6}.
\]
\item If \(Z_5>0\), with no marginal hypothesis, then
\[
 \|p_Y\|_\infty\leq\frac{5e}{\sqrt2}\,Z_5^{-1/10}
                         \leq5\sqrt{2e}\,Z_5^{-1/10}.
\]
\end{enumerate}
These are respectively the three-volume, fractional three-volume,
and unconditional five-volume criteria.
\end{corollary}
\begin{proof}
Take \((k,q)=(3,1),(3,2h),(5,2)\), respectively.  In the last
case every inclusion probability is at most one, so the gate
\(\alpha_i^{(5)}\leq q/2\) is automatic.  The strict inequality
\(k>2q\) holds in all three cases.
\end{proof}

\begin{remark}[Choice of a minor law]\label{four:weighted-volume}
The canonical squared-minor probabilities can be replaced by any
probability law \(v_{S,T}\) supported on nonzero \(k\times k\)
minors.  If its row marginals satisfy the same gate, the bound
\eqref{four:volume-density-bound} acquires the factor
\[
 \exp\!\left(\frac{D(v\Vert w)}{2k}\right),\qquad
 D(v\Vert w)=\sum_{S,T}v_{S,T}\log\frac{v_{S,T}}{w_{S,T}}.
\]
Indeed, apply weighted arithmetic--geometric mean to
\(\sum v_{S,T}
 [\det(M_{S,T})^2/v_{S,T}]
 \prod_{i\in S}a_i\prod_{j\notin T}b_j\).
The determinant lower bound has \(Z_k e^{-D(v\Vert w)}\) in
place of \(Z_k\).  All subsequent steps are identical.  This
form allows the marginal constraint to be traded against an
explicit entropy cost.
\end{remark}

\subsection{Exact marginal factors and entropy optimization}

\begin{theorem}[Exact weighted-volume factors]\label{four:exact-marginal}
Let \(w_{S,T}\) be any probability law on the nonzero
\(k\times k\) minors of a physical cross block \(M\), and let
\(\alpha_i,\beta_j\) be its row and column inclusion probabilities.
Set
\[
 G(w)=\exp\left(\sum_{S,T}w_{S,T}
                 \log\frac{\det(M_{S,T})^2}{w_{S,T}}\right).
\]
For \(q\geq1\) with \(\alpha_i\leq q/2\), define
\[
 R_q(w)=
 \prod_i\left(1-\frac{\alpha_i}{q}\right)^{-(q-\alpha_i)/2}
 \prod_j(1-\beta_j)^{-(1-\beta_j)/2}.
\]
All zero-exponent factors have their continuous value one.
Then, in the setting of Lemma~\ref{four:decoupling},
\begin{equation}\label{four:exact-marginal-decay}
 |\varphi_Y(t)|^{2q}
 \leq R_q(w)q^{-k/2}\pi^kG(w)^{-1/2}|t|^{-k}.
\end{equation}
If \(k>2q\), it follows that
\begin{equation}\label{four:exact-marginal-density}
 \|p_Y\|_\infty
 \leq\frac{k}{(k-2q)\sqrt q}\,
                  R_q(w)^{1/k}G(w)^{-1/(2k)}.
\end{equation}
The density is continuous and vanishes at infinity.
In particular, for \(k=3,q=1\), writing \(R=R_1\),
\begin{equation}\label{four:q1-exact}
 \|p_Y\|_\infty\leq3R(w)^{1/3}G(w)^{-1/6},\qquad
 \|\varphi_Y\|_1\leq6\pi R(w)^{1/3}G(w)^{-1/6}.
\end{equation}
\end{theorem}
\begin{proof}
In the determinant expansion
\eqref{four:determinant-expansion}, apply arithmetic--geometric
mean with the law \(w\) to obtain
\[
 \det(D_b+M^{\mathsf T}D_aM)
 \geq G(w)\prod_i a_i^{\alpha_i}
                   \prod_jb_j^{1-\beta_j}.
\]
The argument for saturated marginals in the proof of
Lemma~\ref{four:volume} applies without change.
Thus its Brascamp--Lieb constant is
\eqref{four:volume-bl-constant} with \(G(w)\) in place of \(Z_k\).
Use the exact scalar estimate \eqref{four:volume-hy} and leave
the row and column factors unrelaxed.
The powers of \(\alpha_i\) cancel, and
\(\sum_i\alpha_i=k\); the result is
\eqref{four:exact-marginal-decay}.
Its exponent is \(k/(2q)\) and its crossover scale is
\(\pi q^{-1/2}R_q(w)^{1/k}G(w)^{-1/(2k)}\).
Lemma~\ref{four:inversion} proves
\eqref{four:exact-marginal-density} and the last assertion.
\end{proof}

The entropy inequality gives \(R_q(w)\leq e^k\), recovering
the coefficient in Lemma~\ref{four:volume}.  For the canonical squared-minor law
\(p_{S,T}=\det(M_{S,T})^2/Z_k\), the identity
\[
 G(w)=Z_k e^{-D(w\Vert p)}
\]
shows that \(p\) maximizes \(G\) without marginal constraints.
It does not assert that \(p\) minimizes the whole expression
in \eqref{four:exact-marginal-density}; the marginal factor
depends on the law as well.

\begin{proposition}[Balanced six-row and six-column blocks]
\label{four:balanced-six}
If a law on \(3\times3\) minors of a six-row, six-column
physical cross block has every row and column marginal \(1/2\),
then
\[
 \|p_Y\|_\infty\leq6G(w)^{-1/6}.
\]
For example, let \(u=6^{-1/2}(1,\ldots,1)^{\mathsf T}\) and
\[
 M=sI_6+(L-s)uu^{\mathsf T},\qquad L>s>0.
\]
With arbitrary self functions on the two banks,
\begin{equation}\label{four:balanced-six-bound}
 \|p_Y\|_\infty
 \leq6[10s^4(L^2+s^2)]^{-1/6}.
\end{equation}
This example need not satisfy the sufficient projection condition
of Corollary~\ref{four:tail}.
\end{proposition}
\begin{proof}
Each of the twelve marginal factors equals \(2^{1/4}\),
so \(R(w)=8\).  Formula \eqref{four:q1-exact} proves the
first assertion.
The displayed matrix has singular values \(L,s,s,s,s,s\).
Simultaneous coordinate-permutation symmetry makes its canonical
row and column marginal vectors constant; since each sums to
three, all entries are \(1/2\).  Moreover
\[
 Z_3=\binom52 L^2s^4+\binom53s^6
     =10s^4(L^2+s^2).
\]
The canonical law has \(G=Z_3\), proving the bound.
For the final assertion, any projection with all diagonal entries
at most \(1/2\) has rank at most three and is therefore proper.
For a unit \(v\) in its orthogonal complement,
\(\|M^{\mathsf T}v\|\geq s\); hence
\(\|(I-P)M\|_{\rm op}\geq s\).
Here \(T_3=4s^2\), so the sufficient requirement
\(\|(I-P)M\|\leq\sqrt{T_3}/15=2s/15\) fails.
\end{proof}

We next optimize the determinant entropy subject to a row
marginal constraint.  All optimization sets are finite
dimensional, and all their coordinates refer to original rows.
For a row \(k\)-set \(S\), define
\[
 a_S=\det(M_SM_S^{\mathsf T})
     =\sum_{|T|=k}\det(M_{S,T})^2,\qquad
 \mathcal B=\{S:a_S>0\},\qquad Z_k=\sum_{S\in\mathcal B}a_S.
\]
For a probability law \(\mu\) on \(\mathcal B\), let
\[
 \alpha_i(\mu)=\sum_{S\ni i}\mu_S,\qquad
 G(\mu)=\exp\left(\sum_S\mu_S\log(a_S/\mu_S)\right).
\]
Terms with \(\mu_S=0\) are interpreted by continuity.

\begin{lemma}[Reduction to row entropy]\label{four:row-entropy}
For fixed row law \(\mu\), the lift
\begin{equation}\label{four:canonical-lift}
 w_{S,T}=\mu_S\frac{\det(M_{S,T})^2}{a_S}
\end{equation}
has the same row marginals and maximizes \(G(w)\) among all
minor laws with row law \(\mu\).  Its value is \(G(\mu)\).
If \(p_S=a_S/Z_k\), then
\begin{equation}\label{four:row-kl}
 G(\mu)=Z_k e^{-D(\mu\Vert p)}.
\end{equation}
\end{lemma}
\begin{proof}
For a general lift let \(\nu_S\) be the conditional column
law given \(S\), and let
\(\pi_S(T)=\det(M_{S,T})^2/a_S\).
Expanding logarithms gives
\[
 \log G(w)
 =\sum_S\mu_S\log(a_S/\mu_S)
                  -\sum_S\mu_SD(\nu_S\Vert\pi_S).
\]
Relative entropy is nonnegative, as follows from
\(\log x\leq x-1\), with equality for identical laws.
Thus the maximum is attained by \(\nu_S=\pi_S\), which is
\eqref{four:canonical-lift}.  Expanding
\(D(\mu\Vert p)\) proves \eqref{four:row-kl}.
\end{proof}

\begin{theorem}[Gibbs dual and concentration on a bounded row set]
\label{four:gibbs}
Assume \(Z_k>0\).  For \(0<h<1\), let
\[
 \mathcal F_h=\{\mu:\alpha_i(\mu)\leq h\text{ for every }i\},
 \qquad G_h(M)=\max_{\mu\in\mathcal F_h}G(\mu),
\]
with \(G_h(M)=0\) if \(\mathcal F_h\) is empty.  Then
\begin{equation}\label{four:gibbs-dual}
 G_h(M)=\inf_{\lambda_i\geq0}
  \exp\left(h\sum_i\lambda_i\right)
       \sum_{S\in\mathcal B}a_S
                    \exp\left(-\sum_{i\in S}\lambda_i\right).
\end{equation}
If a finite vector \(\lambda\) attains the infimum, the
corresponding maximizing law is
\begin{equation}\label{four:gibbs-law}
 \mu_\lambda(S)=
 \frac{a_S e^{-\sum_{i\in S}\lambda_i}}
      {\sum_Ra_R e^{-\sum_{i\in R}\lambda_i}},
\end{equation}
and \(\lambda_i>0\) implies \(\alpha_i(\mu_\lambda)=h\).
Regardless of dual attainment, there exists a set \(H\) of rows
such that
\begin{equation}\label{four:gibbs-delete}
 |H|<k/h,\qquad
 e_k(M_{[m]\setminus H}^{\mathsf T}M_{[m]\setminus H})
                         \leq G_h(M).
\end{equation}
\end{theorem}
\begin{proof}
Let \(\Delta\) be the probability simplex on \(\mathcal B\),
and define its continuous concave entropy function
\[
 F(\mu)=\sum_S\mu_S\log(a_S/\mu_S).
\]
For \(\lambda_i\geq0\), put
\[
 \mathcal L(\mu,\lambda)
 =F(\mu)+\sum_i\lambda_i(h-\alpha_i(\mu)).
\]
The Gibbs variational identity states
\begin{equation}\label{four:gibbs-variational}
 \max_{\mu\in\Delta}\mathcal L(\mu,\lambda)
 =h\sum_i\lambda_i+
         \log\sum_S a_S e^{-\sum_{i\in S}\lambda_i}.
\end{equation}
Indeed, subtracting the right side from
\(\mathcal L(\mu,\lambda)\) gives
\(-D(\mu\Vert\mu_\lambda)\).

We justify the dual equality also when constraints are attained
only on the boundary or are infeasible.
Restrict first to \(0\leq\lambda_i\leq N\).
The compact convex minimax theorem \citep{Sion1958}, applied to
the continuous concave-affine function \(\mathcal L\), gives
\[
 \min_{\lambda\in[0,N]^m}\max_{\mu\in\Delta}\mathcal L(\mu,\lambda)
 =\max_{\mu\in\Delta}
   \left[F(\mu)-N\sum_i(\alpha_i(\mu)-h)_+\right].
\]
If \(\mathcal F_h\ne\varnothing\), the right side decreases,
as \(N\) increases, to \(\max_{\mathcal F_h}F\).
To see this explicitly, take maximizing laws along a sequence
\(N\to\infty\) and a convergent subsequence in the compact
simplex.  Since \(F\) is bounded and a feasible law gives a
fixed finite lower bound, their constraint violations tend
to zero.  The limiting law is feasible.  Dropping the
nonpositive penalty and using continuity of \(F\) proves that
the limiting value is at most the constrained maximum; the
reverse inequality holds for every \(N\).
If \(\mathcal F_h\) is empty, the continuous sum of positive
constraint violations has a strictly positive minimum on
\(\Delta\), so these maxima tend to \(-\infty\).
Taking \(N\to\infty\) and using
\eqref{four:gibbs-variational} proves the logarithm of
\eqref{four:gibbs-dual}, including the value \(-\infty\)
in the infeasible case.

If the dual minimizer is finite, differentiation of its
logarithmic objective gives \(h-\alpha_i(\mu_\lambda)\).
Optimality on the nonnegative orthant therefore implies
\(\alpha_i(\mu_\lambda)\leq h\), and equality if
\(\lambda_i>0\).  Substitution in \(\mathcal L\) shows that
this Gibbs law attains the primal maximum.

It remains to prove \eqref{four:gibbs-delete}.
If \(G_h=Z_k\), take \(H=\varnothing\).
Otherwise choose, for arbitrarily small \(\varepsilon>0\),
a nonnegative vector \(\lambda\) whose dual objective is
at most \(G_h+\varepsilon<Z_k\).
Then \(\Lambda=\sum_i\lambda_i>0\).
Put \(\omega_i=\lambda_i/\Lambda\) and
\[
 H=\{i:\omega_i>h/k\}.
\]
Since \(\sum_i\omega_i=1\), necessarily \(|H|<k/h\).
For every \(k\)-set \(S\) avoiding \(H\),
\(\sum_{i\in S}\lambda_i\leq h\Lambda\).
The corresponding term in the dual objective is therefore
at least \(a_S\).  Thus
\[
 \sum_{S:S\cap H=\varnothing}a_S\leq G_h+\varepsilon.
\]
The left side is the residual \(k\)-volume in
\eqref{four:gibbs-delete}.  There are only finitely many
possible sets \(H\); along a sequence
\(\varepsilon\downarrow0\), one repeats infinitely often.
Passing to that subsequence proves the desired bound,
including when \(G_h=0\).
\end{proof}

\begin{corollary}[A three-volume or bounded-row alternative]
\label{four:gibbs-alternative}
Let \(1/2\leq h<3/4\) and \(g>0\).
For a physical cross block \(M\) with \(Z_3>0\), either
\[
 \|p_Y\|_\infty
 \leq\frac{3e}{(3-4h)\sqrt{2h}}\,g^{-1/6},
\]
with a continuous density vanishing at infinity, or there is
a row set \(H\) with \(|H|<3/h\) for which the residual
third volume is less than \(g\).
In particular \(h=3/5\) requires at most four rows in the
second alternative.
\end{corollary}
\begin{proof}
If \(G_h\geq g\), compactness gives an optimizing feasible law.
Lift it by Lemma~\ref{four:row-entropy} and apply
Theorem~\ref{four:exact-marginal} with \(q=2h\) and
\(R_q\leq e^3\).  If \(G_h<g\), use
\eqref{four:gibbs-delete}.  At \(h=3/5\), the strict bound
\(|H|<5\) is equivalent to \(|H|\leq4\).
\end{proof}

\begin{proposition}[A single active row constraint]
\label{four:one-hub-entropy}
Fix a distinguished row \(i_0\), and write
\[
 Z_1=\sum_{S\ni i_0}a_S,\qquad
 Z_0=\sum_{S\not\ni i_0}a_S.
\]
Suppose \(Z_0,Z_1>0\), \(p=Z_1/(Z_0+Z_1)>h\), and the law
\[
 \mu(S)=
 \begin{cases}
  h\,a_S/Z_1,&i_0\in S,\\
  (1-h)a_S/Z_0,&i_0\notin S
 \end{cases}
\]
satisfies all other row constraints.  Then
\begin{equation}\label{four:one-hub-formula}
 G_h(M)=\left(\frac{Z_1}{h}\right)^h
         \left(\frac{Z_0}{1-h}\right)^{1-h}
       =(Z_0+Z_1)e^{-D_{\rm Ber}(h\Vert p)}.
\end{equation}
A dual optimizer has all multipliers zero except
\[
 e^{-\lambda_{i_0}}=\frac{hZ_0}{(1-h)Z_1}.
\]
\end{proposition}
\begin{proof}
If a row law gives total mass \(x\) to sets containing \(i_0\),
the entropy chain rule bounds its logarithmic value by
\[
 x\log(Z_1/x)+(1-x)\log(Z_0/(1-x)),
\]
with equality exactly when the conditional laws in the two
classes are canonical.  The derivative of this expression is
\(\log[Z_1(1-x)/(Z_0x)]\), which is positive for \(x<p\).
Under \(x\leq h<p\), its maximum is therefore attained at \(h\).
The displayed law attains this value and, by assumption, is
feasible for every other constraint.  This proves
\eqref{four:one-hub-formula}.  Substitution in
\eqref{four:gibbs-law} verifies the stated multiplier and
complementary slackness.
\end{proof}

\subsection{Regularized determinants and ridge leverage}

\begin{lemma}[A graph inequality with a free regularizer]
\label{four:graph}
Let \(R\) be a real \(m\times k\) matrix with columns \(r_j\).
Let \(D_1,\ldots,D_m\) be independent differences of independent
copies of variables with densities bounded by one, and let
\(f_1,\ldots,f_k\) be further probability densities bounded by one.
Fix \(\alpha>0\), \(a\in\mathbb R\), and \(1\leq q\leq2\).  Put
\[
 H=\alpha I_m+a^2RR^{\mathsf T},\qquad
 c_i=\alpha e_i^{\mathsf T}H^{-1}e_i,\qquad
 d_j=a^2r_j^{\mathsf T}H^{-1}r_j,\qquad
 \theta=\sum_jd_j=\sum_i(1-c_i).
\]
If \(d_j\leq q/2\) for every \(j\), then, with
\(L_q=\pi e\sqrt{2/q}\),
\begin{equation}\label{four:graph-bound}
 \mathbb E\prod_j|\widehat f_j(a\langle r_j,D\rangle)|^q
 \leq
 \det(I_m+a^2RR^{\mathsf T}/\alpha)^{-1/2}
          (L_q/\sqrt\alpha)^\theta .
\end{equation}
\end{lemma}
\begin{proof}
For \(a=0\), both the expectation and the determinant are one,
and \(\theta=0\).  Assume \(a\ne0\).
The vectors
\[
 u_i=\sqrt\alpha H^{-1/2}e_i,\qquad
 v_j=aH^{-1/2}r_j
\]
satisfy
\(\sum_i u_iu_i^{\mathsf T}+\sum_jv_jv_j^{\mathsf T}=I_m\).
Their squared lengths are \(c_i,d_j\), respectively; thus
\(0<c_i\leq1\), \(0\leq d_j\leq1\).
The trace identity gives the displayed expression for \(\theta\).
Zero \(v_j\) may be omitted.

Let \(g_i\) be the density of \(D_i\), and put
\(h_j=|\widehat f_j|^q\).  Change variables \(x=H^{-1/2}z\)
in the expectation integral and normalize the nonzero vectors
\(u_i,v_j\) to unit vectors.  The geometric Brascamp--Lieb
inequality gives
\begin{equation}\label{four:graph-bl}
 \det(H)^{-1/2}
 \prod_i\left[\alpha^{c_i/2}c_i^{-c_i/2}\|g_i\|_{1/c_i}\right]
 \prod_j\left[d_j^{-d_j/2}\|h_j\|_{1/d_j}\right].
\end{equation}
For example, the argument of \(g_i\) is
\(\sqrt{c_i/\alpha}\) times the relevant unit linear form, which
accounts for its factor \(\alpha^{c_i/2}c_i^{-c_i/2}\).

Define
\[
 G(c)=2^c(1+c)^{-(1+c)/2},\qquad
 H_q(d)=q^{q/2-d}(q-d)^{-(q-d)/2}.
\]
Lemma~\ref{four:scalar-norms} cancels the powers of \(c_i,d_j\)
in \eqref{four:graph-bl}, leaving
\begin{equation}\label{four:graph-after-scalar}
 \det(H)^{-1/2}\alpha^{(m-\theta)/2}(2\pi)^\theta
                     \prod_iG(c_i)\prod_jH_q(d_j).
\end{equation}
The logarithms of these scalar functions are concave:
\[
 (\log G)''(c)=-\frac1{2(1+c)},\qquad
 (\log H_q)''(d)=-\frac1{2(q-d)}.
\]
The tangent to \(\log G\) at \(c=1\), and the tangent to
\(\log H_q\) at \(d=0\), therefore give
\[
 \log G(c)\leq\frac{1-\log2}{2}(1-c),\qquad
 \log H_q(d)\leq\frac{1-\log q}{2}d.
\]
Their endpoint values are \(G(1)=H_q(0)=1\).
The product in \eqref{four:graph-after-scalar} is consequently
at most \(\exp\{[1-\tfrac12\log(2q)]\theta\}\).
Factor \(\alpha^m\) from \(\det H\).  The remaining scalar
multiplier is
\[
 \left(\frac{2\pi e}{\sqrt{2q\alpha}}\right)^\theta
 =(L_q/\sqrt\alpha)^\theta,
\]
which proves \eqref{four:graph-bound}.
\end{proof}

\begin{theorem}[Regularized cross-block determinant]
\label{four:regularized}
Let \(Y\) have the form in Lemma~\ref{four:decoupling}, with cross
block \(B\).  Let \(1\leq q\leq2\).  Suppose that either \(q=2\)
or the projection \(P_{\operatorname{range}(B)}\), acting on the
row-index space, satisfies
\[
 \max_i(P_{\operatorname{range}(B)})_{ii}\leq q/2.
\]
Then
\begin{equation}\label{four:regularized-bound}
 |\varphi_Y(t)|
 \leq\det\left(I+\frac{2qt^2}{\pi^2e^2}BB^{\mathsf T}\right)^{-1/(4q)}.
\end{equation}
In particular, without any leverage hypothesis,
\begin{equation}\label{four:unrestricted-determinant}
 |\varphi_Y(t)|
 \leq\det\left(I+\frac{4t^2}{\pi^2e^2}BB^{\mathsf T}\right)^{-1/8}.
\end{equation}
If \(B\) has at least \(r>2q\) singular values at least \(\sigma>0\),
the indicated hypothesis implies a continuous density with
\begin{align}
 \|p_Y\|_\infty
 &\leq
 \frac{e\sqrt\pi}{2\sqrt{2q}\,\sigma}
 \frac{\Gamma(r/(4q)-1/2)}{\Gamma(r/(4q))}
 \label{four:regularized-gamma}\\
 &\leq\frac{e}{\sqrt{2q}\,\sigma}\frac r{r-2q}.
 \label{four:regularized-simple}
\end{align}
All these statements allow arbitrary functions in the two self blocks,
and conditioning on additional original coordinates.
\end{theorem}
\begin{proof}
Apply \eqref{four:decoupling-q} and then Lemma~\ref{four:graph}
with \(R=B^{\mathsf T}\), \(a=2t\), and
\(\alpha=L_q^2\).  The graph weights on the Fourier factors are
the diagonal entries of
\[
 (2t)^2B(L_q^2I+(2t)^2B^{\mathsf T}B)^{-1}B^{\mathsf T}.
\]
This matrix is bounded between zero and
\(P_{\operatorname{range}(B)}\), as is seen by singular-value
decomposition.  Its diagonal entries are always at most one,
so \(q=2\) needs no assumption.  The multiplier in
\eqref{four:graph-bound} is one for this choice of \(\alpha\).
Taking the \(2q\)-th root and using equality of the nonzero spectra
of \(BB^{\mathsf T}\) and \(B^{\mathsf T}B\) proves
\eqref{four:regularized-bound}.

Discarding determinant factors other than the \(r\) indicated ones
gives
\[
 |\varphi_Y(t)|
 \leq\left(1+\frac{2q\sigma^2t^2}{\pi^2e^2}\right)^{-r/(4q)}.
\]
The integral
\[
 \int_{\mathbb R}(1+x^2)^{-a}\,dx
 =\sqrt\pi\,\frac{\Gamma(a-1/2)}{\Gamma(a)},\qquad a>1/2,
\]
follows from \(u=x^2\) and the beta integral.  Fourier inversion
now proves \eqref{four:regularized-gamma}.
Bounding the same envelope by
\(\min\{1,(\pi e/(\sqrt{2q}\sigma|t|))^{r/(2q)}\}\)
and applying Lemma~\ref{four:inversion} proves the simpler bound.
The uniformity in self functions and outside coordinates is
inherited from Lemma~\ref{four:decoupling}.
\end{proof}

\begin{theorem}[Ridge-leverage criterion]\label{four:ridge}
For a real cross block \(B\) and \(\gamma>0\), define
\[
 K_\gamma=BB^{\mathsf T}(\gamma^2I+BB^{\mathsf T})^{-1},\quad
 \theta_\gamma=\operatorname{tr}K_\gamma,\quad
 D_\gamma=\det(I+BB^{\mathsf T}/\gamma^2).
\]
Suppose, for some \(1\leq q\leq2\), that
\begin{equation}\label{four:ridge-gate}
 \max_i(K_\gamma)_{ii}\leq q/2,\qquad
 \theta_\gamma>2q.
\end{equation}
Then \(Y\) in Lemma~\ref{four:decoupling} has a continuous density
vanishing at infinity and
\begin{align}
 \|p_Y\|_\infty
 &\leq\frac{e}{\sqrt{2q}\,\gamma}
       \frac{\theta_\gamma}{\theta_\gamma-2q}
       D_\gamma^{-1/(2\theta_\gamma)}
 \label{four:ridge-density}\\
 &\leq\frac{\sqrt e}{\sqrt{2q}\,\gamma}
       \frac{\theta_\gamma}{\theta_\gamma-2q}.
 \label{four:ridge-density-simple}
\end{align}
The two coordinate sets may be interchanged to use the ridge
diagonal on the other side.
\end{theorem}
\begin{proof}
For \(t\ne0\), use the same graph argument as in
Theorem~\ref{four:regularized}, but choose
\(\alpha=4\gamma^2t^2\).  The graph weights \(d_i\) are now
exactly \((K_\gamma)_{ii}\), independent of \(t\); their sum
is \(\theta_\gamma\).  The first condition in
\eqref{four:ridge-gate} makes the scalar Fourier exponents
admissible.  Keeping the multiplier in
\eqref{four:graph-bound}, and taking the \(2q\)-th root, gives
\[
 |\varphi_Y(t)|\leq D_\gamma^{-1/(4q)}
           \left(\frac{L_q}{2\gamma|t|}\right)^{\theta_\gamma/(2q)}.
\]
The power is greater than one.  Its crossover scale is
\[
 a=\frac{L_q}{2\gamma}D_\gamma^{-1/(2\theta_\gamma)}.
\]
Lemma~\ref{four:inversion} and
\(L_q/(2\pi)=e/\sqrt{2q}\) prove
\eqref{four:ridge-density}.
If \(x_j=s_j(B)^2/\gamma^2\), then
\[
 \log D_\gamma=\sum_j\log(1+x_j)
 \geq\sum_j\frac{x_j}{1+x_j}=\theta_\gamma,
\]
because the difference between the summands has derivative
\(x/(1+x)^2\) and vanishes at zero.  Therefore
\(D_\gamma^{-1/(2\theta_\gamma)}\leq e^{-1/2}\), which gives
\eqref{four:ridge-density-simple}.
\end{proof}

\begin{corollary}[Three cross singular values]\label{four:ridge-three}
Suppose \(s_3(B)\geq\sigma>0\), set
\(\gamma=\sigma/\sqrt6\), and assume
\(\max_i(K_\gamma)_{ii}\leq1/2\).  Then
\begin{equation}\label{four:ridge-three-bound}
 \|p_Y\|_\infty\leq\frac{9\sqrt{3e}}{2\sigma}.
\end{equation}
\end{corollary}
\begin{proof}
The first three singular values give
\(\theta_\gamma\geq3/(1+1/6)=18/7\).
Apply \eqref{four:ridge-density-simple} with \(q=1\).
The function \(\theta/(\theta-2)\) decreases for \(\theta>2\),
and its value at \(18/7\) is \(9/2\).
\end{proof}

\begin{lemma}[Stability of ridge leverage]\label{four:ridge-stability}
Let \(B=B_0+E\), \(\|E\|_{\rm op}\leq\eta\), and let \(P\)
be the orthogonal projection onto \(\operatorname{range}(B_0)\).
For every \(\gamma>0\), writing \(\ell_i=P_{ii}\),
\begin{equation}\label{four:ridge-stability-bound}
 (K_\gamma)_{ii}
 \leq\left(\sqrt{\ell_i}
          +\frac{\eta}{\gamma}\sqrt{1-\ell_i}\right)^2.
\end{equation}
In particular, if \(\ell_i\leq\ell\) for all \(i\), the weaker
bound \((\sqrt\ell+\eta/\gamma)^2\) holds.
\end{lemma}
\begin{proof}
The identity
\[
 K_\gamma
 =B(\gamma^2I+B^{\mathsf T}B)^{-1}B^{\mathsf T}
\]
shows \(0\leq K_\gamma\leq I\).
For \(v\perp\operatorname{range}(B_0)\), \(B_0^{\mathsf T}v=0\);
hence
\[
 v^{\mathsf T}K_\gamma v
 \leq\gamma^{-2}\|B^{\mathsf T}v\|^2
 \leq(\eta/\gamma)^2\|v\|^2.
\]
Decompose \(e_i=Pe_i+(I-P)e_i\) and use the triangle inequality
after applying \(K_\gamma^{1/2}\).  Its norm on the first
component is at most \(\sqrt{\ell_i}\), and the preceding
bound on the second is at most
\((\eta/\gamma)\sqrt{1-\ell_i}\).
\end{proof}

\subsection{A spectral-window consequence}

We use the two-part frame partition consequence of
\citet[Corollary 1.5]{MSS2015}: if
\(\sum_i u_iu_i^{\mathsf T}=I_r\) and
\(\|u_i\|^2\leq\varepsilon\), the indices admit a partition
\(I,J\) such that both partial frame operators have norm at most
\((1/\sqrt2+\sqrt\varepsilon)^2\).
The following corollary applies this partition to the range of a
quadratic matrix.

\begin{corollary}[Diffuse full range]\label{four:full-range}
Let \(A\) be symmetric, \(s_3(A)\geq s>0\), and let \(P_A\)
be the projection onto its full range.  If
\(\max_i(P_A)_{ii}\leq1/32\), then, for arbitrary affine terms
and arbitrary additional separable measurable phases,
\begin{equation}\label{four:full-range-bound}
 \|p_{X^{\mathsf T}AX+\sum_iV_i(X_i)}\|_\infty
 \leq\frac{96e}{7\sqrt2\,s}.
\end{equation}
The density is continuous and vanishes at infinity.
\end{corollary}
\begin{proof}
Write \(A=U\Lambda U^{\mathsf T}\), where \(U\) has orthonormal
columns spanning its full range.  The rows of \(U\) form a
tight frame with squared lengths at most \(1/32\).
The cited partition theorem gives
\[
 \frac7{32}I\leq M_I:=U_I^{\mathsf T}U_I,\,
                         M_J:=U_J^{\mathsf T}U_J
                 \leq\frac{25}{32}I.
\]
Indeed \((1/\sqrt2+\sqrt{1/32})^2=25/32\), and the
two operators sum to \(I\).
The nonzero singular values of
\(B=A_{I,J}=U_I\Lambda U_J^{\mathsf T}\) are those of
\(M_I^{1/2}\Lambda M_J^{1/2}\), since
\(U_IM_I^{-1/2}\) and \(U_JM_J^{-1/2}\) have orthonormal
columns.  Multiplication by the two square roots gives
\(s_3(B)\geq(7/32)s\).
Both restricted frames have full column rank, and \(\Lambda\)
is invertible, so \(\operatorname{range}(B)=\operatorname{range}(U_I)\).
Its projection has diagonal
\[
 u_i^{\mathsf T}M_I^{-1}u_i
 \leq\frac{1/32}{7/32}=\frac17<\frac12.
\]
Apply \eqref{four:regularized-simple} with \(q=1\), \(r=3\),
and \(\sigma=7s/32\).
\end{proof}

\begin{theorem}[A diffuse spectral window]\label{four:window}
Let \(A\) be symmetric with \(s_3(A)\geq s>0\), and let \(P\)
be its spectral projection onto the eigenspaces with
\(|\lambda|\geq s/64\).  If
\begin{equation}\label{four:window-hypothesis}
 \max_iP_{ii}\leq1/32,
\end{equation}
then \(X^{\mathsf T}AX+\sum_iV_i(X_i)\), for arbitrary real
measurable \(V_i\) finite almost surely, has a continuous density
vanishing at infinity and
\begin{equation}\label{four:window-bound}
 \|p\|_\infty\leq\frac{288\sqrt{3e}}{13s}<\frac{64}{s}.
\end{equation}
The number, signs, eigenvectors, and nonzero sizes of the smaller
spectral directions are unrestricted.
\end{theorem}
\begin{proof}
Write \(A=A_H+E\), where
\(A_H=U\Lambda_HU^{\mathsf T}\) is the indicated spectral
truncation.  Then \(\|E\|_{\rm op}\leq\eta=s/64\).
Apply the same frame partition as in
Corollary~\ref{four:full-range}.  For
\(B_0=(A_H)_{I,J}\), its proof gives
\[
 s_3(B_0)\geq\frac{7s}{32},\qquad
 \max_i(P_{\operatorname{range}(B_0)})_{ii}\leq\frac17.
\]
The actual cross block is \(B=B_0+E_{I,J}\), and
\(\|E_{I,J}\|_{\rm op}\leq\eta\).  The singular-value
perturbation inequality gives
\[
 s_3(B)\geq\frac{7s}{32}-\frac{s}{64}
          =\sigma:=\frac{13s}{64}.
\]
For completeness, that inequality follows by restricting to the
top-three right-singular subspace of \(B_0\), where
\(\|Bx\|\geq(\frac{7s}{32}-\eta)\|x\|\), and applying the
singular-value min--max principle.

Set \(\gamma=\sigma/\sqrt6\).
Lemma~\ref{four:ridge-stability} yields
\[
 \max_i(K_\gamma)_{ii}
 \leq\left(\frac1{\sqrt7}+\frac{\sqrt6}{13}\right)^2<\frac12.
\]
For example \(1/\sqrt7<2/5\) and \(\sqrt6/13<1/5\), so
the square is less than \(9/25\).
Corollary~\ref{four:ridge-three} proves
\[
 \|p\|_\infty\leq\frac{9\sqrt{3e}}{2\sigma}
              =\frac{288\sqrt{3e}}{13s}.
\]
All small spectral directions have remained in the actual cross
block \(B\); none was removed from the random quadratic expression.
All self terms are permitted by decoupling.  This proves the
claim including arbitrary separable perturbations.
\end{proof}

\subsection{Stability of volume marginals and distributed tails}

\begin{proposition}[Spectral formula for volume marginals]
\label{four:volume-operator}
Let \(M\) have rank at least \(k\), with singular-value
decomposition \(M=U\operatorname{diag}(s_j)V^{\mathsf T}\).
Put \(\lambda_j=s_j^2\), \(Z_k=e_k(\lambda)\), and
\[
 p_j=\frac{\lambda_je_{k-1}(\lambda_{\ne j})}{Z_k},
 \qquad W_k=U\operatorname{diag}(p_j)U^{\mathsf T}.
\]
Then the row marginals in \eqref{four:volume-marginals} satisfy
\((W_k)_{ii}=\alpha_i\), and
\begin{equation}\label{four:volume-operator-bound}
 0\leq W_k\leq I,\qquad
 W_k\leq\frac{e_{k-1}(\lambda)}{Z_k}MM^{\mathsf T}.
\end{equation}
Writing \(T_k=\sum_{j\geq k}\lambda_j>0\), we also have
\begin{equation}\label{four:volume-tail-operator}
 W_k\leq\frac{k}{T_k}MM^{\mathsf T}
       \leq\frac{k}{s_k(M)^2}MM^{\mathsf T}.
\end{equation}
\end{proposition}
\begin{proof}
Multiply row \(i\) of \(M\) by \(\sqrt{1+t}\).
Each squared \(k\)-minor containing that row is multiplied by
\(1+t\), while the other minors are unchanged.  Thus the
coefficient of \(t\) in the new \(Z_k\) is \(Z_k\alpha_i\).
Let \(C=MM^{\mathsf T}\).  Using Sylvester's determinant
identity and the rank-one determinant formula, the generating
polynomial of this row scaling is
\[
 \det(I+zC)
 \left[1+tz\,e_i^{\mathsf T}(I+zC)^{-1}Ce_i\right].
\]
Its coefficient of \(tz^k\), after diagonalizing \(C\), is
\(\sum_jU_{ij}^2\lambda_je_{k-1}(\lambda_{\ne j})\).
Comparison proves \((W_k)_{ii}=\alpha_i\).  This polynomial
argument also covers repeated singular values and zero modes.

The decomposition
\[
 Z_k=\lambda_je_{k-1}(\lambda_{\ne j})
                         +e_k(\lambda_{\ne j})
\]
gives \(0\leq p_j\leq1\), while
\(e_{k-1}(\lambda_{\ne j})\leq e_{k-1}(\lambda)\).
These observations prove \eqref{four:volume-operator-bound}.
Finally,
\[
 kZ_k=\sum_{|S|=k-1}\left(\prod_{i\in S}\lambda_i\right)
                           \sum_{j\notin S}\lambda_j.
\]
Removing any \(k-1\) indices removes at most the sum of the
largest \(k-1\) squared singular values.  Hence every inner
sum is at least \(T_k\), and
\(kZ_k\geq T_ke_{k-1}(\lambda)\).
This proves \eqref{four:volume-tail-operator}.
\end{proof}

\begin{lemma}[Projection control of volume marginals]
\label{four:volume-stability}
Let \(P\) be any orthogonal projection on the row-index space,
\(\ell_i=P_{ii}\), and
\(\varepsilon=\|(I-P)M\|_{\rm op}\).  With the notation of
Proposition~\ref{four:volume-operator},
\begin{equation}\label{four:volume-stability-tail}
 \sqrt{\alpha_i}
 \leq\sqrt{\ell_i}
      +\sqrt{\frac{k}{T_k}}\,\varepsilon\sqrt{1-\ell_i}.
\end{equation}
In particular for \(k=3\), \(\kappa=s_3(M)>0\),
\begin{equation}\label{four:volume-stability-three}
 \sqrt{\alpha_i}
 \leq\sqrt{\ell_i}
      +\frac{\sqrt3\,\varepsilon}{\kappa}\sqrt{1-\ell_i}.
\end{equation}
\end{lemma}
\begin{proof}
Use the triangle inequality after applying \(W_k^{1/2}\) to
\(e_i=Pe_i+(I-P)e_i\).  The bound \(W_k\leq I\) controls
the first component by \(\sqrt{\ell_i}\).
The first inequality in \eqref{four:volume-tail-operator}
controls the second by
\[
 \sqrt{k/T_k}\,\|M^{\mathsf T}(I-P)e_i\|
 \leq\sqrt{k/T_k}\,\varepsilon\sqrt{1-\ell_i}.
\]
Since \(\|W_k^{1/2}e_i\|=\sqrt{\alpha_i}\), this proves the
first assertion.  The bound \(T_3\geq\kappa^2\) proves the second.
\end{proof}

\begin{theorem}[A cross block near a diffuse range]
\label{four:near-diffuse}
Suppose a physical cross block of a quadratic form satisfies
\[
 M=M_0+E,\qquad s_3(M_0)\geq a>0,\qquad
 \max_i(P_{\operatorname{range}(M_0)})_{ii}\leq\tfrac12,
 \qquad \|E\|_{\rm op}\leq a/16.
\]
Then, with arbitrary self blocks, outside couplings, and affine
terms, the quadratic form has a continuous density vanishing
at infinity and
\begin{equation}\label{four:near-diffuse-bound}
 \|p\|_\infty\leq\frac{24\sqrt3\,e}{5a}<\frac{23}{a}.
\end{equation}
No rank bound on either \(M_0\) or \(E\) is required.
\end{theorem}
\begin{proof}
Let \(P=P_{\operatorname{range}(M_0)}\).
Then \(\varepsilon=\|(I-P)M\|\leq a/16\), and singular-value
perturbation gives \(\kappa=s_3(M)\geq15a/16\).
Lemma~\ref{four:volume-stability} gives
\[
 \sqrt{\alpha_i}\leq\sqrt{\ell_i}
                    +\frac{\sqrt3}{15}\sqrt{1-\ell_i}.
\]
For \(0\leq\ell\leq1/2\) and \(0\leq d\leq1\), the function
\(\sqrt\ell+d\sqrt{1-\ell}\) is increasing in \(\ell\).
Consequently
\[
 \alpha_i\leq\frac{(1+\sqrt3/15)^2}{2}<\frac23.
\]
For a rational check, replacing \(\sqrt3\) by \(2\) gives
the still smaller-than-\(2/3\) bound \(289/450\).
Theorem~\ref{four:volume-density} with \(k=3\), \(q=4/3\)
therefore gives
\[
 \|p\|_\infty\leq\frac{9\sqrt3\,e}{2}\,Z_3^{-1/6}.
\]
Since \(Z_3\geq s_1(M)^2s_2(M)^2s_3(M)^2\geq\kappa^6\),
the claimed constant follows from \(\kappa\geq15a/16\).
\end{proof}

\begin{corollary}[Distributed spectral tails]\label{four:tail}
Let \(M\) be a physical cross block of rank at least three and put
\[
 T_3=\sum_{j\geq3}s_j(M)^2.
\]
Suppose a projection \(P\) on its row-index space satisfies
\[
 \max_iP_{ii}\leq\tfrac12,\qquad
 \|(I-P)M\|_{\rm op}\leq\sqrt{T_3}/15.
\]
Then
\begin{align}
 \|p\|_\infty
 &\leq\frac{9\sqrt3\,e}{2}\,e_3(M^{\mathsf T}M)^{-1/6}
 \label{four:tail-bound}\\
 &\leq\frac{9\sqrt3\,e}{2}
           [s_1(M)^2s_2(M)^2T_3]^{-1/6}.
 \nonumber
\end{align}
If \(s_2(M)\geq a>0\) and \(T_3\geq b^2>0\), the latter
bound is \((9\sqrt3\,e/2)a^{-2/3}b^{-1/3}\).
In particular, when \(P\) projects onto the top two left singular
directions and its diagonal is at most \(1/2\), the condition
\(T_3/s_3(M)^2\geq225\) suffices.
\end{corollary}
\begin{proof}
Use \eqref{four:volume-stability-tail} with \(k=3\).
The same marginal estimate as in
Theorem~\ref{four:near-diffuse} gives \(\alpha_i<2/3\).
The fractional three-volume theorem proves the first bound.
The terms of \(e_3\) containing the first two indices sum to
\(s_1^2s_2^2T_3\), proving the second.  For the indicated
projection, \(\|(I-P)M\|_{\rm op}=s_3(M)\).
\end{proof}

\subsection{Small balls for an arbitrary nonzero quadratic form}

\begin{lemma}[Quarter-power concentration]\label{four:quarter}
Let \(S\ne0\) be any real symmetric matrix and let
\[
 R=X^{\mathsf T}SX+b^{\mathsf T}X+c.
\]
For every \(\rho>0\),
\begin{equation}\label{four:quarter-bound}
 \sup_w\mathbb P(|R-w|<\rho)
 \leq\min\left\{1,5\left(\frac{\rho}{\|S\|_{\rm op}}\right)^{1/4}\right\}.
\end{equation}
In particular \(R\) has no atoms.  The same concentration estimate
holds with a closed interval.
\end{lemma}
\begin{proof}
Put \(L=\|S\|_{\rm op}>0\) and \(m=\max_i|S_{ii}|\).
First suppose \(m\geq L/2\), and condition on all coordinates
except one attaining \(m\).
A real quadratic \(ax^2+dx+e\) with \(a\ne0\) has a preimage
of any interval of length \(2\rho\) of total length at most
\(2\sqrt{2\rho/|a|}\).  Indeed, complete the square and note
that for an interval \([u,u+\delta]\) the set where
\((x-x_0)^2\) belongs to that interval has length at most
\(2\sqrt\delta\): on the nonnegative half-line this is
\(2(\sqrt{u+\delta}-\sqrt u)\), with the evident truncation
when \(u<0\).
The density cap therefore bounds the conditional probability by
\(4\sqrt{\rho/L}\).
For \(\rho\leq L\), this is at most
\(4(\rho/L)^{1/4}\); for \(\rho>L\), the bound one suffices.

Now suppose \(m<L/2\).  Let \(S_0=S-\operatorname{diag}(S)\),
so \(\|S_0\|_{\rm op}\geq L/2\).
For independent uniform signs on the diagonal of a matrix \(E\),
define \(C_E=(S_0-ES_0E)/2\).
The zero diagonal of \(S_0\) implies
\(\mathbb EC_E=S_0/2\).  Convexity of the operator norm shows
that some sign choice satisfies
\(\|C_E\|_{\rm op}\geq\|S_0\|_{\rm op}/2\geq L/4\).
For the cut into positive and negative signs,
\[
 C_E=\begin{pmatrix}0&B\\B^{\mathsf T}&0\end{pmatrix},
 \qquad \|B\|_{\rm op}=\|C_E\|_{\rm op}\geq L/4 .
\]
This is a cut in the original coordinates.
The unrestricted determinant bound
\eqref{four:unrestricted-determinant}, retaining only the
largest singular value of \(B\), gives
\begin{equation}\label{four:quarter-decay}
 |\varphi_R(t)|
 \leq\left(1+\frac{4\|B\|_{\rm op}^2t^2}{\pi^2e^2}\right)^{-1/8}
 \leq\left(\frac{2\pi e}{L|t|}\right)^{1/4}.
\end{equation}
The affine terms are allowed throughout the decoupling argument.

For the concentration constant, let \(v>0\).  Then
\(\mathbf1_{\{|x|<\rho\}}\leq e^v e^{-vx^2/\rho^2}\).
Gaussian Fourier inversion and Fubini give
\[
 \mathbb P(|R-w|<\rho)
 \leq\frac{e^v\rho}{2\sqrt{\pi v}}
      \int_{\mathbb R}e^{-\rho^2t^2/(4v)}
                         |\varphi_R(t)|\,dt.
\]
For \(0<\alpha<1\),
\[
 \int_{\mathbb R}e^{-at^2}|t|^{-\alpha}\,dt
 =a^{(\alpha-1)/2}\Gamma((1-\alpha)/2),
\]
by the substitution \(u=at^2\).
Use \eqref{four:quarter-decay}, \(\alpha=1/4\), and \(v=1/8\).
The coefficient of \((\rho/L)^{1/4}\) is
\[
 2^{3/8}\pi^{-1/4}e^{3/8}\Gamma(3/8)<5.
\]
For an elementary verification,
\(\Gamma(3/8)\leq8/3+e^{-1}<73/24\), obtained by splitting
its defining integral at one.  Also
\(2^{3/8}<4/3\), \(\pi^{-1/4}<4/5\), and \(e^{3/8}<3/2\).
The product of these four upper bounds is \(73/15<5\).
This proves \eqref{four:quarter-bound}.
Letting \(\rho\downarrow0\) proves absence of atoms.
The closed-interval version follows by applying the result with
\(\rho+\varepsilon\) and then letting \(\varepsilon\downarrow0\).
\end{proof}

\begin{corollary}[Logarithmic and fractional potentials]
\label{four:log-potential}
In the setting of Lemma~\ref{four:quarter}, let
\(L=\|S\|_{\rm op}\).  For every \(a>0\), put
\(M=5(a/L)^{1/4}\).  Uniformly in \(w\),
\begin{align}
 \mathbb E\log_+\frac a{|R-w|}
 &\leq
 \begin{cases}
  4M,&M\leq1,\\
  4+4\log M,&M\geq1,
 \end{cases}
 \label{four:log-potential-bound}\\
 &\leq20(a/L)^{1/4}.
 \nonumber
\end{align}
For every \(0<\beta<1/4\), moreover,
\begin{equation}\label{four:negative-moment}
 \sup_w\mathbb E|R-w|^{-\beta}
 \leq\frac{5^{4\beta}}{1-4\beta}\,L^{-\beta}.
\end{equation}
\end{corollary}
\begin{proof}
Tonelli's theorem and the quarter-power bound give
\[
 \mathbb E\log_+\frac a{|R-w|}
 =\int_0^\infty
      \mathbb P(|R-w|<ae^{-u})\,du
 \leq\int_0^\infty\min\{1,Me^{-u/4}\}\,du.
\]
If \(M\leq1\), this integral is \(4M\).
If \(M>1\), split at \(u=4\log M\) to obtain \(4+4\log M\).
Both cases are at most \(4M\).
For the negative moment use
\[
 \mathbb E|R-w|^{-\beta}
 =\beta\int_0^\infty r^{-\beta-1}
                       \mathbb P(|R-w|<r)\,dr.
\]
Split at \(r_0=L/5^4\).  The probability is at most
\((r/r_0)^{1/4}\) for \(r<r_0\) and at most one otherwise.
The two integrals sum to
\(r_0^{-\beta}/(1-4\beta)\), giving
\eqref{four:negative-moment}.
\end{proof}

The spectral-window, volume, and ridge criteria above are
quantitative sufficient hypotheses.  A lower bound for the third
singular value of the full quadratic matrix alone does not verify
their coordinate cut or marginal conditions.  Their usefulness is
that the estimates retain arbitrary self interactions and, where
specified, all additional spectral directions without any assumption
that a remainder is independent of the principal part.

\section{Fractional potentials of rank-two forms}
\label{sec:rank-two}

An indefinite rank-two form can have an unbounded density,
but its half-power potential admits a uniform determinant
bound, including for affine terms outside the range of its
quadratic matrix.  Such terms arise under conditioning on
the original coordinates.

\begin{lemma}[A rank-one affine small-ball estimate]
\label{rt:rank-one}
Let \(|u|=1\), \(\lambda\ne0\), and let
\[
 R=\lambda\langle u,X\rangle^2+b^{\mathsf T}X+d.
\]
For every \(w\in\mathbb R\) and \(\rho>0\),
\[
 \mathbb P\{|R-w|<\rho\}
 \leq\min\{1,8\sqrt{\rho/|\lambda|}\}.
\]
\end{lemma}
\begin{proof}
If \(u_i^2\geq1/2\) for some original coordinate, conditioning
on the others leaves a quadratic of leading coefficient
\(\lambda u_i^2\).  Its inverse image of an interval of
length \(2\rho\) has total length at most
\(2\sqrt{2\rho/|\lambda u_i^2|}\leq4\sqrt{\rho/|\lambda|}\).

Otherwise Theorem~\ref{four:diffuse}, with one spectral
direction and the given affine phase, gives
\[
 |\phi_R(t)|\leq\sqrt{e\pi}\,|\lambda t|^{-1/2}.
\]
For \(v>0\), Gaussian Fourier inversion and
\(\boldsymbol1_{\{|z|<\rho\}}\leq e^v e^{-vz^2/\rho^2}\)
imply
\[
 \mathbb P\{|R-w|<\rho\}
 \leq
 \frac{e^v\rho}{2\sqrt{\pi v}}
 \int_{\mathbb R}e^{-\rho^2t^2/(4v)}|\phi_R(t)|\,dt
 \leq
 \frac{e^v\sqrt e\,\Gamma(1/4)}
 {\sqrt2\,v^{1/4}}\sqrt{\rho/|\lambda|}.
\]
Take \(v=1/4\).  The resulting constant
\(e^{3/4}\Gamma(1/4)\) is less than \(8\).
An elementary upper bound sufficient to verify this
numerical inequality is
\[
 \Gamma(1/4)\leq
 4-\frac45+\frac29-\frac2{39}+\frac1{102}+e^{-1}.
\]
Indeed, on \([0,1]\) use the degree-four Taylor upper
bound for \(e^{-x}\), and on \([1,\infty)\) use
\(x^{-3/4}\leq1\).  Together with
\(e^{-1}<3/8\) and \(e^{3/4}<17/8\), the displayed
expression gives a product less than \(8\).
\end{proof}

Define an even function on \(\mathbb R\) by
\[
 \Psi(a)=
 \begin{cases}
 1+\log(1/|a|),&0<|a|\leq1,\\
 |a|^{-1/4},&|a|\geq1,\\
 +\infty,&a=0.
 \end{cases}
\]

\begin{lemma}[A translated quadratic kernel]
\label{rt:kernel}
For \(a\ne0\) and every probability density \(f\leq1\),
\[
 \int_{\mathbb R}f(x)|a-x^2|^{-1/2}\,dx\leq5\Psi(a).
\]
More generally, a probability density bounded by \(K\)
satisfies
\[
 \int_{\mathbb R}f(x)|a-x^2|^{-1/2}\,dx
 \leq5K\Psi(K^2a).
\]
\end{lemma}
\begin{proof}
Let \(m_a(s)\) be the measure of
\(\{x:|a-x^2|^{-1/2}>s\}\).
The integral is at most
\(\int_0^\infty\min\{1,m_a(s)\}\,ds\).
For \(a>0\), put \(\varepsilon=s^{-2}\).
If \(\varepsilon\geq a\), then
\(m_a(s)=2\sqrt{a+\varepsilon}\leq2\sqrt2/s\).
If \(\varepsilon<a\), then
\[
 m_a(s)
 =2\bigl(\sqrt{a+\varepsilon}-\sqrt{a-\varepsilon}\bigr)
 \leq\frac4{\sqrt a\,s^2}.
\]
For \(0<a\leq1\), split the integral at \(1\) and
\(a^{-1/2}\).  It is at most
\(5+\sqrt2\log(1/a)\leq5\Psi(a)\).
For \(a\geq1\), splitting at \(2a^{-1/4}\) gives
at most \(4a^{-1/4}\).

If \(a<0\), then \(m_a(s)\leq2/s\), and it vanishes
when \(s\geq|a|^{-1/2}\).
For \(-1\leq a<0\), split at \(1\) to obtain
\(1+\log(1/|a|)\).  For \(a\leq-1\), the integrand
kernel itself is at most \(|a|^{-1/2}\leq|a|^{-1/4}\).
This proves the first assertion.  For the second,
substitute \(y=Kx\); the rescaled density is bounded
by one and the kernel gains the factor \(K\).
\end{proof}

\begin{lemma}[Averaging the kernel]
\label{rt:kernel-average}
If \(\mathbb P\{|W|<\rho\}\leq L\sqrt\rho\) for all
\(\rho>0\), then
\[
 \mathbb E\Psi(W)\leq4\sqrt L.
\]
\end{lemma}
\begin{proof}
The hypothesis excludes an atom at zero.
For \(0<s<1\), the event \(\Psi(W)>s\) implies
\(|W|<s^{-4}\), so its probability is at most \(Ls^{-2}\).
For \(s\geq1\), it is at most \(Le^{(1-s)/2}\).
Since
\(\sup_{s\geq1}s^2e^{(1-s)/2}=16e^{-3/2}<4\),
both ranges are bounded by \(\min\{1,4L/s^2\}\).
Its integral over \(s>0\) is \(4\sqrt L\).
\end{proof}

\begin{theorem}[Rank-two half-power potential]
\label{rt:half}
Let \(U^{\mathsf T}U=I_2\), let \(\Lambda\) be an
invertible real symmetric \(2\)-by-\(2\) matrix, and
put \(D=|\det\Lambda|\).
For every original-coordinate affine term,
\[
 R=(U^{\mathsf T}X)^{\mathsf T}\Lambda(U^{\mathsf T}X)
       +b^{\mathsf T}X+c_0,
\]
one has
\begin{equation}
 \sup_w\mathbb E|R-w|^{-1/2}
 \leq40\,2^{3/4}D^{-1/4}<68D^{-1/4}.
 \label{eq:rank-two-half}
\end{equation}
\end{theorem}
\begin{proof}
Suppose first that every row leverage of \(U\) is at
most \(1/2\).
Theorem~\ref{four:diffuse} gives
\[
 |\phi_R(t)|\leq\min\{1,K/|t|\},
 \qquad K=e\pi D^{-1/2}.
\]
The Gaussian representation
\[
 |z|^{-1/2}
 =\frac1{\Gamma(1/4)}
   \int_0^\infty v^{-3/4}e^{-vz^2}\,dv
\]
and Gaussian Fourier inversion give
\[
 \mathbb E|R-w|^{-1/2}
 \leq(2\pi)^{-1/2}
       \int_{\mathbb R}|t|^{-1/2}|\phi_R(t)|\,dt.
\]
To justify the order of integration, first write the
Gaussian formula at each \(v>0\), take absolute values,
and then apply Tonelli to the nonnegative integrand.
The integral in \(v\) is
\(\sqrt2\,\Gamma(1/4)|t|^{-1/2}\) before its fixed
prefactor, giving the stated coefficient.
Splitting the last integral at \(|t|=K\) gives
\[
 \mathbb E|R-w|^{-1/2}
 \leq4\sqrt{2e}\,D^{-1/4}.
\]

Otherwise an original coordinate \(i\) has leverage
\(h>1/2\).  Rotate only the two-dimensional spectral
frame so that its row is \((\sqrt h,0)\), and write
\[
 Y_1=\sqrt h\,X_i+T,\qquad Y_2=Z,\qquad
 \Lambda=\begin{pmatrix}\beta&\gamma\\\gamma&\delta\end{pmatrix}.
\]
The variables \(T,Z\) depend on \(X_{-i}\), and \(Z\)
is a unit linear form in those original coordinates.
Put \(k=h^{-1/2}\), \(d=b_i/\sqrt h\), and
\(L_0=b_{-i}^{\mathsf T}X_{-i}-dT+c_0\).
Then
\[
 R=\beta Y_1^2+2\gamma Y_1Z+\delta Z^2+dY_1+L_0.
\]
Conditional on \(X_{-i}\), the density of \(Y_1\) is
bounded by \(k\); its translation is unrestricted.

If \(\beta\ne0\), complete the square:
\[
 R=\beta\left(Y_1+\frac{\gamma Z}{\beta}
                    +\frac d{2\beta}\right)^2+R_0,
 \qquad
 R_0=\frac{\det\Lambda}{\beta}Z^2
       -\frac{\gamma d}{\beta}Z-\frac{d^2}{4\beta}+L_0 .
\]
This is a rank-one quadratic in the remaining original
coordinates plus an arbitrary affine term.
Lemma~\ref{rt:kernel}, conditionally, yields
\[
 \mathbb E[|R-w|^{-1/2}\mid X_{-i}]
 \leq5k|\beta|^{-1/2}\Psi(W),
 \qquad W=k^2(w-R_0)/\beta .
\]
The magnitude of the rank-one coefficient of \(W\) is
\(\lambda_0=k^2D/\beta^2\).
Lemmas~\ref{rt:rank-one} and \ref{rt:kernel-average}
therefore give
\[
 \mathbb E|R-w|^{-1/2}
 \leq20\sqrt8\,k|\beta|^{-1/2}\lambda_0^{-1/4}
 =20\sqrt8\,\sqrt k\,D^{-1/4}
 \leq40\,2^{3/4}D^{-1/4}.
\]
All powers of \(\beta\) cancel.

If \(\beta=0\), then \(D=\gamma^2>0\), and the
conditional polynomial is affine in \(X_i\), of slope
\(\sqrt h(2\gamma Z+d)\).
For any probability density bounded by \(H\),
layer cake gives
\[
 \int f(x)|x-a|^{-1/2}\,dx\leq2\sqrt{2H}.
\]
Apply this first to \(X_i\), then to \(Z\), whose
density is at most \(\sqrt2\) by
Lemma~\ref{four:projection}.  The result is
\[
 \mathbb E|R-w|^{-1/2}
 \leq4\,2^{3/4}h^{-1/4}D^{-1/4}
 \leq8D^{-1/4}.
\]
The exceptional value \(2\gamma Z+d=0\) has probability
zero.  These cases exhaust the possibilities.
\end{proof}

\begin{lemma}[A conditional two-coordinate potential]
\label{rt:conditional}
Suppose \(Z\) has density at most \(H\), and the
conditional density of \(S\) given \(Z\) is at most
\(K\).  For an invertible real symmetric matrix \(L\),
and arbitrary deterministic affine terms,
\[
 \sup_w\mathbb E\left|
 (S,Z)L(S,Z)^{\mathsf T}+\alpha S+\beta Z+\gamma-w
 \right|^{-1/2}
 \leq10\sqrt{15}\sqrt{KH}\,|\det L|^{-1/4}.
\]
\end{lemma}
\begin{proof}
An invertible deterministic quadratic matrix permits
removing its affine term by translating \(S,Z\);
both density caps are preserved.
Write
\(L=\bigl(\begin{smallmatrix}a&b\\b&d\end{smallmatrix}\bigr)\).
For \(a\ne0\), complete the square in \(S\).
Lemma~\ref{rt:kernel} reduces the conditional bound to
\[
 5K|a|^{-1/2}\Psi(A+\epsilon\delta_0Z^2),
 \qquad
 \delta_0=K^2|\det L|/a^2,\quad \epsilon\in\{-1,1\}.
\]
The elementary sublevel estimate
\(|\{z:|z^2+A|<\rho\}|\leq4\sqrt\rho\)
shows that the function
\(v(z)=\Psi(A+\epsilon\delta_0z^2)\) has weak \(L^2\)
norm at most \(\sqrt{15}\delta_0^{-1/4}\).
Indeed, for levels \(0<s<1\),
\(s^2|\{v>s\}|\leq4\delta_0^{-1/2}\).
For \(s\geq1\), the bound is
\(4\delta_0^{-1/2}s^2e^{(1-s)/2}
\leq64e^{-3/2}\delta_0^{-1/2}
<15\delta_0^{-1/2}\).
For any probability density \(g\leq H\), layer cake
gives
\[
 \int gv\leq
 \int_0^\infty\min\{1,H\|v\|_{2,\infty}^2/s^2\}\,ds
 =2\sqrt H\,\|v\|_{2,\infty}.
\]
Substitution cancels \(a\) and proves the asserted
constant.  For \(a=0\), condition first on \(Z\) and
use the affine negative-half-moment estimate from the
last proof; it gives the smaller constant
\(4\sqrt2\sqrt{KH}\,|\det L|^{-1/4}\).
\end{proof}

\begin{corollary}[Affine terms in the range]
\label{rt:range}
In Theorem~\ref{rt:half}, if \(b\in\operatorname{range}(U)\),
then
\[
 \sup_w\mathbb E|R-w|^{-1/2}
 \leq10\sqrt{30}\,D^{-1/4}.
\]
\end{corollary}
\begin{proof}
A deterministic translation of the original coordinates
removes the affine term.  In the diffuse case the
bound \(4\sqrt{2e}D^{-1/4}\) obtained above suffices.
In the heavy-coordinate notation of that proof,
the conditional density of \(Y_1\) given \(Z\) is at
most \(h^{-1/2}\): average the conditional bound given
all \(X_{-i}\).  The marginal density of \(Z\) is at
most \(\sqrt2\).
Lemma~\ref{rt:conditional} gives
\[
 10\sqrt{15}(2/h)^{1/4}D^{-1/4}
 \leq10\sqrt{30}\,D^{-1/4}.
\]
\end{proof}

\begin{theorem}[An exact-rank-three Schur estimate]
\label{rt:rank-three-schur}
Suppose an original-coordinate principal block has
representation
\[
 A_{II}=VMV^{\mathsf T},\qquad
 V^{\mathsf T}V=I_3,\qquad \det M\ne0.
\]
Choose \(i\in I\) with \(a=A_{ii}\ne0\), and write
\(h=|V^{\mathsf T}e_i|^2\).
Then the full ambient quadratic form, with arbitrary
other blocks and affine terms, satisfies
\[
 \|p_Q\|_\infty
 \leq68\bigl(h|a|\,|\det M|\bigr)^{-1/4}.
\]
For a full exact-rank-three matrix with all nonzero
eigenvalues of modulus at least one, this implies
\(\|p_Q\|_\infty\leq68/\sqrt{|a|}\).
\end{theorem}
\begin{proof}
Condition first on the original coordinates outside
\(I\).  They create only an arbitrary affine term in
the independent coordinates of \(I\).
Completing the square in \(X_i\) gives a remainder
with Schur matrix
\[
 S=A_{I\setminus\{i\},I\setminus\{i\}}
 -a^{-1}A_{I\setminus\{i\},i}A_{i,I\setminus\{i\}}.
\]
The square kernel gives
\[
 p_Q(t)\leq |a|^{-1/2}\mathbb E|t-R|^{-1/2}.
\]
Put \(v=V^{\mathsf T}e_i\), choose a \(3\)-by-\(2\)
orthonormal frame \(T\) for \(v^\perp\), and let
\(V_-\) be \(V\) with row \(i\) removed.
Then \(W=V_-T\) satisfies
\[
 W^{\mathsf T}W=T^{\mathsf T}(I-vv^{\mathsf T})T=I_2.
\]
The symmetric matrix
\(H=M-Mvv^{\mathsf T}M/a\) has kernel
\(\operatorname{span}(v)\).  Writing
\(N=T^{\mathsf T}HT\) gives \(S=WNW^{\mathsf T}\).
In a spectral-space basis with \(v=\sqrt h\,e_1\),
write
\[
 M=\begin{pmatrix}\alpha&q^{\mathsf T}\\q&G\end{pmatrix},
 \qquad a=h\alpha.
\]
Then \(N=G-qq^{\mathsf T}/\alpha\), up to an
orthogonal change of its two-dimensional basis, and
\[
 \det N=\det M/\alpha=h\det M/a.
\]
Theorem~\ref{rt:half}, with arbitrary induced affine
terms allowed, proves the first bound.
In the normalized full-rank-three case,
\(|\det M|\geq\|M\|\) and \(|a|\leq h\|M\|\);
hence \(h|a||\det M|\geq |a|^2\), proving the last
assertion.
\end{proof}

\section{Arbitrary quadratic forms under Gaussian mixture inputs}\label{sec:gaussian-mixture}

The Gaussian determinant formula and weighted arithmetic--geometric mean yield a dimension-free estimate for product laws with controlled Gaussian-mixture component peaks. The result applies to quadratic forms of arbitrary rank and signature, provided the hypothesis holds for every coordinate. For related Gaussian quadratic-form and regularity estimates, see \citet{APSSS2017,GNSU2019,HMP2025}.

\begin{theorem}[Gaussian mixtures with controlled average component peaks]\label{gm:main}
Let
\[
 X_i=M_i+\sqrt{T_i}\,G_i,\qquad 1\le i\le n,
\]
where the pairs $(M_i,T_i)$ are independent across coordinates, $0<T_i<\infty$ almost surely, and the $G_i$ are independent standard Gaussian variables independent of all pairs. Suppose
\begin{equation}\label{gm:moment}
 \mathbb E T_i^{-1/2}\le\sqrt{2\pi}\qquad(1\le i\le n).
\end{equation}
Let $A$ be any real symmetric matrix, let $b,c$ be arbitrary, and let $3\le k\le\rank A$. Put
\[
 Z_k=e_k(A^2)=\sum_{|I|=|J|=k}\det(A_{IJ})^2>0,
 \qquad
 I_k=\int_0^\infty(1+u^{2k})^{-1/4}\,du.
\]
Then $Q=X^TAX+b^TX+c$ has a bounded continuous density satisfying
\begin{equation}\label{gm:density}
 \|p_Q\|_\infty\le I_k Z_k^{-1/(2k)},\qquad
 I_k=\frac1{2k}\mathrm B\!\left(\frac1{2k},\frac14-\frac1{2k}\right)
 <\frac{k}{k-2}.
\end{equation}
In particular, $s_3(A)\ge1$ implies $\|p_Q\|_\infty<3$ in this input class.
\end{theorem}

\begin{proof}
Condition first on all the mixture variables $(M_i,T_i)$. With $D=\diag(T_1,\ldots,T_n)$ and $M=(M_1,\ldots,M_n)$, write
\[
 Q=G^T B G+\ell^TG+c',\qquad
 B=D^{1/2}AD^{1/2},\quad
 \ell=D^{1/2}(2AM+b).
\]
Orthogonally diagonalizing $B$ is legitimate for the standard Gaussian vector $G$. If its eigenvalues are $\mu_j$ and the rotated linear coefficients are $\beta_j$, direct one-dimensional Gaussian integration gives
\begin{align*}
 |\varphi_{Q\mid M,D}(t)|
 &=\prod_j(1+4t^2\mu_j^2)^{-1/4}
   \exp\left[-\sum_j\frac{t^2\beta_j^2}{2(1+4t^2\mu_j^2)}\right]\\
 &\le\det(I+4t^2B^2)^{-1/4}.
\end{align*}
Terms with $\mu_j=0$ obey the same formula and contribute only a Gaussian linear factor. Expanding the determinant in elementary symmetric functions yields
\[
 |\varphi_{Q\mid M,D}(t)|
 \le[1+(2t)^{2k}W_k]^{-1/4},\qquad W_k=e_k(B^2)>0.
\]
This bound is integrable because $k>2$. Fourier inversion and the substitution $u=2tW_k^{1/(2k)}$ give
\begin{equation}\label{gm:conditional}
 \|p_{Q\mid M,D}\|_\infty
 \le\frac1{2\pi}\int_{\mathbb R}[1+(2t)^{2k}W_k]^{-1/4}\,dt
 =\frac{I_k}{2\pi}W_k^{-1/(2k)}.
\end{equation}

It remains to average the random spectral factor. For $T_I=\prod_{i\in I}T_i$, Cauchy--Binet gives the exact identity
\begin{equation}\label{gm:minors}
 W_k=\sum_{|I|=|J|=k}\det(A_{IJ})^2T_I T_J.
\end{equation}
The sum is over ordered pairs of coordinate subsets; intersections are allowed. On the nonzero minors, define probability weights
\[
 w_{IJ}=\frac{\det(A_{IJ})^2}{Z_k},\qquad
 \alpha_i=\sum_{I,J}w_{IJ}(\mathbf1_{\{i\in I\}}+\mathbf1_{\{i\in J\}}).
\]
Then $0\le\alpha_i\le2$ and $\sum_i\alpha_i=2k$. Weighted arithmetic--geometric mean in \eqref{gm:minors} gives
\[
 W_k\ge Z_k\prod_iT_i^{\alpha_i},\qquad
 W_k^{-1/(2k)}\le Z_k^{-1/(2k)}\prod_iT_i^{-\alpha_i/(2k)}.
\]
Independence across coordinates and concavity of $s\mapsto s^{\alpha_i/k}$ give
\begin{align*}
 \mathbb E W_k^{-1/(2k)}
 &\le Z_k^{-1/(2k)}\prod_i\mathbb E(T_i^{-1/2})^{\alpha_i/k}\\
 &\le Z_k^{-1/(2k)}\prod_i(\mathbb E T_i^{-1/2})^{\alpha_i/k}\\
 &\le2\pi Z_k^{-1/(2k)}.
\end{align*}
Here $\alpha_i/k\le2/k<1$, and zero exponents are omitted. This estimate also proves that the expectation of the conditional Fourier $L^1$ bound is finite. Hence
$\int|\varphi_Q(t)|\,dt\le\mathbb E\int|\varphi_{Q\mid M,D}(t)|\,dt<\infty$
by conditional expectation, the triangle inequality, and Tonelli. Fourier inversion therefore supplies a continuous density, and averaging \eqref{gm:conditional} proves \eqref{gm:density}.

The beta-integral formula follows from $v=u^{2k}$. Finally,
\[
 I_k<\int_0^1 1\,du+\int_1^\infty u^{-k/2}\,du
 =1+\frac2{k-2}=\frac{k}{k-2}.
\]
If $s_3(A)\ge1$, then $Z_3\ge s_1(A)^2s_2(A)^2s_3(A)^2\ge1$, which gives the final assertion.
\end{proof}

\begin{corollary}[Centered scale mixtures]\label{gm:centered}
Conjecture~\ref{conj:main} holds, with a constant smaller than $3$, for all independent coordinate laws that are Gaussian scale mixtures about deterministic coordinate centers and have densities bounded by one. There is no lower bound on the component variances and no moment assumption on the coordinate variables themselves.
\end{corollary}
\begin{proof}
If $M_i=m_i$ is deterministic, the mixture density is
\[
 f_i(x)=\mathbb E\left[(2\pi T_i)^{-1/2}
 e^{-(x-m_i)^2/(2T_i)}\right].
\]
It is maximized at $m_i$, and its essential supremum is
$(2\pi)^{-1/2}\mathbb E T_i^{-1/2}$, also when this is infinite. The latter assertion follows by monotone convergence as $x\to m_i$ along either side. Thus the density cap is equivalent to \eqref{gm:moment}, and Theorem~\ref{gm:main} applies.
\end{proof}

\begin{remark}[Scope of the mixture condition]\label{gm:scope}
For random component centers, \eqref{gm:moment} still implies a cap-one marginal density, but the converse can fail. Narrow components placed far apart can have a bounded final mixture density while their average component peak is large. The theorem therefore does not extend to all bounded densities by an arbitrary Gaussian-mixture approximation. Independent centered Laplace laws, after scaling to the stated cap, are examples within the centered scale-mixture class. The result also differs from the adapted-hub theorem: here the entire input vector satisfies the mixture hypothesis, whereas the adapted theorem allows a conditional hub law outside that class.
\end{remark}

\begin{remark}[Unequal component-peak bounds]
The same proof retains the factor
\[
 \frac{I_k}{2\pi}Z_k^{-1/(2k)}
 \prod_i(\mathbb E T_i^{-1/2})^{\alpha_i/k}.
\]
In particular, a common bound $\mathbb E T_i^{-1/2}\le\sqrt{2\pi}K$ gives $I_kK^2Z_k^{-1/(2k)}$. The powers sum to two, as required by quadratic scaling. The weighted-minor argument allows the conditional Gaussian eigenvectors to depend on the mixture variables.
\end{remark}

\section{Semidefinite forms and spectral spread}
\label{sec:semidefinite}

Throughout this section, $X=(X_1,\ldots,X_n)$ has independent coordinates
with Lebesgue densities $0\leq f_i\leq1$.  We write
\[
 Q=X^TAX+b^TX+c,\qquad \varphi_Q(t)=\mathbb E e^{itQ},
\]
where $A$ is real and symmetric.  The vector $b$ and scalar $c$ are
arbitrary.  Density bounds refer to the essential supremum.  Eigenvalues
of a positive semidefinite matrix are numbered in decreasing order,
including zeros.

\begin{theorem}[Semidefinite forms]\label{sd:main}
If $A$ is positive semidefinite and $\lambda_3(A)>0$, then $Q$ has a
density satisfying
\begin{equation}\label{sd:main-bound}
 \|p_Q\|_\infty\leq \frac{17100}{\lambda_3(A)}.
\end{equation}
The same conclusion holds for negative semidefinite $A$, with
$\lambda_3(-A)$ in the denominator.  The bound is independent of dimension
and rank, and allows a linear term outside the range of $A$.
\end{theorem}

\begin{remark}[Density at the lower endpoint]\label{sd:endpoint-zero}
For the homogeneous form $Q=X^TAX$ with $A\succeq0$ and
$\lambda_3(A)>0$, only three positive eigendirections are needed to
show that the density vanishes \emph{at zero}.  Let $u_1,u_2,u_3$
be orthonormal eigenvectors with positive eigenvalues
$\lambda_1,\lambda_2,\lambda_3$, and put
$Y_j=\langle u_j,X\rangle$.  Lemma~\ref{four:projection} bounds the
joint density of $(Y_1,Y_2,Y_3)$ by $2^{3/2}$.  Positivity of the
discarded spectral terms gives
\[
 Q\geq\sum_{j=1}^3\lambda_jY_j^2,
 \qquad
 \mathbb P(Q\leq t)
 \leq \frac{8\sqrt2\,\pi}{3\sqrt{\lambda_1\lambda_2\lambda_3}}
       t^{3/2}\quad(t\geq0).
\]
Thus the distribution function has right derivative zero at zero;
the continuous density furnished by
Theorem~\ref{fg:fixed-matrix-continuity} satisfies $p_Q(0)=0$.
The comparison concerns the lower tail of a nonnegative homogeneous
form.  It does not identify the law of $Q$ after spectral truncation
or prove a density bound at positive levels, and it does not apply at
zero when an affine term moves zero away from the lower endpoint.
\end{remark}

\subsection{Concentration and conditioning}
For a real random variable $Z$, define
\[
 L_Z(u)=\sup_s\mathbb P(|Z-s|\leq u),\qquad
 \alpha_1(Z)=\sup_{u>0}\frac{L_Z(u)}{\sqrt u},\qquad
 \alpha_2(Z)=\sup_{0<u\leq1}
       \frac{L_Z(u)}{u(1+\log(1/u))}.
\]

\begin{lemma}[Fourier concentration]\label{sd:fejer}
For every $u>0$,
\begin{equation}\label{sd:fejer-bound}
 L_Z(u)\leq\frac{u}{2\sin^2 1}
              \int_{|t|\leq2/u}|\varphi_Z(t)|\,dt.
\end{equation}
If $|\varphi_Z(t)|\leq B|t|^{-1/2}$ for $t\ne0$, then
$\alpha_1(Z)\leq2\sqrt2 B/\sin^2 1$.
If $|\varphi_Z(t)|\leq\min(1,B/|t|)$ and $B\geq2$, then
$\alpha_2(Z)\leq B/\sin^2 1$.
\end{lemma}

\begin{proof}
The function $K(x)=(\sin x/x)^2$, with $K(0)=1$, is at least
$\sin^2 1$ on $[-1,1]$.  Its Fourier transform, with kernel $e^{-itx}$,
is $\pi(1-|t|/2)_+$.  Fourier inversion and Fubini's theorem therefore give
\[
 \mathbb P(|Z-s|\leq u)
 \leq\frac{\mathbb E K((Z-s)/u)}{\sin^2 1}
 \leq\frac{u}{2\sin^2 1}
            \int_{|t|\leq2/u}|\varphi_Z(t)|\,dt.
\]
Integrating $B|t|^{-1/2}$ proves the first consequence.
For $0<u\leq\min(1,2/B)$, integration of $\min(1,B/|t|)$ gives
\[
 L_Z(u)\leq\frac{Bu}{\sin^2 1}
       \left(1+\log\frac{2}{Bu}\right)
 \leq\frac{Bu}{\sin^2 1}(1+\log(1/u)).
\]
For $2/B<u\leq1$, the last inequality follows from $L_Z(u)\leq1$
and $Bu>2$.
\end{proof}

\begin{lemma}[Square-root profiles]\label{sd:composition}
Suppose that, for every nonnegative Borel function $h$,
\begin{equation}\label{sd:profile}
 \mathbb E h(Q)\leq
 K\,\mathbb E\int_{\mathbb R}h(y)|y-\xi|^{-1/2}\,dy.
\end{equation}
Then $Q$ has a density.  Moreover:
\begin{enumerate}
\item $\alpha_1(Q)\leq4K$;
\item if $\alpha_1(\xi)\leq K_1$, then
$\alpha_2(Q)\leq4\sqrt2 K\max(K_1,1)$;
\item if $\alpha_2(\xi)\leq K_2$ and $K_2\geq1$, then
\begin{equation}\label{sd:composition-density}
 \|p_Q\|_\infty\leq3K\sqrt{K_2(3+\log K_2)}.
\end{equation}
\end{enumerate}
\end{lemma}

\begin{proof}
Tonelli's theorem shows that the right side of \eqref{sd:profile}
vanishes when $h$ is the indicator of a Lebesgue null set.  Thus $Q$
has a density, with
$p_Q(y)\leq K\mathbb E|y-\xi|^{-1/2}$ almost everywhere.
The integral of $|y-\xi|^{-1/2}$ over an interval of length $2u$
is at most $4\sqrt u$, proving the first assertion.

For any $d\in\mathbb R$,
\begin{equation}\label{sd:interval-kernel}
 \int_{-u}^u|y+d|^{-1/2}\,dy
 \leq4\sqrt2\,u\max(|d|,u)^{-1/2}.
\end{equation}
For $|d|\leq2u$ this follows from the preceding integral bound.
For $|d|>2u$, use $|y+d|\geq|d|/2$.
When $0<u\leq1$, layer cake gives
\begin{align*}
 \mathbb E\max(|s-\xi|,u)^{-1/2}
 &=\int_0^{u^{-1/2}}\mathbb P(|s-\xi|<t^{-2})\,dt\\
 &\leq1+K_1\int_1^{u^{-1/2}}\frac{dt}{t}\\
 &\leq\max(K_1,1)(1+\log(1/u)).
\end{align*}
Combine this with \eqref{sd:profile} and \eqref{sd:interval-kernel}.

For the last assertion, put $V=\sqrt{K_2(3+\log K_2)}\geq1$.
For every $s$,
\begin{align*}
 \mathbb E|s-\xi|^{-1/2}
 &\leq V+K_2\int_V^\infty t^{-2}(1+2\log t)\,dt\\
 &=V+K_2(3+2\log V)/V.
\end{align*}
Writing $a=3+\log K_2\geq3$ gives
$3+2\log V=a+\log a\leq2a$.
The last display is at most $3V$, proving the claim.
\end{proof}

\begin{lemma}[Positive pivots]\label{sd:pivot}
Suppose $A$ is positive semidefinite and $a=a_{ii}>0$.
Completing the square gives
$Q=a(X_i-m(X_{-i}))^2+R(X_{-i})$, with $m$ affine and $R$ quadratic.
The matrix of $R$ is the positive semidefinite Schur complement
\[
 S=A_{-i,-i}-a^{-1}A_{-i,i}A_{i,-i}.
\]
The profile \eqref{sd:profile} holds with $K=a^{-1/2}$ and $\xi=R$.
For $k\geq2$,
\begin{equation}\label{sd:schur-eigenvalue}
 \lambda_{k-1}(S)\geq\lambda_k(A).
\end{equation}
If $A$ has rank $r$, then $S$ has rank $r-1$.
\end{lemma}

\begin{proof}
For fixed $X_{-i}$, the density of $a(X_i-m)^2$ at $w>0$ is
\[
 \frac{f_i(m+\sqrt{w/a})+f_i(m-\sqrt{w/a})}{2\sqrt{aw}}
 \leq(aw)^{-1/2}.
\]
Integration and averaging prove the profile.  No independence between
$m$ and $R$ is needed.  Writing
$A=\left(\begin{smallmatrix}a&v^T\\v&D\end{smallmatrix}\right)$, the identity
\[
 (x,y)^TA(x,y)=a(x+a^{-1}v^Ty)^2+y^TSy
\]
proves positivity of $S$ and the rank assertion.
A $k$-dimensional subspace where the form of $A$ is at least
$\lambda_k(A)$ times the squared norm intersects $ax+v^Ty=0$
in dimension at least $k-1$.  Projection onto $y$ is injective on
this intersection, and on its image
$y^TSy\geq\lambda_k(A)(x^2+|y|^2)\geq\lambda_k(A)|y|^2$.
The variational principle proves \eqref{sd:schur-eigenvalue}.
\end{proof}

\subsection{The decoupling normalization}
For a $p\times q$ matrix $M$ of rank at least $k$, set
\[
 Z_k=e_k(M^TM)
 =\sum_{\substack{S\subseteq[p],\,T\subseteq[q]\\|S|=|T|=k}}
       \det(M_{S,T})^2,\qquad
 w_{S,T}=\det(M_{S,T})^2/Z_k.
\]
Let $\alpha_i=\sum_{S\ni i,T}w_{S,T}$ and
$\beta_j=\sum_{S,T\ni j}w_{S,T}$ be the row and column inclusion
probabilities.  If $M=A_{I,J}$ for disjoint coordinate sets $I,J$
and $\alpha_i\leq1/2$ for every $i\in I$, Lemma~\ref{four:volume} gives
\begin{equation}\label{sd:volume-decay}
 |\varphi_Q(t)|^2\leq
 e^{k/2}(2\pi)^k|t|^{-k}Z_k^{-1/2}.
\end{equation}
Arbitrary quadratic terms inside $I$ and $J$, and arbitrary linear
terms, are allowed.  For $k=3$, Fourier inversion gives
\begin{equation}\label{sd:volume-density}
 \|p_Q\|_\infty\leq6\sqrt e\,Z_3^{-1/6}.
\end{equation}
Indeed, with $B=(2\pi)^{3/2}e^{3/4}Z_3^{-1/4}$,
\[
 (2\pi)^{-1}\int_{\mathbb R}\min(1,B|t|^{-3/2})\,dt
 =3B^{2/3}/\pi=6\sqrt e\,Z_3^{-1/6}.
\]
The analytic input in \eqref{sd:volume-decay} is the integrated
Fourier estimate of Lemma~\ref{four:lp} and the Gaussian
Brascamp--Lieb theorem \citep{BL1976,Lieb1990}.  No pointwise decay
assumption on an individual characteristic function is used.

\subsection{A partition in Gram coordinates}
For positive semidefinite $A$ of rank $r$, choose
\[
 A=BB^T,\qquad B^TB=\Lambda
 =\operatorname{diag}(\lambda_1,\ldots,\lambda_r)>0.
\]
Write $b_i^T$ for row $i$ of $B$, so that $|b_i|^2=a_{ii}$.
For independent equiprobable signs $\varepsilon_i$, define
\[
 I=\{i:\varepsilon_i=1\},\quad J=\{i:\varepsilon_i=-1\},\quad
 N=\tfrac12\sum_i\varepsilon_i b_ib_i^T.
\]
Then $H_I=B_I^TB_I=\Lambda/2+N$ and
$H_J=B_J^TB_J=\Lambda/2-N$.

\begin{lemma}[Partition identity]\label{sd:partition-identity}
For every $g\in\mathbb R^r$,
\begin{align}
 \langle B_Ig,A_{I,J}B_Jg\rangle
 &=|\Lambda g|^2/4-|Ng|^2,\label{sd:gram-identity}\\
 |B_Ig|\,|B_Jg|&\leq g^T\Lambda g/2.\label{sd:gram-product}
\end{align}
If $K\subseteq[r]$ has cardinality $k$, $\lambda_m\geq\theta>0$
for $m\in K$, and
$T_K=\sum_{m\in K}|Ne_m|^2/\lambda_m^2\leq1/8$, then
$s_k(A_{I,J})\geq\theta/4$.
If $\max_i a_{ii}\leq d$, then
\[
 \mathbb ET_K\leq kd/(4\theta),\qquad
 \mathbb P(T_K>1/8)\leq2kd/\theta.
\]
\end{lemma}

\begin{proof}
The scalar product is $g^TH_IH_Jg$.  Expanding $H_IH_J$ gives
$\Lambda^2/4-N^2+(N\Lambda-\Lambda N)/2$, whose last term is
skew symmetric.  This proves \eqref{sd:gram-identity}.
Equation~\eqref{sd:gram-product} follows from $H_I+H_J=\Lambda$.
For $g$ in the span $G$ of $\{e_m:m\in K\}$,
Cauchy--Schwarz gives $|Ng|\leq|\Lambda g|T_K^{1/2}$.
Thus the scalar product is at least $|\Lambda g|^2/8$.
It is positive for $g\ne0$, so $B_J$ is injective on $G$.
For $x=B_Jg\ne0$,
\[
 \frac{|A_{I,J}x|}{|x|}
 \geq\frac{\langle B_Ig,A_{I,J}B_Jg\rangle}{|B_Ig|\,|B_Jg|}
 \geq\frac{|\Lambda g|^2}{4g^T\Lambda g}\geq\theta/4.
\]
The singular-value variational principle on the $k$-dimensional
subspace $B_JG$ proves the assertion.  Finally,
\[
 \mathbb E|Ne_m|^2
 =\tfrac14\sum_i b_i(m)^2|b_i|^2
 \leq\tfrac d4\sum_i b_i(m)^2=d\lambda_m/4.
\]
Sum over $m$ and apply Markov's inequality.
\end{proof}

\subsection{Volume marginals and ridge leverage}

\begin{lemma}[Volume marginals]\label{sd:ridge}
Let $G=MM^T$ have eigenvalues $\sigma_1\geq\sigma_2\geq\cdots\geq0$
and an orthonormal eigenbasis $(u_m)$.  If $\operatorname{rank}M\geq k$,
then
\begin{equation}\label{sd:marginal-spectrum}
 \alpha_i=\sum_m p_m u_m(i)^2,\qquad
 p_m=\frac{\sigma_m e_{k-1}(\sigma_{-m})}{e_k(\sigma)},
\end{equation}
where $e_0=1$ and $\sigma_{-m}$ is the list with entry $m$ removed.
Moreover $p_m\leq\min(1,k\sigma_m/\sigma_k)$.
If $0<\kappa\leq\sigma_k/k$ and $m_i=M^Te_i$, then
\begin{equation}\label{sd:ridge-bound}
 \alpha_i\leq2m_i^T(M^TM+\kappa I)^{-1}m_i.
\end{equation}
\end{lemma}

\begin{proof}
By Cauchy--Binet, the numerator of $\alpha_i$ is
$\sum_{S\ni i,\,|S|=k}\det G_{S,S}$.
Differentiate $\det(I+tD(z)G)$ with respect to $z_i$ at
$z=(1,\ldots,1)$, where $D(z)=\operatorname{diag}(z)$, and take
the coefficient of $t^k$.  The principal-minor expansion gives
this numerator, whereas spectral decomposition gives the
coefficient in
\[
 \sum_m u_m(i)^2\,t\sigma_m
                     \prod_{\ell\ne m}(1+t\sigma_\ell).
\]
This proves \eqref{sd:marginal-spectrum}.
The identity
$e_k(\sigma)=\sigma_me_{k-1}(\sigma_{-m})+e_k(\sigma_{-m})$
gives $p_m\leq1$.  For $m\leq k$ it also gives
$p_m\leq k\sigma_m/\sigma_k$.  Otherwise the list
$\mu=\sigma_{-m}$ contains $k$ entries at least $\sigma_k$, and
\[
 k e_k(\mu)
 =\sum_{|S|=k-1}\left(\prod_{j\in S}\mu_j\right)
                       \sum_{j\notin S}\mu_j
 \geq\sigma_k e_{k-1}(\mu).
\]
Indeed, each $(k-1)$-element set omits one of those $k$ entries.
This proves the remaining bound for $p_m$.
The inequality $\min(1,x)\leq2x/(1+x)$ now yields
$p_m\leq2\sigma_m/(\sigma_m+\kappa)$.
The spectral decomposition of
$M(M^TM+\kappa I)^{-1}M^T$ proves \eqref{sd:ridge-bound}.
\end{proof}

\begin{lemma}[Gram comparison]\label{sd:ridge-comparison}
For $M=A_{I,J}=B_IB_J^T$, $i\in I$, and $\kappa>0$,
\begin{equation}\label{sd:ridge-gram}
 m_i^T(M^TM+\kappa I)^{-1}m_i
 \leq b_i^T(H_I+\kappa\Lambda^{-1})^{-1}b_i .
\end{equation}
\end{lemma}

\begin{proof}
The identity
$(B_JH_IB_J^T+\kappa I)B_J=B_J(H_IH_J+\kappa I)$ gives
\[
 m_i^T(M^TM+\kappa I)^{-1}m_i
 =b_i^TH_J(H_IH_J+\kappa I)^{-1}b_i .
\]
If $H_J>0$, the middle matrix is
$(H_I+\kappa H_J^{-1})^{-1}$.
Since $H_J\leq\Lambda$, order reversal under inversion gives
\eqref{sd:ridge-gram}.  If $H_J$ is singular, replace it in this
matrix identity by $H_J+\varepsilon I\leq\Lambda+\varepsilon I$
and let $\varepsilon\downarrow0$.
Continuity at a positive semidefinite matrix $T$ follows from the
expression
$T^{1/2}(T^{1/2}H_IT^{1/2}+\kappa I)^{-1}T^{1/2}$.
\end{proof}

\begin{proposition}[A semidefinite partition]\label{sd:psd-partition}
For $k\in\{1,2,3\}$, let $\theta=\lambda_k(A)>0$, with $A$
positive semidefinite, and put
$\gamma_1=128$, $\gamma_2=300$, $\gamma_3=640$.
If $\max_i a_{ii}\leq\theta/\gamma_k$, some partition
$[n]=I\sqcup J$ satisfies
\[
 s_k(A_{I,J})\geq\theta/4,\qquad
 \max_{i\in I}\alpha_i
 \leq(64k/\theta)\max_i a_{ii}\leq1/2.
\]
Uniformly in the linear and constant terms, this implies
\begin{align}
 k=1:\quad&
 |\varphi_Q(t)|\leq2e^{1/4}\sqrt{2\pi}\,
                         (\theta|t|)^{-1/2},\label{sd:psd-k1}\\
 k=2:\quad&
 |\varphi_Q(t)|\leq8\pi\sqrt e\,(\theta|t|)^{-1},
                         \label{sd:psd-k2}\\
 k=3:\quad&
 \|p_Q\|_\infty\leq24\sqrt e/\theta.\label{sd:psd-k3}
\end{align}
\end{proposition}

\begin{proof}
Put $d=\max_i a_{ii}$.  Split the Gram-coordinate space into its
first $k$ coordinates, denoted by the subscript $b$, and the
remaining coordinates, denoted by $s$.  These subscripts denote
blocks, not the linear coefficient of $Q$.  Set
\[
 \kappa=\theta^2/(16k),\qquad \rho=\kappa/\theta=\theta/(16k),
 \qquad \ell_i=|\Lambda_b^{-1/2}b_i^b|^2.
\]
Then $\sum_i\ell_i=k$ and $\ell_i\leq d/\theta$.
Consider the events
\begin{align*}
 \mathcal S:\quad&
 \sum_{m=1}^k |Ne_m|^2/\lambda_m^2\leq1/8,\\
 \mathcal E_1:\quad&
 \|\Lambda_b^{-1/2}N_{bb}\Lambda_b^{-1/2}\|_F^2\leq1/64,\\
 \mathcal E_2:\quad&
 (2/\rho)\|N_{sb}\Lambda_b^{-1/2}\|_F^2\leq1/8.
\end{align*}
An empty $s$ block contributes zero.
On $\mathcal S$, Lemma~\ref{sd:partition-identity} gives
$s_k(M)\geq\theta/4$, so $\kappa\leq s_k(M)^2/k$.
Lemmas~\ref{sd:ridge} and \ref{sd:ridge-comparison} give
$\alpha_i\leq2b_i^T(H_I+\kappa\Lambda^{-1})^{-1}b_i$.

For $T>0$, completing a square gives the block inequality
\begin{equation}\label{sd:block-comparison}
 \begin{pmatrix}P&R\\R^T&T\end{pmatrix}
 \geq
 \begin{pmatrix}P-2RT^{-1}R^T&0\\0&T/2\end{pmatrix}.
\end{equation}
The quadratic form of the difference is
$2|T^{-1/2}R^Tx+T^{1/2}y/2|^2$.
As $\Lambda_s\leq\theta I$ and $(H_I)_{ss}\geq0$,
\eqref{sd:block-comparison}, with $T=\rho I$, yields
\[
 H_I+\kappa\Lambda^{-1}
 \geq
 \left(\Lambda_b/2+N_{bb}-(2/\rho)N_{bs}N_{sb}\right)
       \oplus(\rho/2)I_s.
\]
On $\mathcal E_1$, $N_{bb}\geq-\Lambda_b/8$.
On $\mathcal E_2$, conjugating by $\Lambda_b^{-1/2}$ and bounding
operator norm by trace gives
$(2/\rho)N_{bs}N_{sb}\leq\Lambda_b/8$.
Therefore
$H_I+\kappa\Lambda^{-1}\geq\Lambda_b/4\oplus(\rho/2)I_s$ and
\[
 \alpha_i\leq8\ell_i+(4/\rho)|b_i^s|^2
 \leq(64k/\theta)|b_i|^2=(64k/\theta)a_{ii}.
\]

It remains to show that these events can hold simultaneously.
Lemma~\ref{sd:partition-identity} gives
$\mathbb P(\mathcal S^c)\leq2kd/\theta$.
Independence and centering of the signs also give
\begin{align*}
 \mathbb E\|\Lambda_b^{-1/2}N_{bb}\Lambda_b^{-1/2}\|_F^2
 &=\tfrac14\sum_i\ell_i^2\leq kd/(4\theta),\\
 \mathbb E\|N_{sb}\Lambda_b^{-1/2}\|_F^2
 &=\tfrac14\sum_i|b_i^s|^2\ell_i\leq kd/4.
\end{align*}
Markov's inequality gives failure bounds $16kd/\theta$ and
$64k^2d/\theta$ for $\mathcal E_1$ and $\mathcal E_2$.
The total is at most
\[
 (18k+64k^2)/\gamma_k
 =
 \begin{cases}82/128,&k=1,\\292/300,&k=2,\\630/640,&k=3,\end{cases}
\]
which is strictly below one.  The marginal bound is at most
$1/2$ for each stated $\gamma_k$.
Since $Z_k\geq\prod_{j=1}^k s_j(M)^2\geq(\theta/4)^{2k}$,
\eqref{sd:volume-decay}--\eqref{sd:volume-density} prove
\eqref{sd:psd-k1}--\eqref{sd:psd-k3}.
\end{proof}

\subsection{Completion of the semidefinite proof}

\begin{proposition}[The first two levels]\label{sd:psd-levels}
For positive semidefinite $A$ and arbitrary $b,c$:
\begin{enumerate}
\item if $\lambda_1(A)\geq1$, then
$\alpha_1(Q)\leq K_1:=4\sqrt{128}<45.26$;
\item if $\lambda_2(A)\geq1$, then
$\alpha_2(Q)\leq K_2:=4\sqrt{600}\,K_1<4435$.
\end{enumerate}
\end{proposition}

\begin{proof}
If some $a_{ii}\geq1/128$, Lemmas~\ref{sd:pivot} and
\ref{sd:composition}(1) give the first bound.
Otherwise Proposition~\ref{sd:psd-partition}, with $k=1$,
and Lemma~\ref{sd:fejer} give
\[
 \alpha_1(Q)\leq
 (2\sqrt2/\sin^2 1)\,2e^{1/4}\sqrt{2\pi}<25.72<K_1.
\]
For the second assertion, a coefficient $a_{ii}\geq1/300$ leaves
a positive semidefinite Schur complement with an eigenvalue at least
one.  Apply the first assertion to its quadratic form, which is the
random center in the profile.  Lemma~\ref{sd:composition}(2)
gives $\alpha_2(Q)\leq4\sqrt2\sqrt{300}\,K_1=K_2$.
If there is no such coefficient,
\eqref{sd:psd-k2} and Lemma~\ref{sd:fejer} give
$\alpha_2(Q)\leq8\pi\sqrt e/\sin^2 1<58.53<K_2$.
\end{proof}

\begin{proof}[Proof of Theorem~\ref{sd:main}]
Suppose first that $\lambda_3(A)\geq1$.  A coefficient
$a_{ii}\geq1/640$ leaves a positive semidefinite Schur complement
with at least two eigenvalues at least one.
Apply Proposition~\ref{sd:psd-levels}(2) to its quadratic form and
Lemma~\ref{sd:composition}(3) to obtain
\[
 \|p_Q\|_\infty\leq
 3\sqrt{640}\sqrt{K_2(3+\log K_2)}<17100.
\]
Here $K_2=4\sqrt{600}\,4\sqrt{128}$ is used without rounding.
If every diagonal coefficient is below $1/640$,
Proposition~\ref{sd:psd-partition}, with $k=3$, instead gives
$\|p_Q\|_\infty\leq24\sqrt e$.
For $\lambda_3(A)=\theta>0$, apply this argument to $Q/\theta$
and use density scaling.  Applying the result to $-Q$ proves
the negative semidefinite case.
\end{proof}

\subsection{Indefinite forms with a small absolute spectral diagonal}
For a real symmetric $A$, write $|A|=(A^2)^{1/2}$ and let
$|\lambda|_{(1)}(A)\geq|\lambda|_{(2)}(A)\geq\cdots$
be its eigenvalue moduli.  In particular $(|A|)_{ii}\geq|a_{ii}|$.

\begin{theorem}[Absolute spectral spread]\label{sd:indefinite}
Let $k\in\{1,2,3\}$, set $\theta=|\lambda|_{(k)}(A)>0$, and put
$\gamma'_1=130$, $\gamma'_2=360$, $\gamma'_3=680$.
If
\begin{equation}\label{sd:absolute-spread}
 \max_i(|A|)_{ii}\leq\theta/\gamma'_k,
\end{equation}
some partition $[n]=I\sqcup J$ satisfies
$s_k(A_{I,J})\geq23\theta/80$ and
$\max_{i\in I}\alpha_i\leq1/2$.
Consequently, for every $b,c$,
\begin{equation}\label{sd:indefinite-fourier}
 |\varphi_Q(t)|\leq
 e^{k/4}(2\pi)^{k/2}(23\theta|t|/80)^{-k/2}.
\end{equation}
For $k=3$, $Q$ has a density and
\begin{equation}\label{sd:indefinite-density}
 \|p_Q\|_\infty\leq \frac{480\sqrt e}{23\theta}
 <\frac{35}{\theta}.
\end{equation}
\end{theorem}

\begin{proof}
On the range of $A$, factor
\[
 A=B\Sigma B^T,\qquad B^TB=\Lambda>0,\qquad
 \Sigma=\operatorname{diag}(\pm1),
\]
where the entries of $\Lambda$ are the nonzero eigenvalue moduli
in decreasing order.  Thus $BB^T=|A|$ and
$|b_i|^2=(|A|)_{ii}\leq d:=\max_i(|A|)_{ii}$.
Define the random partition and $N$ as before.  Then
$H_I=\Lambda/2+N$, $H_J=\Lambda/2-N$, and
$M=A_{I,J}=B_I\Sigma B_J^T$.
Let $G$ be the first $k$ Gram coordinates.
Decompose $g\in G$ as $g=g_++g_-$ according to the signs in $\Sigma$.
Direct multiplication gives
\begin{align}
 \langle B_I\Sigma g,MB_Jg\rangle
 &=|\Lambda g|^2/4-g_-^T\Lambda Ng_+
      -g_+^T\Lambda Ng_-
      -(\Sigma N\Sigma g)^TNg,\label{sd:signed-identity}\\
 |B_I\Sigma g|^2+|B_Jg|^2
 &=g^T\Lambda g-4g_+^TNg_-.\label{sd:signed-norm}
\end{align}
Suppose
\begin{equation}\label{sd:signed-events}
 T:=\sum_{m=1}^k |Ne_m|^2/\lambda_m^2\leq\eta^2,\qquad
 \|\Lambda_b^{-1/2}N_{bb}\Lambda_b^{-1/2}\|_F\leq\eta',
\end{equation}
where $b$ denotes the first $k$ Gram coordinates.
Then $|Ng|\leq\eta|\Lambda g|$ for $g\in G$.
The sum of the two middle terms in \eqref{sd:signed-identity}
has modulus at most $\eta|\Lambda g|^2$, by applying this bound
separately to $g_+$ and $g_-$ and then using $2ab\leq a^2+b^2$.
The last term has modulus at most $2\eta^2|\Lambda g|^2$
(the smaller coefficient $\eta^2$ would also suffice).
Since
$|g_+^TNg_-|\leq
\eta'|\Lambda^{1/2}g_+|\,|\Lambda^{1/2}g_-|$,
\eqref{sd:signed-norm} gives
\[
 |B_I\Sigma g|\,|B_Jg|
 \leq(1/2+\eta')g^T\Lambda g.
\]
If $1/4-\eta-2\eta^2>0$, positivity of the scalar product makes
$B_J$ injective on $G$.  Testing against $B_I\Sigma g$ and applying
the singular-value variational principle gives
\begin{equation}\label{sd:signed-sv}
 s_k(M)\geq
 \theta\,\frac{1/4-\eta-2\eta^2}{1/2+\eta'}.
\end{equation}
Take $\eta=1/16$ and $\eta'=1/8$; the last ratio is $q=23/80$.

The ridge comparison carries over to the signed factorization.
The calculation of Lemma~\ref{sd:ridge-comparison}, with
$\Sigma H_I\Sigma$ and $\Sigma b_i$, gives
\begin{align*}
 m_i^T(M^TM+\kappa I)^{-1}m_i
 &\leq(\Sigma b_i)^T
    (\Sigma H_I\Sigma+\kappa\Lambda^{-1})^{-1}(\Sigma b_i)\\
 &=b_i^T(H_I+\kappa\Lambda^{-1})^{-1}b_i.
\end{align*}
Here $\Sigma$ commutes with $\Lambda$; the same limiting argument
handles a singular $H_J$.
Set $\kappa=q^2\theta^2/k$ and $\rho=\kappa/\theta=q^2\theta/k$,
and additionally require
\[
 (2/\rho)\|N_{sb}\Lambda_b^{-1/2}\|_F^2\leq1/8.
\]
On these events, $\kappa\leq s_k(M)^2/k$, and the block comparison
in Proposition~\ref{sd:psd-partition} gives
$H_I+\kappa\Lambda^{-1}\geq\Lambda_b/4\oplus(\rho/2)I_s$.
Thus
\[
 \alpha_i\leq8\ell_i+(4/\rho)|b_i^s|^2
       \leq\frac{4k}{q^2\theta}(|A|)_{ii}.
\]

All second moments used in Proposition~\ref{sd:psd-partition}
depend only on $|b_i|^2\leq d$ and are unchanged.
Markov's inequality bounds the failure probabilities of the
three events by
$64kd/\theta$, $16kd/\theta$, and $4k^2d/(q^2\theta)$.
Their sum is at most
\[
 \frac{80k+(25600/529)k^2}{\gamma'_k}<1,\qquad k=1,2,3.
\]
The three ratios are less than $0.988$, $0.983$, and $0.994$.
Also $4k/(q^2\gamma'_k)<1/2$ in every case.
A good partition therefore exists.  Since $Z_k\geq(q\theta)^{2k}$,
\eqref{sd:volume-decay}--\eqref{sd:volume-density}
prove the theorem.
\end{proof}

\begin{remark}\label{sd:scope}
Condition \eqref{sd:absolute-spread} concerns the diagonal of
$|A|$, not merely that of $A$.  The positive and negative spectral
contributions to $a_{ii}$ can cancel for an indefinite matrix.
Such cancellation is absent from $(|A|)_{ii}$.
Thus Theorem~\ref{sd:indefinite} does not settle the unrestricted
indefinite case.
\end{remark}

\subsection{A smaller constant for a separated nonzero spectrum}
The following argument gives a smaller constant when every nonzero
eigenvalue, rather than just the third, is bounded away from zero.

\begin{lemma}[A spread spectral frame]\label{sd:frame}
Let $R$ be an $n\times r$ matrix with $R^TR=I_r$ and let
$\Lambda$ be a nonsingular real diagonal $r\times r$ matrix.
If every row $r_i^T$ of $R$ satisfies $|r_i|^2\leq1/2$, then
for $Q=X^TR\Lambda R^TX+b^TX+c$,
\begin{equation}\label{sd:frame-fourier}
 |\varphi_Q(t)|
 \leq(4\pi)^{r/2}|t|^{-r/2}|\det\Lambda|^{-1/2}.
\end{equation}
\end{lemma}

\begin{proof}
Put $c_i=|r_i|^2$ and $u_i=r_i/\sqrt{c_i}$ when $c_i>0$.
Then $\sum_i c_i u_i u_i^T=I_r$ and $\sum_i c_i=r$.
The geometric Brascamp--Lieb inequality \citep{BL1976,Lieb1990}
applied to
$F_i(s)=|\varphi_{X_i}(\sqrt{c_i}s+tb_i)|^{1/c_i}$ gives
\begin{align*}
 \int_{\mathbb R^r}
  \prod_i|\varphi_{X_i}(\langle r_i,u\rangle+tb_i)|\,du
 &\leq\prod_{i:c_i>0}\left(\int_{\mathbb R}F_i(s)\,ds\right)^{c_i}\\
 &\leq(4\pi)^r.
\end{align*}
The last inequality is Lemma~\ref{four:lp} with exponent
$1/c_i\geq2$, after changing variables.  Rows with $c_i=0$
contribute factors of modulus at most one.

For a justified quadratic Fourier inversion, put $Y=R^TX$ and
regularize the quadratic exponential as
\[
 e^{itY^T\Lambda Y-\varepsilon|Y|^2},\qquad \varepsilon>0.
\]
The Fourier transform with kernel $e^{-iu^Ty}$ of this complex
Gaussian has modulus at most
$\pi^{r/2}|t|^{-r/2}|\det\Lambda|^{-1/2}$.
Fourier inversion and independence therefore bound the
regularized expectation by
\[
 (4\pi|t|)^{-r/2}|\det\Lambda|^{-1/2}
 \int_{\mathbb R^r}
   \prod_i|\varphi_{X_i}(\langle r_i,u\rangle+tb_i)|\,du.
\]
For $\varepsilon>0$ the Gaussian transform is integrable, which
justifies the initial use of Fubini.  The preceding
Brascamp--Lieb bound supplies an integrable majorant for the
product after multiplication by the uniformly bounded Gaussian
Fourier factor.  Let $\varepsilon\downarrow0$ to prove
\eqref{sd:frame-fourier}.
\end{proof}

\begin{theorem}[Separated nonzero semidefinite spectrum]\label{sd:separated}
Suppose $A$ is positive semidefinite of rank $r\geq3$ and every
nonzero eigenvalue is at least $\theta>0$.  Then
\begin{equation}\label{sd:separated-bound}
 \|p_Q\|_\infty\leq130/\theta
\end{equation}
for all $b,c$.  The same assertion holds if all nonzero
eigenvalues are at most $-\theta$.
\end{theorem}

\begin{proof}
Scale so that $\theta=1$, and write $A=R\Lambda R^T$ with
$R^TR=I_r$ and $\Lambda\geq I_r$.  Then $|r_i|^2\leq a_{ii}$.
Use a positive pivot if some $a_{ii}\geq1/2$.
Otherwise Lemma~\ref{sd:frame} applies, and
\[
 |\varphi_Q(t)|
 \leq\min(1,(4\pi/|t|)^{r/2})
 \leq\min(1,(4\pi/|t|)^{j/2})\quad(r\geq j).
\]
For rank at least one, a pivot gives $\alpha_1(Q)\leq4\sqrt2<15$.
Without a pivot, Lemma~\ref{sd:fejer} gives
$\alpha_1(Q)\leq2\sqrt2\sqrt{4\pi}/\sin^2 1<15$.
Thus $\alpha_1\leq15$ throughout this spectral class.

For rank at least two, a pivot leaves a Schur complement of rank
one less with all nonzero eigenvalues at least one.
The first level and Lemma~\ref{sd:composition}(2) give
$\alpha_2(Q)\leq4\sqrt2\sqrt2\cdot15=120$.
Without a pivot, Lemma~\ref{sd:fejer} gives
$\alpha_2(Q)\leq4\pi/\sin^2 1<18$.

For rank at least three, a pivot and the second level give
\[
 \|p_Q\|_\infty\leq3\sqrt2\sqrt{120(3+\log120)}<130.
\]
Without a pivot, integrating
$\min(1,(4\pi/|t|)^{3/2})$ gives $\|p_Q\|_\infty\leq12$.
Scaling and reflection complete the proof.
\end{proof}

\subsection{Sharper bounds in definite rank three}

\begin{theorem}[Exact definite rank three]\label{sd:rank-three}
Suppose $A$ has rank three and its nonzero eigenvalues are all
at least $\theta>0$, or all at most $-\theta$.  Then, for every
$b\in\mathbb R^n$ and $c\in\mathbb R$,
\begin{equation}\label{sd:rank-three-bound}
 \|p_{X^TAX+b^TX+c}\|_\infty\leq58/\theta.
\end{equation}
If in addition $b\in\operatorname{range}(A)$, the stronger bound
\begin{equation}\label{sd:rank-three-range}
 \|p_{X^TAX+b^TX+c}\|_\infty\leq3e/\theta
\end{equation}
holds.
\end{theorem}

\begin{proof}
Reflection and scaling reduce the assertion to $A\geq0$,
$\operatorname{rank}(A)=3$, and $\lambda_3(A)\geq1$.
Write $A=R\Lambda R^T$ with $R^TR=I_3$ and $\Lambda\geq I_3$.
If every row $r_i^T$ of $R$ satisfies $|r_i|^2\leq1/2$,
Theorem~\ref{four:diffuse}, with the individual linear terms as
separable phases, gives
\[
 \|p_Q\|_\infty\leq3e\,(\det\Lambda)^{-1/3}\leq3e.
 \tag*{}
\]
This proves both assertions in this case.

Otherwise choose $i$ with $|r_i|^2>1/2$.  Then
$a=a_{ii}\geq|r_i|^2>1/2$.
Lemma~\ref{sd:pivot} writes
$Q=a(X_i-m(X_{-i}))^2+T(X_{-i})$.
The matrix of $T$ has rank two, with both nonzero eigenvalues
at least one; its linear term may be arbitrary.
Theorem~\ref{thm:rank-two-affine} therefore gives
$\alpha_2(T)\leq46$.  The conditional square density yields
\[
 \|p_Q\|_\infty
 \leq \sqrt2\sup_s\mathbb E|s-T|^{-1/2}.
 \tag*{}
\]
For any $V\geq1$, the layer-cake calculation used in
Lemma~\ref{sd:composition} gives
\[
 \mathbb E|s-T|^{-1/2}
 \leq V+\frac{46(3+2\log V)}{V}.
 \tag*{}
\]
Taking $V=20$ proves \eqref{sd:rank-three-bound}, since
\[
 \sqrt2\left(20+\frac{46(3+2\log20)}{20}\right)
 <57.531<58.
 \tag*{}
\]

Now suppose $b\in\operatorname{range}(A)$.  Choose $z$ with
$2Az=b$ and replace $X$ by $X+z$, absorbing the resulting
constant into $c$.  Translation preserves independence and all
coordinate density caps, so it is enough to treat $b=0$.
In the heavy-row case the Schur complement is a positive
semidefinite matrix of rank two with nonzero eigenvalues
$\mu_1,\mu_2\geq1$.  Thus
\[
 T=Y^T\operatorname{diag}(\mu_1,\mu_2)Y+c',
 \qquad Y=U^TX_{-i},\quad U^TU=I_2.
 \tag*{}
\]
Lemma~\ref{four:projection} bounds the joint density of $Y$ by
$2$.  Lemma~\ref{geom:definite} consequently gives
$\|p_T\|_\infty\leq2\pi$.
Keep the support restriction in the square density formula:
\[
 p_Q(s)\leq a^{-1/2}
      \mathbb E\bigl[(s-T)^{-1/2}{\bf1}_{\{T<s\}}\bigr]
 \quad\hbox{for almost every }s.
 \tag*{}
\]
For any real random variable with density bounded by $M$,
layer cake gives
\[
 \mathbb E\bigl[(s-T)^{-1/2}{\bf1}_{\{T<s\}}\bigr]
 =\int_0^\infty\mathbb P(s-t^{-2}<T<s)\,dt
 \leq\int_0^\infty\min(1,Mt^{-2})\,dt=2\sqrt M.
 \tag*{}
\]
It follows that $\|p_Q\|_\infty\leq4\sqrt\pi<3e$.
Together with the diffuse case, this proves
\eqref{sd:rank-three-range}.  Restoring the scale $\theta$
completes the proof.
\end{proof}

\section{The first two levels without a coordinate pivot}
\label{sec:lower-levels}

Let $X=(X_1,\ldots,X_n)$ have independent coordinates with densities
bounded by one, and write
\[
 Q=X^{\mathsf T}MX+b^{\mathsf T}X+c_0,
 \qquad \varphi_Q(t)=\mathbb E e^{itQ}.
\]
Here $M$ is real symmetric, and all estimates are uniform in $b,c_0$.
The definition of a $c$-pivot is given in
Section~\ref{sec:structure}.  In particular, absence of a $c$-pivot,
with $0<c\leq1$, implies
\begin{equation}\label{lower:entry-cap}
 |m_{ii}|<c,\qquad m_{ij}^2<c+c^2\leq2c\quad(i\ne j).
\end{equation}
Set
\[
 O=M-\operatorname{diag}M,\qquad
 r_i^2=\sum_{j\ne i}m_{ij}^2,\qquad
 L_Q(u)=\sup_s\mathbb P(|Q-s|\leq u).
\]
The truncated half-power potential is
\[
 \mathcal H(Q)=\sup_s
 \mathbb E\bigl[|Q-s|^{-1/2};\ |Q-s|\leq1\bigr].
\]

\begin{theorem}[One large eigenvalue]\label{lower:one}
Suppose that $M$ has no $c$-pivot, where $c\leq2\cdot10^{-6}$,
and $\|M\|\geq1$.  Then
\[
 L_Q(u)\leq25\sqrt u\qquad(u>0).
\]
In particular,
\[
 \sup_s\mathbb E\log_+\frac1{|Q-s|}
 \leq2+2\log25<8.44.
\]
\end{theorem}

\begin{theorem}[Two large eigenvalues]\label{lower:two}
Put $c_*=10^{-15}$.  If $M$ has no $c_*$-pivot and has at least
two eigenvalues of modulus at least one, then
\[
 \mathcal H(Q)\leq10^6.
\]
\end{theorem}

The second theorem is a negative-moment estimate.  Its statement does
not require an estimate of order $u(1+\log(1/u))$ for $L_Q(u)$.
The small pivot threshold is used to make the separate cases below
quantitative.

\subsection{Fourier concentration and volume marginals}

Put $\kappa_0=4\sin^2(1/2)>0.91939$.
The Fej\'er kernel identity
\[
 \left(\frac{\sin(x/2)}{x/2}\right)^2
 =\int_{-1}^1(1-|\tau|)e^{i\tau x}\,d\tau
\]
gives, for any real random variable $Z$,
\begin{equation}\label{lower:fejer}
 L_Z(u)\leq\kappa_0^{-1}u
             \int_{|t|\leq1/u}|\varphi_Z(t)|\,dt.
\end{equation}
Indeed, the kernel is at least $\kappa_0$ on $[-1,1]$;
take its expectation at $(Z-s)/u$, use Fubini, and then take absolute
values.  In particular,
\begin{equation}\label{lower:power-concentration}
 |\varphi_Z(t)|\leq B|t|^{-\beta},\quad 0<\beta<1
 \quad\Longrightarrow\quad
 L_Z(u)\leq\frac{2B}{\kappa_0(1-\beta)}u^\beta.
\end{equation}

For disjoint coordinate sets $I,J$, let $A=M_{I,J}$.  Write
\[
 Z_k(A)=\sum_{\substack{|S|=|T|=k}}
                   \det(A_{S,T})^2=e_k(A^{\mathsf T}A),
 \qquad
 \alpha_i^{(k)}=
 \frac{\sum_{S\ni i,T}\det(A_{S,T})^2}{Z_k(A)}.
\]
The sums are over row and column subsets of $A$ of the indicated
cardinalities.  Lemma~\ref{four:volume}, with $1\leq q\leq2$,
gives
\begin{equation}\label{lower:volume-q}
 |\varphi_Q(t)|\leq
 \left[e^{k/2}(2\pi)^kq^{-k/2}Z_k(A)^{-1/2}\right]^{1/(2q)}
 |t|^{-k/(2q)},
 \qquad \max_i\alpha_i^{(k)}\leq q/2.
\end{equation}
This estimate allows arbitrary quadratic terms within $I$ and $J$.
Other coordinates may be conditioned on, and the orientation of $A$
may be reversed.  These facts follow from the decoupling step in
that lemma; the resulting bounds are uniform in the conditioned
affine terms.

\begin{lemma}[First-volume estimate]\label{lower:first-volume}
Suppose $A=M_{I,J}\ne0$, and let $1/2\leq h\leq1$.
If every row satisfies $|A_{i,\cdot}|^2\leq h\|A\|_F^2$, then
\begin{equation}\label{lower:first-volume-bound}
 |\varphi_Q(t)|\leq
 \left[2\pi e^{1/2}(2h)^{-1/2}\|A\|_F^{-1}\right]^{1/(4h)}
 |t|^{-1/(4h)}.
\end{equation}
The same assertion holds with columns in place of rows.  When $h=1$
there is no additional marginal hypothesis.
\end{lemma}
\begin{proof}
For $k=1$, the minors are the entries of $A$, so that
$Z_1(A)=\|A\|_F^2$ and
$\alpha_i^{(1)}=|A_{i,\cdot}|^2/\|A\|_F^2$.
Apply \eqref{lower:volume-q} with $q=2h$.
\end{proof}

\begin{lemma}[Choice of orientation]\label{lower:orientation}
Let $Z=\|A\|_F^2>0$ and $1/2<h<1$.  If
$A_{ij}^2\leq(2h-1)Z$ for every entry, then either every row share
is at most $h$, or every column share is at most $h$.
\end{lemma}
\begin{proof}
If a row $i$ and a column $j$ both had shares greater than $h$, then
\[
 Z\geq |A_{i,\cdot}|^2+|A_{\cdot,j}|^2-A_{ij}^2
   >2hZ-A_{ij}^2,
\]
contrary to the entry hypothesis.
\end{proof}

Combining these two lemmas with \eqref{lower:power-concentration}
at $h=3/4$ gives the following useful consequence:
\begin{equation}\label{lower:third-power}
 \begin{split}
 \|M_{I,J}\|_F\geq\rho,\qquad
 m_{ij}^2\leq\tfrac12\|M_{I,J}\|_F^2
 \quad\Longrightarrow\quad
 L_Q(u)\leq K(\rho)u^{1/3},\\
 K(\rho)=\frac3{\kappa_0}
       \left(\frac{2\pi\sqrt{2e/3}}{\rho}\right)^{1/3}
 <3.264\left(\frac{8.46}{\rho}\right)^{1/3}.
 \end{split}
\end{equation}

\subsection{Proof at the first level}

\begin{proof}[Proof of Theorem~\ref{lower:one}]
Let $F^2=\|O\|_F^2$.  The triangle inequality gives
$F\geq\|O\|\geq1-c$.  Put $\varepsilon=1/100$ and call an
index heavy if $r_i^2>\varepsilon F^2$.  There are fewer than
$1/\varepsilon$ heavy indices.  Place all of them in $J$, and place
each remaining index independently in $I$ or $J$, with equal
probabilities.  Set $Z=\|M_{I,J}\|_F^2$.

The squared entries between heavy indices have total at most
$2c/\varepsilon^2$, by \eqref{lower:entry-cap}.  Every other
off-diagonal squared entry contributes to $\mathbb EZ$ with
probability at least $1/4$, with the two orientations counted in
$F^2$.  Consequently,
\[
 \mathbb EZ\geq\frac14(F^2-2c/\varepsilon^2)
          \geq\frac{19}{80}F^2.
\]
The last step uses $c\leq\varepsilon^2F^2/40$, which follows from
$c\leq2\cdot10^{-6}$ and $F\geq1-c$.

Changing the coin at a light index $i$ changes $Z$ by at most
$r_i^2$.  Conditional variance for one fair coin, followed by
variance tensorization, gives
\[
 \operatorname{Var}Z
 \leq\frac14\sum_{i\ {
m light}}r_i^4
 \leq\frac\varepsilon4F^4.
\]
Chebyshev's inequality therefore gives
\[
 \mathbb P\{Z<\tfrac12\mathbb EZ\}
 \leq\frac{\varepsilon F^4}{(\mathbb EZ)^2}
 \leq\frac{6400\varepsilon}{361}<0.18.
\]
Choose a realization with $Z\geq\mathbb EZ/2>F^2/9$.
All its rows are light, so
\[
 \frac{|M_{i,J}|^2}{Z}\leq9\varepsilon=0.09<1/2.
\]
Lemma~\ref{lower:first-volume}, with $h=1/2$, now yields
\[
 |\varphi_Q(t)|
 \leq\left(\frac{6\pi\sqrt e}{F}\right)^{1/2}|t|^{-1/2}
 <5.6|t|^{-1/2}.
\]
Equation~\eqref{lower:power-concentration} gives
$L_Q(u)\leq(22.4/\kappa_0)\sqrt u<25\sqrt u$.
Finally, layer cake gives
\[
 \mathbb E\log_+\frac1{|Q-s|}
 =\int_0^\infty\mathbb P(|Q-s|<e^{-v})\,dv
 \leq\int_0^\infty\min(1,25e^{-v/2})\,dv
 =2+2\log25.
\]
\end{proof}

With only $\max_i|m_{ii}|\leq1/2$ and $\|M\|\geq1$, a fair
partition satisfies $\mathbb E\|M_{I,J}\|_F^2=\|O\|_F^2/4\geq1/16$.
The unrestricted $h=1$ case gives a fourth-root estimate;
Lemma~\ref{four:quarter} provides such an estimate for every nonzero
quadratic matrix.  In particular it already gives a
uniform logarithmic potential.  The argument above upgrades the
concentration exponent when the principal pivots are small.

\subsection{A discriminant estimate}

Fix a coordinate $h$, and write $X'=X_{-h}$ and
\[
 a=m_{hh},\qquad w=M_{-h,h},\qquad
 L=w^{\mathsf T}X'+b_h/2,\qquad
 Q=aX_h^2+2LX_h+R'(X').
\]
For a target level $s$, set
\begin{equation}\label{lower:delta}
 \Delta_s=L^2+a(s-R').
\end{equation}
Its quadratic matrix is
\begin{equation}\label{lower:minor-matrix}
 N^{(h)}=ww^{\mathsf T}-aM_{-h,-h},
 \qquad
 N^{(h)}_{ij}=m_{hi}m_{hj}-m_{hh}m_{ij}.
\end{equation}
The diagonal entries are negatives of the principal two-coordinate
minors through $h$.

\begin{lemma}[Conditional sublevel bound]\label{lower:conditional}
For every $u>0$,
\begin{equation}\label{lower:conditional-bound}
 \mathbb P(|Q-s|\leq u\mid X')
 \leq\min\left\{1,\frac{2\sqrt2\sqrt u}{\sqrt{|a|}},
                     \frac{4u}{\sqrt{|\Delta_s|}}\right\}.
\end{equation}
When $a=0$, omit the middle expression.
\end{lemma}
\begin{proof}
If $a\ne0$, complete the square:
\[
 Q=a(X_h+L/a)^2+\xi,
 \qquad \xi=R'-L^2/a,
 \qquad |a|\,|s-\xi|=|\Delta_s|.
\]
Apply the one-dimensional quadratic sublevel estimate in
Lemma~\ref{geom:square}.  If $a=0$, the conditional polynomial is
affine with slope $2L$ and has sublevel probability at most
$\min(1,u/|L|)$.  Since $\Delta_s=L^2$, this implies the stated
weaker bound, including the case $L=0$.
\end{proof}

\begin{lemma}[Integrated discriminant bound]\label{lower:integrated}
With the middle term omitted for $a=0$, define
\[
 I(d,a)=\int_1^\infty
 \min\left\{1,\frac{2\sqrt2}{v\sqrt{|a|}},
                  \frac4{v^2\sqrt{|d|}}\right\}\,dv.
\]
Then
\begin{align}
 I(d,a)&\leq4|d|^{-1/4},\label{lower:quarter-discriminant}\\
 \mathbb E[|Q-s|^{-1/2};\ |Q-s|\leq1]
       &\leq1+\mathbb E I(\Delta_s,a).
       \label{lower:riesz-discriminant}
\end{align}
If $0<|a|\leq1$, one also has
\begin{equation}\label{lower:log-discriminant}
 I(d,a)\leq\frac{2\sqrt2}{\sqrt{|a|}}
       \left(2+\frac12\log_+\frac{a^2}{4|d|}\right).
\end{equation}
\end{lemma}
\begin{proof}
Put $D=4|d|^{-1/2}$.  Dropping the middle term, the integral is at
most $D$ if $D\leq1$, and at most $2\sqrt D$ if $D>1$, by splitting
at $\sqrt D$.  Both bounds imply \eqref{lower:quarter-discriminant}.
For \eqref{lower:log-discriminant}, put
$A_0=2\sqrt2|a|^{-1/2}\geq1$ and $V_0=D/A_0$.
If $V_0\leq A_0$, integrate $1$ to $A_0$ and $D/v^2$ afterwards;
the result is at most $2A_0$.  If $V_0>A_0$, integrate $1$ to
$A_0$, $A_0/v$ to $V_0$, and $D/v^2$ afterwards.  The result is at
most $A_0(2+\log(V_0/A_0))$.
Since $V_0/A_0=|a|/(2\sqrt{|d|})$, this proves the claim.
Finally, layer cake and Lemma~\ref{lower:conditional} give
\[
 \mathbb E[|Q-s|^{-1/2};\ |Q-s|\leq1]
 \leq1+\int_1^\infty\mathbb P(|Q-s|<v^{-2})\,dv
 \leq1+\mathbb E I(\Delta_s,a).
\]
\end{proof}

\begin{theorem}[A coordinate discriminant]\label{lower:discriminant}
Let $0<\nu\leq1$.  Suppose $M$ has no $c$-pivot, where
$c\leq\nu/9$, and some coordinate $h$ satisfies
$\|N^{(h)}_{\rm off}\|_F\geq\nu$.  Then
\begin{equation}\label{lower:discriminant-result}
 \mathcal H(Q)\leq5+40\left(\frac{16.92}{\nu}\right)^{1/3}.
\end{equation}
No spectral hypothesis is needed for this assertion.
\end{theorem}
\begin{proof}
By \eqref{lower:entry-cap} and $c\leq1/9$,
\[
 |N^{(h)}_{ii}|<c,
 \qquad |N^{(h)}_{ij}|
 \leq2c+c\sqrt{2c}\leq3c.
\]
A random partition of the coordinates other than $h$ yields a
realization satisfying
\[
 \|N^{(h)}_{I,J}\|_F^2\geq\nu^2/4,
 \qquad |N^{(h)}_{ij}|^2\leq9c^2\leq\nu^2/9
                          \leq\tfrac12\|N^{(h)}_{I,J}\|_F^2.
\]
Apply \eqref{lower:third-power} to $\Delta_s$, whose affine terms
are arbitrary.  Uniformly in $s$ and in a second target $z$,
\[
 \mathbb P(|\Delta_s-z|\leq r)\leq K_\Delta r^{1/3},
 \qquad
 K_\Delta=K(\nu/2)
 <3.264(16.92/\nu)^{1/3}.
\]
In particular $\Delta_s$ has no atom at zero, and layer cake gives
\[
 \mathbb E|\Delta_s|^{-1/4}
 \leq1+\int_1^\infty K_\Delta v^{-4/3}\,dv
 =1+3K_\Delta.
\]
Equations~\eqref{lower:quarter-discriminant} and
\eqref{lower:riesz-discriminant} now give
$\mathcal H(Q)\leq5+12K_\Delta$, which implies
\eqref{lower:discriminant-result}.
\end{proof}

\begin{lemma}[Two tests for a discriminant]\label{lower:tests}
If $M$ has no $c$-pivot, then every unit eigenpair $(\mu,v)$ and
every coordinate $h$ satisfy
\begin{align}
 \|N^{(h)}_{\rm off}\|_F
 &\geq |\mu|\bigl(|\mu|v(h)^2-|m_{hh}|\bigr)-c,
       \label{lower:eigen-test}\\
 \|N^{(h)}_{\rm off}\|_F
 &\geq r_h^2-|m_{hh}|\,\|M\|-c.
       \label{lower:row-test}
\end{align}
\end{lemma}
\begin{proof}
Extend $N^{(h)}$ by a zero $h$th row and column.  Then
$N^{(h)}=m_hm_h^{\mathsf T}-m_{hh}M$, where $m_h=Me_h$.
Its diagonal norm is at most $c$, so
\[
 \|N^{(h)}_{\rm off}\|_F\geq\|N^{(h)}\|-c.
\]
The Rayleigh quotient on $v$ is
$\mu^2v(h)^2-m_{hh}\mu$, proving \eqref{lower:eigen-test}.
Also $\|N^{(h)}\|\geq|m_h|^2-|m_{hh}|\|M\|$, and
$|m_h|^2\geq r_h^2$, proving \eqref{lower:row-test}.
\end{proof}

\subsection{Two-volume marginals and random partitions}

\begin{lemma}[Second-volume marginals]\label{lower:two-marginals}
Let $A$ have rank at least two, with singular values
$s_1\geq s_2\geq\cdots$ and left singular vectors $u_m$.
Put $G=AA^{\mathsf T}$, $\lambda_m=s_m^2$ and $Z_2=e_2(G)$.
Then
\begin{equation}\label{lower:marginal-formula}
 \alpha_i^{(2)}
 =\frac{G_{ii}\operatorname{tr}G-(G^2)_{ii}}{e_2(G)}
 =\sum_m p_m u_m(i)^2,
 \quad
 p_m=\frac{\lambda_m\sum_{\ell\ne m}\lambda_\ell}{e_2(\lambda)}
 \leq1.
\end{equation}
Consequently,
\begin{align}
 \alpha_i^{(2)}&\leq2|A_{i,\cdot}|^2/s_2^2,
                          \label{lower:row-marginal}\\
 \alpha_i^{(2)}&\leq u_1(i)^2+
 |(A-s_1u_1v_1^{\mathsf T})_{i,\cdot}|^2
 \left(\frac1T+\frac1S\right),
                          \label{lower:residual-marginal}\\
 \qquad S&=s_1^2,\qquad T=\sum_{m\geq2}s_m^2.
 \nonumber
\end{align}
There are identical formulas for the column marginals.
\end{lemma}
\begin{proof}
Cauchy--Binet gives
$\sum_T\det(A_{S,T})^2=\det G_{S,S}$.
Summing the two-coordinate minors that contain $i$ gives
\[
 \sum_{j\ne i}(G_{ii}G_{jj}-G_{ij}^2)
 =G_{ii}\operatorname{tr}G-(G^2)_{ii}.
\]
Spectral decomposition gives \eqref{lower:marginal-formula}.
The bound $p_m\leq1$ follows because its numerator is the sum of
the terms of $e_2(\lambda)$ that contain $\lambda_m$.
Moreover,
\[
 e_2(\lambda)=\tfrac12\sum_m\lambda_m
                         \sum_{\ell\ne m}\lambda_\ell
 \geq\tfrac12\lambda_2\sum_m\lambda_m,
\]
which proves \eqref{lower:row-marginal}.  Since $e_2\geq ST$,
for $m\geq2$ one has $p_m\leq\lambda_m(S+T)/(ST)$.
Separating the term $m=1$ proves \eqref{lower:residual-marginal}.
\end{proof}

Let $(\mu_1,v_1),(\mu_2,v_2)$ be orthonormal eigenpairs of $M$,
write $V=[v_1\ v_2]$, and put $\Lambda=\operatorname{diag}(\mu_1,\mu_2)$.
For independent fair signs $\epsilon_i$, put
\[
 E=\operatorname{diag}(\epsilon_i),\quad
 I=\{i:\epsilon_i=1\},\quad J=I^c,\quad
 \Sigma=V^{\mathsf T}EV,\quad B=V_I^{\mathsf T}M_{I,J}V_J.
\]
Writing $D=\operatorname{diag}M$, direct expansion gives
\begin{equation}\label{lower:partition-expansion}
 4B=\Lambda+[\Sigma,\Lambda]-V^{\mathsf T}DV-C,
 \qquad C_{kl}=\sum_{i\ne j}\epsilon_i\epsilon_j
                         v_k(i)m_{ij}v_l(j).
\end{equation}
Indeed $4B=V^{\mathsf T}(I+E)M(I-E)V$, and $MV=V\Lambda$.
If $M_3=M-V\Lambda V^{\mathsf T}$, the alternative expansion is
\begin{equation}\label{lower:remainder-expansion}
 4B=\Lambda+[\Sigma,\Lambda]-\Sigma\Lambda\Sigma
       -V^{\mathsf T}\operatorname{diag}(M_3)V-C^{(3)},
\end{equation}
where $C^{(3)}$ is defined as $C$ with $M_3$ in place of $M$.
This follows from $M_3V=0$.

The elementary second-moment identities used below are
\begin{align}
 \mathbb E\Sigma_{kl}^2
 &=\sum_i v_k(i)^2v_l(i)^2,\label{lower:sigma-variance}\\
 \mathbb E(C^{(H)}_{kl})^2
 &\leq2\sum_{i\ne j}v_k(i)^2H_{ij}^2v_l(j)^2
 \leq2\max_i v_k(i)^2\|H\|^2.
 \label{lower:chaos-variance}
\end{align}
For the second identity only equal or reversed index pairs survive
the sign average; bound each cross product by half the sum of its
two squares.  Finally,
\[
 [\Sigma,\Lambda]=\begin{pmatrix}0&s\\-s&0\end{pmatrix},
 \qquad s=(\mu_2-\mu_1)\Sigma_{12},
\]
Thus the matrix
$P=\left(\begin{smallmatrix}\mu_1&s\\-s&\mu_2\end{smallmatrix}\right)$
has determinant $\mu_1\mu_2+s^2$.  Since $\|V_I\|,\|V_J\|\leq1$,
the singular values satisfy $s_k(M_{I,J})\geq s_k(B)$ for $k=1,2$.

\begin{lemma}[Comparable eigenvalues and light rows]\label{lower:balanced}
Suppose
\[
 |m_{ii}|\leq c,\quad
 1\leq|\mu_2|\leq|\mu_1|\leq K,\quad
 r_i\leq\rho,\quad
 c\leq\frac1{8(K+2)},\quad
 \rho\leq\frac3{64(K+2)}.
\]
With probability at least $11/16$, a fair partition has
\[
 s_2(M_{I,J})\geq\sigma_A:=\frac3{32(K+2)},
 \qquad \max_i\alpha_i^{(2)}\leq1/2.
\]
Consequently,
\begin{equation}\label{lower:balanced-decay}
 |\varphi_Q(t)|\leq2\pi\sqrt e\,\sigma_A^{-1}|t|^{-1}.
\end{equation}
\end{lemma}
\begin{proof}
The eigenvector equation gives
\[
 |v_k(i)|\leq\frac{r_i}{|\mu_k-m_{ii}|}
            \leq\frac\rho{1-c}\qquad(k=1,2).
\]
Equation~\eqref{lower:chaos-variance}, using the off-diagonal row
bound before the last inequality, gives
\[
 \mathbb E\|C\|_F^2\leq\frac{8\rho^4}{(1-c)^2}\leq16\rho^4.
\]
Also $\mathbb E\Sigma_{12}^2\leq\rho^2/(1-c)^2$.
Except on a set of probability at most $1/16+1/4=5/16$, both
\[
 |\Sigma_{12}|\leq4\rho/(1-c),\qquad \|C\|_F\leq8\rho^2
\]
hold.  On this event,
\[
 |s|\leq8K\rho/(1-c)\leq1/2,
 \quad |\det P|\geq3/4,\quad \|P\|\leq K+1/2.
\]
Thus $s_2(P)\geq3/[4(K+2)]$.  Using
\eqref{lower:partition-expansion},
\[
 s_2(4B)\geq\frac3{4(K+2)}-c-8\rho^2
             \geq\frac3{8(K+2)}.
\]
Division by four gives the claimed singular-value lower bound.
Every row of $M_{I,J}$ has norm at most $\rho$, and hence
\eqref{lower:row-marginal} gives
$\alpha_i^{(2)}\leq2\rho^2/\sigma_A^2\leq1/2$.
Finally $Z_2\geq s_2^4\geq\sigma_A^4$, so
\eqref{lower:volume-q} with $k=2,q=1$ proves
\eqref{lower:balanced-decay}.
\end{proof}

\begin{lemma}[A dominant eigenvalue]\label{lower:dominant}
Order the eigenvalues of $M$ by decreasing modulus, and suppose
\[
 |\mu_2|=1,\quad K=|\mu_1|\geq10^4,\quad
 |m_{ii}|\leq c,\quad |m_{ij}|\leq\sqrt{2c},\quad c\leq10^{-6}.
\]
Assume, for all $i$, that
$K v_1(i)^2\leq\eta$ and $v_2(i)^2\leq\eta$, where
$\eta\leq6\cdot10^{-5}$.  Some partition and orientation satisfy
\[
 Z_2(M_{I,J})\geq1,
 \qquad \max_i\alpha_i^{(2)}\leq0.6.
\]
In particular,
\[
 |\varphi_Q(t)|\leq6.51|t|^{-5/6},\qquad
 L_Q(u)\leq85u^{5/6},\qquad \mathcal H(Q)<129.
\]
\end{lemma}
\begin{proof}
We first find a partition with two singular directions.
Since $\|M_3\|\leq1$, equations
\eqref{lower:sigma-variance}--\eqref{lower:chaos-variance} show
that, with probability at least $53/64$, all four bounds
\begin{equation}\label{lower:dominant-event}
 \begin{gathered}
 |\Sigma_{11}|,|\Sigma_{12}|\leq8\sqrt{\eta/K},
 \qquad |\Sigma_{22}|\leq8\sqrt\eta,\\
 \|C^{(3)}\|_F\leq8\sqrt\eta
 \end{gathered}
\end{equation}
hold.  The failure probabilities for the three entries of $\Sigma$
are at most $1/64$ each.  The bound
$\mathbb E\|C^{(3)}\|_F^2\leq4\eta/K+4\eta\leq8\eta$
gives an additional failure probability at most $8/64$.

On this event,
\[
 \|\Sigma\Lambda\Sigma\|_F
 \leq K(\Sigma_{11}^2+\Sigma_{12}^2)
               +\Sigma_{12}^2+\Sigma_{22}^2\leq256\eta.
\]
Moreover,
\[
 \|V^{\mathsf T}\operatorname{diag}(M_3)V\|
 \leq c+2\eta,
 \qquad
 s^2\leq64\eta K(1+1/K)^2\leq65\eta K.
\]
It follows that $|\det P|\geq0.996K$, $\|P\|\leq1.001K$,
and $s_2(P)\geq0.995$.  The remaining terms in
\eqref{lower:remainder-expansion} have norm at most
\[
 256\eta+c+2\eta+8\sqrt\eta<0.078.
\]
Hence, with $A=M_{I,J}$,
\begin{equation}\label{lower:dominant-singular}
 s_2(A)\geq0.229,
 \qquad s_1(A)\geq0.249K.
\end{equation}

Next write
\[
 A=\mu_1xy^{\mathsf T}+F,
 \qquad x=(v_1)_I,\quad y=(v_1)_J,\quad \|F\|\leq1.
\]
Equation~\eqref{lower:dominant-event} implies
$|x|^2,|y|^2\geq0.4996$.  Put
$\sigma_0=K|x||y|$, so that
$0.4996K\leq\sigma_0\leq K/2$.
Absorbing the sign of $\mu_1$ into one normalized vector, write
$\widehat x=x/|x|$, $\widehat y=\pm y/|y|$.
Then
\[
 AA^{\mathsf T}=\sigma_0^2\widehat x\widehat x^{\mathsf T}+P',
 \qquad \|P'\|\leq2\sigma_0+1.
\]
If $u$ is a top left singular vector of $A$, projecting its
eigenvector equation perpendicular to $\widehat x$ gives
\[
 \sin\angle(u,\widehat x)
 \leq\frac{2\sigma_0+1}{(\sigma_0-1)^2}
 \leq\frac{2.01}{\sigma_0}.
\]
Here $s_1(A)\geq\sigma_0-1$, and $\sigma_0\geq4996$.
The identical bound holds for the top right singular vector $v$
relative to $\widehat y$.

Let $A_0=A-s_1uv^{\mathsf T}$.  For each applicable coordinate,
\[
 |u_i|,|v_j|\leq R_K:=
      \sqrt{\frac\eta{0.4996K}}+\frac{2.01}{0.4996K},
 \qquad R_K^2<3\cdot10^{-7}.
\]
Using $A_0=A(I-vv^{\mathsf T})$ and its transposed version gives
\[
 |(A_0)_{i,\cdot}|,|(A_0)_{\cdot,j}|
 \leq1+2.01\sqrt{\frac\eta{0.4996K}}<1.001.
\]
For the entries, $s_1\leq\sigma_0+1\leq K/2+1$ yields
\[
 |(A_0)_{ij}|\leq\sqrt{2c}+(K/2+1)R_K^2<0.003.
\]
The last bound follows by substituting $c\leq10^{-6}$,
$\eta\leq6\cdot10^{-5}$ and $K\geq10^4$; after expansion,
$(K/2+1)R_K^2$ is decreasing in $K$.
By \eqref{lower:dominant-singular},
$T=\|A_0\|_F^2\geq0.229^2$.
Since $0.003^2<0.02\cdot0.229^2$, Lemma~\ref{lower:orientation}
with $h=0.51$ gives an orientation in which every residual row
has squared norm at most $0.51T$.

Equation~\eqref{lower:residual-marginal} therefore gives
\[
 \alpha_i^{(2)}
 \leq3\cdot10^{-7}+0.51+
            \frac{1.001^2}{(0.249K)^2}<0.6.
\]
Also $Z_2\geq s_1^2s_2^2\geq(0.249K)^2\,0.229^2>1$.
Apply \eqref{lower:volume-q} with $k=2,q=6/5$:
\[
 |\varphi_Q(t)|\leq
 \left[\frac{5e(2\pi)^2}{6}\right]^{5/12}|t|^{-5/6}
 <6.51|t|^{-5/6}.
\]
Equation~\eqref{lower:power-concentration} gives $L_Q(u)\leq85u^{5/6}$.
Finally,
\[
 \mathcal H(Q)\leq1+\int_1^\infty85v^{-5/3}\,dv
 =128.5<129.
\]
\end{proof}

\subsection{Completion of the second-level estimate}

\begin{proof}[Proof of Theorem~\ref{lower:two}]
Order the eigenvalues by modulus and let $r=|\mu_2|\geq1$.
Scaling by $r$ preserves the absence of a $c_*$-pivot.  More
precisely,
\[
 \mathcal H(Q)\leq r^{-1/2}\mathcal H(Q/r),
\]
because the event $|Q-s|\leq1$ becomes
$|Q/r-s/r|\leq1/r\leq1$.
It is therefore enough to assume $|\mu_2|=1$ and put
$K=|\mu_1|=\|M\|$.

Suppose first that $K\leq10^4$, and put
\[
 \rho_A=\frac3{64\cdot10002},
 \qquad \nu_A=\rho_A^2/2.
\]
If some $r_h>\rho_A$, equation~\eqref{lower:row-test} gives
\[
 \|N^{(h)}_{\rm off}\|_F
 \geq\rho_A^2-10^{-15}(10^4+1)\geq\nu_A.
\]
Here $\nu_A>1.098\cdot10^{-11}$ and $10^{-15}<\nu_A/9$.
Theorem~\ref{lower:discriminant} consequently gives
\[
 \mathcal H(Q)\leq5+40(16.92/\nu_A)^{1/3}<4.7\cdot10^5.
\]

If all $r_i\leq\rho_A$, apply Lemma~\ref{lower:balanced} with
$K=10^4$.  It gives $|\varphi_Q(t)|\leq\min(1,B_0/|t|)$, where
\[
 B_0=2\pi\sqrt e\,\frac{32\cdot10002}{3}<1.106\cdot10^6.
\]
For $u\leq1/B_0$, equation~\eqref{lower:fejer} implies
\[
 L_Q(u)\leq\frac{2B_0u}{\kappa_0}
                   \left(1+\log\frac1{B_0u}\right).
\]
For larger $u$ use $L_Q(u)\leq1$.  Splitting the half-power
integral at $v=\sqrt{B_0}$ yields
\begin{align*}
 \mathcal H(Q)
 &\leq\sqrt{B_0}+
 \frac{2B_0}{\kappa_0}\int_{\sqrt{B_0}}^\infty
 v^{-2}\left(1+2\log\frac v{\sqrt{B_0}}\right)dv\\
 &=\left(1+\frac6{\kappa_0}\right)\sqrt{B_0}<8000.
\end{align*}

Now suppose $K>10^4$ and put $\nu_B=5\cdot10^{-5}$.
If some $h$ and $k\in\{1,2\}$ satisfy
\[
 |\mu_k|\bigl(|\mu_k|v_k(h)^2-c_*\bigr)
                      \geq\nu_B+2c_*,
\]
equation~\eqref{lower:eigen-test} gives
$\|N^{(h)}_{\rm off}\|_F\geq\nu_B$.
Theorem~\ref{lower:discriminant} gives $\mathcal H(Q)<3000$.
Otherwise, for every $i$,
\[
 K v_1(i)^2<c_*+\frac{\nu_B+2c_*}{K}\leq6\cdot10^{-5},
 \qquad v_2(i)^2<\nu_B+3c_*\leq6\cdot10^{-5}.
\]
The entry bounds \eqref{lower:entry-cap} permit an application of
Lemma~\ref{lower:dominant}, which gives $\mathcal H(Q)<129$.
Every case is bounded by $10^6$, as asserted.
\end{proof}

Theorems~\ref{lower:one} and \ref{lower:two} provide the first two
integrability estimates used in the coordinate-pivot reduction.
Lemma~\ref{lem:schur-count} preserves the required number of large
eigenvalues at each pivot.  The reduction still requires a third-level
estimate for matrices without a fixed small pivot.  The estimates above
allow all affine terms produced by conditioning.

\section{Spectral reductions and finite coordinate cores}
\label{sec:structure}

Throughout this section \(X=(X_1,\ldots,X_n)\) has independent
coordinates with probability densities bounded by one.  Write
\[
 Q_A(X)=X^{\mathsf T}AX+b^{\mathsf T}X+c_0,
 \qquad
 N_A(t)=\#\{j:|\lambda_j(A)|\geq t\}.
\]
All assertions are uniform in the affine terms.  For \(0<c\leq1\),
a \emph{\(c\)-pivot} is either a diagonal entry of modulus at least
\(c\), or a principal \(2\)-by-\(2\) submatrix whose determinant has
modulus at least \(c\).  If there is no \(c\)-pivot, then
\begin{equation}
 |a_{ii}|<c,\qquad |a_{ij}|^2<c+c^2\leq2c .
 \label{eq:no-pivot-entries}
\end{equation}

\subsection{Schur complements and logarithmic potentials}

\begin{lemma}[Spectral count under a coordinate pivot]
\label{lem:schur-count}
Let \(B\) be an invertible principal submatrix of a real symmetric
matrix \(A\), of order \(k\), and let \(S=A/B\) be its Schur
complement.  For every \(t>0\),
\[
 N_S(t)\geq N_A(t)-k.
\]
\end{lemma}
\begin{proof}
First suppose \(A\) is invertible.  Then \(S\) is invertible and
\(S^{-1}\) is the complementary principal submatrix of \(A^{-1}\).
Interlacing bounds the number of eigenvalues of \(S^{-1}\) in
\((t^{-1},\infty)\) by the corresponding number for \(A^{-1}\).
It gives the analogous bound in \((-\infty,-t^{-1})\).
Consequently, at most \(n-N_A(t)\) of the \(n-k\) eigenvalues of
\(S^{-1}\) have modulus greater than \(t^{-1}\).  This is the
asserted inequality.  For singular \(A\), apply the argument to
\(A+\varepsilon I\) along values for which this matrix is invertible,
at threshold \(t-|\varepsilon|\).  The pivot remains invertible
for sufficiently small \(\varepsilon\); its Schur complement
converges to \(S\).  Continuity of the ordered eigenvalues gives
the result.
\end{proof}


Define
\[
 {\cal L}(Y)=\sup_s\mathbb E\log_+\frac1{|Y-s|},
 \qquad
 {\cal H}(Y)=\sup_s\mathbb E\left[
 |Y-s|^{-1/2}\boldsymbol 1_{\{|Y-s|\leq1\}}\right].
\]

\begin{lemma}[A global logarithmic bound]
\label{lem:global-log}
If \(A\ne0\), then
\[
 {\cal L}(Q_A)\leq
 4+\log_+\frac{5^4}{\|A\|}.
\]
In particular, \(N_A(1)\geq1\) implies
\({\cal L}(Q_A)\leq L_0:=4+4\log5\).
\end{lemma}
\begin{proof}
The quarter-power estimate of Lemma~\ref{four:quarter} and the
layer-cake formula give
\[
 \mathbb E\log_+\frac1{|Q_A-s|}
 \leq\int_0^\infty
 \min\{1,5\|A\|^{-1/4}e^{-u/4}\}\,du.
\]
The integral is \(4+4\log K\) when \(K=5\|A\|^{-1/4}\geq1\),
and \(4K\leq4\) otherwise.
\end{proof}

\begin{lemma}[Half-power integrability at a pivot]
\label{lem:pivot-half}
If \(N_A(1)\geq2\) and \(A\) has a \(c\)-pivot, then
\[
 {\cal H}(Q_A)\leq44c^{-1/2}.
\]
\end{lemma}
\begin{proof}
Suppose first that \(a=a_{hh}\) satisfies \(|a|\geq c\).
Completing the square in the original coordinate \(X_h\) writes
\(Q_A=a(X_h-m)^2+K\), where \(m\) is affine in the remaining
coordinates and the quadratic matrix of \(K\) is \(A/a\).
Lemma~\ref{lem:schur-count} gives \(N_{A/a}(1)\geq1\).
The conditional density estimate of Lemma~\ref{geom:square}
therefore yields
\[
 {\cal H}(Q_A)
 \leq |a|^{-1/2}\sup_s\mathbb E J(s-K),
 \qquad
 J(d)=\int_{-1}^1 |x|^{-1/2}|x-d|^{-1/2}\,dx .
\]
For \(d\ne0\),
\begin{equation}
 J(d)\leq8+3\log_+\frac1{|d|}.
 \label{eq:root-convolution}
\end{equation}
Indeed, on \(|x|\leq2|d|\) substitute \(x=|d|u\).
The resulting integral is at most
\[
 \int_{-2}^2 |u|^{-1/2}|u-1|^{-1/2}\,du
 =2\operatorname{arsinh}\sqrt2+\pi
   +2\operatorname{arcosh}\sqrt2<8.
\]
On \(2|d|<|x|\leq1\), the integrand is at most
\(\sqrt2/|x|\), giving the remaining term in
\eqref{eq:root-convolution}.  Lemma~\ref{lem:global-log}
bounds this case by \(c^{-1/2}(8+3L_0)<40c^{-1/2}\).

Next suppose that a principal two-coordinate block \(B\) is
definite with determinant \(D\geq c\).  Its conditional density
is at most \(\pi D^{-1/2}\); integrating
\(|x-s|^{-1/2}\boldsymbol1_{\{|x-s|\leq1\}}\) gives
\({\cal H}(Q_A)\leq4\pi c^{-1/2}\).

Finally suppose \(\det B=-D\), \(D\geq c\).
Lemma~\ref{geom:pair} bounds the conditional density by
\[
 D^{-1/2}\left(5.2+
       \log_+\frac{16\sqrt D}{|y-K|}\right).
\]
For every real \(a\),
\[
 \int_{-1}^1 |x|^{-1/2}\log_+\frac1{|x-a|}\,dx\leq8.
\]
To see this without a regularity assumption, apply layer cake
to both factors.  Their superlevel sets are intervals, and the
overlap of two intervals of prescribed lengths is maximal when
their centres coincide.  The integral is therefore at most its
value at \(a=0\), namely
\(2\int_0^1x^{-1/2}\log(1/x)\,dx=8\).
It follows that
\[
 {\cal H}(Q_A)\leq
 D^{-1/2}\left(28.8+4\log_+(16\sqrt D)\right)
 \leq c^{-1/2}(28.8+4\log16)+4/e
 <44c^{-1/2}.
\]
All conditionings are on original independent coordinates.
\end{proof}

\subsection{Gaussian padding and reduction to small pivots}

\begin{theorem}[Reduction to the three-eigenvalue no-pivot case]
\label{thm:no-pivot-equivalence}
Fix \(0<c\leq1\).  Suppose that
\begin{equation}
 \|p_{Q_A}\|_\infty\leq C_3
 \quad\text{whenever }N_A(1)\geq3
 \text{ and \(A\) has no \(c\)-pivot}.
 \label{eq:no-pivot-hypothesis}
\end{equation}
Then the bounded-density conjecture holds for all real symmetric
matrices with \(N_A(1)\geq3\).  One admissible bound is
\[
 \max\left\{C_3,\ c^{-1/2}(1+G_2),\ 20c^{-1/2}\right\},
 \qquad
 G_2=\max\{2\sqrt{2\pi e}\,C_3,\ 44c^{-1/2}\}.
\]
The converse implication is immediate.
\end{theorem}
\begin{proof}
First consider a matrix \(M\) with \(N_M(1)\geq2\) and no
\(c\)-pivot.  Choose \(N>1/c\), put
\(u=N^{-1/2}(1,\ldots,1)\), and let \(G_1,\ldots,G_N\) be
independent standard Gaussian variables, independent of \(X\).
The variable \(Z=\langle u,G\rangle\) is standard Gaussian.
The matrices \(M\oplus uu^{\mathsf T}\) and
\(M\oplus(-uu^{\mathsf T})\) have at least three eigenvalues
of modulus at least one and no \(c\)-pivot: all new diagonal
entries have modulus \(1/N<c\), all two-coordinate minors
inside the new block vanish, and a mixed minor is the product
of two diagonal entries of modulus less than \(c\).
The original Gaussian coordinate densities are bounded by one.
Thus \eqref{eq:no-pivot-hypothesis} applies to \(Q_M\pm Z^2\).

The density of \(Z^2\) is
\[
 k(t)=\frac{e^{-t/2}}{\sqrt{2\pi t}}\boldsymbol1_{\{t>0\}}.
\]
For \(0<|t|\leq1\),
\[
 k(t)+k(-t)\geq(2\pi e)^{-1/2}|t|^{-1/2}.
\]
Adding the two convolution bounds gives
\({\cal H}(Q_M)\leq2\sqrt{2\pi e}\,C_3\).
The estimate holds for every centre \(s\), rather than only
almost every centre: the symmetrized kernel
\((2\pi|t|)^{-1/2}e^{-|t|/2}\), with value \(+\infty\) at zero,
is lower semicontinuous.  Its convolution has the same property,
so an almost-everywhere upper bound extends everywhere.
Lemma~\ref{four:quarter} also ensures that \(Q_M\) has no atoms.
Together with Lemma~\ref{lem:pivot-half}, this proves
\begin{equation}
 N_M(1)\geq2\quad\Longrightarrow\quad {\cal H}(Q_M)\leq G_2
 \label{eq:global-half-conditional}
\end{equation}
for every matrix \(M\).

Now let \(N_A(1)\geq3\).  The case with no pivot is the
hypothesis.  At a one-coordinate pivot \(a\), the square
kernel and Lemma~\ref{lem:schur-count} give
\[
 p_{Q_A}(s)\leq |a|^{-1/2}\mathbb E|s-K|^{-1/2}
 \leq c^{-1/2}(1+G_2),
\]
using \eqref{eq:global-half-conditional} for the Schur
complement.  A definite two-coordinate pivot gives
\(\pi c^{-1/2}\).  At an indefinite pivot of determinant
\(-D\), the centre has at least one large eigenvalue, so
Lemmas~\ref{geom:pair} and \ref{lem:global-log} give
\[
 p_{Q_A}(s)\leq
 D^{-1/2}\bigl(5.2+\log_+(16\sqrt D)+L_0\bigr)
 \leq20c^{-1/2}.
\]
This exhausts the cases.
\end{proof}

\begin{remark}
The reduction assumes \eqref{eq:no-pivot-hypothesis} and requires
neither regularity of the original coordinate densities nor moment
bounds.
\end{remark}

\subsection{Finite-dimensional bounds and projected gradients}

\begin{theorem}[Slicing by a projected gradient]
\label{thm:projected-gradient-slicing}
Let \(E\) be a three-dimensional invariant subspace of \(A\)
on which every eigenvalue has modulus at least \(\sigma>0\).
Put
\[
 \ell_i=|P_Ee_i|^2,\qquad {\cal R}(E)=\sum_i\sqrt{\ell_i}.
\]
Then
\begin{equation}
 \|p_{Q_A}\|_\infty
 \leq\frac{2\sqrt{2\pi}}{\sigma}{\cal R}(E)
 \leq\frac{2\sqrt{6\pi n}}{\sigma}.
 \label{eq:projected-slicing}
\end{equation}
For \(n=3\), the bound improves to \(6\sqrt\pi/\sigma\).
All eigenvalues outside \(E\) are unrestricted.
\end{theorem}
\begin{proof}
Write \(v(x)=\nabla Q_A(x)\) and \(g(x)=P_Ev(x)\).
There is a unique \(z_0\in E\) with
\(2(A|_E)z_0+P_Eb=0\), and hence
\[
 g(x)=2(A|_E)(P_Ex-z_0),\qquad
 |g(x)|\geq2\sigma|P_Ex-z_0|.
\]
The set \(g=0\) has codimension three and has probability zero.
The identities \(v^{\mathsf T}g=|g|^2\) and \(g\in E\) imply
\begin{equation}
 1\leq\sum_i\frac{|\partial_iQ_A|\,|g_i|}{|g|^2},
 \qquad |g_i|\leq\sqrt{\ell_i}|g|.
 \label{eq:projected-field-partition}
\end{equation}
The absolute values in this inequality are necessary; individual
terms of \(v^{\mathsf T}g\) need not be nonnegative.

Let \(\psi\) be a nonnegative bounded compactly supported
Borel function.  Multiply \eqref{eq:projected-field-partition}
by \(\psi(Q_A(x))\prod_jf_j(x_j)\), integrate, and apply
Tonelli.  For the \(i\)-th summand, fix all coordinates except
\(x_i=t\).  On the monotonicity intervals of the resulting
nonconstant quadratic or linear polynomial, one-dimensional
change of variables gives
\[
 \int \psi(Q_A)f_i(t)
       \frac{|\partial_iQ_A|\,|g_i|}{|g|^2}\,dt
 =
 \int\psi(s)
 \sum_{\substack{t:Q_A(t,x_{-i})=s\\ \partial_iQ_A\ne0}}
 f_i(t)\frac{|g_i(t,x_{-i})|}{|g(t,x_{-i})|^2}\,ds .
\]
Critical values form a null set.  If the fibre polynomial is
constant, the left side is zero and the fibre contributes zero.

Choose a two-dimensional subspace \(E_i\subset E\cap e_i^\perp\).
It exists because \(\dim E=3\).  The projection
\(P_{E_i}x\) is independent of the coordinate \(t=x_i\), and
\[
 |g(x)|\geq2\sigma|P_{E_i}x-P_{E_i}z_0|.
\]
There are at most two noncritical roots, and \(f_i\leq1\).
The root sum is consequently at most
\[
 \frac{\sqrt{\ell_i}}
 {\sigma|P_{E_i}x_{-i}-P_{E_i}z_0|}.
\]
The vector \(P_{E_i}X\) is an orthonormal two-dimensional
projection of the remaining original product coordinates.
Its density is at most \(2\), by
\citet[Theorem~1.1]{LPP2016}.
For a two-dimensional probability density bounded by \(M\),
layer cake gives
\[
 \mathbb E|Y-z|^{-1}
 \leq\int_0^\infty\min\{1,\pi M/u^2\}\,du
 =2\sqrt{\pi M}.
\]
The \(i\)-th summand is therefore bounded by
\((2\sqrt{2\pi}/\sigma)\sqrt{\ell_i}\int\psi\).
Summing proves domination of the output law by the first
constant in \eqref{eq:projected-slicing} times Lebesgue
measure.  This establishes absolute continuity as well as
the essential density bound, without a pointwise density
assumption on a level set.

Finally \(\sum_i\ell_i=3\), so
\({\cal R}(E)\leq\sqrt{3n}\).
When \(n=3\), \(E_i\) is the full original coordinate
complement, whose density is at most one; also
\({\cal R}(E)=3\).  This gives \(6\sqrt\pi/\sigma\).
\end{proof}

\begin{corollary}[The finite-dimensional spectral problem]
\label{thm:finite-dimension}
If \(A\) is \(n\)-by-\(n\) and \(N_A(\lambda)\geq3\), then
\[
 \|p_{Q_A}\|_\infty\leq2\sqrt{6\pi n}/\lambda.
\]
\end{corollary}
\begin{proof}
Choose \(E\) to be spanned by three corresponding eigenvectors
and apply Theorem~\ref{thm:projected-gradient-slicing}.
\end{proof}

\begin{corollary}[A finite coordinate core]
\label{cor:finite-core}
Suppose a principal submatrix \(A_{KK}\), with \(|K|\leq N\),
has three eigenvalues of modulus at least \(\lambda>0\).
Then
\[
 \|p_{Q_A}\|_\infty\leq2\sqrt{6\pi N}/\lambda,
\]
irrespective of the ambient dimension and all other blocks of \(A\).
\end{corollary}
\begin{proof}
Condition on \(X_{K^c}\).  The remaining variables retain their
independent densities, the quadratic matrix is \(A_{KK}\),
and the conditioned values affect only its affine terms.
Apply Corollary~\ref{thm:finite-dimension} and average.
\end{proof}

\begin{theorem}[A localized spectral core with a dense perturbation]
\label{thm:physical-core-stability}
Let \(H\) contain \(m\geq3\) original coordinates, and suppose
\[
 A_{HH}=U_H\Lambda U_H^{\mathsf T}+T_H,\qquad
 U^{\mathsf T}U=I_3,\qquad
 U_H^{\mathsf T}U_H\geq aI_3,
\]
where \(\Lambda\) is symmetric,
\(s_{\min}(\Lambda)\geq s>0\), and
\(\|T_H\|\leq\varepsilon<as\).
Then every choice of the other blocks of \(A\) satisfies
\[
 \|p_{Q_A}\|_\infty\leq
 \frac{2\sqrt{6\pi m}}{as-\varepsilon}.
\]
\end{theorem}
\begin{proof}
Put \(G=U_H^{\mathsf T}U_H\) and
\(V=U_HG^{-1/2}\).  The columns of \(V\) are orthonormal.
On their span, the rank-three matrix
\(U_H\Lambda U_H^{\mathsf T}\) is represented by
\(K=G^{1/2}\Lambda G^{1/2}\), for which
\[
 |Kz|\geq\sqrt a\,|\Lambda G^{1/2}z|
 \geq as|z|.
\]
For \(x\in\operatorname{range}(V)\), the triangle inequality
gives \(|A_{HH}x|\geq(as-\varepsilon)|x|\).
The variational characterization of singular values yields
\(s_3(A_{HH})\geq as-\varepsilon\).
Now apply Corollary~\ref{cor:finite-core}.
\end{proof}

The perturbation in Theorem~\ref{thm:physical-core-stability}
is retained in the actual physical matrix.  Its small operator
norm certifies the spectrum on the core; the proof makes no
assertion that \(\varepsilon|X|^2\) is small.

\begin{proposition}[Small balls of a projected gradient]
\label{prop:projected-gradient}
Let \(u_1,\ldots,u_r\) be orthonormal eigenvectors of \(A\),
with nonzero eigenvalues \(\lambda_1,\ldots,\lambda_r\),
and put \(U=(u_1,\ldots,u_r)\).
The vector \(Y=U^{\mathsf T}\nabla Q_A(X)\) has density at most
\[
 K=\frac{2^{-r/2}}{\prod_{j=1}^r|\lambda_j|}.
\]
Writing \(v_r\) for the volume of the unit ball in
\(\mathbb R^r\), for every \(z\in\mathbb R^r\), \(a>0\),
and \(0<p<r\),
\[
 \mathbb P\{|Y-z|\leq a\}\leq\min\{1,Kv_ra^r\},
 \qquad
 \mathbb E|Y-z|^{-p}
 \leq\frac r{r-p}(Kv_r)^{p/r}.
\]
\end{proposition}
\begin{proof}
Since
\[
 U^{\mathsf T}\nabla Q_A
 =2\operatorname{diag}(\lambda_1,\ldots,\lambda_r)U^{\mathsf T}X
   +U^{\mathsf T}b,
\]
the projection-density bound of Lemma~\ref{four:projection}
and change of variables give \(K\).
Integrating this density over a ball gives the first estimate.
The second follows by integrating the first estimate in the
layer-cake identity
\(\mathbb E|Y-z|^{-p}
=\int_0^\infty\mathbb P\{|Y-z|<t^{-1/p}\}\,dt\).
\end{proof}

The projected-gradient estimate controls neighbourhoods of
critical points.  It does not, by itself, bound the density of
\(Q_A\): the level-set integral also contains the product
density, which is allowed to be rough.

\subsection{An affine estimate for exact rank two}

\begin{lemma}[A square with an arbitrary linear term]
\label{lem:rank-one-affine}
For \(|u|=1\), \(\mu\ne0\), and arbitrary \(b,c_0\),
\[
 \sup_s\mathbb P\{|\mu\langle u,X\rangle^2+
 b^{\mathsf T}X+c_0-s|\leq a\}
 \leq8\sqrt{a/|\mu|}.
\]
\end{lemma}
\begin{proof}
This is Lemma~\ref{rt:rank-one}; its Gaussian smoothing
argument includes arbitrary original-coordinate affine terms.
The closed-interval version follows by decreasing the radius
to \(a\) from above.
\end{proof}

\begin{theorem}[Exact rank two]
\label{thm:rank-two-affine}
Suppose \(A=RMR^{\mathsf T}\), where \(R^{\mathsf T}R=I_2\),
\(M\) is real symmetric, and \(|\det M|\geq1\).
Then, for \(0<a\leq1\),
\[
 \sup_s\mathbb P\{|Q_A-s|\leq a\}
 \leq32\sqrt2\,a(1+\log(1/a))\leq46a(1+\log(1/a)).
\]
Moreover, Theorem~\ref{rt:half} gives the stronger fractional-potential bound
\(\sup_s\mathbb E|Q_A-s|^{-1/2}<68\).
\end{theorem}
\begin{proof}
Let \(\rho_i=R^{\mathsf T}e_i\).
If every \(|\rho_i|^2\leq1/2\), Fresnel linearization in the
two spectral directions and the geometric
Brascamp--Lieb estimate give
\(|\phi_{Q_A}(t)|\leq4\pi/(|t|\sqrt{|\det M|})\).
The Fej\'er kernel
\((\sin(x/2)/(x/2))^2\) is greater than \(9/10\) on \([-1,1]\).
Its triangular Fourier representation yields
\[
 \mathbb P\{|Y-s|\leq a\}
 \leq\frac{10}{9}a\int_{|t|\leq a^{-1}}|\phi_Y(t)|\,dt.
\]
Using \(|\phi|\leq1\) on \(|t|\leq1\) gives the asserted
bound with constant less than \(28\).

Otherwise choose \(h\) with \(l=|\rho_h|^2>1/2\).
Rotate the two-dimensional frame so that
\(\rho_h=\sqrt l\,e_2\), and write
\[
 M=\begin{pmatrix}s&q\\q&a_0\end{pmatrix}.
\]
The first column \(r_1\) of \(R\), restricted to \(h^c\), is
a unit vector; the second restricted column \(r_2\) is
orthogonal to it.  As a quadratic in \(x=X_h\), write
\(Q_A=ax^2+2Lx+T\), and define
\(\Delta=L^2+a(s_0-T)\), for a target level \(s_0\).
Its quadratic matrix is
\[
 cc^{\mathsf T}-aA_{h^c,h^c}
 =-l\det(M)\,r_1r_1^{\mathsf T},
 \qquad
 c=A_{h^c,h}=\sqrt l(qr_1+a_0r_2),\quad a=la_0.
\]
This identity follows by expanding both sides; the coefficients
of \(r_2r_2^{\mathsf T}\) and
\(r_1r_2^{\mathsf T}+r_2r_1^{\mathsf T}\) cancel.
The other terms of \(\Delta\) are affine.
Lemma~\ref{lem:rank-one-affine} thus gives
\(\mathbb P\{|\Delta-t|\leq u\}\leq8\sqrt2\sqrt u\).

The one-variable quadratic sublevel estimate, including the
linear case \(a=0\), bounds the conditional probability by
\(\min\{1,4a_{\!s}/\sqrt{|\Delta|}\}\), where \(a_{\!s}\)
is the small-ball radius.  Hence
\[
 \mathbb P\{|Q_A-s_0|\leq a_{\!s}\}
 \leq\int_0^1\min\{1,32\sqrt2\,a_{\!s}/\theta\}\,d\theta
 \leq32\sqrt2\,a_{\!s}(1+\log(1/a_{\!s})).
\]
The potential estimate is Theorem~\ref{rt:half}, since
\(|\det M|\geq1\).
\end{proof}

\section{Off-diagonal rank and quantitative skeletons}
\label{sec:skeletons}

For a real symmetric matrix \(A\), let
\[
 \rho(A)=\max_{S\cap T=\varnothing}\operatorname{rank}A_{ST},
 \qquad
 \operatorname{mr}(A)=\min_{D\ {\rm diagonal}}\operatorname{rank}(A-D).
\]
The sets \(S,T\) in this definition need not have the same size.

\begin{theorem}[Exact completion by a diagonal]
\label{thm:diagonal-completion}
For every real symmetric \(A\),
\[
 \operatorname{mr}(A)\leq3\rho(A).
\]
For every positive integer \(r\), examples satisfy
\(\rho(A)=r\) and \(\operatorname{mr}(A)=2r\).
If \(\rho(A)\leq1\), then \(\operatorname{mr}(A)\leq2\).
\end{theorem}
\begin{proof}
Put \(r=\rho(A)\).  The case \(r=0\) means that \(A\) is
diagonal.  Otherwise choose disjoint \(r\)-sets \(S,T\)
with \(A_{ST}\) invertible, and put \(C=A_{ST}^{-1}\),
\(K=S\cup T\), \(L=K^c\).
For \(l\in L\), write \(x_l=A_{lT}\), \(y_l=A_{lS}\),
and collect these rows in matrices \(X,Y\).
Every off-diagonal \((r+1)\)-minor vanishes; adjoining
distinct \(l,m\in L\) to the selected minor gives
\[
 a_{lm}=x_lCy_m^{\mathsf T}.
\]
Symmetry therefore implies that the span of the vectors
\((y_l,x_l)\in\mathbb R^{2r}\) is isotropic for the
nondegenerate alternating form
\[
 \omega((y,x),(y',x'))=xC(y')^{\mathsf T}-x'Cy^{\mathsf T}.
\]
An isotropic subspace has dimension at most \(r\):
it is contained in its orthogonal complement for this
nondegenerate form, and the two dimensions sum to \(2r\).
The matrix \(P=XCY^{\mathsf T}\) is symmetric because its
off-diagonal entries agree with those of \(A_{LL}\).
Choose \(D_K=0\) and
\(D_L=\operatorname{diag}(A_{LL}-P)\).
Each \(L\)-row of \(A-D\) is a fixed linear image of
\((y_l,x_l)\), so these rows span at most \(r\) dimensions.
The \(2r\) rows in \(K\) add at most \(2r\), proving
\(\operatorname{mr}(A)\leq3r\).

A star with centre \(h\), zero leaf--leaf entries and at
least two nonzero edges has \(\rho=1\) and
\(\operatorname{mr}=2\).  Its off-diagonal matrix already
has rank two.  A symmetric rank-one completion would have
the form \(\pm vv^{\mathsf T}\); two nonzero centre--leaf
entries would then force a nonzero entry between those
leaves.  Direct sums of \(r\) such stars have the stated
parameters, since both quantities add over direct sums.

It remains to prove the improved upper bound for \(\rho\leq1\).
Consider the support graph of the off-diagonal entries.
For four distinct indices the hypothesis reads
\(a_{ik}a_{jl}=a_{il}a_{jk}\).
Two disjoint edges consequently force all six edges on their
four endpoints.  A vertex adjacent to a clique of order at
least three is adjacent to every vertex of that clique:
pair its incident edge with edges between the other clique
vertices and use the same identity.
Thus, if disjoint edges occur, all nonzero edges lie in a
single complete graph with at least four vertices.
On this graph the numbers
\[
 c_i=\frac{a_{ij}a_{ik}}{a_{jk}}
\]
are independent of the distinct choices \(j,k\ne i\).
The four-index identities give \(c_ic_j=a_{ij}^2>0\);
all \(c_i\) have a common sign.  The edge signs obey the
same identities, so choices of signs for
\(v_i=\pm\sqrt{|c_i|}\) give \(a_{ij}=\pm v_iv_j\)
with a common overall sign.  This supplies a rank-one
diagonal completion.  If there are no disjoint edges,
the graph is a star, a triangle, or a subgraph of one.
A star has the rank-two completion already described.
For a triangle, subtracting any eigenvalue from the three
diagonal entries makes the corresponding block singular,
giving rank at most two.
\end{proof}

For \(r\geq1\), define
\[
 V_r(A)=
 \sum_{\substack{|S|=|T|=r\\S\cap T=\varnothing}}
 \det(A_{ST})^2,\qquad V_0(A)=1.
\]
The sum is over ordered pairs of subsets; each subset is
listed in its natural order.

\begin{lemma}[Volume-sampled skeleton inequality]
\label{lem:skeleton}
Suppose \(V_r(A)>0\).  Sample disjoint \(r\)-sets \(S,T\)
with probability \(\det(A_{ST})^2/V_r(A)\), omitting zero
weights, and put
\[
 P^{ST}=A_{\cdot T}A_{ST}^{-1}A_{S\cdot}.
\]
Then
\begin{equation}
 \mathbb E_{S,T}
 \sum_{\substack{l,m\notin S\cup T\\l\ne m}}
 (A-P^{ST})_{lm}^2
 \leq(r+1)^2\frac{V_{r+1}(A)}{V_r(A)}.
 \label{eq:skeleton-inequality}
\end{equation}
\end{lemma}
\begin{proof}
For every sampled nonsingular seed and every allowed \(l,m\),
the bordered determinant formula gives
\[
 \det(A_{S+l,T+m})^2
 =\det(A_{ST})^2(A-P^{ST})_{lm}^2.
\]
Sum over the seeds and \(l,m\).  A disjoint pair of
\((r+1)\)-sets is counted at most \((r+1)^2\) times,
according to which row and column are distinguished.
Some of these deletions may leave a singular seed and are
then absent from the sum.  Division by \(V_r(A)\) proves
\eqref{eq:skeleton-inequality}.
\end{proof}

\begin{remark}[Why equality is not automatic]
\label{rem:skeleton-zero-seeds}
Take
\[
 A=\begin{pmatrix}
 0&0&1&0\\0&0&0&1\\1&0&0&0\\0&1&0&0
 \end{pmatrix},\qquad r=1.
\]
Then \(V_1=V_2=4\).  Every nonzero seed has squared
off-diagonal residual sum \(2\), so the left side of
\eqref{eq:skeleton-inequality} is \(2\), whereas the
right side is \(4\).  Singular seeds may have nonsingular
extensions.  Thus the valid general statement is an
inequality, which suffices for the following consequence.
\end{remark}

\begin{theorem}[A robust low-rank completion]
\label{thm:robust-skeleton}
Suppose \(V_r(A)>0\) and
\(V_{r+1}(A)\leq\kappa V_r(A)\).
There is a decomposition
\[
 A=D+R+E,\qquad
 \operatorname{rank}R\leq4r,\qquad
 \|E\|_{\mathrm F}\leq3(r+1)\sqrt\kappa,
\]
where \(D\) is diagonal and \(E\) is symmetric, has zero
diagonal, and vanishes on \(2r\) selected rows and columns.
\end{theorem}
\begin{proof}
Choose a seed from Lemma~\ref{lem:skeleton} for which the
off-diagonal residual on \(L=(S\cup T)^c\) has squared
Frobenius norm at most \((r+1)^2\kappa\).
Write
\[
 P=A_{LT}A_{ST}^{-1}A_{SL},\qquad
 F=\operatorname{offdiag}(A_{LL}-P).
\]
Since \(A_{LL}\) is symmetric,
\(\|P-P^{\mathsf T}\|_{\mathrm F}\leq2\|F\|_{\mathrm F}\).
Thus \(P_{\rm sym}=(P+P^{\mathsf T})/2\) satisfies
\(\|P_{\rm sym}-P\|_{\mathrm F}\leq\|F\|_{\mathrm F}\).
Let \(P_r\) retain the \(r\) largest eigenvalues in modulus
of \(P_{\rm sym}\).  The best rank-\(r\) approximation
property and \(\operatorname{rank}P\leq r\) imply
\[
 \|P_{\rm sym}-P_r\|_{\mathrm F}\leq\|F\|_{\mathrm F}.
\]
All columns of \(P_r\) lie in
\(\operatorname{col}[A_{LS},A_{LT}]\), a space of
dimension at most \(2r\), because this contains the range
of \(P_{\rm sym}\).

With \(K=S\cup T\), set
\[
 R=\begin{pmatrix}A_{KK}&A_{KL}\\A_{LK}&P_r\end{pmatrix},
 \quad
 E_{LL}=\operatorname{offdiag}(A_{LL}-P_r),
 \quad E_{K,\cdot}=E_{\cdot,K}=0,
\]
and let \(D=A-R-E\).  The \(L\)-row block of \(R\) has
rank at most \(2r\), and its \(K\)-row block adds at most
\(2r\), so \(\operatorname{rank}R\leq4r\).
The matrix \(D\) is diagonal and
\[
 \|E\|_{\mathrm F}
 \leq\|F\|_{\mathrm F}+\|P-P_r\|_{\mathrm F}
 \leq3\|F\|_{\mathrm F}
 \leq3(r+1)\sqrt\kappa.
\]
\end{proof}

\begin{corollary}
\label{cor:tail-skeleton}
If every coordinate partition \(I,J\) satisfies
\(\sum_{j\geq r+1}s_j(A_{IJ})^2\leq\delta^2\), and
\(V_r(A)>0\), then Theorem~\ref{thm:robust-skeleton}
applies with \(\kappa=4\delta^2\).  In particular
\(\|E\|_{\mathrm F}\leq6(r+1)\delta\).
\end{corollary}
\begin{proof}
For nonincreasing nonnegative numbers \(x_j\),
\[
 e_{r+1}(x)\leq
 \left(\sum_{j\geq r+1}x_j\right)e_r(x).
\]
Assign each \((r+1)\)-set to its largest index to verify
this inequality term by term.
For a random fair coordinate partition, Cauchy--Binet gives
\(\mathbb Ee_r(A_{IJ}A_{IJ}^{\mathsf T})=4^{-r}V_r(A)\).
Average the preceding inequality with \(x_j=s_j(A_{IJ})^2\).
\end{proof}

The conclusion of Corollary~\ref{cor:tail-skeleton} uses a
bound on the entire squared singular-value tail.  A bound
only on \(s_{r+1}\) yields a dimension factor through that
tail and is not the same hypothesis.

\section{Grouped projections and projected crossing blocks}
\label{sec:grouped-crossing}

The following arguments retain arbitrary phases inside each
coordinate group.  The projected coordinates need not be
independent.  This distinction is useful when many original
coordinates share a small number of coupling directions.

\begin{lemma}[A phase-weighted projection]
\label{lem:phase-projection}
Let a random vector \(W\) have density \(\rho\leq M\), and
let \(Z\) be any random variable on the same probability
space with \(|Z|\leq1\).  The complex measure
\(\mathbb E[Z;W\in dw]\) has a density \(a(w)\) satisfying
\[
 |a(w)|\leq\rho(w),\qquad
 \|a\|_1\leq1,\qquad \|a\|_2^2\leq M.
\]
\end{lemma}
\begin{proof}
Disintegration gives
\(a(w)=\rho(w)\mathbb E[Z\mid W=w]\) almost everywhere.
Its modulus is at most \(\rho\), and
\(\int|a|^2\leq\int\rho^2\leq M\int\rho=M\).
\end{proof}

\begin{theorem}[An exact-rank crossing block]
\label{thm:exact-crossblock}
Let \(X\) and \(Y\) be independent random vectors and
\(B=U\Sigma V^{\mathsf T}\) have exact rank \(r\), where
\(U,V\) have orthonormal columns and
\(\Sigma=\operatorname{diag}(s_1,\ldots,s_r)>0\).
Suppose \(U^{\mathsf T}X\) and \(V^{\mathsf T}Y\) have
densities bounded by \(M_X,M_Y\).
For arbitrary measurable real functions \(F,G\), put
\[
 Q=F(X)+2X^{\mathsf T}BY+G(Y).
\]
Then, for \(t\ne0\),
\[
 |\phi_Q(t)|\leq
 \min\left\{1,
 \left(\frac{\pi}{|t|}\right)^{r/2}
 \frac{\sqrt{M_XM_Y}}{\sqrt{s_1\cdots s_r}}\right\}.
\]
If \(r>2\), its density satisfies
\[
 \|p_Q\|_\infty\leq
 \frac r{r-2}
 \left(\frac{M_XM_Y}{s_1\cdots s_r}\right)^{1/r}.
\]
For product inputs with coordinate densities bounded by one,
one may take \(M_X,M_Y\leq2^{r/2}\).
\end{theorem}
\begin{proof}
Push \(e^{itF(X)}d\mathbb P_X\) onto \(U^{\mathsf T}X\),
obtaining a complex density \(a_t\), and similarly obtain
\(b_t\) from \(e^{itG(Y)}d\mathbb P_Y\).
Lemma~\ref{lem:phase-projection} gives
\(\|a_t\|_2^2\leq M_X\) and
\(\|b_t\|_2^2\leq M_Y\).
Independence and Fubini give
\[
 \phi_Q(t)=\int_{\mathbb R^r}
       \widehat a_t(2t\Sigma z)b_t(z)\,dz .
\]
By Plancherel,
\[
 \int|\widehat a_t(2t\Sigma z)|^2\,dz
 =\left(\frac{\pi}{|t|}\right)^r
   (s_1\cdots s_r)^{-1}\|a_t\|_2^2 .
\]
Cauchy--Schwarz proves the characteristic-function bound.
For \(\alpha>1\),
\[
 \frac1{2\pi}\int_{\mathbb R}\min\{1,K|t|^{-\alpha}\}\,dt
 =\frac{\alpha}{\pi(\alpha-1)}K^{1/\alpha}.
\]
Take \(\alpha=r/2\).  The last assertion is the projection
bound of \citet{LPP2016}.
\end{proof}

We next allow the light crossing matrix to have arbitrary
rank.  For a matrix \(C\) with rows \(c_i\), let
\[
 Z_\ell(C)=e_\ell(CC^{\mathsf T})
 =\sum_{|S|=\ell}\det(C_SC_S^{\mathsf T}).
\]
When \(Z_\ell(C)>0\), its \(\ell\)-volume row marginals are
\[
 \alpha_i=\sum_{S\ni i}
 \frac{\det(C_SC_S^{\mathsf T})}{Z_\ell(C)}.
\]
They satisfy \(\sum_i\alpha_i=\ell\).

\begin{lemma}[A weighted-subset Fourier estimate]
\label{lem:weighted-subset}
Let \(D\in\mathbb R^m\) have independent coordinate
densities bounded by one.  Let \(M\) be a \(k\)-by-\(m\)
matrix of row rank \(k\), with product of singular values
\(S_M\).  Write \(\Pi\) for projection onto its row space.
For a matrix \(A_0\) with rows \(a_i\), put
\(C=A_0(I-\Pi)\).
Suppose \(Z_\ell(C)>0\) and every \(\ell\)-volume row
marginal is at most \(1/2\).
Let \(\psi_i\) be characteristic functions of scalar
probability densities bounded by one, and let \(\kappa\)
be any nonnegative probability density on \(\mathbb R^k\).
Then
\begin{equation}
 \mathbb E\left[\kappa(MD)
     \prod_i|\psi_i(2t\langle a_i,D\rangle)|\right]
 \leq
 \frac{2^{(k+\ell)/2}(2\pi)^\ell e^{\ell/2}}
 {S_M\sqrt{Z_\ell(C)}\,|t|^\ell}.
 \label{eq:weighted-subset}
\end{equation}
When \(\ell=1\), the factor \(e^{\ell/2}\) may be omitted.
\end{lemma}
\begin{proof}
Set
\[
 w_S=\frac{\det(C_SC_S^{\mathsf T})}{Z_\ell(C)}.
\]
Discard factors with \(\alpha_i=0\), which only increases
the left side.  Weighted H\"older for the finite measure
\(\kappa(Md)d\mathbb P_D(d)\) bounds the left side by
\(\prod_{S:w_S>0}I_S^{w_S}\), where
\[
 I_S=\mathbb E\left[\kappa(MD)
       \prod_{i\in S}
       |\psi_i(2t\langle a_i,D\rangle)|^{1/\alpha_i}\right].
\]
Indeed, the exponent of each factor after taking this
weighted product is
\(\sum_{S\ni i}w_S/\alpha_i=1\).

For \(i\in S\), write \(a_i=c_i+\lambda_iM\).
The linear image
\((MD,(\langle c_i,D\rangle)_{i\in S})\) has density
bounded by
\[
 \frac{2^{(k+\ell)/2}}
 {S_M\sqrt{\det(C_SC_S^{\mathsf T})}}.
\]
To verify the determinant, the rows of \(C\) are
perpendicular to the rows of \(M\), so the Gram determinant
of the combined map is
\(S_M^2\det(C_SC_S^{\mathsf T})\).
Factor this map through an orthonormal projection and
apply Lemma~\ref{four:projection}.
For each fixed \(x=MD\), shift the remaining coordinates
by \(\lambda_i x\), integrate \(\kappa\), and use
Lemma~\ref{four:lp} at exponent \(1/\alpha_i\geq2\).
The result is
\[
 I_S\leq
 \frac{2^{(k+\ell)/2}(2\pi)^\ell}
 {S_M|t|^\ell}
 \frac{\prod_{i\in S}\sqrt{\alpha_i}}
 {\sqrt{\det(C_SC_S^{\mathsf T})}}.
\]
The weighted product of the last factors is
\[
 Z_\ell(C)^{-1/2}
 \exp\left\{\frac12\left(
       H(w)+\sum_i\alpha_i\log\alpha_i\right)\right\},
\]
where \(H(w)=-\sum_Sw_S\log w_S\).
Entropy subadditivity for the indicator coordinates of
the random subset \(S\) gives
\[
 H(w)\leq\sum_i\left[-\alpha_i\log\alpha_i
 -(1-\alpha_i)\log(1-\alpha_i)\right].
\]
This also follows directly from nonnegativity of relative
entropy against the independent Bernoulli distribution
with parameters \(\alpha_i\).
Since \(-(1-a)\log(1-a)\leq a\) and
\(\sum_i\alpha_i=\ell\), the exponential is at most
\(e^{\ell/2}\).  For \(\ell=1\), \(w_{\{i\}}=\alpha_i\),
so the entropy expression cancels exactly.
\end{proof}

\begin{theorem}[Grouped hubs with an arbitrary-rank crossing remainder]
\label{thm:grouped-crossing}
Let \(X_H\) and \(X'\) be independent, with \(X'\) having
independent coordinate densities bounded by one.
Let \(U\) have \(k\) orthonormal columns, and suppose
\(W=U^{\mathsf T}X_H\) has density at most \(M_H\).
Consider
\[
 Q=F(X_H)+2\langle W,MX'+\beta\rangle+R(X'),
\]
where \(F\) is arbitrary measurable, \(M\) has row rank
\(k\), and \(R\) is an arbitrary quadratic polynomial
with matrix \(A_R\).
Split the light coordinates into \(I,J\), assume
\(M_J\) has row rank \(k\), and write \(S\) for the
product of its singular values.  Define
\[
 C=(A_R)_{IJ}(I-\Pi_{\operatorname{row}(M_J)}).
\]
Suppose \(Z_\ell(C)>0\) and its \(\ell\)-volume row
marginals are at most \(1/2\).  Put \(d=k+\ell\).
For \(t\ne0\),
\[
 |\phi_Q(t)|\leq\min\{1,K|t|^{-d/2}\},
 \quad
 K=\frac{M_H^{1/2}2^{(k+3\ell)/4}\pi^{d/2}e^{\ell/4}}
 {S^{1/2}Z_\ell(C)^{1/4}}.
\]
If \(d>2\), then
\begin{equation}
 \|p_Q\|_\infty\leq
 \frac d{d-2}
 2^{(k+3\ell)/(2d)}e^{\ell/(2d)}
 \frac{M_H^{1/d}}{S^{1/d}Z_\ell(C)^{1/(2d)}}.
 \label{eq:grouped-crossing-density}
\end{equation}
For \(\ell=1\), omit the exponential factors.
\end{theorem}
\begin{proof}
Push \(e^{itF(X_H)}d\mathbb P_H\) onto \(W\), obtaining
a complex density \(a_t\) with
\(\|a_t\|_2^2\leq M_H\), and put
\(h_t(l)=\widehat a_t(2t(l+\beta))\).
Plancherel gives
\[
 \int|h_t(l)|^2\,dl\leq
 M_H(\pi/|t|)^k.
\]
Let \(\mu_t(l)\) be the density of the complex measure
\(\mathbb E[e^{itR(X')};MX'\in dl]\).
The ordinary projection density of \(MX'\) is bounded,
so Lemma~\ref{lem:phase-projection} implies \(\mu_t\in L^2\).
Independence yields
\[
 \phi_Q(t)=\int h_t(l)\mu_t(l)\,dl,\qquad
 |\phi_Q(t)|^2\leq M_H(\pi/|t|)^k\|\mu_t\|_2^2.
 \label{eq:grouped-pairing}
\]

To bound the last norm, take a nonnegative even compactly
supported approximate identity \(k_\varepsilon\) on
\(\mathbb R^k\) and write
\[
 m_\varepsilon(l)=
 \mathbb E[e^{itR(X')}k_\varepsilon(l-MX')].
\]
Conditional Cauchy--Schwarz in \(X_I\), followed by
integration in \(l\) and an independent copy \(X_J^*\),
gives
\[
 \|m_\varepsilon\|_2^2
 \leq
 \mathbb E\left[
 e^{it(R(X_I,X_J)-R(X_I,X_J^*))}
 \kappa_\varepsilon(M_JD_J)\right],
\]
where \(D_J=X_J-X_J^*\) and
\(\kappa_\varepsilon=k_\varepsilon*k_\varepsilon\)
is a nonnegative probability density.
Every term purely in \(I\) cancels.  Terms purely in
\(J\) give a phase of modulus one, while the cross terms
give \(2\langle X_I,(A_R)_{IJ}D_J\rangle\).
Averaging over \(X_I\) and taking absolute values bounds
the last expression by
\[
 \mathbb E\left[
 \kappa_\varepsilon(M_JD_J)
 \prod_{i\in I}
 |\widehat f_i(2t\langle(A_R)_{iJ},D_J\rangle)|
 \right].
\]
The coordinates of \(D_J\) are independent with densities
\(f_j*\widetilde f_j\leq1\).
Apply Lemma~\ref{lem:weighted-subset}, uniformly in
\(\varepsilon\), and pass to the limit by \(L^2\)
approximate-identity convergence.  This yields
\[
 \|\mu_t\|_2^2\leq
 \frac{2^{(k+\ell)/2}(2\pi)^\ell e^{\ell/2}}
 {S\sqrt{Z_\ell(C)}\,|t|^\ell}.
\]
Combining with the hub Plancherel estimate proves the
characteristic-function bound.  The integration formula
in Theorem~\ref{thm:exact-crossblock} gives
\eqref{eq:grouped-crossing-density}.
\end{proof}

\begin{corollary}[One hub and a spread second crossing volume]
\label{cor:one-hub-arbitrary-crossing}
For one original hub, let
\[
 Q=F(X_h)+2X_h(\langle w,X'\rangle+\beta)+R(X').
\]
For a split \(I,J\), suppose \(w_J\ne0\), and let
\(C=(A_R)_{IJ}(I-\Pi_{w_J})\).
If \(Z_2(C)>0\) and every second-volume row marginal is
at most \(1/2\), then
\[
 \|p_Q\|_\infty
 \leq3\,2^{7/6}e^{1/3}
       (|w_J|^2Z_2(C))^{-1/6}
 <9.4(|w_J|^2Z_2(C))^{-1/6}.
\]
There is no restriction on \(\operatorname{rank}C\).
\end{corollary}
\begin{proof}
Use Theorem~\ref{thm:grouped-crossing} with
\(k=1,\ell=2\), and \(M_H=1\).
\end{proof}

\begin{corollary}[Many coordinates with two exact hub directions]
\label{cor:many-two-hubs}
Suppose a hub--light cross block has exact rank two and
write it as \(UM\), with \(U^{\mathsf T}U=I_2\).
Every internal hub quadratic, and even every measurable
internal hub function, is allowed.
If a light split has \(\operatorname{rank}M_J=2\),
singular-value product \(S\), and projected crossing
matrix \(C\) with
\[
 F_0=\|C\|_{\mathrm F}>0,\qquad
 \|c_i\|^2\leq F_0^2/2,
\]
then product hub and light laws satisfy
\[
 \|p_Q\|_\infty
 \leq3\,2^{7/6}(SF_0)^{-1/3}<6.74(SF_0)^{-1/3}.
\]
\end{corollary}
\begin{proof}
The projection \(U^{\mathsf T}X_H\) has density at most
two.  Apply Theorem~\ref{thm:grouped-crossing} with
\(k=2,\ell=1\), \(M_H=2\), and \(Z_1(C)=F_0^2\).
\end{proof}

\section{Concentrated directions and uniform bounds}
\label{sec:structural-boundaries}

\begin{lemma}[Few concentrated directions in exact finite rank]
\label{lem:heavy-isotropic}
Let \(A=RMR^{\mathsf T}\), with \(R^{\mathsf T}R=I_r\)
and every eigenvalue of \(M\) of modulus at least one.
Assume \(A\) has no \(c\)-pivot.  Let \(H\) be any set of
original coordinates and
\[
 G_H=R_H^{\mathsf T}R_H,\qquad
 \tau=|H|\sqrt{2c}.
\]
If the eigenvalues of \(G_H\) are \(g_1\geq\cdots\geq g_r\),
then \(g_k\leq\tau\) whenever \(k>r/2\).
If \(M\) is definite, every row leverage
\(|R^{\mathsf T}e_i|^2\) is less than \(c\).
\end{lemma}
\begin{proof}
By \eqref{eq:no-pivot-entries},
\(\|A_{HH}\|\leq\|A_{HH}\|_{\mathrm F}\leq\tau\).
Rotate the frame so that \(G_H\) is diagonal.
If \(g_k>\tau\), the compression of \(M\) to its first
\(k\) frame coordinates has norm at most \(\tau/g_k<1\).
Indeed, writing \(R_H=VG_H^{1/2}\) on its support
expresses \(A_{HH}\) as an isometric copy of
\(G_H^{1/2}MG_H^{1/2}\), and division by the first
\(k\) eigenvalues of \(G_H^{1/2}\) gives this norm bound.
Interlacing then forces \(M\) to have at least \(k\)
eigenvalues greater than \(-1\), hence at least \(1\),
and at least \(k\) eigenvalues less than \(1\), hence at
most \(-1\).  These are disjoint sets, so \(2k\leq r\).
For definite \(M\), its eigenvalue hypothesis gives
\(|a_{ii}|\geq|R^{\mathsf T}e_i|^2\); use the
no-pivot diagonal bound.
\end{proof}

For example, if \(H=\{i:|R^{\mathsf T}e_i|^2\geq\delta\}\),
then \(|H|\leq r/\delta\).  In exact rank three,
Lemma~\ref{lem:heavy-isotropic} permits only one strongly
concentrated frame direction when \(c\) is small relative
to \(\delta\).  It makes no assertion about a decomposition
with an arbitrary additional diagonal: that diagonal may
cancel the diagonal entries of the low-rank component.

\begin{proposition}[A normalized exact-rank-three transition]
\label{prop:rank-three-epsilon}
Let \(w,v,p\) be orthonormal vectors on the light
coordinates, write
\(L=\langle w,X'\rangle\), \(V=\langle v,X'\rangle\),
\(P=\langle p,X'\rangle\), and put
\[
 Q_\varepsilon=2X_hL+V^2+2\varepsilon LP.
\]
For every real \(\varepsilon\), its quadratic matrix has
exact rank three and nonzero eigenvalues
\[
 \sqrt{1+\varepsilon^2},\qquad
 -\sqrt{1+\varepsilon^2},\qquad 1.
\]
For every fixed \(c>0\) and \(0\leq\varepsilon\leq1\),
there are arbitrarily large spread frames for which this
matrix has no \(c\)-pivot.
If the light frame entries all have modulus \(m^{-1/2}\),
then every principal submatrix of order at most \(N\)
has third singular value at most \(N/m\).
\end{proposition}
\begin{proof}
In the orthonormal basis \((e_h,w,v,p)\), the matrix is
\[
 \begin{pmatrix}
 0&1&0&0\\
 1&0&0&\varepsilon\\
 0&0&1&0\\
 0&\varepsilon&0&0
 \end{pmatrix}.
\]
On the plane spanned by \(w\) and
\((e_h+\varepsilon p)/\sqrt{1+\varepsilon^2}\), it is
\(\bigl(\begin{smallmatrix}0&\sqrt{1+\varepsilon^2}\\
\sqrt{1+\varepsilon^2}&0\end{smallmatrix}\bigr)\).
The \(v\)-direction has eigenvalue one and the remaining
orthogonal direction has eigenvalue zero.

Take three orthogonal Walsh columns on \(m\) coordinates,
normalized to entries of modulus \(m^{-1/2}\), with
\(m>6/c\).  The hub diagonal is zero, every hub--leaf
minor has modulus \(1/m<c\), and light-block entries
have modulus at most \(3/m\).
Their diagonals are less than \(c\) and their principal
two-coordinate determinants have modulus at most
\(18/m^2<c\), after enlarging \(m\) if necessary.
This proves the no-pivot assertion.

For any coordinate set \(K\), the principal matrix is
the sum of a matrix of rank at most two,
\[
 \bigl((e_h+\varepsilon p)w^{\mathsf T}
       +w(e_h+\varepsilon p)^{\mathsf T}\bigr)_{KK},
\]
and \(v_Kv_K^{\mathsf T}\), of norm
\(|v_K|^2\leq |K|/m\).
The singular-value approximation inequality gives
\(s_3(A_{KK})\leq |K|/m\).
\end{proof}

The family in Proposition~\ref{prop:rank-three-epsilon}
has a third singular value of one for every \(\varepsilon\).
The endpoint \(\varepsilon=0\) and every fixed nonzero
\(\varepsilon\) are covered by results in
Section~\ref{sec:hubs}, but the available bounds for
arbitrary rough laws need not stay bounded as
\(\varepsilon\) tends to zero.
Corollary~\ref{cor:hub-loglog-density} gives a much smaller
loss than the direct leakage estimate.
Theorem~\ref{thm:grouped-crossing} also retains a dependence
on this parameter: in the one-hub split its projected
crossing matrix is
\[
 C=v_I(v_J^\perp)^{\mathsf T}
       +\varepsilon w_I(p_J^\perp)^{\mathsf T},
\]
so \(Z_2(C)\leq4\varepsilon^2\) for
\(0\leq\varepsilon\leq1\).
The finite-core theorem is likewise insufficient uniformly,
by the last assertion of Proposition~\ref{prop:rank-three-epsilon}.
These are limitations of the respective estimates, not
counterexamples to a bounded density.

\begin{proposition}[Two hubs do not suffice with an arbitrary remainder]
\label{prop:two-hub-counterexample}
For every \(0<c\leq1\), there is a real symmetric matrix
with no \(c\)-pivot, exactly two eigenvalues of modulus
at least one, and a two-row hub--light block whose
singular values are both one, such that the corresponding
quadratic form under independent standard Gaussian
coordinates has unbounded density.
\end{proposition}
\begin{proof}
Choose \(0<b<c/2\), put
\(B=\operatorname{diag}(b,-b)\), and let \(M\) have
two orthonormal, sufficiently spread rows.
Set
\[
 A=\begin{pmatrix}
 B&M\\ M^{\mathsf T}&M^{\mathsf T}B^{-1}M
 \end{pmatrix}.
\]
Writing \(L=MX'\),
\[
 Q_A=(X_H+B^{-1}L)^{\mathsf T}
       B(X_H+B^{-1}L).
\]
The matrix has rank two.  On the two mutually orthogonal
hub--partner planes its nonzero eigenvalues are
\(b+b^{-1}\) and \(-b-b^{-1}\).
By taking the two rows of \(M\) with entries of modulus
at most \(\eta\), where
\(\eta<c/2\) and \(2\eta^2/b<c/2\), every entry of
\(A\) has modulus less than \(c/2\).  Thus every
diagonal entry is smaller than \(c\) and every principal
two-coordinate determinant has modulus less than
\(c^2/2<c\).

For Gaussian inputs, \(Y=X_H+B^{-1}L\) is a nondegenerate
Gaussian vector whose density is bounded below by some
\(\gamma>0\) on a neighbourhood of the origin.
Use the change of variables
\(u=\sqrt b(Y_1-Y_2)\),
\(v=\sqrt b(Y_1+Y_2)\).
Then \(Q_A=uv\) and the Jacobian is \(2b\).
Choose \(r>0\) so that the square \(|u|,|v|\leq r\)
maps into that neighbourhood.  For \(0<a\leq r^2\),
\[
 \mathbb P\{|Q_A|\leq a\}
 \geq\frac{\gamma}{2b}
 \left|\{(u,v):|u|,|v|\leq r,\ |uv|\leq a\}\right|
 =\frac{2\gamma}{b}a(1+\log(r^2/a)).
\]
The ratio to \(2a\) is unbounded as \(a\downarrow0\).
\end{proof}

Proposition~\ref{prop:two-hub-counterexample} has only two
large eigenvalues.  It shows why a two-hub bound must use
additional information about the remainder, such as that
in Theorem~\ref{thm:grouped-crossing}; it does not refute
the three-eigenvalue conjecture.

The structural statements above do not imply that every
no-pivot matrix admits a split for which the quantities in
\eqref{eq:grouped-crossing-density} are uniformly nondegenerate.
In particular, the grouped hub block must have exact rank;
discarding the smaller singular directions does not cover a
nearly rank-two block.  A dimension-independent bound requires
uniform control of these limiting configurations and of the
one-hub transition above.

\section{Quadratic forms with hub coordinates}\label{sec:hubs}

Let $X_0,X_1,\ldots,X_m$ be independent real random variables with
probability densities bounded by one, and write $X'=(X_1,\ldots,X_m)$.
For orthonormal vectors $w,v\in\mathbb R^m$ put
$L=\langle w,X'\rangle$ and $V=\langle v,X'\rangle$.
The basic family in this section is
\begin{equation}\label{eq:hub-normal}
 Q=aX_0^2+2X_0(L+\beta)+\lambda V^2
       +2L\langle q,X'\rangle+\sum_{i=1}^m\phi_i(X_i),
 \qquad \phi_i(x)=d_ix^2+b_ix+c_i,
\end{equation}
where $\lambda\ne0$.  Unless a stronger condition is stated, assume
\begin{equation}\label{eq:hub-frame}
 w_i^2+v_i^2\le\tfrac12\qquad(1\le i\le m).
\end{equation}
The coefficients of the one-coordinate polynomials are unrestricted.
We use the Fourier convention
$\widehat\mu(t)=\int e^{itx}\,\mu(dx)$, and all logarithms are natural.

The one-coordinate term $X_0$ is a hub: its coupling to the remaining
coordinates is carried by the single linear statistic $L$.  The estimates
below distinguish the independent hub law from the dependence between
$L$, $V$, and the other linear statistics of $X'$.

\subsection{Projection and Fourier estimates}

We use the following consequence of the density bound for projections of
product measures \citep{LPP2016,RV2015}: if $T$ is an $r\times m$
matrix of rank $r$, then
\begin{equation}\label{eq:hub-projection}
 \|p_{TX'}\|_\infty\le
       \frac{2^{r/2}}{\sqrt{\det(TT^T)}}.
\end{equation}
The bound for an orthonormal projection is $2^{r/2}$.  No independence
between the projected coordinates is asserted.

\begin{lemma}[Modulated frame integral]\label{lem:hub-frame-integral}
Let $f\le1$ be a probability density and let $g$ be a complex function
with $|g|\le f$.  For $p\ge2$,
\begin{equation}\label{eq:hub-scalar-lp}
 \int_{\mathbb R}|\widehat g(t)|^p\,dt\le\frac{4\pi}{\sqrt p}.
\end{equation}
Let $e_1,\ldots,e_r$ be orthonormal vectors in $\mathbb R^m$, set
$u_i=(e_{1i},\ldots,e_{ri})$, and assume $|u_i|^2\le1/2$.
For complex functions $g_i$ dominated by the coordinate densities and
arbitrary real shifts $s_i$,
\begin{equation}\label{eq:hub-frame-int}
 \int_{\mathbb R^r}\prod_i
   |\widehat g_i(\langle u_i,z\rangle+s_i)|\,dz\le(4\pi)^r.
\end{equation}
\end{lemma}

\begin{proof}
For an integer $k\ge1$, Plancherel and convolution give
\[
 \int|\widehat g|^{2k}
 =2\pi(g*\widetilde g)^{*k}(0)
 \le2\pi(f*\widetilde f)^{*k}(0),
 \qquad \widetilde g(x)=\overline{g(-x)}.
\]
The density on the right is that of a sum of $2k$ independent variables
whose densities are bounded by one.  Applying
\eqref{eq:hub-projection} to this one-dimensional linear map gives a
bound $1/\sqrt k$ for that density.  Since $|\widehat g|\le1$, choose
$k=\lfloor p/2\rfloor$.  For $p\ge4$, $k\ge p/4$, and for
$2\le p<4$ Plancherel gives the required bound directly.

Set $c_i=|u_i|^2$ and omit the zero rows.  The vectors
$u_i/\sqrt{c_i}$ with weights $c_i$ form geometric
Brascamp--Lieb data, since $\sum_i u_i u_i^T=I_r$.
The geometric Brascamp--Lieb inequality
\citep{BL1976,Ball1986,Barthe1998}, including arbitrary translations of
the one-dimensional arguments, bounds the integral in
\eqref{eq:hub-frame-int} by
\[
 \prod_i\left(c_i^{-1/2}
          \int|\widehat g_i(t)|^{1/c_i}\,dt\right)^{c_i}
 \le(4\pi)^{\sum_i c_i}=(4\pi)^r.
\]
Here $1/c_i\ge2$, so \eqref{eq:hub-scalar-lp} applies.
\end{proof}

For $\ell\in\mathbb R$ define the frozen polynomial
\[
 Y_\ell=\lambda V^2+2\ell\langle q,X'\rangle
                         +\sum_i\phi_i(X_i).
\]
For each fixed $\ell_0$, let $\mu_{\ell_0,\ell}$ denote the
unnormalized slice of the joint law of $(L,Y_{\ell_0})$ at $L=\ell$.
Its total mass is $\rho_L(\ell)$, not one.

\begin{lemma}[Frozen slices]\label{lem:hub-frozen}
For every fixed $\ell_0$, for almost every $\ell$, and every $t\ne0$,
\begin{equation}\label{eq:hub-slice-fourier}
 |\widehat\mu_{\ell_0,\ell}(t)|
 \le\min\left\{\rho_L(\ell),
                   \frac{4\sqrt\pi}{\sqrt{|t\lambda|}}\right\}.
\end{equation}
In particular, for every interval $I$,
\begin{equation}\label{eq:hub-slice-interval}
 \mu_{\ell_0,\ell}(I)
       \le2\pi^2\sqrt{\frac{|I|}{|\lambda|}}.
\end{equation}
\end{lemma}

\begin{proof}
Write $Y_{\ell_0}=\lambda V^2+\sum_i(d_iX_i^2+c_i'X_i)+c'$.
For fixed $t\ne0$ the Fourier transform in $L$ of the complex measure
$\mathbb E[e^{itY_{\ell_0}};L\in\cdot]$ is
\[
 \Psi(z,t)=\mathbb E e^{izL+itY_{\ell_0}}.
\]
The Fresnel identity, first with Gaussian damping, gives
\[
 |\Psi(z,t)|\le(4\pi|t\lambda|)^{-1/2}
  \int\prod_i|\widehat g_i(zw_i+uv_i+tc_i')|\,du,
 \qquad g_i(x)=f_i(x)e^{itd_ix^2}.
\]
Lemma~\ref{lem:hub-frame-integral} bounds the integral in $(z,u)$ by
$(4\pi)^2$.  Fourier inversion in $z$ therefore gives the second term
in \eqref{eq:hub-slice-fourier}.  The first follows from the total mass
of the slice.  Disintegration first identifies this transform for almost
every pair $(\ell,t)$; continuity in $t$ identifies it for every $t$ on
a common full-measure set of $\ell$.

If $I$ has center $y_0$ and length $2b$, the function
\[
 K(y)=\left(\frac{\sin(\pi(y-y_0)/(4b))}
                         {\pi(y-y_0)/(4b)}\right)^2
\]
satisfies $K\ge(8/\pi^2)\mathbf1_I$; its Fourier transform is supported
on $|t|\le\pi/(2b)$ and has modulus at most $4b$.
Parseval and \eqref{eq:hub-slice-fourier} give
\[
 \mu(I)\le\frac{\pi^2}{8}\frac{4b}{2\pi}
   \int_{|t|\le\pi/(2b)}\frac{4\sqrt\pi}{\sqrt{|t\lambda|}}\,dt
 =2\pi^2\sqrt{\frac{2b}{|\lambda|}}.
\]
All the Fourier manipulations follow first with a positive Gaussian
damping; the displayed integrable bounds justify its removal.
\end{proof}

\begin{remark}[Freezing and exceptional slices]\label{rem:hub-freezing}
An estimate valid for every fixed $\ell_0$ and almost every $\ell$ cannot
be substituted on the diagonal $\ell_0=\ell$ without justification.
We use the following discretization whenever the slice parameter varies.
Partition the real line into cells $C_j$ of width $h$, choose
$\ell_j\in C_j$, and replace $Y_L$ by $Y_{\ell_j}$ on
$\{L\in C_j\}$.  Then
$Y_{\ell_j}-Y_L=2(\ell_j-L)\langle q,X'\rangle\to0$ almost surely.
For a bounded kernel $G(\ell,y)$ continuous off the null set
$\ell=-\beta$, dominated convergence gives
\begin{equation}\label{eq:hub-freezing-limit}
 \mathbb E G(L,Y_L)
 =\lim_{h\downarrow0}\sum_j\int_{C_j}\int
       G(\ell_j,y)\,\mu_{\ell_j,\ell}(dy)\,d\ell.
\end{equation}
For each partition only countably many frozen parameters occur, so the
slice bounds apply simultaneously.  The analogous construction with
rectangular cells applies to several hub partners.  Riemann sums used
below can first be restricted to bounded sets; tails are controlled by
the total mass.  This also justifies the formulas for arbitrary
unbounded coordinate laws.
\end{remark}

\subsection{Real-variable estimates for one hub}

\begin{lemma}[A profile criterion]\label{lem:hub-profile}
Assume $a=0$.  Suppose that for every fixed $\ell_0$, almost every
$\ell$, and every Borel set $S$,
\[
 \mu_{\ell_0,\ell}(S)\le K|S|^\theta,
 \qquad 0<\theta\le1.
\]
Then
\begin{equation}\label{eq:hub-profile-bound}
 \|p_Q\|_\infty\le\frac12+\frac{2^\theta K}{\theta}.
\end{equation}
More generally, a concave nondecreasing profile $\omega$ with
$\omega(0)=0$ gives
$\|p_Q\|_\infty\le\frac12+\int_0^2\omega(u)\,du/u$ whenever
that integral is finite.
\end{lemma}

\begin{proof}
Fix a level $s$ and a window radius $\delta>0$.  With $t=\ell+\beta$ put
\[
 G(\ell,y)=\int f_0(x)
            \mathbf1_{\{|2tx+y-s|\le\delta\}}\,dx.
\]
For $t\ne0$, $0\le G\le\delta/|t|$, and $\int G(\ell,y)\,dy=2\delta$.
If $M=\delta/|t|$, layer cake and Jensen's inequality give
\[
 \int G\,d\mu
  \le K\int_0^M |\{G>u\}|^\theta\,du
  \le KM(2\delta/M)^\theta
  =K\delta\,2^\theta|t|^{\theta-1}.
\]
For $|t|\ge1$ use $G\le\delta$ and the slice mass
$\rho_L(\ell)$, whose integral is one.  Apply
\eqref{eq:hub-freezing-limit}, integrate the displayed bound over
$|t|<1$, divide by $2\delta$, and let $\delta\downarrow0$.
This proves \eqref{eq:hub-profile-bound}.  The same layer-cake argument
gives $M\omega(2\delta/M)$ for a concave profile and yields the final
assertion.  The resulting uniform interval bounds imply absolute
continuity and the stated essential density bound.
\end{proof}

\begin{theorem}[An interval-supported hub]\label{thm:hub-interval}
Let $a=0$ and suppose $X_0$ is supported in an interval of length $D$.
Then every form \eqref{eq:hub-normal} satisfies
\begin{equation}\label{eq:hub-interval-bound}
 \|p_Q\|_\infty\le
 \frac32(4\sqrt2\pi^2)^{2/3}(D/|\lambda|)^{1/3}
 <22(D/|\lambda|)^{1/3}.
\end{equation}
\end{theorem}

\begin{proof}
Let $J$ be the supporting interval and replace $f_0$ by $\mathbf1_J$.
For $t=\ell+\beta\ne0$ the resulting window kernel is
\[
 G(\ell,y)=\frac{|(s-y+[-\delta,\delta])\cap2tJ|}{2|t|}.
\]
Three estimates follow by Fubini and
Lemma~\ref{lem:hub-frozen}:
\begin{align*}
 \int G\,d\mu&\le (\delta/|t|)\rho_L(\ell),\\
 \int G\,d\mu&\le (\delta/|t|)2\pi^2
                  \sqrt{(2|t|D+2\delta)/|\lambda|},\\
 \int G\,d\mu&\le 2\pi^2D\sqrt{2\delta/|\lambda|}.
\end{align*}
Use the first on $|t|\ge T$ and the minimum of the other two on
$|t|<T$, with the freezing procedure of
Remark~\ref{rem:hub-freezing}.  After division by $2\delta$, the range
$|t|<\delta$ tends to zero.  On $\delta\le|t|<T$, dominated
convergence applies with an integrable multiple of $|t|^{-1/2}$.
It follows that, uniformly in $s$,
\[
 \limsup_{\delta\downarrow0}
 \frac{\mathbb P(|Q-s|\le\delta)}{2\delta}
 \le\frac1{2T}+4\sqrt2\pi^2\sqrt{DT/|\lambda|}.
\]
Choose $T=(4\sqrt2\pi^2)^{-2/3}(D/|\lambda|)^{-1/3}$.
\end{proof}

The next result needs no spread condition.  Its only use of the product
law is the two-dimensional projection bound.

\begin{theorem}[A transverse square with an independent hub]
\label{thm:hub-pure}
Let $w,v$ be orthonormal and let $v_c,F_c$ be arbitrary measurable
real functions.  Set
\[
 Q=2X_0(L+\beta)+\lambda(V-v_c(L))^2+F_c(L),
 \qquad \mu=|\lambda|>0.
\]
Then
\begin{equation}\label{eq:hub-pure-bound}
 \|p_Q\|_\infty\le
 \begin{cases}
 2\sqrt2+\sqrt2\log(1/3+4\sqrt2/\mu),&0<\mu\le6\sqrt2,\\
 (192/\mu)^{1/3},&\mu\ge6\sqrt2.
 \end{cases}
\end{equation}
\end{theorem}

\begin{proof}
For fixed $s$, integrating out $X_0$ gives
\[
 p_Q(s)=\int\rho_{L,V}(\ell,v)
 \frac{f_0((s-\lambda(v-v_c(\ell))^2-F_c(\ell))/(2t))}
      {2|t|}\,d\ell\,dv,
 \qquad t=\ell+\beta.
\]
Denote the fraction by $K(\ell,v)$.  It is at most $1/(2|t|)$.
Changing variables on the two branches of the square and using
$f_0\le1$, $\int f_0=1$, gives
\[
 \int K(\ell,v)\,dv\le\sqrt2(\mu|t|)^{-1/2}.
\]
Indeed the remaining integral is a one-sided negative half moment of
$f_0$, bounded by $\int_0^1u^{-1/2}\,du=2$.
Since $\rho_{L,V}\le2$ and $\rho_L\le\sqrt2$, the contribution at $t$
is bounded by $g(t)/(2|t|)$ for some nonnegative function satisfying
\[
 \int g\le1,\qquad
 g(t)\le M(t):=\min\{\sqrt2,4\sqrt2\mu^{-1/2}\sqrt{|t|}\}.
\]
The weight $1/(2|t|)$ decreases with $|t|$.  Thus its maximal integral
under these constraints is obtained by taking $g=M$ on
$[-\tau,\tau]$, where $\int_{-\tau}^{\tau}M=1$.
The envelope switches at $t_* =\mu/16$, and its mass on
$[-t_*,t_*]$ is $\sqrt2\mu/12$.
If this mass is at most one, the inner contribution is $2\sqrt2$ and
the outer contribution is $\sqrt2\log(\tau/t_*)$, with
$\tau/t_*=1/3+4\sqrt2/\mu$.
Otherwise
$\tau=(3\sqrt\mu/(16\sqrt2))^{2/3}$, and the integral is
$8\sqrt2\mu^{-1/2}\sqrt\tau=(192/\mu)^{1/3}$.
The argument can equivalently be applied to interval probabilities,
so no prior existence of $p_Q$ is required.
\end{proof}

\subsection{A nonzero transverse cross term}

Write
\[
 q=q_ww+q_vv+\varepsilon e,\qquad
 e\perp w,v,\quad |e|=1,\quad\varepsilon>0,
\]
and set $P=\langle e,X'\rangle$.

\begin{theorem}[Nonzero leakage]\label{thm:hub-leak}
Assume $w_i^2+v_i^2+e_i^2\le1/2$ for every light coordinate.  For
\eqref{eq:hub-normal}, including an arbitrary hub diagonal $a$,
\begin{align}
 |\mathbb E e^{itQ}|&\le
 4\sqrt2\pi^{3/2}(\varepsilon|\lambda|)^{-1/2}|t|^{-3/2},
 \label{eq:hub-leak-fourier}\\
 \|p_Q\|_\infty&\le
 3\,2^{5/3}(\varepsilon|\lambda|)^{-1/3}.
 \label{eq:hub-leak-density}
\end{align}
\end{theorem}

\begin{proof}
Complete the square in $V$:
$V'=V+(q_v/\lambda)L$ and
$\sigma=2q_w-q_v^2/\lambda$.
After moving $-2\varepsilon\beta P$ into the linear terms, write
\[
 Q=aX_0^2+2X_0(L+\beta)+\sigma L^2+
   2\varepsilon(L+\beta)P+W,
 \quad
 W=\lambda V'^2-2\varepsilon\beta P+\sum_i\phi_i(X_i).
\]
Let $g(\omega)=\int f_0(x)e^{ita x^2+i\omega x}\,dx$ and
\[
 H(\ell,p)=g(2t(\ell+\beta))
                   e^{it\sigma\ell^2+2it\varepsilon(\ell+\beta)p}.
\]
Let $n(\ell,p)$ be the density of
$\mathbb E[e^{itW};(L,P)\in\cdot]$.  Its Fourier transform is
integrable by the Fresnel identity and the three-dimensional frame
integral.  Fourier transformation of $H$ in $p$ pins
$\ell+\beta=-\gamma/(2t\varepsilon)$; consequently
\[
 |\widehat H(z,\gamma)|=
      \frac{\pi}{|t|\varepsilon}|g(-\gamma/\varepsilon)|.
\]
This identity is interpreted with Gaussian damping, or as a tempered
distribution followed by the indicated integrable pairing.
Put $G(\gamma)=\int|\widehat n(z,\gamma)|\,dz$.
Fresnel linearization in $V'$ and translation of the $z$ variable give
\[
 \|G\|_\infty\le(4\pi|t\lambda|)^{-1/2}(4\pi)^2,
 \qquad
 \|G\|_1\le(4\pi|t\lambda|)^{-1/2}(4\pi)^3.
\]
Thus $\|G\|_2\le(4\pi|t\lambda|)^{-1/2}(4\pi)^{5/2}$.
Plancherel gives
$\int|g(\gamma/\varepsilon)|^2\,d\gamma\le2\pi\varepsilon$.
Parseval and Cauchy--Schwarz yield
\[
 |\mathbb E e^{itQ}|
 \le\frac{1}{4\pi|t|\varepsilon}
       (2\pi\varepsilon)^{1/2}
       (4\pi|t\lambda|)^{-1/2}(4\pi)^{5/2},
\]
which is \eqref{eq:hub-leak-fourier}.
Finally, $\int_{\mathbb R}\min(1,B|t|^{-3/2})\,dt=6B^{2/3}$.
Fourier inversion gives \eqref{eq:hub-leak-density}.
\end{proof}

\begin{example}[A frozen-hub obstruction]\label{ex:hub-frozen-obstruction}
The nonzero-leakage estimate cannot be made uniform by bounding every
conditional density obtained on fixing $X_0=x$.  Fix $x=r>0$, let the
light vector have law $X'=-(r/\varepsilon)e+G$, where the coordinates
of $G$ are independent standard Gaussians, and take no other terms.
Then the conditional quadratic form is
\[
 \lambda\langle v,G\rangle^2+
       2\varepsilon\langle w,G\rangle\langle e,G\rangle.
\]
For $\lambda>0$, let $Z_1,Z_2,Z_3$ be its three independent standard
Gaussian coordinates.  The density of $2Z_1Z_2$ is positive and
continuous on any compact interval avoiding zero.  In the convolution
formula for $\lambda Z_3^2+2\varepsilon Z_1Z_2$ at zero, restrict
$\lambda Z_3^2$ to $[\varepsilon,2\varepsilon]$.
The square density is bounded below by a constant times
$(\lambda u)^{-1/2}$ there when $\varepsilon\le\lambda$; the second
density is bounded below by a constant times $\varepsilon^{-1}$.
Integration gives a lower bound
$c(\lambda\varepsilon)^{-1/2}$.
All input densities are bounded by one.  This is an obstruction to
uniform frozen-hub estimates, not a counterexample to the full quadratic
form after averaging the independent hub.
\end{example}

\subsection{Near-linear small-ball bounds and symmetric-window recovery}

\begin{theorem}[Double-logarithmic small balls]\label{thm:hub-loglog-smallball}
Assume $a=0$ and \eqref{eq:hub-frame}.  No restriction is imposed on
$q$ beyond \eqref{eq:hub-normal}.  For every $s\in\mathbb R$ and
$\delta>0$,
\begin{equation}\label{eq:hub-loglog-smallball}
 \frac{\mathbb P(|Q-s|\le\delta)}{2\delta}
 \le A_\delta:=\frac12+2\sqrt2+
                 \sqrt2\log_+(2B_\delta^2),
\end{equation}
where
\begin{equation}\label{eq:hub-bdelta}
 B_\delta=\frac{4+\sqrt{4/\pi+32\log_+(|\lambda|/\delta)}}
                     {2\sqrt{|\lambda|}}.
\end{equation}
\end{theorem}

\begin{proof}
Use the window kernel in the proof of Lemma~\ref{lem:hub-profile}.
For fixed $t=\ell+\beta\ne0$, it is
$G(\ell,\cdot)=\mathbf1_{[-\delta,\delta]}*g$, where
$g(u)=f_0((s-u)/(2t))/(2|t|)$.  Hence
\[
 \widehat G(\eta)=\frac{2\sin(\eta\delta)}{\eta}\widehat g(\eta),
 \qquad \int|\widehat g|^2\le\frac\pi{|t|}.
\]
For every applicable frozen slice, Lemma~\ref{lem:hub-frozen} gives
\[
 \int_{|\eta|\le1/\delta}|\widehat\mu(\eta)|^2\,d\eta
 \le\frac{4+32\pi\log_+(|\lambda|/\delta)}{|\lambda|}.
\]
To see this, use $|\widehat\mu|^2\le2$ below $1/|\lambda|$
and $|\widehat\mu|^2\le16\pi/(|\eta\lambda|)$ above it.
Parseval and Cauchy--Schwarz on the regions below and above
$1/\delta$ now give
\begin{align*}
 \int G\,d\mu
 &\le\frac1{2\pi}\left(
 2\delta\int_{|\eta|\le1/\delta}|\widehat g||\widehat\mu|
 +\int_{|\eta|>1/\delta}\frac2{|\eta|}
                          |\widehat g||\widehat\mu|\right)\\
 &\le\delta(|\lambda t|)^{-1/2}
          \left(\sqrt{4/\pi+32\log_+(|\lambda|/\delta)}+4\right).
\end{align*}
The high-frequency calculation uses
$\int_{|\eta|>1/\delta}|\eta|^{-3}\,d\eta=\delta^2$.
The alternative bound is
$\int G\,d\mu\le\sqrt2\delta/|t|$.
For $|t|\ge1$ retain the slice mass and obtain a contribution at most
$\delta$.  Formula~\eqref{eq:hub-freezing-limit} therefore implies
\[
 \frac{\mathbb P(|Q-s|\le\delta)}{2\delta}
 \le\frac12+\int_{|t|<1}
       \min\left\{\frac{\sqrt2}{2|t|},
                         B_\delta|t|^{-1/2}\right\}\,dt.
\]
If $2B_\delta^2\ge1$, split at
$|t|=1/(2B_\delta^2)$; the integral is
$2\sqrt2+\sqrt2\log(2B_\delta^2)$.
Otherwise it equals $4B_\delta\le2\sqrt2$.
\end{proof}

\begin{lemma}[Recovery from a symmetric window]\label{lem:hub-window-recovery}
Suppose a real random variable $Z$ satisfies
$|\mathbb E e^{itZ}|\le B|t|^{-3/2}$ for $t\ne0$.
Then it has a continuous bounded density and, for every $h>0$,
\begin{equation}\label{eq:hub-window-recovery}
 \|p_Z\|_\infty\le
 \sup_s\frac{\mathbb P(|Z-s|\le h)}{2h}
       +\frac{2\sqrt{2\pi}}{3\pi}B\sqrt h.
\end{equation}
\end{lemma}

\begin{proof}
The integrable characteristic function gives the density by Fourier
inversion.  Let $U_h=\mathbf1_{[-h,h]}/(2h)$.
Its Fourier transform is $\sin(ht)/(ht)$, which is at most one.
Consequently
\[
 |p_Z(s)-(p_Z*U_h)(s)|
 \le\frac{B\sqrt h}{\pi}
   \int_0^\infty u^{-3/2}\left(1-\frac{\sin u}{u}\right)\,du.
\]
Since $\sin u/u=\int_0^1\cos(ru)\,dr$, Tonelli and an integration by
parts give
\begin{align*}
 \int_0^\infty u^{-3/2}\left(1-\frac{\sin u}{u}\right)\,du
 &=\int_0^1\int_0^\infty u^{-3/2}(1-\cos(ru))\,du\,dr\\
 &=\sqrt{2\pi}\int_0^1\sqrt r\,dr
 =\frac23\sqrt{2\pi}.
\end{align*}
For completeness, the inner integral at $r=1$ is
$2\int_0^\infty u^{-1/2}\sin u\,du=\sqrt{2\pi}$;
this identity follows by exponential damping and the Gamma integral.
Finally $(p_Z*U_h)(s)=\mathbb P(|Z-s|\le h)/(2h)$.
\end{proof}

\begin{corollary}[Improved double-logarithmic density estimate]
\label{cor:hub-loglog-density}
Under the assumptions of Theorem~\ref{thm:hub-leak}, with $a=0$,
\begin{equation}\label{eq:hub-loglog-recovery}
 \|p_Q\|_\infty\le A_h+
       \frac{16\pi}{3}\sqrt{\frac{h}{\varepsilon|\lambda|}}
       \qquad(h>0).
\end{equation}
In particular,
\begin{equation}\label{eq:hub-loglog-density}
 \|p_Q\|_\infty\le\frac12+2\sqrt2+
       \sqrt2\log_+(2B_\varepsilon^2)
                +\frac{16\pi}{3\sqrt{|\lambda|}}.
\end{equation}
For fixed $|\lambda|\ge1$, this is
$O(1+\log\log(1/\varepsilon))$ as $\varepsilon\downarrow0$.
\end{corollary}

\begin{proof}
Apply Lemma~\ref{lem:hub-window-recovery} with the constant in
\eqref{eq:hub-leak-fourier}, and then
Theorem~\ref{thm:hub-loglog-smallball}.
The recovery constant simplifies to
$(2\sqrt{2\pi}/(3\pi))4\sqrt2\pi^{3/2}=16\pi/3$.
\end{proof}

\begin{remark}\label{rem:hub-endpoint-open}
The estimates above hold for each positive $\varepsilon$.
Their constants are not bounded as $\varepsilon\downarrow0$.
Thus they do not establish a uniform three-eigenvalue theorem.
Even though the family at exactly zero leakage is bounded in several
important normal forms, continuity for fixed input laws does not supply
uniformity over all bounded product densities.
\end{remark}

\subsection{Localization and regularity of the hub law}

\begin{proposition}[A localized leakage direction]\label{prop:hub-localized}
Consider
\[
 Q=aX_0^2+2X_0(L+\beta)+\lambda V^2
                          +2\varepsilon LP+\sum_i\phi_i(X_i).
\]
Suppose the coefficient vector of $P$ is supported on $T$, let
$I=T^c$, and put
$G_I=\sum_{i\in I}(w_i,v_i)^T(w_i,v_i)$.
If $G_I$ is invertible and
$(w_i,v_i)G_I^{-1}(w_i,v_i)^T\le1/2$ on $I$, then
\[
 |\mathbb E e^{itQ}|\le
 4\pi^{3/2}(|\lambda|\det G_I)^{-1/2}|t|^{-3/2},
 \qquad
 \|p_Q\|_\infty\le
 3\,2^{4/3}(|\lambda|\det G_I)^{-1/3}.
\]
The estimate is uniform in $\varepsilon$, $a$, and all other displayed
coefficients.
\end{proposition}

\begin{proof}
For fixed $t$, integrating the hub and using $f_0\le1$ bounds the
characteristic function by $(2|t|)^{-1}\int|\widehat n(z)|\,dz$,
where $\widehat n(z)$ is the light expectation with the additional
phase $zL$.  Fresnel linearization of $\lambda V^2$ reduces its
integral to $(4\pi|t\lambda|)^{-1/2}J$, where
\[
 J=\iint\left|\mathbb E e^{izL+iuV+2it\varepsilon LP
                                     +it\sum_i\phi_i(X_i)}\right|dz\,du.
\]
Disintegrate the complex light measure over $P=p$.  On that slice the
cross term translates $z$ by $2t\varepsilon p$.  The triangle
inequality followed by this translation bounds $J$ by the averaged
conditional frame integral.  The coordinates in $I$ are independent
of those carrying $P$; the latter push forward to a complex measure of
total variation at most one.  Therefore
\[
 J\le\iint\prod_{i\in I}
        |\widehat g_i(zw_i+uv_i+tb_i)|\,dz\,du
 \le(4\pi)^2(\det G_I)^{-1/2}.
\]
The last inequality follows by the change of variables associated with
$G_I$ and Lemma~\ref{lem:hub-frame-integral}.
Multiplication of the three factors proves the characteristic-function
bound; its integral proves the density estimate.
\end{proof}

\begin{proposition}[A square on independent coordinates]
\label{prop:hub-free-square}
Let $I$ be a set of light coordinates on which $w$ vanishes and let
$X_R$ denote the others.  Suppose
\[
 Q=2X_0(L+\beta)+\lambda V^2+2Vg(X_R)+F(X_R),
 \qquad \omega=|v_I|>0,
\]
where $g,F$ are arbitrary measurable functions.  Then
\[
 \|p_Q\|_\infty\le\frac12+
             \frac{8\sqrt2}{\omega\sqrt{|\lambda|}}.
\]
\end{proposition}

\begin{proof}
Write $V=V_I+V_R$.  The variable $V_I$ is independent of $X_R$ and
has density at most $\sqrt2/\omega$ by
\eqref{eq:hub-projection}.  Given $X_R$, the remainder is
$\lambda(V_I+c_1)^2+c_2$.  For every Borel set $S$,
\[
 |\{u:\lambda u^2\in S\}|\le2\sqrt{|S|/|\lambda|},
\]
which follows by placing a set of fixed measure at the origin in the
one-sided density of the square.  Thus the conditional probability is
at most $2\sqrt2\omega^{-1}\sqrt{|S|/|\lambda|}$.
Since $L$ is a function of $X_R$ and $\rho_L\le\sqrt2$, its
unnormalized slices have profile
$4\omega^{-1}\sqrt{|S|/|\lambda|}$.
Lemma~\ref{lem:hub-profile} finishes the proof.
\end{proof}

\begin{proposition}[A block carrying only the hub partner]
\label{prop:hub-free-partner}
Assume $a=0$.  Suppose a set $I$ of light coordinates satisfies
$v_i=d_i=0$, $q_i=\kappa w_i$, and $b_i=\mu w_i$ on $I$.
Put $\omega=|w_I|>0$.  Then
\begin{align*}
 \|p_Q\|_\infty
 &\le\frac12+\frac{\sqrt2}{\omega}
 \left[\frac72\left(2+\log(1+4\pi^4/|\lambda|)\right)
                          +\frac12\log(1+4|\kappa|)\right].
\end{align*}
If $\kappa=0$, the sharper bound is
\[
 \|p_Q\|_\infty\le\frac12+\frac{\sqrt2}{\omega}
                        \left(2+\log(1+2\pi^4/|\lambda|)\right).
\]
\end{proposition}

\begin{proof}
Let $R=I^c$, $L_I=\langle w_I,X_I\rangle$ and $L_R=L-L_I$.
Collecting terms gives
$Q=2X_0t+2\kappa L^2+2LP+W$, where $t=L+\beta$ and $P,W$
depend only on $X_R$.  The value of the hub at level $s$ has the form
\[
 x_* =\frac{\alpha}{2t}+c_0-\kappa t,
 \quad
 \alpha=s-2\kappa\beta^2+2\beta P-W,
 \quad c_0=2\kappa\beta-P.
\]
Conditional on $X_R$, the variable $t$ has density at most
$\sqrt2/\omega$.  On $|t|\ge1$ its contribution to $p_Q$ is at most
$1/2$.  Thus the remaining contribution is bounded by
\[
 \frac{\sqrt2}{\omega}\mathbb E_{X_R}
 \int_{0<|t|<1}\widetilde f_0\!\left(\frac\alpha{2t}-\kappa t\right)
                                          \frac{dt}{2|t|},
 \qquad\widetilde f_0(x)=f_0(x+c_0).
\]
For $\kappa=0$, the substitution $x=\alpha/(2t)$ and the bathtub
principle bound the inner integral by $\log(1+1/|\alpha|)$.
For $\kappa\ne0$, set $\sigma=2\kappa$ and
$\tau=\min(1,\sqrt{|\alpha|/(2|\sigma|)})$.
For $0<|t|<\tau$, the map
$x(t)=\alpha/(2t)-\sigma t/2$ is monotone on each half-line,
its two images are disjoint, $|x(t)|\ge|\alpha|/(4\tau)$, and
$dt/|t|\le3\,|dx|/|x|$.
The bathtub principle bounds this part by
$3\log(1+2\tau/|\alpha|)$; the range $\tau\le|t|<1$ contributes at
most $\log(1/\tau)$.
Together these are at most
\[
 \frac72\log(1+2/|\alpha|)+\frac12\log(1+2|\sigma|).
\]
The variable $\alpha$ is a quadratic polynomial on $X_R$ consisting of
$-\lambda V^2$ and one-coordinate terms: the $L$-cross terms have
already been put into $P$.  One-dimensional Fresnel linearization and
Lemma~\ref{lem:hub-frame-integral} give
$|\mathbb E e^{it\alpha}|\le2\sqrt{\pi/|t\lambda|}$.
The kernel in Lemma~\ref{lem:hub-frozen} therefore gives
$\mathbb P(|\alpha-u|\le r)\le\pi^2\sqrt{2r/|\lambda|}$.
For $c=2\pi^2/\sqrt{|\lambda|}$, layer cake implies
\[
 \mathbb E\log(1+2/|\alpha|)
 \le\int_0^\infty\min\{1,c(e^u-1)^{-1/2}\}\,du
 \le\log(1+c^2)+2.
\]
Indeed, above $u=\log(1+c^2)$ the integral equals
$2c\arctan(1/c)<2$.
Replacing $2$ by $1$ in the logarithm gives
$2+\log(1+2\pi^4/|\lambda|)$, proving both assertions.
\end{proof}

\begin{proposition}[A hub density of bounded variation]
\label{prop:hub-bv}
Assume $a=0$ and let the total variation of $f_0$ be $V_0<\infty$.
For $A=4\sqrt{\pi/|\lambda|}$,
\begin{equation}\label{eq:hub-bv-fourier}
 |\mathbb E e^{itQ}|\le V_0|t|^{-3/2}
       \left[A\left(1+\log_+(2|t|^{3/2}/V_0)\right)+\frac12\right].
\end{equation}
In particular, for fixed $|\lambda|\ge1$,
$\|p_Q\|_\infty\le C V_0^{2/3}$ with an absolute constant $C$,
uniformly in every light diagonal, linear term, and cross term.
\end{proposition}

\begin{proof}
Integration by parts against the distributional derivative of $f_0$
gives $|\widehat f_0(z)|\le\min(1,V_0/|z|)$.
Freeze $L$ and apply \eqref{eq:hub-slice-fourier}.  With $u=L+\beta$,
\[
 |\mathbb E e^{itQ}|\le
 \int\min\{1,V_0/(2|tu|)\}
       \min\{\rho_L(u-\beta),A|t|^{-1/2}\}\,du.
\]
On $|u|\le T=|t|^{1/2}$, use the second entry in the second minimum.
The integral of the first minimum is at most
$(V_0/|t|)(1+\log_+(2|t|T/V_0))$.
On $|u|>T$, retain $\rho_L$ and use its unit mass, obtaining
$V_0/(2|t|^{3/2})$.  This proves
\eqref{eq:hub-bv-fourier}.
Substitute $t=V_0^{2/3}u$ into the integral of the minimum of one and
the right side.  The resulting integral is finite and uniformly
bounded for $A\le4\sqrt\pi$, proving the density estimate.
No numerical quadrature is needed for the existence of $C$.
\end{proof}

\subsection{Bounded light supports and a fixed transverse projection}

\begin{proposition}[A maximal-density criterion]
\label{prop:hub-maximal-projection}
Let $E$ be a fixed $k$-dimensional subspace orthogonal to $w,v$,
and write $Z$ for the orthonormal coordinates of the projection of
$X'$ onto $E$.  Let
\[
 Q=2X_0(L+\beta)+\lambda V^2+2Vg(L,Z)+F(L,Z),
\]
where $g,F$ are measurable.  Put
\[
 M_k=\mathop{\rm ess\,sup}_{\ell}
            \int_{\mathbb R^k}\mathop{\rm ess\,sup}_u
                       \rho_{L,V,Z}(\ell,u,z)\,dz.
\]
If $M_k<\infty$, then
\[
 \|p_Q\|_\infty\le\frac12+8M_k|\lambda|^{-1/2}.
\]
This result needs no leverage condition.
\end{proposition}

\begin{proof}
On the slice $L=\ell$ complete the square in $V$:
$\lambda(V-V_0(z))^2+F_0(z)$.
The slab $|V-V_0(z)|\le r$ has slice mass at most $2rM_k$.
Outside the slab, the two inverse branches each have derivative
at most $(2|\lambda|r)^{-1}$ in the output variable, so the
restricted output density is at most $M_k/(|\lambda|r)$.
For every Borel set $S$, choose
$r=\sqrt{|S|/(2|\lambda|)}$ to obtain the profile
$2\sqrt2M_k\sqrt{|S|/|\lambda|}$.
Lemma~\ref{lem:hub-profile} gives the result; freezing is unnecessary
here because the remainder is already a function of $(L,V,Z)$.
\end{proof}

\begin{proposition}[Bounded light supports]\label{prop:hub-light-intervals}
Suppose every light coordinate is supported in an interval of length
at most $K$.  Then the quantity in
Proposition~\ref{prop:hub-maximal-projection} satisfies
\[
 M_k\le2^{(k+2)/2}K^k
  \int_{\mathbb R^k}e^{-2(|y|-k-2)_+^2}\,dy.
\]
Moreover, for \eqref{eq:hub-normal} with $a=d_i=0$, arbitrary $q$
and arbitrary linear terms,
\begin{equation}\label{eq:hub-light-interval-bound}
 \|p_Q\|_\infty\le\frac12+
 8\sqrt2\left[2^{3/2}\left(6+\sqrt{\pi/2}\right)
                         \frac K{|\lambda|}\right]^{1/2}
 <\frac12+52\sqrt{K/|\lambda|}.
\end{equation}
The hub law is arbitrary subject to its density cap.
\end{proposition}

\begin{proof}
Tilt the product law by $e^{\langle\theta,Z\rangle}$.
Each tilted coordinate density is at most $e^{|\theta|K}$ times
its original cap, because the tilting factor varies by at most
that amount on its supporting interval.  The $(k+2)$-dimensional
projection-density bound for the tilted product law is therefore
$2^{(k+2)/2}e^{(k+2)|\theta|K}$.
Independence and the bounded-variable estimate of
\citet{Hoeffding1963} give
\[
 \mathbb E e^{\langle\theta,Z-\mathbb EZ\rangle}
                 \le e^{|\theta|^2K^2/8}.
\]
Undoing the tilt and optimizing $\theta$ in the direction
$z-\mathbb EZ$ yields
\[
 \rho_{L,V,Z}(\ell,u,z)
 \le2^{(k+2)/2}
  \exp\left[-2\frac{(|z-\mathbb EZ|-(k+2)K)_+^2}{K^2}\right].
\]
The optimization uses
$|\theta|=4(|z-\mathbb EZ|-(k+2)K)_+/K^2$.
Integration over $z$ proves the first assertion.

For the second assertion freeze $L$ as in
Remark~\ref{rem:hub-freezing}.  With no light diagonal, every frozen
remainder has the form
$\lambda(V-V_0)^2+\gamma B+c$, where $B$ is a unit linear
projection orthogonal to $w,v$; its direction may depend on the
frozen parameter.  The preceding tilt bound is uniform in that
direction and gives
\[
 \int\sup_u\rho_{L,V,B}(\ell,u,b)\,db
 \le M:=2^{3/2}\left(6+\sqrt{\pi/2}\right)K.
\]
The slab $|V-V_0|\le r$ has slice mass at most $4r$, using
$\rho_{L,V}\le2$.  Off the slab the output density is at most
$M/(|\lambda|r)$.  Optimizing
$4r+M|S|/(|\lambda|r)$ gives the Borel profile
$4\sqrt{M|S|/|\lambda|}$.
Lemma~\ref{lem:hub-profile} proves
\eqref{eq:hub-light-interval-bound}.  If $\gamma=0$, the direct
square profile is at least as strong, since $K\ge1$.
\end{proof}

\begin{corollary}[Window-mixture light laws]\label{cor:hub-window-lights}
Suppose each light law is of the form $\mathcal U[0,1]*\nu_i$.
For \eqref{eq:hub-normal} with $a=d_i=0$, every independent cap-one
hub law satisfies
\[
 \|p_Q\|_\infty\le\frac12+52|\lambda|^{-1/2}.
\]
\end{corollary}

\begin{proof}
Represent $X_i=U_i+Z_i$, with independent uniform $U_i$ on $[0,1]$,
and condition on the shifts $Z_i$.  The light coordinates are now
supported on unit intervals.  The shifts change only the linear
coefficients and the constant of the quadratic form, which are
unrestricted in Proposition~\ref{prop:hub-light-intervals}.
Apply that proposition and average.
\end{proof}

\subsection{Several hub coordinates}

\begin{theorem}[Three or more hubs]\label{thm:hub-many}
Let $X_H\in\mathbb R^k$ and $X'$ be independent, with
$\int f_H^2\le1$, and suppose $L=MX'\in\mathbb R^k$ has density
bounded by $D$.  For any real measurable $F,R$ and $\beta\in\mathbb R^k$,
\[
 Q=F(X_H)+2\langle X_H,L+\beta\rangle+R(X')
\]
satisfies
\[
 |\mathbb E e^{itQ}|\le
       \min\{1,D^{1/2}(\pi/|t|)^{k/2}\}.
\]
If $k>2$, then
\begin{equation}\label{eq:hub-many-density}
 \|p_Q\|_\infty\le\frac{k}{k-2}D^{1/k}.
\end{equation}
For independent cap-one light coordinates and singular values
$s_1,\ldots,s_k>0$ of $M$, this is at most
$\frac{k}{k-2}\sqrt2(s_1\cdots s_k)^{-1/k}$.
\end{theorem}

\begin{proof}
Let
$h_t(\ell)=\int f_H(x)e^{itF(x)+2it\langle x,\ell+\beta\rangle}\,dx$.
Plancherel gives
\[
 \int_{\mathbb R^k}|h_t(\ell)|^2\,d\ell
   =(\pi/|t|)^k\int f_H^2\le(\pi/|t|)^k.
\]
After conditioning on $X'$,
\[
 |\mathbb E e^{itQ}|\le\mathbb E|h_t(L)|
 \le\|\rho_L\|_2\|h_t\|_2
 \le D^{1/2}(\pi/|t|)^{k/2}.
\]
Integrating the minimum with one gives
\eqref{eq:hub-many-density}.  The final assertion follows from
\eqref{eq:hub-projection}.
\end{proof}

\begin{theorem}[Two hubs and a transverse square]\label{thm:hub-two-square}
Let $X_H\in\mathbb R^2$ have density $f_H\le1$ and be independent
of the cap-one product vector $X'$.  Let $M$ have rank two, singular
values $s_1,s_2$, and orthonormal row-space basis $e_1,e_2$.
Suppose $v\perp e_1,e_2$, $|v|=1$, and
$e_{1i}^2+e_{2i}^2+v_i^2\le1/2$.
Put $L=MX'$ and $V=\langle v,X'\rangle$.  For every real symmetric
$2\times2$ matrix $B$, vectors $q_1,q_2$, and coordinate quadratics,
\[
 Q=X_H^TBX_H+2\langle X_H,L+\beta\rangle+\lambda V^2
                   +2\sum_{j=1}^2L_j\langle q_j,X'\rangle
                   +\sum_i\phi_i(X_i)
\]
has
\[
 \|p_Q\|_\infty\le
 5\,2^{6/5}(s_1s_2)^{-2/5}|\lambda|^{-1/5}
 <11.5(s_1s_2)^{-2/5}|\lambda|^{-1/5}.
\]
\end{theorem}

\begin{proof}
Freeze the two partners on rectangular cells.  Plancherel gives
\[
 \int|h_t(\ell)|^2\,d^2\ell\le\pi^2/t^2
\]
for the hub factor.
Fresnel linearization of the transverse square and the
three-dimensional frame bound give for the frozen slice
\[
 |\widehat\mu_{\ell_0,\ell}(t)|
 \le\min\left\{\rho_L(\ell),
            \frac{8\sqrt\pi}{s_1s_2\sqrt{|t\lambda|}}\right\}.
\]
The factor $(s_1s_2)^{-1}$ is the Jacobian converting the frequency
of $L$ to that of its orthonormal row-space coordinates.
For each fixed $t$, Cauchy--Schwarz on the cells and passage to the
Riemann-sum limit give
\[
 |\mathbb E e^{itQ}|\le K|t|^{-5/4},
 \quad K=\pi\left(\frac{8\sqrt\pi}{s_1s_2}\right)^{1/2}
                              |\lambda|^{-1/4}.
\]
Indeed, the squared slice factor integrates to at most its displayed
uniform bound times $\int\rho_L=1$.
The characteristic function is also at most one.  Its integral is
$10K^{4/5}$, and division by $2\pi$ gives the stated constant.
\end{proof}

\begin{theorem}[Two hubs with an interval slice profile]
\label{thm:hub-two-profile}
Let $X_H\in\mathbb R^2$ have density at most one, be independent of
$X'$, and let $L=MX'$ have a density.  Let $Y_\ell$ be a family of
real functions of $X'$ such that $Y_{\ell'}\to Y_\ell$ pointwise as
$\ell'\to\ell$.  Suppose its frozen slices satisfy
\[
 \mu_{\ell_0,\ell}(I)\le K|I|^b
 \quad(0<b\le1)
\]
for every fixed $\ell_0$, almost every $\ell$, and every interval $I$.
Then, for every real symmetric $B$ and $\beta\in\mathbb R^2$,
\[
 Q=X_H^TBX_H+2\langle X_H,L+\beta\rangle+Y_L
\]
has
\[
 \|p_Q\|_\infty\le\frac1\pi+
             \frac{\sqrt{10K}}{1-2^{-b/2}}.
\]
\end{theorem}

\begin{proof}
First restrict the light measure to $\{\rho_L(L)\le N\}$.
The resulting subprobability measure has bounded partner density, and
all the assumed slice inequalities persist with the same constants.
The bounds below are independent of $N$; monotone convergence at the
end removes the restriction.  Use rectangular cells $C_j$ and frozen
values $\ell_j$.  Define the
finite positive measure
\[
 \nu_h=\sum_j\int_{C_j}
          \mu_{\ell_j,\ell}*\widetilde\mu_{\ell_j,\ell}\,d^2\ell.
\]
It has nonnegative Fourier transform
$\widehat\nu_h(t)=\sum_j\int_{C_j}|\widehat\mu_{\ell_j,\ell}(t)|^2
d^2\ell$ and satisfies
$\nu_h([-r,r])\le K(2r)^b$, since each slice has mass
$\rho_L(\ell)$ and $\int\rho_L=1$.

For any positive measure $\nu$ with nonnegative transform and
$\nu([-r,r])\le K'r^b$, the Fej\'er kernel gives
\[
 \int_{-T}^T\widehat\nu(t)\,dt
 \le4T\int\left(\frac{\sin(Tx)}{Tx}\right)^2\nu(dx)
 \le20K'T^{1-b}.
\]
For the last inequality split at $|x|=1/T$ and into the dyadic shells
$2^j/T<|x|\le2^{j+1}/T$; use
$\min(1,(Tx)^{-2})$ and
$1+2^b/(1-2^{b-2})\le5$.

Let $\Omega_A=\{A\le|t|\le2A\}$.  Cauchy--Schwarz in the partners
and then in $t$, followed by Fatou and the freezing limit, give
\[
 \int_{\Omega_A}|\mathbb E e^{itQ}|\,dt
 \le \left(\int_{\Omega_A}\int|h_t(\ell)|^2d^2\ell\,dt\right)^{1/2}
       \liminf_{h\downarrow0}
         \left(\int_{\Omega_A}\widehat\nu_h(t)dt\right)^{1/2}.
\]
This can first be done on finitely many cells and then exhausted; all
slice-mass estimates remain unchanged.  The first factor is at most
$\pi A^{-1/2}$ by hub Plancherel, and the second is at most
$(20\cdot2^bK(2A)^{1-b})^{1/2}$.
Thus the integral is at most $\pi\sqrt{40K}A^{-b/2}$.
Sum over $A=2^j$, $j\ge0$, add the bound two on $|t|<1$, and divide
by $2\pi$.
\end{proof}

\begin{proposition}[A projected crossing remainder]
\label{prop:hub-crossing}
Let $k=1$ or $2$, let $X_H\in\mathbb R^k$ have density at most one,
and let $X'$ be an independent cap-one product vector.  Consider
\[
 Q=X_H^TBX_H+2\langle X_H,MX'+\beta\rangle+R(X'),
 \qquad R(x)=x^TA_Rx+b^Tx+c.
\]
Split the light coordinates as $I\cup J$.  Suppose $M_J$ has rank
$k$, with product $S$ of its singular values, and put
\[
 C=(A_R)_{I,J}(I-\Pi),\qquad F=\|C\|_F>0,
\]
where $\Pi$ projects onto the row space of $M_J$.
If every row $c_i$ of $C$ satisfies $|c_i|^2\le F^2/2$, then
\begin{equation}\label{eq:hub-crossing-fourier}
 |\mathbb E e^{itQ}|\le
 2^{(k+3)/4}\pi^{(k+1)/2}(SF)^{-1/2}|t|^{-(k+1)/2}.
\end{equation}
For $k=2$ this gives
\[
 \|p_Q\|_\infty\le3\,2^{5/6}(SF)^{-1/3}<5.35(SF)^{-1/3}.
\]
For $k=1$, $S=|w_J|$, it gives
\[
 \mathbb P(|Q-s|\le\delta)\le
 \frac{\pi^2\delta}{\sqrt{SF}}
 \left(1+\log_+\frac{\sqrt{SF}}{4\delta}\right).
\]
\end{proposition}

\begin{proof}
Let $\mu_t(\ell)$ be the density of
$\mathbb E[e^{itR(X')};MX'\in d\ell]$.
The hub factor and Plancherel imply
\[
 |\mathbb E e^{itQ}|^2\le(\pi/|t|)^k\int|\mu_t(\ell)|^2d\ell.
\]
Convolve $\mu_t$ with an even nonnegative normalized approximate
identity $k_h$.  Conditional Cauchy--Schwarz in $X_I$, an independent
copy $X_J^*$, and integration over $\ell$ yield
\[
 \int|\mu_t*k_h|^2\le
 \mathbb E\left[\kappa_h(M_JD_J)
       \prod_{i\in I}|\widehat f_i(2t\langle a_i,D_J\rangle)|\right],
\]
where $D_J=X_J-X_J^*$, $\kappa_h=k_h*k_h$, and $a_i$ is the
$i$th row of $(A_R)_{I,J}$.  Terms of $R$ internal to $I$ cancel;
terms internal to $J$ contribute only a phase of modulus one.

Put $\theta_i=|c_i|^2/F^2$ and discard the zero rows, whose factors
are at most one.  H\"older with weights $\theta_i$ bounds the last
expectation by
\[
 \prod_i\left(\mathbb E\left[\kappa_h(M_JD_J)
 |\widehat f_i(2t\langle a_i,D_J\rangle)|^{1/\theta_i}\right]
                                                   \right)^{\theta_i}.
\]
The coordinates of $D_J$ are independent with densities bounded by
one.  Since $c_i\perp\operatorname{row}(M_J)$, the density of
$(M_JD_J,\langle c_i,D_J\rangle)$ is at most
$2^{(k+1)/2}/(S|c_i|)$.
Write $a_i=c_i+\alpha_iM_J$ and translate the last variable for each
value of $M_JD_J$.  Applying \eqref{eq:hub-scalar-lp} gives a bound
\[
 \frac{2^{(k+1)/2}}{S|c_i|}\frac{4\pi\sqrt{\theta_i}}{2|t|}
 =\frac{2^{(k+3)/2}\pi}{SF|t|}
\]
for every factor.  It is independent of $h$.
Pass to the $L^2$ approximate-identity limit and combine with the
hub factor to obtain \eqref{eq:hub-crossing-fourier}.
For $k=2$, integrate its minimum with one.
For $k=1$, put $K_1=2\pi/\sqrt{SF}$ and apply the kernel used in
Lemma~\ref{lem:hub-frozen} to
$\min(1,K_1/|t|)$.  The integral over
$|t|\le\pi/(2\delta)$ is at most
$2K_1(1+\log_+(\pi/(2K_1\delta)))$, which gives the assertion.
\end{proof}

\begin{remark}
The projected-crossing estimate permits arbitrary terms internal to
either light block.  Its extension to several volume directions and
arbitrary-rank crossing matrices is given in
Theorem~\ref{thm:grouped-crossing}; exact rank of the crossing matrix
is not needed there.  The nondegeneracy and volume-marginal assumptions
remain essential hypotheses of that theorem.
\end{remark}

\begin{corollary}[A logarithmic bound with a hub diagonal]
\label{cor:hub-diagonal-log}
Suppose Theorem~\ref{thm:hub-leak} and the $k=1$ case of
Proposition~\ref{prop:hub-crossing} both apply.  Then
\[
 \|p_Q\|_\infty\le
 \frac{2}{\sqrt{SF}}
 \left(3+\log_+\frac{4(SF)^{3/2}}{\varepsilon|\lambda|}\right).
\]
\end{corollary}

\begin{proof}
The characteristic function is bounded by
$\min(1,K_1|t|^{-1},K_2|t|^{-3/2})$, where
$K_1=2\pi/\sqrt{SF}$ and
$K_2=4\sqrt2\pi^{3/2}/\sqrt{\varepsilon|\lambda|}$.
Splitting the positive half-line at $K_1$ and $(K_2/K_1)^2$ gives
an integral at most
$K_1[3+\log_+(K_2^2/K_1^3)]$.
The same bound holds when the latter split precedes the former by
using just the first and third entries of the minimum.
Since $K_2^2/K_1^3=4(SF)^{3/2}/(\varepsilon|\lambda|)$,
division of the two-sided integral by $2\pi$ proves the claim.
\end{proof}

\begin{proposition}[A spread linear hub]\label{prop:hub-spread}
Let $X_H$ and $X'$ be independent cap-one product vectors, let
$U=\langle c,X_H\rangle$, and assume
$c_h^2\le|c|^2/2$ for every hub coordinate.  For
\[
 Q=2U(L+\beta)+\sum_{h\in H}\phi_h(X_h)
       +\lambda V^2+2L\langle q,X'\rangle+\sum_i\phi_i(X_i),
\]
with \eqref{eq:hub-frame},
\[
 |\mathbb E e^{itQ}|\le
 8\pi^{3/2}|c|^{-1}|\lambda|^{-1/2}|t|^{-3/2},
 \qquad
 \|p_Q\|_\infty\le12(|c|^2|\lambda|)^{-1/3}.
\]
\end{proposition}

\begin{proof}
The conditional hub expectation is
\[
 H_t(u)=\prod_{h\in H}
  \widehat g_h(2tuc_h+tb_h),\qquad
 g_h(x)=f_h(x)e^{itd_hx^2}.
\]
H\"older with exponents $p_h=|c|^2/c_h^2\ge2$ and
\eqref{eq:hub-scalar-lp} gives
$\int|H_t(u)|\,du\le2\pi/(|t||c|)$.
Freeze $L$ and pair this factor with the uniform slice-transform bound
$4\sqrt\pi/\sqrt{|t\lambda|}$ from
Lemma~\ref{lem:hub-frozen}.  Their product is the claimed
characteristic-function bound.  Integrating the minimum with one gives
the density bound.
\end{proof}

\subsection{An adapted-hub variational formulation}

The entangled case admits a variational formulation in terms of a
restricted conditional law.  Its dependence structure is absent from
arbitrary bounded two-dimensional densities.

\begin{proposition}[Adapted-hub variational reduction]
\label{prop:hub-adapted}
Let $a=0$ and $q=q_0+\varepsilon e$, where
$q_0\in\operatorname{span}\{w,v\}$ and $P=\langle e,X'\rangle$.
Absorb $-2\varepsilon\beta P$ into the light linear terms and write
\[
 Q=2Y(L+\beta)+R_0(X'),\qquad Y=X_0+\varepsilon P.
\]
Then the conditional density of $Y$ given $X'$ is at most one and
depends on $X'$ only through $P$.
For a fixed level $s$, put
$x_*=(s-R_0(X'))/(2(L+\beta))$.
If $W_s(x,p)$ is the joint density of $(x_*,P)$ under the
nonnegative weight $1/(2|L+\beta|)$, then the supremum at $s$ over
all conditional hub densities $f(y\mid p)\le1$ is
\begin{equation}\label{eq:hub-adapted-functional}
 \Phi(s)=\int_{\mathbb R}
       \sup_{|A|=1}\int_A W_s(x,p)\,dx\,dp.
\end{equation}
In particular, $p_Q(s)\le\Phi(s)$ for the independent-hub form.
The same assertion can be formulated for window probabilities, without
assuming that the weighted pushforward has a density.
\end{proposition}

\begin{proof}
The identity
$2X_0(L+\beta)+2\varepsilon LP
 =2(L+\beta)(X_0+\varepsilon P)-2\varepsilon\beta P$
proves the first assertion.  Its conditional density is
$f_0(y-\varepsilon P)$.
Given $X'$, the map from $Y$ to $Q$ is affine with derivative
$2(L+\beta)$; hence
\[
 p_Q(s)=\iint f(x\mid p)W_s(x,p)\,dx\,dp.
\]
For each $p$, maximizing a nonnegative linear functional subject to
$0\le f(\cdot\mid p)\le1$ and $\int f(\cdot\mid p)=1$
gives the integral over a set of measure one on which $W_s$ is largest.
Measurable approximate maximizers can be chosen by thresholding
$W_s$ and filling a level set; monotone approximation then gives
\eqref{eq:hub-adapted-functional}.  The window version follows by
applying the same conditional maximization to the nonnegative
conditional window kernel before taking the density limit.
\end{proof}

\begin{conjecture}[A sufficient adapted-hub estimate]
\label{conj:hub-adapted}
Let $X'$ have independent coordinate densities bounded by one, and let
$w,v,e$ be orthonormal with $w_i^2+v_i^2+e_i^2\le1/2$.
Set $L=\langle w,X'\rangle$, $V=\langle v,X'\rangle$, and
$P=\langle e,X'\rangle$.
If the conditional density of $Y$ given $X'$ is $g(y\mid P)\le1$
of unit mass, then $2YL+V^2$ has density bounded by an absolute constant.
Equivalently, its functional \eqref{eq:hub-adapted-functional} is
bounded uniformly in the level.
\end{conjecture}

The conjecture is stronger than the estimate for a translated
independent hub, because it permits an arbitrary probability density
depending on $P$.  This conditional estimate remains open.  The
preceding estimates also do not cover arbitrary quadratic twists on
the orthogonal complement of $\operatorname{span}\{w,v\}$ or yield an
exhaustive quantitative reduction of all three-eigenvalue matrices
to the stated hub configurations.

\subsection{Adapted hubs for Gaussian mixtures}

The adapted-hub estimate can be proved for a class allowing arbitrarily
small Gaussian component variances.  The resulting theorem concerns the
entangled core $2YL+\lambda V^2$; it does not assert an estimate for
arbitrary light diagonal terms.

\begin{lemma}[A conditional Gaussian bound]\label{lem:hub-gaussian}
Suppose $(L,V,P)$ is a Gaussian vector with covariance matrix at least
$hI_3$, where $h>0$.  Let the conditional law of $Y$ given $(L,V,P)$
have density $g(y\mid P)\le1$ and unit integral in $y$.
For $|\lambda|\ge1$ and arbitrary finite real measurable $B,C$, set
\[
 Q=2YL+\lambda V^2+B(Y,P)V+C(Y,P).
\]
Then
\begin{equation}\label{eq:hub-gaussian-bound}
 \|p_Q\|_\infty\le4(2\pi)^{-3/4}h^{-3/4}.
\end{equation}
If $h\ge(2\pi)^{-1}$, the sharper numerical bound is
$\|p_Q\|_\infty\le3$.
The means of the Gaussian vector are arbitrary.
\end{lemma}

\begin{proof}
Condition on $P$.  The conditional covariance of $(L,V)$ is at least
$hI_2$.  Indeed, for the full covariance matrix $M\ge hI_3$, its
Schur complement satisfies
\[
 a^TK a=\inf_b(a,b)^TM(a,b)\ge h|a|^2.
\]
Regress $L$ on $V$: $L=\mu+\alpha V+W$, where $W$ is an independent
Gaussian with variance at least $h$, and $V$ also has variance at
least $h$.  Given $P$, the variable $Y$ is independent of these
Gaussian residuals.  For fixed $Y=y$, the $W$ factor and the
one-dimensional quadratic Gaussian integral imply
\[
 |\mathbb E[e^{itQ}\mid P,Y=y]|
 \le e^{-2hy^2t^2}(1+4h^2t^2)^{-1/4}.
\]
The added terms change only a linear phase in $V$ and a constant.
Completing the Gaussian square shows that the modulus contributed by
the linear phase is at most one, for every conditional mean.
Using both $g\le1$ and $\int g=1$ gives
\begin{equation}\label{eq:hub-gaussian-envelope}
 |\mathbb E e^{itQ}|\le
 (1+4h^2t^2)^{-1/4}
       \min\left\{1,\frac{\sqrt\pi}{\sqrt{2h}|t|}\right\}.
\end{equation}
This integrable bound is independent of $P$.
Set $a=\sqrt{2\pi h}$ and $u=2h|t|$.  Fourier inversion gives
\[
 \|p_Q\|_\infty\le a^{-2}\int_0^\infty
             (1+u^2)^{-1/4}\min(1,a/u)\,du.
\]
Bound $(1+u^2)^{-1/4}$ by $u^{-1/2}$ and split at $u=a$.
The integral is at most $4\sqrt a$, proving
\eqref{eq:hub-gaussian-bound}.
For $h=(2\pi)^{-1}$, $a=1$; splitting at one and using one below
that point and $u^{-1/2}$ above it gives the sharper value $1+2=3$.
The envelope decreases with $h$, so this holds for larger $h$ too.
\end{proof}

\begin{corollary}[A Gaussian covariance floor]\label{cor:hub-gaussian-floor}
Let $X'=Z+G$, where $Z$ is an arbitrary random vector and $G$ is an
independent Gaussian vector with covariance at least $(2\pi)^{-1}I$.
Let $L,V,P$ be orthonormal linear projections and let $Y$ have a
conditional density at most one depending on $X'$ only through $P$.
Then every $Q$ in Lemma~\ref{lem:hub-gaussian} has density at most
three.  No spread hypothesis on the frame is needed.
\end{corollary}

\begin{proof}
Generate the conditional hub after $X'$, using independent auxiliary
randomness.  Conditional on $Z$, the vector $(L,V,P)$ is Gaussian with
the required covariance floor and with arbitrary mean.  Its hub kernel
is still $g(y\mid P)$.  Apply Lemma~\ref{lem:hub-gaussian} and average.
\end{proof}

\begin{lemma}[A uniform partition of fixed-dimensional frames]
\label{lem:hub-frame-partition}
For each positive integer $d$ there is $\eta_d>0$ such that every
finite family $r_i\in\mathbb R^d$ satisfying
\[
 \sum_i r_i r_i^T=I_d,\qquad |r_i|^2\le\tfrac12
\]
can be partitioned into $A,B$ with
\begin{equation}\label{eq:hub-frame-partition}
 \sum_{i\in A}r_i r_i^T\ge\eta_d I_d,
 \qquad
 \sum_{i\in B}r_i r_i^T\ge\eta_d I_d.
\end{equation}
The constant does not depend on the number of vectors.
\end{lemma}

\begin{proof}
We first recall the finite vector-matroid base-packing criterion:
a family of rank $d$ contains two disjoint bases if and only if
$|E\setminus F|\ge2(d-\operatorname{rank}F)$ for every subfamily $F$
\citep{Edmonds1965}.  In a Parseval frame as above, the energy
orthogonal to a subspace of codimension $k$ is $k$.
Since a single vector contributes at most $1/2$, at least $2k$
vectors lie outside that subspace.  Thus two disjoint bases exist
for each fixed frame.

Suppose the uniform assertion were false.  Choose frames whose best
lower bound in \eqref{eq:hub-frame-partition} tends to zero.
Sort each frame in decreasing order of vector norm, pad with zero
vectors, and take a diagonal subsequence so that its $i$th vector
converges to $r_i$ for each fixed $i$.
Because the total squared norm is $d$, every tail vector beyond
index $N$ has squared norm at most $d/(N+1)$.
The limiting atoms satisfy $\sum_i r_i r_i^T\le I_d$.
Put
\[
 S=I_d-\sum_i r_i r_i^T\ge0,\qquad K=\ker S.
\]
The projections of the limiting atoms onto $K$ form a Parseval frame
for $K$, each with squared norm at most $1/2$.
They contain two disjoint finite bases of $K$.
To justify this for the possibly infinite family, choose a finite
initial segment whose covariance on $K$ is at least
$(1-\epsilon)I_K$, with $\epsilon<1/(2d)$.
Every codimension-$k$ subspace of $K$ then has outside energy greater
than $k-1/2$, so the number of outside vectors is at least $2k$.
The finite base-packing criterion applies to this initial segment.

Let $A_0,B_0$ be the indices of these disjoint bases, taking both
empty if $K=\{0\}$.  The matrices
\[
 H_A=\sum_{i\in A_0}r_i r_i^T+S/2,
 \qquad
 H_B=\sum_{i\in B_0}r_i r_i^T+S/2
\]
are positive definite.  A null vector would belong to $K$ and be
orthogonal to a basis of $K$.  Let $\gamma>0$ be the smaller of their
least eigenvalues.  Choose $N$ containing both bases and large enough
that $d/(2\sqrt{N+1})<\gamma/4$.

In each approximating frame assign the indicated basis vectors to
their prescribed sides, and assign the other first $N$ vectors
arbitrarily.  Partition the tail independently by fair coin flips.
If $T$ is the tail covariance and $T_A$ its contribution to side $A$,
\[
 \mathbb E\|T_A-T/2\|_F^2
 =\tfrac14\sum_{i>N}|r_i^{\rm app}|^4
 \le\frac{d^2}{4(N+1)}.
\]
There is consequently a tail partition with operator error at most
$d/(2\sqrt{N+1})$ on both sides.  For fixed $N$, the tail covariance
converges to $I_d-\sum_{i\le N}r_i r_i^T\ge S$.
The prescribed basis covariances converge as well.  Thus sufficiently
late approximating frames have both partition covariances at least
$\gamma I_d/2$, a contradiction.
\end{proof}

\begin{theorem}[An adapted hub for Gaussian component mixtures]
\label{thm:hub-gaussian-mixtures}
Let
\[
 X_i=M_i+\sqrt{T_i}\,G_i,\qquad T_i>0,
\]
where the pairs $(M_i,T_i)$ are independent across coordinates and
the standard Gaussian variables $G_i$ are independent of all these
pairs.  Suppose
\begin{equation}\label{eq:hub-mixture-peak}
 \mathbb E T_i^{-1/2}\le\sqrt{2\pi}\qquad\text{for every }i.
\end{equation}
Let $w,v,e$ be orthonormal with
$w_i^2+v_i^2+e_i^2\le1/2$, put
$(L,V,P)=(\langle w,X'\rangle,\langle v,X'\rangle,
\langle e,X'\rangle)$, and let $Y$ have conditional density
$g(y\mid P)\le1$ of unit mass.
For $|\lambda|\ge1$ and arbitrary finite real measurable $B,C$,
\[
 Q=2YL+\lambda V^2+B(Y,P)V+C(Y,P)
\]
has density bounded by
\begin{equation}\label{eq:hub-mixture-bound}
 4\,2^{-3/4}3^{3/2}\eta_3^{-3/4},
\end{equation}
where $\eta_3$ is the constant in
Lemma~\ref{lem:hub-frame-partition}.
In particular this gives a uniform bound for
\[
 aX_0^2+2X_0(L+\beta)+\lambda V^2+2\varepsilon LP
\]
for every independent cap-one hub law, every $a,\beta,\varepsilon$,
and every light law satisfying \eqref{eq:hub-mixture-peak}.
\end{theorem}

\begin{proof}
Put $r_i=(w_i,v_i,e_i)^T$ and choose a partition $A,B$ from
Lemma~\ref{lem:hub-frame-partition}.  Its two covariance sums satisfy
$S_A\ge\eta_3I_3$ and $S_B\ge\eta_3I_3$.
Realize $Y$ after $X'$ using fresh auxiliary randomness and the
prescribed kernel $g(y\mid P)$.  This preserves the joint law relevant
to $Q$ and ensures that the same kernel remains valid after
conditioning on the latent mixture variables.
Conditional on the mixture variables, the covariance of $(L,V,P)$ is
\[
 M=\sum_iT_i r_i r_i^T=M_A+M_B.
\]
If $h=\lambda_{\min}(M)$, Lemma~\ref{lem:hub-gaussian} bounds the
conditional density by $4(2\pi)^{-3/4}h^{-3/4}$, uniformly in all
conditional means and in $P$.  It remains to average this bound.

For side $A$, set $R_A=(S_A^{-1/2}r_i)_{i\in A}$, so that
$R_AR_A^T=I_3$.  For any positive diagonal matrix $D$, compression
satisfies
\begin{equation}\label{eq:hub-compression}
 (R_ADR_A^T)^{-1/2}\le R_AD^{-1/2}R_A^T.
\end{equation}
Indeed the corresponding resolvent inequality follows from a Schur
complement, and integrating it against
$\pi^{-1}s^{-1/2}ds$ proves \eqref{eq:hub-compression}, using
$x^{-1/2}=\pi^{-1}\int_0^\infty s^{-1/2}(x+s)^{-1}\,ds$.
Since
$M_A=S_A^{1/2}[R_A\operatorname{diag}(T_i)R_A^T]S_A^{1/2}$,
taking traces and using \eqref{eq:hub-mixture-peak} gives
\[
 \mathbb E\lambda_{\min}(M_A)^{-1/2}
 \le\eta_3^{-1/2}\sum_{i\in A}
            |S_A^{-1/2}r_i|^2\mathbb ET_i^{-1/2}
 \le3\sqrt{2\pi}\eta_3^{-1/2}.
\]
The same bound holds for $M_B$.
The positive random variables
$a_0=\lambda_{\min}(M_A)$ and $b_0=\lambda_{\min}(M_B)$ are
independent, and $h\ge a_0+b_0\ge2\sqrt{a_0b_0}$.
Therefore, by independence and the monotonicity of positive moments,
\[
 \mathbb Eh^{-3/4}
 \le2^{-3/4}(\mathbb Ea_0^{-1/2})^{3/4}
                 (\mathbb Eb_0^{-1/2})^{3/4}
 \le2^{-3/4}(3\sqrt{2\pi}\eta_3^{-1/2})^{3/2}.
\]
Combining the two inequalities yields
\eqref{eq:hub-mixture-bound}.

For the final assertion, set $Y=X_0+\varepsilon P$.
Then its conditional density is $f_0(y-\varepsilon P)\le1$, and
\[
 aX_0^2+2X_0(L+\beta)+\lambda V^2+2\varepsilon LP
 =2YL+\lambda V^2+a(Y-\varepsilon P)^2
                         +2\beta(Y-\varepsilon P),
\]
which is in the stated class with $B=0$.
\end{proof}

\begin{remark}[Scope of the mixture condition]
Condition \eqref{eq:hub-mixture-peak} bounds the average peak of the
Gaussian components and implies that each marginal density is at most
one.  For centered Gaussian scale mixtures, meaning that each $M_i$
is deterministic, it is equivalent to the cap-one assumption: the
maximum is attained at $M_i$ and equals
$(2\pi)^{-1/2}\mathbb ET_i^{-1/2}$.
Thus Theorem~\ref{thm:hub-gaussian-mixtures} covers every cap-one
centered Gaussian scale mixture, even when its component variances
have no positive lower bound.  A mixture of many narrow components
separated in location can have a bounded final density while violating
\eqref{eq:hub-mixture-peak}; such laws are not covered.
\end{remark}

\begin{corollary}[Laplace light coordinates]\label{cor:hub-laplace}
Let the light coordinates have independent translated Laplace
densities $e^{-2|x-m_i|}$.  For the adapted forms of
Theorem~\ref{thm:hub-gaussian-mixtures}, no spread condition is needed,
and
\[
 \|p_Q\|_\infty\le12\pi^{-3/4}\Gamma(1/4)<19.
\]
\end{corollary}

\begin{proof}
The Laplace law has the representation
$m_i+\sqrt{T_i}G_i$ with $T_i$ exponential of rate two:
its characteristic function is
$\mathbb E e^{-t^2T_i/2}=4/(4+t^2)$.
For the covariance $M=\sum_iT_i r_i r_i^T$, the same resolvent
compression argument, now with exponent $3/4$, gives
\[
 \mathbb E\lambda_{\min}(M)^{-3/4}
 \le\mathbb E\operatorname{tr}M^{-3/4}
 \le\sum_i|r_i|^2\mathbb ET_i^{-3/4}
 =3\,2^{3/4}\Gamma(1/4).
\]
Average \eqref{eq:hub-gaussian-bound}.  No partition of the frame is
needed because the full negative $3/4$ moment exists here.
\end{proof}

\subsection{A limitation of density-cap arguments for adapted hubs}

\begin{proposition}[Conditional information is not controlled by marginal caps]
\label{prop:hub-cap-obstruction}
For every sufficiently large integer $m$, there is a joint law
$(L,V,P)$ whose joint density and all its coordinate marginal
densities are bounded by one, and a conditional hub density
$g(y\mid P)\le1$ of unit mass, such that
\[
 \|p_{2YL+V^2}\|_\infty\ge\frac14\log m-\frac12\log9.
\]
Here coordinate marginals include all one- and two-coordinate
projections of this three-dimensional vector.  The constructed joint
law is not asserted to be an orthogonal projection of a product law.
\end{proposition}

\begin{proof}
Let $\delta=1/m$.  Partition $[0,1]\times[1,2]$ into squares
\[
 C_{jk}=[j\delta,(j+1)\delta]\times
                    [1+k\delta,1+(k+1)\delta],\qquad0\le j,k<m.
\]
Assign to each square a distinct disjoint unit interval $I_{jk}$
on the $P$ axis.  Put joint density one on
$\bigcup_{j,k}C_{jk}\times I_{jk}$ and zero elsewhere.
Its total integral is $m^2\delta^2=1$.
The $(L,V)$ marginal is uniform on the rectangle; the $(L,P)$ and
$(V,P)$ densities are $\delta$ on their supports; and the $P$ density
is $\delta^2$.  The one-dimensional $L,V$ densities are one.

Write $(\ell_j,v_k)$ for the center of $C_{jk}$.
When $P\in I_{jk}$, let $Y$ be uniform on the unit interval centered
at $-v_k^2/(2\ell_j)$.  This kernel depends only on $P$.
For cells whose lower $L$ endpoint is at least $8\sqrt\delta$,
the root $-v^2/(2\ell)$ stays in this unit interval throughout the
cell.  The mean-value theorem bounds its distance from the center
root by
\[
 \frac\delta{\ell_{\min}}+
             \frac\delta{\ell_{\min}^2}
 \le\frac{\sqrt\delta}{8}+\frac1{64}<\frac12.
\]
The same remains true for the root $(s-v^2)/(2\ell)$ when
$|s|\le\sqrt\delta$, since the added displacement is at most
$1/16$.
Let $b_m$ be the first grid boundary at least $8\sqrt\delta$;
then $b_m\le9\sqrt\delta$ for large $m$.
For every such level $s$, the conditional density formula in $Y$
therefore gives
\[
 p_{2YL+V^2}(s)\ge
 \int_{b_m}^1\int_1^2\frac{dv\,d\ell}{2\ell}
 =\frac12\log(1/b_m)
 \ge\frac14\log m-\frac12\log9.
\]
This holds on a nonempty interval of levels, and hence proves an
essential-supremum lower bound.
\end{proof}

This example shows why the adapted problem requires more than caps
on the displayed joint and marginal densities: $P$ can encode a fine
$(L,V)$ cell and allow the hub to choose a separate center for that
cell.  It is an obstruction to a reduced proof strategy, not a
counterexample to the bounded-density conjecture for independent
coordinates.  The mixture theorem above uses both the conditional
Gaussian structure and the independence of two frame blocks, and
therefore retains information absent from this example.

\subsection{Collisions on shells and a logarithmic support dependence}

\begin{lemma}[A shell collision estimate]\label{lem:hub-shell}
Let $w,v,b$ be orthonormal, put $L=w\cdot X'$, $V=v\cdot X'$,
$B=b\cdot X'$, and let
$Y=\lambda(V-v_0)^2+\gamma B$, where $\lambda\ne0$.
For almost every $\ell$, let $h_\ell$ be the unnormalized density
of $Y$ on the slice $L=\ell$, restricted to
$r\le |V-v_0|\le K$, where $0<r<K$.  Then
\[
 \int_{\mathbb R}h_\ell(y)^2\,dy
 \le \frac8{|\lambda|}\log(K/r).
\]
No bound on the coordinate leverages is required.
\end{lemma}

\begin{proof}
Take an independent copy $(L',V',B')$ and put
$S=V+V'-2v_0$, $D=V-V'$, $E=B-B'$.
The square integral is the collision density at $(\ell,\ell,0)$
of $(L,L',Y-Y')$ for the two restricted measures, and
$Y-Y'=\lambda SD+\gamma E$.
First restrict to pairs on the same side of the fold $v_0$.
Then $2r\le |S|\le2K$.  For fixed $S=s$ put
$n=(\lambda s,\gamma)/\sqrt{\lambda^2s^2+\gamma^2}$.
On the product space of the two copies, the coefficient vectors of
\[
 L,\quad L',\quad (S+2v_0)/\sqrt2,\quad (n_1D+n_2E)/\sqrt2
\]
are orthonormal.  Their joint density is at most four by
\eqref{eq:hub-projection}.  Undoing the two factors $\sqrt2$ bounds
the density in the unscaled coordinates by two.  Restrictions of
the underlying measure can only decrease it.  Thus the contribution
of the same-side pairs is at most
\[
 \int_{2r\le |s|\le2K}
       \frac{2\,ds}{\sqrt{\lambda^2s^2+\gamma^2}}
 \le\frac4{|\lambda|}\log(K/r).
\]
For opposite-side pairs fix $D=d$ instead.  Now
$2r\le |D|\le2K$, and the four orthogonal directions are
$L,L',D/\sqrt2,(n_1S+n_2E)/\sqrt2$, with
$n=(\lambda d,\gamma)/\sqrt{\lambda^2d^2+\gamma^2}$.
The same calculation gives the remaining four units.
These collision identities can first be applied to mollified
densities and bounded restrictions; disintegration and monotone
convergence then give the asserted almost-everywhere statement.
\end{proof}

\begin{theorem}[Logarithmic dependence on bounded light supports]
\label{thm:hub-shell-support}
Suppose each light coordinate has density at most one and is
supported in an interval of length at most $K$.  In the normal form
at the start of this section suppose $a=0$ and the light diagonal
coefficients vanish.  Then, uniformly in all cross terms, affine
terms, and the hub law,
\[
 \|p_Q\|_\infty
 \le \frac12+\frac{C}{\sqrt{|\lambda|}}\,
              \sqrt{1+\log(1+K^2|\lambda|)}
\]
for an absolute constant $C$.
\end{theorem}

\begin{proof}
We have $K\ge1$.  On every frozen slice the light polynomial,
after completing the square, has the form in
Lemma~\ref{lem:hub-shell}, up to a constant; its direction $b$,
coefficient $\gamma$, and fold location may depend on the frozen
value.  A missing $b$ direction can be supplied by an extra,
independent coordinate whose coefficient in $Y$ is zero.
The exponential-tilt proof of
Proposition~\ref{prop:hub-light-intervals}, applied just in direction
$v$, gives
\[
 \rho_{L,V}(\ell,u)
 \le 2\exp\{-2(|u-\mathbb EV|-2K)_+^2/K^2\}.
\]
Consequently the slice mass on
$|V-\mathbb EV|>2K+t$ is at most
$\sqrt{2\pi}K e^{-2t^2/K^2}$.

Put $W=2K+t$ and $D_0=|v_0-\mathbb EV|$.
On the remaining part of the slice and outside the fold slab of
radius $r\le W$, the values of $|V-v_0|$ lie between
$\max(r,D_0-W)$ and $D_0+W$.  The ratio of these endpoints is at
most $3W/r$: if $D_0\le2W$ use the upper bound $3W$, and otherwise
use the ratio bound three.  The fold slab has slice mass at most
$4r$, since $\rho_{L,V}\le2$.  Cauchy--Schwarz and
Lemma~\ref{lem:hub-shell} therefore give, for every Borel set $S$,
\[
 \mu_{\ell_0,\ell}(S)
 \le4r+\sqrt{\frac8{|\lambda|}\log(3W/r)}\,\sqrt{|S|}
                  +\sqrt{2\pi}K e^{-2t^2/K^2}.
\]
Here and below harmless changes in the polynomial's additive
constant do not affect the estimate.

Write $x=|S|$ and $\mu=|\lambda|$.  If $x<\mu$, take
\[
 r=\sqrt{x/\mu},\qquad
 t=K\sqrt{\tfrac12\log_+
                       (\sqrt{2\pi}K\sqrt{\mu/x})}.
\]
Then $r<1\le K\le W$, and the last term is at most
$\sqrt{x/\mu}$.  For an absolute $C_0$ this proves
\[
 \mu_{\ell_0,\ell}(S)
 \le\frac{C_0}{\sqrt\mu}
       \sqrt{x\log\!\left(e+\frac{1+K^2\mu}{x}\right)}.
\]
Indeed $\log(3W/r)\le
C[1+\log(K\sqrt{\mu/x})]$, since the extra square root of a logarithm
in $W$ contributes at most the same expression.  If $x\ge\mu$,
the displayed bound follows instead from the total slice mass
$\rho_L(\ell)\le\sqrt2$, after enlarging $C_0$.

This last displayed expression defines an increasing concave
profile for all $x>0$, with value zero at zero.
To check concavity, $x\mapsto x\log(e+A/x)$ is increasing and
concave for every $A>0$, as direct differentiation shows;
composition with the increasing concave square root preserves both
properties.  Finally, substituting $x=2e^{-u}$ shows
\[
 \int_0^2\sqrt{x\log(e+A/x)}\,\frac{dx}{x}
 \le C_1\sqrt{1+\log(1+A)}.
\]
For example use
$\log(e+(A/2)e^u)\le u+\log(e+A/2)$ and integrate
$e^{-u/2}\sqrt{u+\log(e+A/2)}$.
Lemma~\ref{lem:hub-profile} proves the theorem.
\end{proof}

\begin{proposition}[A weighted collision criterion]
\label{prop:hub-weighted-collision}
In the notation of Lemma~\ref{lem:hub-shell}, set
\[
 G=\sqrt{\lambda^2
       (|V-v_0|+|V'-v_0|)^2+\gamma^2},
 \qquad Z=(Y-Y')/G.
\]
Suppose that, uniformly in the frozen parameters, the
disintegrated collision density of $(L,L',Z)$ at
$(\ell,\ell,0)$ is at most $C_0$ for almost every $\ell$.
More explicitly, in each of the two branch classes used in the proof
of Lemma~\ref{lem:hub-shell}, assume that the nonnegative density
$F_\ell(s)$ in the remaining shell coordinate at $Z=0$
satisfies $\int F_\ell(s)\,ds\le C_0$.
Then the density restricted to $|V-v_0|\ge r$ satisfies
\[
 \int h_\ell(y)^2\,dy
 \le\frac8{|\lambda|}
              \log\!\left(1+\frac{C_0}{8r}\right).
\]
In particular, if this hypothesis holds uniformly and
$|\lambda|\ge1$, the corresponding zero-diagonal one-hub forms
have uniformly bounded densities, with a constant depending only
on $C_0$.
\end{proposition}

\begin{proof}
In the same-side class use $s=S$ and in the opposite-side class
use $s=D$.  The orthogonal projection calculation above gives
$0\le F_\ell(s)\le2$, even though its direction depends on $s$.
The collision integral in either class is at most
$|\lambda|^{-1}\int_{|s|\ge2r}F_\ell(s)|s|^{-1}\,ds$.
The maximum under the constraints $F\le2$ and $\int F\le C_0$
is obtained by filling the intervals closest to the two endpoints
$\pm2r$.  Its value is
$4|\lambda|^{-1}\log(1+C_0/(8r))$.
Adding the two classes proves the first assertion.
The fold slab again has mass at most $4r$.  Hence
\[
 \mu_{\ell_0,\ell}(S)
 \le4r+\sqrt{\frac{8|S|}{|\lambda|}
                           \log(1+C_0/(8r))}.
\]
For $|\lambda|\ge1$, choose $r=\sqrt{|S|}$ near zero and use the
trivial slice mass bound away from zero.  An increasing concave
majorant is $C(C_0)\sqrt{x\log(e+(1+C_0)^2/x)}$.
Its integral against $dx/x$ on $(0,2)$ is finite, so
Lemma~\ref{lem:hub-profile} applies.
\end{proof}

The weighted collision hypothesis is an additional nonlinear
condition.  The flat projection density bound proves the
pointwise bound $F_\ell\le2$ but does not prove its integral
bound.  Thus the proposition gives a sufficient condition whose
validity for arbitrary product densities remains open.

\subsection{Averaging the hub quadratic phase}

\begin{lemma}[A parabola extension estimate]\label{lem:hub-strichartz}
If $0\le f\le1$ and $\int f=1$, and
$g(s,\tau)=\int f(x)e^{isx^2+i\tau x}\,dx$, then
\[
 \iint_{\mathbb R^2}|g(s,\tau)|^6\,ds\,d\tau
 \le \frac{4\pi^3}{\sqrt3}.
\]
The same bound holds if $f$ in the definition of $g$ is replaced
by a complex function dominated in absolute value by $f$.
\end{lemma}

\begin{proof}
The cube $g^3$ is the Fourier transform of the pushforward of
$f(x_1)f(x_2)f(x_3)\,dx$ under
$T(x)=(x_1+x_2+x_3,x_1^2+x_2^2+x_3^2)$.
If $r$ is the distance from $x$ to the line spanned by
$(1,1,1)$, the coarea Jacobian is $2\sqrt3\,r$.
Each nonempty regular fiber is a circle of radius $r$.
Its image density $m$ is consequently at most
$\pi/\sqrt3$, since the product of the three densities is at
most one.  It has integral one.  Plancherel gives
\[
 \iint |g|^6=(2\pi)^2\int m^2
 \le(2\pi)^2(\pi/\sqrt3)\int m.
\]
For a complex amplitude, the absolute value of its pushforward
is bounded by the same nonnegative density $m$, so the argument
is unchanged.  This coarea proof is the elementary
one-dimensional parabola case of the Fourier extension estimates
associated with \citet{Strichartz1977,Foschi2007,HZ2006}.
\end{proof}

\begin{theorem}[A degenerate hub with a quadratic partner]
\label{thm:hub-degenerate}
Let $l,u$ be linearly independent, let
$L=l\cdot X'$, $U=u\cdot X'$, and put
$G=\operatorname{Gram}(l,u)$.
Assume $(l_i,u_i)G^{-1}(l_i,u_i)^T\le1/2$ for every $i$.
For independent coordinates with density at most one, consider
\[
 Q=aH^2+2H(L+\beta)+\sigma L^2+\lambda U^2
                                      +\sum_i\phi_i(X_i),
\]
where the $\phi_i$ are arbitrary quadratic polynomials.
If $\sigma\lambda\ne0$ and $a\sigma\ne1$, define
\[
 \begin{split}
 C_6&=(4\pi^3/\sqrt3)^{1/6},\\
 B_0&=(4\pi)^{5/6}C_6
 |\lambda|^{-1/2}|\sigma|^{-1/6}|1-a\sigma|^{-1/6}
               (\det G)^{-5/12}|u|^{-1/6}.
 \end{split}
\]
Then, for every $T>0$,
\begin{equation}\label{eq:hub-strichartz-block}
 \int_T^{2T}|\mathbb E e^{itQ}|\,dt\le B_0T^{-1/6},
\end{equation}
and
\[
 \|p_Q\|_\infty
 \le\frac7\pi
 \left(\frac{B_0}{6(1-2^{-1/6})}\right)^{6/7}.
\]
\end{theorem}

\begin{proof}
The leverage assumption also implies
$u_i^2/|u|^2\le1/2$.  Absorb all linear hub terms into $\beta$.
Write $\phi_i(x)=d_ix^2+b_ix+c_i$ and
$\psi_i^{(t)}(\theta)=\int f_i(x)e^{itd_ix^2+i\theta x}\,dx$.
For a fixed hub value $x$ use the exact identity
\[
 ax^2+2xL+\sigma L^2
   =\sigma(L+x/\sigma)^2+(a-1/\sigma)x^2.
\]
Two Fresnel linearizations, with Gaussian damping followed by
dominated convergence, therefore give
\[
 |\varphi_Q(t)|\le
 \frac1{4\pi|t|\sqrt{|\sigma\lambda|}}
 \int_{\mathbb R}
 |g_H(t(a-1/\sigma),2t\beta+\zeta/\sigma)|K_t(\zeta)\,d\zeta,
\]
where
$K_t(\zeta)=\int\prod_i
 |\psi_i^{(t)}(\zeta l_i+\xi u_i+tb_i)|\,d\xi$.
Lemma~\ref{lem:hub-frame-integral}, first in the $u$ direction
and then in both directions, yields
\[
 \|K_t\|_\infty\le4\pi/|u|,
 \quad \|K_t\|_1\le(4\pi)^2(\det G)^{-1/2},
 \quad
 \|K_t\|_{6/5}\le
 (4\pi)^{11/6}(\det G)^{-5/12}|u|^{-1/6}.
\]
Apply H\"older with exponents $6$ and $6/5$ on
$[T,2T]\times\mathbb R$.  The map
\[
 (t,\zeta)\longmapsto
 (t(a-1/\sigma),\,2t\beta+\zeta/\sigma)
\]
has Jacobian $|1-a\sigma|/\sigma^2$.
Lemma~\ref{lem:hub-strichartz} bounds the first factor by
$C_6(\sigma^2/|1-a\sigma|)^{1/6}$.
The second factor is at most
$T^{5/6}(4\pi)^{11/6}(\det G)^{-5/12}|u|^{-1/6}$.
Multiplying these factors and $1/(4\pi T\sqrt{|\sigma\lambda|})$
proves \eqref{eq:hub-strichartz-block}.
Since $Q$ is real, the negative-frequency bound is identical.
The full Fourier integral is at most
\[
 2T_0+\frac{2B_0T_0^{-1/6}}{1-2^{-1/6}}.
\]
Its minimum occurs at
$T_0=[B_0/(6(1-2^{-1/6}))]^{6/7}$.
Fourier inversion proves the density estimate.
\end{proof}

\begin{theorem}[Uniformity through zero partner curvature]
\label{thm:hub-degenerate-uniform}
In Theorem~\ref{thm:hub-degenerate}, assume $a=0$.
There is a finite bound depending only on
$|\lambda|$, $\det G$, and $|u|$, uniform in every
$\sigma\in\mathbb R$, in the affine terms, and in the light
diagonal terms.  One explicit bound is
\[
 \|p_Q\|_\infty\le
 1+\frac{13}{4}C_H+
 C_{\mathrm{av}}|\lambda|^{-1/2}
                         (\det G)^{-5/12}|u|^{-1/6},
 \qquad
 C_H=3\,2^{4/3}(|\lambda|\det G)^{-1/3},
\]
where
\[
 C_{\mathrm{av}}=
 \frac{(4\pi)^{5/6}C_6}{\pi}
 \left\{\frac2{1-2^{-1/6}}
             +\frac{2}{3(2^{11/6}-1)}\right\}.
\]
\end{theorem}

\begin{proof}
When $\sigma=0$, Proposition~\ref{prop:hub-localized} with empty
localization set, after changing the two frame coordinates,
gives the bound $C_H$.  Suppose $\sigma\ne0$.
Set $Q_1=Q-\sigma L^2$ and
$Q_2=Q-\sigma(L+\beta)^2$.
Each has the preceding zero-curvature form and density at most
$C_H$, since their difference consists only of affine light terms.
Fix $\delta\ge|\sigma|$ and put
$\Lambda=\sqrt{\delta/|\sigma|}\ge1$.
On $|L|\le\Lambda$, comparison with $Q_1$ bounds a window of
radius $\delta$ by $4C_H\delta$.
On $|L|>\Lambda$, $|L+\beta|\ge\Lambda/2$, conditioning on
$X'$ and using the affine hub slope gives probability at most
$\delta/|L+\beta|\le2\delta$.
On the remaining event,
$|Q-Q_2|\le|\sigma|(L+\beta)^2<\delta/4$,
so comparison with $Q_2$ gives at most $(5/2)C_H\delta$.
Consequently
\[
 \sup_s\frac{\mathbb P(|Q-s|\le\delta)}{2\delta}
 \le1+\frac{13}{4}C_H.
\]

To pass to smaller scales, put $h=|\sigma|$.
The preceding theorem gives an integrable characteristic function.
For the uniform probability density $U_h$ on $[-h,h]$,
Fourier inversion and $0\le1-\sin z/z\le\min(2,z^2/6)$ imply
\[
 \|p_Q-p_Q*U_h\|_\infty
 \le\frac1{2\pi}\int |\varphi_Q(t)|
                       (1-\operatorname{sinc}(ht))\,dt.
\]
Use \eqref{eq:hub-strichartz-block} on blocks
$T=2^j/h$, $j\in\mathbb Z$.
On such a block the last factor is at most
$\min(2,(2/3)4^j)$.  The geometric sums over $j\ge0$ and
$j\le-1$ are respectively bounded by
$2/(1-2^{-1/6})$ and $2/[3(2^{11/6}-1)]$.
Thus the error is at most
$B_0h^{1/6}$ times their sum divided by $\pi$.
When $a=0$, the factor $|\sigma|^{-1/6}$ in $B_0$
cancels $h^{1/6}$ exactly.  The estimate for $p_Q*U_h$
already follows from the window bound at $\delta=h$.
Adding the two bounds proves the result.
\end{proof}

\begin{corollary}[A normalized degenerate hub]
\label{cor:hub-degenerate-normalized}
If $w,v$ are orthonormal and $w_i^2+v_i^2\le1/2$, then
\[
 4H(w\cdot X')+2(w\cdot X')^2+2(v\cdot X')^2
                         +\sum_i\phi_i(X_i)+b_HH
\]
has density less than $18$, for every choice of the independent
coordinate laws with density at most one.
\end{corollary}

\begin{proof}
Apply Theorem~\ref{thm:hub-degenerate} with $l=2w$, $u=v$,
$a=0$, $\sigma=1/2$, and $\lambda=2$.
Then $\det G=4$, $|u|=1$, and the exact expression for $B_0$
is less than $7.5$.  The resulting exact density bound is
less than $18$.
\end{proof}

For a quadratic polynomial confined to $\operatorname{span}(w,v)$,
completing the square in $V$ produces
$\lambda'(V+\kappa L+\gamma)^2+\sigma L^2$ and affine terms.
The replacement $u=v+\kappa w$ preserves
$\det\operatorname{Gram}(w,u)=1$ and the relative leverage
$w_i^2+v_i^2$.  Hence the uniform theorem applies whenever the
transverse coefficient $|\lambda'|$ stays bounded below.
For $a\ne0$, the bound retains its factors
$|\sigma|^{-1/6}|1-a\sigma|^{-1/6}$ and diverges as the
corresponding curvatures vanish.

\subsection{Hub diagonals and square-root profiles with random centers}

Here $H$ is independent of $(t,R)$, has density $f_H\le1$, and
\[
 Q=aH^2+2Ht+R.
\]
The random variable $t$ may be a shifted linear projection of the
light coordinates.  We use the following explicitly joint
hypothesis: there are a random center $C$ coupled with $t$ and
a constant $K$ such that for every nonnegative Borel function
$\Phi$,
\begin{equation}\label{eq:hub-sq}
 \mathbb E\Phi(t,R)\le
 K\mathbb E\int_{\mathbb R}
                         \Phi(t,r)|r-C|^{-1/2}\,dr.
\end{equation}
For example, if
$R=\lambda(V-D_0)^2+C$ and $V$ has conditional density at most
$\kappa$ given $(t,D_0,C)$, then \eqref{eq:hub-sq} holds with
$K=\kappa/\sqrt{|\lambda|}$, by the two-root change of variables.
The auxiliary variables used to represent the right-hand side of
\eqref{eq:hub-sq} are sampled independently of $H$.

\begin{lemma}[A scalar discriminant bound]\label{lem:hub-sq-scalar}
For a probability density $0\le f\le1$ and
$\Delta=t^2+a s\ne0$,
\[
 \int_{\mathbb R}f(x)|s-ax^2-2tx|^{-1/2}\,dx
 \le2^{7/4}|\Delta|^{-1/4}.
\]
If $a=0$ and $t\ne0$, the sharper bound is $2|t|^{-1/2}$.
\end{lemma}

\begin{proof}
For $a\ne0$, put $A=|a|$, $d=|\Delta|$, and
$z=A(x+t/a)/\sqrt d$.
The sublevel condition $|s-ax^2-2tx|<u$ becomes
$|z^2-\operatorname{sgn}\Delta|<Au/d$.
If $\Delta>0$, the length of this set in the $z$ variable is
\[
 2\{\sqrt{1+r}-\sqrt{(1-r)_+}\},\qquad r=Au/d,
\]
which is at most $2\sqrt2\,r$.  For $r\le1$ rationalize
the difference and use
$\sqrt{1+r}+\sqrt{1-r}\ge\sqrt2$; for $r\ge1$ use
$\sqrt{1+r}\le\sqrt2\,r$.
If $\Delta<0$, its length is $2\sqrt{(r-1)_+}\le r$,
so the same bound holds.  After returning to $x$,
\[
 |\{x:|s-ax^2-2tx|<u\}|
                      \le 2\sqrt2\,u/\sqrt d.
\]
Layer cake and $f\le1$, $\int f=1$ yield
\[
 \int f|s-ax^2-2tx|^{-1/2}
 \le\int_0^\infty
       \min(1,2\sqrt2\,v^{-2}/\sqrt d)\,dv
 =2^{7/4}d^{-1/4}.
\]
For $a=0$ the sublevel length is exactly $u/|t|$,
and the same integral gives $2|t|^{-1/2}$.
\end{proof}

\begin{theorem}[Composition with a random center]
\label{thm:hub-sq-composition}
Under \eqref{eq:hub-sq}, for almost every $s$,
\begin{equation}\label{eq:hub-sq-composition}
 p_Q(s)\le1+2^{7/4}K\,
                    \mathbb E|t^2+a(s-C)|^{-1/4}.
\end{equation}
If $C=\sigma_c t^2+\eta_1t+\eta_0+C'$ for deterministic
$\sigma_c,\eta_1,\eta_0$, and the conditional density of $t$
given $C'$ is bounded by $D$, then, provided
$D_H=1-a\sigma_c\ne0$,
\[
 \|p_Q\|_\infty
 \le1+2^{13/4}K\sqrt D\,|D_H|^{-1/4}.
\]
At $a=0$ the first bound requires no information about the
random center beyond \eqref{eq:hub-sq}.
\end{theorem}

\begin{proof}
For fixed $(t,r)$ the roots of $ax^2+2tx+r=s$ have derivative
$2|ax+t|$.  Roots with $|ax+t|\ge1$ contribute at most one
in total.  Let
\[
 \mathcal K_s(t,r)=
 \sum_{\substack{x:\,ax^2+2tx+r=s\\|ax+t|<1}}
                         \frac{f_H(x)}{2|ax+t|}.
\]
Changing $r$ at fixed $t$ on each monotonicity branch gives
\[
 \int \mathcal K_s(t,r)\Psi(r)\,dr
 =\int f_H(x)\mathbf1_{\{|ax+t|<1\}}
                    \Psi(s-ax^2-2tx)\,dx.
\]
Apply \eqref{eq:hub-sq} to $\mathcal K_s$, use this identity,
and drop the indicator.  Lemma~\ref{lem:hub-sq-scalar},
with its level parameter replaced by $s-C$, proves
\eqref{eq:hub-sq-composition}.
The formulas also follow by smoothing in $s$, which avoids
exceptional critical values; Tonelli then justifies all the
nonnegative integrals.

Conditionally on $C'$, the discriminant in the second assertion
is a quadratic polynomial in $t$ with leading coefficient $D_H$.
For real roots, factor it and use Cauchy--Schwarz.
For conjugate nonreal roots, its absolute value is at least
$|D_H|(t-r)^2$, where $r$ is their common real part.
In either case
\[
 \mathbb E\bigl(|t^2+a(s-C)|^{-1/4}\mid C'\bigr)
 \le |D_H|^{-1/4}
             \sup_r\mathbb E(|t-r|^{-1/2}\mid C')
 \le2\sqrt2\,\sqrt D\,|D_H|^{-1/4}.
\]
The final bound is the bathtub integral over the interval of
length $1/D$ centered at $r$.  Multiplication of the constants
proves the result.
\end{proof}

\begin{proposition}[Randomness at the degenerate center]
\label{prop:hub-sq-log}
Assume \eqref{eq:hub-sq}, the center representation, and the
conditional density cap $D$ in
Theorem~\ref{thm:hub-sq-composition}, now allowing $D_H=0$.
Suppose $\sigma_c\ne0$.
For a random variable $Z$ write
$\alpha_1(Z)=\sup_{s\in\mathbb R,\,u>0}\mathbb P(|Z-s|\le u)/\sqrt u$.
With the convention that the second argument below is infinite
when $D_H=0$, put
\[
 A_*=\min\left\{\alpha_1(C'),
                    2\sqrt2\sqrt{|\sigma_c|/|D_H|}\right\}.
\]
If $A_*<\infty$, then
\[
 \|p_Q\|_\infty
 \le1+KD|\sigma_c|^{-1/2}
       \{12+2\log_+|\sigma_c|+4\log_+ A_*\}.
\]
In particular, center randomness with a finite
$\alpha_1(C')$ gives a finite bound at $D_H=0$.
\end{proposition}

\begin{proof}
Retain the indicator in the proof of
Theorem~\ref{thm:hub-sq-composition}, and condition on $(H,C')$.
The density cap $D$ reduces the remaining integral to
\[
 D\int_{|ax+\tau|<1}
 |s-c'-\eta_0-ax^2-(2x+\eta_1)\tau-\sigma_c\tau^2|^{-1/2}
                                                        \,d\tau.
\]
The interval has length two.  A quadratic with leading
coefficient $-\sigma_c$ and critical value $h_c$ has integral
over any such interval at most
\begin{equation}\label{eq:hub-quadratic-log-kernel}
 |\sigma_c|^{-1/2}
       \{8+2\log_+(|\sigma_c|/|h_c|)\}.
\end{equation}
To verify this bound, after translating and dividing by
$|\sigma_c|$ write the integrand as $|z^2-c|^{-1/2}$.
On $|z|\le2\sqrt{|c|}$ its integral over the whole region
is at most $\pi+2\operatorname{arcosh}2<6$.
Outside that region it is at most $2/(\sqrt3|z|)$.
On a set of length at most two the integral of this last
majorant is at most
$(4/\sqrt3)\log(1+1/(2\sqrt{|c|}))$.
Since this is at most
$(4/\sqrt3)\log(3/2)+(2/\sqrt3)\log_+(1/|c|)$,
\eqref{eq:hub-quadratic-log-kernel} follows.  The case $c=0$
is interpreted with an infinite right side.

The critical value here is $s-C'-q_c(H)$, where
\[
 q_c(x)=\eta_0+ax^2-(x+\eta_1/2)^2/\sigma_c,
\]
whose quadratic coefficient is $-D_H/\sigma_c$.
The variables $C'$ and $H$ are independent.
Convolution does not increase the concentration function, so
\[
 \mathbb P(|C'+q_c(H)-s|\le u)
 \le\min(1,A_*\sqrt u).
\]
For the bound from $q_c(H)$ use the elementary quadratic
sublevel length $2\sqrt{2u/|D_H/\sigma_c|}$ and $f_H\le1$.
Layer cake now gives
\[
 \mathbb E\log_+\frac1{|s-C'-q_c(H)|}
 \le\int_0^\infty\min(1,A_*e^{-v/2})\,dv
 \le2+2\log_+ A_*.
\]
Insert this and
$\log_+(|\sigma_c|/|h_c|)
\le\log_+|\sigma_c|+\log_+(1/|h_c|)$
into \eqref{eq:hub-quadratic-log-kernel}.
The near-root contribution is the asserted expression after
multiplication by $KD$, and the far roots cost at most one.
\end{proof}

\begin{corollary}[An independent transverse square]
\label{cor:hub-sq-independent}
Suppose $V=v\cdot X_S$ and $t=w\cdot X_T+\beta$ use disjoint
coordinate sets, and $D_0,C$ are arbitrary functions of $X_T$.
For
$Q=aH^2+2Ht+\lambda(V-D_0)^2+C$,
\[
 \|p_Q\|_\infty
 \le\frac{2^{13/4}}{|v|\sqrt{|\lambda|}}\,
        \alpha_1(t^2-aC)^{1/2}.
\]
When $a=0$ one has the unconditional bound
\[
 \|p_Q\|_\infty
 \le 2^{13/4}|\lambda|^{-1/2}|v|^{-1}|w|^{-1/2}.
\]
\end{corollary}

\begin{proof}
Condition first on $(H,X_T)$ and integrate the independent square.
Since $\|p_V\|_\infty\le\sqrt2/|v|$, the resulting density
is bounded by
\[
 \frac{\sqrt2}{|v|\sqrt{|\lambda|}}\,
 \mathbb E|s-aH^2-2Ht-C|^{-1/2}.
\]
There is no far-root term in this argument.
The scalar lemma bounds the expectation by
$2^{7/4}\mathbb E|t^2+a(s-C)|^{-1/4}$.
For any $Z$ with $\alpha_1(Z)<\infty$, layer cake gives
$\mathbb E|Z-s|^{-1/4}\le2\alpha_1(Z)^{1/2}$.
This proves the first displayed bound.
For $a=0$, use instead the scalar estimate
$2|t|^{-1/2}$ and the linear density cap
$\sqrt2/|w|$.  Its bathtub integral is
$2\sqrt2(\sqrt2/|w|)^{1/2}$, which yields the second bound.
\end{proof}

The square-root hypothesis is a statement about the joint law
with the hub partner.  A concentration estimate for the marginal
law of $R$ alone cannot replace it.  For example, if
$R=2Yt$ with an additional independent variable $Y$, then at
$a=0$ the complete form is $2(H+Y)t$; for suitable interval
laws this product has a logarithmically unbounded density.
Likewise, at $C=t^2/a$ and opposite signs of $a$ and $\lambda$,
the form $(aH+t)^2/a+\lambda V^2$ has a hyperbolic
two-dimensional singularity.  These examples have only two
effective quadratic directions and explain the determinant
and center-randomness requirements without contradicting the
rank-three target.

\section{Convolution bounds for independent coordinate blocks}
\label{sec:gluing}

The estimates in this section use disjoint groups of the original
independent coordinates.  A phase carried by one group may be arbitrary,
but the independence required of a separate smoothing term is explicit.

\subsection{Product densities and two-coordinate quadratic forms}

\begin{lemma}[A sharp product-density estimate]\label{fg:product}
Let $U,V$ be independent real random variables with probability
densities $f\leq M$ and $g\leq N$.  The density $q$ of $UV$ satisfies
\begin{equation}\label{fg:product-l2}
 \|q\|_2^2\leq4MN.
\end{equation}
Equality holds for the centered uniform densities of heights $M,N$.
Consequently, for independent cap-one variables $U,V$ and $c\ne0$,
\begin{equation}\label{fg:edge-l2}
 \|p_{2cUV+\alpha U+\beta V+\gamma}\|_2
 \leq\sqrt2\,|c|^{-1/2}
\end{equation}
for all real $\alpha,\beta,\gamma$.
\end{lemma}

\begin{proof}
The change of variables $(x,y)\mapsto(x,xy)$, away from the
null set $x=0$, gives
$q(z)=\int f(x)g(z/x)\,dx/|x|$ for almost every $z$.
All integrands in the following calculation are nonnegative,
so Tonelli's theorem applies even before finiteness is known.
Expanding the square, putting $z=xy$, and then $x=rx'$ gives
\begin{align*}
 \|q\|_2^2
 &=\int_{\mathbb R^2}\frac{f(x)f(x')}{|xx'|}
       \int_{\mathbb R}g(z/x)g(z/x')\,dz\,dx\,dx'\\
 &=\int_{\mathbb R}K_f(r)K_g(r)\,dr,\qquad
 K_f(r)=\int_{\mathbb R}f(x)f(rx)\,dx .
\end{align*}
The mass and height bounds imply
$K_f(r)\leq M\min(1,|r|^{-1})$ and
$K_g(r)\leq N\min(1,|r|^{-1})$ for $r\ne0$.
Since the integral of $\min(1,|r|^{-1})^2$ over the real line
equals $4$, this proves \eqref{fg:product-l2}.
For centered uniform densities both correlation inequalities
are equalities, proving sharpness.

Finally,
\[
 2cUV+\alpha U+\beta V+\gamma
 =2c\left(U+\frac{\beta}{2c}\right)
       \left(V+\frac{\alpha}{2c}\right)
        +\gamma-\frac{\alpha\beta}{2c}.
\]
Translations preserve the caps and independence.
For a scalar dilation by $a\ne0$, the squared $L^2$ norm of a
density is divided by $|a|$.  Apply \eqref{fg:product-l2}
with $M=N=1$ and $a=2c$.
\end{proof}

\begin{theorem}[Two coordinates with an independent $L^2$ remainder]
\label{fg:two-hub}
Let $H=(H_1,H_2)$ and $Y=(Y_1,\ldots,Y_m)$ consist of independent
cap-one original coordinates.  Let $Z$ be independent of $(H,Y)$
and have a density in $L^2(\mathbb R)$.  Suppose
\[
 Q=H^TCH+2H^T(BY+d)+R(Y)+Z,
\]
where $C$ is real symmetric, $B$ is a real $2$-by-$m$ matrix
of rank two, $d\in\mathbb R^2$, and $R$ is an arbitrary finite
measurable function.  Put
\[
 D=\det(BB^T),\qquad P=\|p_{BY+d}\|_2,\qquad R_2=\|p_Z\|_2 .
\]
Then $Q$ has a density in $C_0(\mathbb R)$ and
\begin{equation}\label{fg:two-hub-bound}
 \|p_Q\|_\infty
 \leq3\,2^{-2/3}(PR_2)^{2/3}
 \leq3\,2^{-1/3}D^{-1/6}R_2^{2/3}.
\end{equation}
If $Z=2cUV+\alpha U+\beta V+\gamma$, where $U,V$ are
independent cap-one coordinates independent of $(H,Y)$ and
$c\ne0$, then
\begin{equation}\label{fg:two-hub-edge}
 \|p_Q\|_\infty\leq3D^{-1/6}|c|^{-1/3}.
\end{equation}
\end{theorem}

\begin{proof}
By Lemma~\ref{four:projection}, the density of $S=BY+d$ is at
most $2D^{-1/2}$.  Thus $P\leq\sqrt2 D^{-1/4}$.
For fixed $t$, push the complex measure
$e^{itR(Y)}\,d\mathbb P_Y$ forward under $Y\mapsto S$.
Its total variation is dominated by the law of $S$, so it has
a density $a_t$ with $|a_t|\leq p_S$ almost everywhere.
In particular, $\|a_t\|_2\leq P$.

Let $V=Q-Z$.  The joint density of $H$ has $L^2$ norm at most
one.  Cauchy--Schwarz in the original $H$ coordinates,
followed by two-dimensional Plancherel, yields
\begin{align*}
 |\varphi_V(t)|
 &\leq
 \left(\int_{\mathbb R^2}
       \left|\int_{\mathbb R^2}e^{2it h^Ts}a_t(s)\,ds\right|^2
       dh\right)^{1/2}\\
 &=(2|t|)^{-1}\|\widehat a_t\|_2
   =\frac{\pi}{|t|}\|a_t\|_2
   \leq\frac{\pi P}{|t|}\qquad(t\ne0).
\end{align*}
The factor $e^{it h^TCh}$ has modulus one throughout this
argument.  In particular, no sign condition on $C$ is needed.

Independence of $Z$ gives
$\varphi_Q=\varphi_V\varphi_Z$.  For any $T>0$, Plancherel
and Cauchy--Schwarz give
\begin{align*}
 \int_{\mathbb R}|\varphi_Q(t)|\,dt
 &\leq2T+\pi P
       \int_{|t|>T}|t|^{-1}|\varphi_Z(t)|\,dt\\
 &\leq2T+
       \pi P(2/T)^{1/2}(2\pi)^{1/2}R_2\\
 &=2T+2\pi^{3/2}PR_2T^{-1/2}.
\end{align*}
Take $T=(\pi^{3/2}PR_2/2)^{2/3}$.  The last expression is $6T$.
Fourier inversion therefore gives a $C_0$ density with bound
$6T/(2\pi)=3\,2^{-2/3}(PR_2)^{2/3}$.
Substituting the bound for $P$ proves \eqref{fg:two-hub-bound}.
Equation \eqref{fg:edge-l2} proves \eqref{fg:two-hub-edge}.
\end{proof}

\begin{remark}
The independence of $Z$ is a hypothesis of
Theorem~\ref{fg:two-hub}.  The function $R(Y)$ can depend on the
entire slope bank $Y$, but the proof does not permit replacing
$Z$ by a variable coupled to that bank.  For example, if
$BB^T=\kappa^2 I_2$ and $|c|=\kappa^{-2}$, the bound in
\eqref{fg:two-hub-edge} is $3$, independently of $\kappa$ and
of the measurable phase $R$.
\end{remark}

\subsection{A quadratic pair and an independent square}

\begin{lemma}[Fractional integrals of a capped density]
\label{fg:bathtub-power}
For a probability density $f\leq1$, $a\in\mathbb R$, and
$0<\alpha<1$,
\begin{equation}\label{fg:fractional-bathtub}
 \int_{\mathbb R}f(x)|x-a|^{-\alpha}\,dx
 \leq\frac{2^\alpha}{1-\alpha}.
\end{equation}
\end{lemma}
\begin{proof}
Layer cake and the density cap bound the left side by
\[
 \int_0^\infty\min(1,2t^{-1/\alpha})\,dt
 =\frac{2^\alpha}{1-\alpha}.
\]
The bound is attained by the uniform density on
$[a-1/2,a+1/2]$.
\end{proof}

\begin{lemma}[A pair with one zero quadratic coefficient]
\label{fg:pair-lp}
Let $X,Y$ be independent cap-one variables and let
$b\ne0$.  For every $1<p<\infty$, the variable
\[
 T=aX^2+2bXY+2uX+2vY+c
\]
has a density satisfying
\begin{equation}\label{fg:pair-lp-bound}
 \|p_T\|_p\leq p\,|b|^{-1+1/p}.
\end{equation}
\end{lemma}
\begin{proof}
For almost every $x$, conditioning on $X=x$ leaves an affine
function of $Y$ of nonzero slope $2(bx+v)$.
Writing $\alpha=1-1/p$, its density has $L^p$ norm at most
$|2(bx+v)|^{-\alpha}$, since $\|f_Y\|_p\leq1$.
Minkowski's integral inequality and
Lemma~\ref{fg:bathtub-power} give
\[
 \|p_T\|_p
 \leq (2|b|)^{-\alpha}
      \int f_X(x)|x+v/b|^{-\alpha}\,dx
 \leq p\,|b|^{-\alpha}.
\]
The quadratic term $aX^2$ and the remaining terms only
translate the conditional density.
\end{proof}

\begin{lemma}[An $L^p$ bound for a square]\label{fg:square-lp}
For a cap-one variable $X$, $\gamma\ne0$, and $1<p<2$,
\begin{equation}\label{fg:square-lp-bound}
 \|p_{\gamma X^2+2wX}\|_p^p
 \leq |\gamma|^{1-p}\frac{2^{p-1}}{2-p}.
\end{equation}
In particular,
$\|p_{\gamma X^2+2wX}\|_{3/2}\leq2|\gamma|^{-1/3}$.
\end{lemma}
\begin{proof}
Complete the square and put $m=-w/\gamma$.
The two-root density formula and $ds=2|\gamma|r\,dr$ give
\begin{align*}
 \|p_{\gamma X^2+2wX}\|_p^p
 &=(2|\gamma|)^{1-p}
     \int_0^\infty
      \bigl(f_X(m+r)+f_X(m-r)\bigr)^p r^{1-p}\,dr\\
 &\leq|\gamma|^{1-p}
       \int_{\mathbb R}f_X(x)|x-m|^{-(p-1)}\,dx .
\end{align*}
Here the sum of the two density values is at most two.
Apply \eqref{fg:fractional-bathtub} with $\alpha=p-1$.
\end{proof}

\begin{theorem}[A physical pair and a square]\label{fg:pair-square}
Suppose $A$ has an original-coordinate principal block
\[
 \begin{pmatrix}
  a&b&0\\
  b&0&0\\
  0&0&\gamma
 \end{pmatrix},
 \qquad b\gamma\ne0.
\]
All other blocks of $A$ and all affine terms are arbitrary.
Then $Q=X^TAX+b_0^TX+c_0$ has a density in $C_0(\mathbb R)$
and
\begin{equation}\label{fg:pair-square-bound}
 \|p_Q\|_\infty\leq6|b|^{-2/3}|\gamma|^{-1/3}.
\end{equation}
If $b^2\geq|a|+1$ and $|\gamma|\geq1$, this is at most $6$;
these conditions also guarantee three eigenvalues of modulus
at least one in the displayed principal block.
\end{theorem}
\begin{proof}
Condition on all coordinates outside the displayed block.
The remaining polynomial is a sum of two independent
variables of the forms in Lemmas~\ref{fg:pair-lp} and
\ref{fg:square-lp}, with arbitrary induced affine coefficients.
Young's inequality with exponents $3$ and $3/2$ gives
\[
 \|p_T*p_S\|_\infty
 \leq\|p_T\|_3\|p_S\|_{3/2}
 \leq6|b|^{-2/3}|\gamma|^{-1/3}.
\]
The convolution belongs to $C_0$: approximate its two
factors in their respective $L^p$ norms by continuous
compactly supported functions and use Young's inequality
to obtain uniform convergence of the convolutions.
The bound is uniform in the conditioned coordinates.
Measurable conditional versions may be obtained by taking
limits of centered interval averages of the conditional laws.
Dominated convergence then shows that their average is
continuous and vanishes at infinity, and retains the bound.

The two eigenvalues of the upper-left pair are
$(a\pm\sqrt{a^2+4b^2})/2$.  Both have modulus at least one
exactly when $b^2\geq|a|+1$.
The third eigenvalue of the displayed block is $\gamma$.
\end{proof}

\section{Qualitative regularity and variational formulations}
\label{sec:regularity-normal-forms}

Absolute continuity, continuity of a density, and a bound uniform over
all dimensions and matrices are distinct conclusions.  We first prove
that every fixed quadratic matrix of rank at least three has a $C_0$
density for the product laws under consideration.  The estimates used
for this regularity statement may depend on the matrix.

\subsection{Continuity for rough product densities}

\begin{lemma}[A quadratic integral near its critical point]
\label{fg:critical-ball}
Let $B$ be an invertible real symmetric $d$-by-$d$ matrix,
where $d=3$ or $4$.  Write
$\Lambda=\|B\|$, $\sigma=s_d(B)>0$, and
$q(x)=x^TBx$.  The pushforward under $q$ of Lebesgue measure
restricted to $|x|\leq r$ has a density bounded by
\begin{equation}\label{fg:critical-bound}
 C_d\Lambda^{(d-2)/2}|\det B|^{-1/2}r^{d-2},
 \qquad C_3=2\pi,\quad C_4=4\pi.
\end{equation}
If $F(x)=\prod_{i=1}^d f_i(x_i)$ with probability densities
$f_i\leq1$, the pushforward of
$F(x){\bf1}_{\{|x|>R\}}\,dx$ has density at most
\begin{equation}\label{fg:far-bound}
 \frac{d^{3/2}}{\sigma R}.
\end{equation}
\end{lemma}

\begin{proof}
Use the linear change of variables $z=|B|^{1/2}x$, followed
by an orthogonal change in this deterministic integral.
Its Jacobian is $|\det B|^{-1/2}$, the quadratic is a sum
of positive squares minus a sum of negative squares, and
the image of the ball lies in $|z|\leq\rho=\sqrt\Lambda r$.

For a definite form, polar coordinates give the bound
$|\mathbb S^{d-1}|\rho^{d-2}/2$.
This is $2\pi\rho$ for $d=3$ and $\pi^2\rho^2$ for $d=4$.
For signature $(p,m)$ with $p,m>0$ and $p+m=d$, put
$u=|z_+|^2$ and $v=|z_-|^2$.
The quadratic is $u-v$, while the radial measure is
\[
 \frac{|\mathbb S^{p-1}||\mathbb S^{m-1}|}{4}
       u^{p/2-1}v^{m/2-1}\,du\,dv,
 \qquad u,v\geq0,\quad u+v\leq\rho^2 .
\]
In dimension three, signature $(2,1)$ gives a density at
most $\pi\int_0^{\rho^2}v^{-1/2}\,dv=2\pi\rho$;
the opposite signature follows by reflection of the output.
In dimension four, signature $(2,2)$ gives at most
$\pi^2\rho^2/2$, because the allowed radial interval on
$u-v=t$ has length at most $\rho^2/2$.
For signature $(3,1)$, the density is at most
$2\pi\rho\int_0^{\rho^2}v^{-1/2}\,dv=4\pi\rho^2$.
Reflection treats $(1,3)$.  These estimates prove
\eqref{fg:critical-bound}.

For the second assertion, $|Bx|\geq\sigma R$ when $|x|>R$.
Hence at least one coordinate satisfies
$|\partial_iq(x)|=2|(Bx)_i|\geq2\sigma R/\sqrt d$.
For fixed $x_{-i}$, a quadratic in $x_i$ has at most two
inverse branches.  On the indicated set their inverse
derivatives are each at most $\sqrt d/(2\sigma R)$.
The cap on $f_i$ and change of variables therefore bound
the conditional pushforward density by
$\sqrt d/(\sigma R)$.
Integrating the other probability densities and summing
over the possible $i$ proves \eqref{fg:far-bound}.
\end{proof}

\begin{lemma}[The region with only one large derivative]
\label{fg:thin-tubes}
In the notation of Lemma~\ref{fg:critical-ball}, fix
$0<r<R<\infty$.  Let
\[
 E_{r,R,\delta}
 =\{x:r\leq|x|\leq R,\ 
       \text{at most one }|\partial_iq(x)|\text{ exceeds }\delta\}.
\]
For $0<\delta\leq1$, the pushforward of
$F{\bf1}_{E_{r,R,\delta}}\,dx$ has density at most
\begin{equation}\label{fg:tube-bound}
 C(B,r,R)\delta^{d-2},
\end{equation}
uniformly over all products of cap-one probability densities.
\end{lemma}

\begin{proof}
Partition the set measurably according to an index $i$
maximizing $|(Bx)_i|$, breaking ties by the smallest index.
On this part,
\[
 |\partial_iq|\geq2\sigma r/\sqrt d,\qquad
 |(Bx)_j|\leq\delta/2\quad(j\ne i).
\]
Write $Bx=se_i+v$, where $v_i=0$.
Then $|s|\leq\Lambda R$, $|v|\leq\sqrt{d-1}\delta/2$, and
\[
 x=sB^{-1}e_i+B^{-1}v.
\]
The projection onto the $d-1$ coordinates other than $i$
therefore lies within distance
$\epsilon=\sqrt{d-1}\delta/(2\sigma)$ of a line segment of
length at most $2\Lambda R/\sigma$.
In $\mathbb R^{d-1}$ the volume of this neighborhood is at
most
\[
 \frac{2\Lambda R}{\sigma}\,
          \omega_{d-2}\epsilon^{d-2}
       +\omega_{d-1}\epsilon^{d-1}
 \leq C(B,R)\delta^{d-2},
\]
where $\omega_j$ is the volume of the unit ball in
$\mathbb R^j$.  The estimate also holds when the projected
segment is a point.

For each fixed value of these $d-1$ coordinates, the
inverse-branch argument from the preceding lemma bounds
the one-dimensional pushforward density by
$\sqrt d/(\sigma r)$.
Here all density factors are bounded by one, so integration
over the projected set costs at most its Lebesgue volume.
Summing over $i$ proves \eqref{fg:tube-bound}.
\end{proof}

\begin{lemma}[Continuity in three or four invertible coordinates]
\label{fg:invertible-continuity}
Let $d=3$ or $4$, let $B$ be invertible and symmetric, and let
$X_1,\ldots,X_d$ be independent with densities bounded by one.
For every affine term, $X^TBX+b^TX+c$ has a density in
$C_0(\mathbb R)$.  A height bound independent of the affine
term and of the coordinate laws is
\begin{equation}\label{fg:fixed-block-height}
 M_B=C_d\|B\|^{(d-2)/2}|\det B|^{-1/2}
                +d^{3/2}/s_d(B).
\end{equation}
For fixed $B,b,c$, convergence of all coordinate densities
in $L^1$ implies uniform convergence of these continuous
output densities.
\end{lemma}

\begin{proof}
Translation of the original independent coordinates removes
the affine term, up to a constant shift of the output.
It suffices to consider $q(x)=x^TBx$.
Splitting into the unit ball and its complement in
Lemma~\ref{fg:critical-ball} immediately gives
\eqref{fg:fixed-block-height}.

We give the regularity argument in detail to allow rough,
disconnected, and unbounded coordinate supports.
For a product density $F$ and measurable set $E$, denote by
$p_{F;E}$ the density of the restricted pushforward
$q_*(F{\bf1}_E\,dx)$.
These restrictions are absolutely continuous because they
are dominated by the full pushforward, which is absolutely
continuous by Lemma~\ref{geom:ac}.

Let $F=\prod_i f_i$ and $G=\prod_i g_i$ be two products of
cap-one probability densities.  On the set
\[
 \mathcal G_\delta
 =\{x:\text{at least two }|\partial_iq(x)|\text{ exceed }\delta\},
\]
the telescoping identity
\[
 F-G=\sum_{i=1}^d(f_i-g_i)
              \prod_{k<i}f_k\prod_{k>i}g_k
\]
and coordinatewise change of variables give
\begin{equation}\label{fg:good-stability}
 \|p_{F;\mathcal G_\delta}-p_{G;\mathcal G_\delta}\|_\infty
 \leq\frac{2(d-1)}{\delta}
             \sum_{i=1}^d\|f_i-g_i\|_1 .
\end{equation}
Indeed, for the summand containing $f_i-g_i$, every point
of $\mathcal G_\delta$ has some $j\ne i$ with
$|\partial_jq|>\delta$.
Sum over those $d-1$ choices.  On each choice, solve for
$x_j$ along the at most two inverse branches.
The $j$th density factor is at most one and the inverse
derivative is at most $\delta^{-1}$.
Integration of the remaining factors gives
$\|f_i-g_i\|_1$, since all other factors have mass one.
The same estimate holds with an arbitrary measurable subset
of $\mathcal G_\delta$ in place of that set.

Fix a desired error $\eta>0$.
Choose $r$ sufficiently small and $R$ sufficiently large
that \eqref{fg:critical-bound} and \eqref{fg:far-bound} are
each less than $\eta$.
For these fixed $r,R$, choose $\delta$ sufficiently small
that \eqref{fg:tube-bound} is less than $\eta$.
The complement of $\mathcal G_\delta$ is covered by these
three regions.  Thus both $F$ and $G$ contribute at most
$3\eta$ in density there.  Equation
\eqref{fg:good-stability} gives the uniform estimate
\begin{equation}\label{fg:rough-stability}
 \|p_F-p_G\|_\infty
 \leq6\eta+\frac{2(d-1)}{\delta}
                    \sum_{i=1}^d\|f_i-g_i\|_1.
\end{equation}
The parameters depend only on $B,\eta$, not on the
coordinate densities.

Each probability density $f_i\leq1$ can be approximated in
$L^1$ by smooth compactly supported probability densities
with the same cap.  One explicit construction first
truncates to a large bounded interval, places the missing
mass with height one on a disjoint interval, and then
convolves with a nonnegative smooth compactly supported
approximate identity.  Truncation and the added mass cost
at most twice the discarded mass in $L^1$; convolution
preserves the total mass and the cap.

For a smooth compactly supported product density, its
pushforward under $q$ belongs to $C_0$.
To see this without assuming regularity at the critical
value, remove a ball of radius $r$ about the origin by a
smooth cutoff.  On the compact support that remains, the
gradient of $q$ never vanishes.  Cover that support by
finitely many coordinate charts with one nonzero partial
derivative and take a smooth partition of unity.
In each chart the coordinates $(q,x_{-i})$ form a smooth
change of variables.  The transformed integrand is smooth
with compact support, and integration in $x_{-i}$ gives a
smooth compactly supported function of $q$.
The removed part has density bounded by a constant times
$r^{d-2}$, by \eqref{fg:critical-bound}.
Letting $r\downarrow0$ shows that the full density is a
uniform limit of functions in $C_0$, and hence is in $C_0$.

Apply \eqref{fg:rough-stability} to the smooth product
approximations.  Their output densities are uniformly
Cauchy and therefore converge uniformly to a function in
$C_0$.  The input products converge in $L^1$, so their
pushforward measures converge in total variation.
The limiting continuous function is consequently a
density of $q(X)$.  Finally, the same inequality, first
choosing $\eta$ and then the input $L^1$ error, proves the
stability assertion.
\end{proof}

\begin{theorem}[Every fixed rank-three-or-higher matrix]
\label{fg:fixed-matrix-continuity}
Let $A$ be any real symmetric matrix of rank at least three.
For independent cap-one coordinates and arbitrary affine
terms, $Q=X^TAX+b^TX+c$ has a bounded density in
$C_0(\mathbb R)$.
\end{theorem}

\begin{proof}
There is an invertible original-coordinate principal
submatrix $B=A_{II}$ of order $d=3$ or $4$.
For completeness, symmetric elimination proves this as
follows.  If a nonzero diagonal entry is available, use it
as an invertible one-coordinate pivot.  If all diagonal
entries vanish but the remaining matrix is nonzero, a
nonzero off-diagonal entry gives the invertible block
$\left(\begin{smallmatrix}0&a\\a&0\end{smallmatrix}\right)$.
An invertible pivot reduces rank by its order, through its
Schur complement.  Continue until the total pivot order
first reaches three.  The rank assumption permits doing
so, and the total is either three or four.
The corresponding principal submatrix of the original
matrix is invertible, since its determinant is the product
of the successive nonzero pivot determinants.

Condition on $Y=X_{I^c}$.  The conditional quadratic matrix
in the coordinates $X_I$ is the fixed invertible matrix
$B$; all conditioning effects are affine.
Lemma~\ref{fg:invertible-continuity} supplies for every $y$
a conditional density $p_y\in C_0$ bounded by $M_B$.
There is a jointly measurable choice: for each $(y,t)$
take the limit of
\[
 \frac{m}{2}\,
 \mathbb P\bigl(|X_I^TBX_I+b(y)^TX_I+c(y)-t|<1/m\bigr)
 \quad\text{as }m\to\infty.
\]
The limit exists at every $t$ because the conditional
continuous density exists.
Its average $p(t)=\mathbb E p_Y(t)$ is a density of $Q$.
The bound $M_B$ permits dominated convergence both as
$t$ approaches a finite point and as $|t|\to\infty$.
Thus $p$ is continuous, vanishes at infinity, and is bounded.
\end{proof}

\begin{remark}
The constant $M_B$ in this proof can deteriorate when the
selected principal block becomes ill-conditioned.
Theorem~\ref{fg:fixed-matrix-continuity} therefore gives a
qualitative regularity statement for every fixed matrix;
the dimension-independent constant in
Conjecture~\ref{conj:main} remains a separate question.
\end{remark}

\subsection{Recovery from one fixed smoothing kernel}

Let $\mathfrak C$ be the optimal constant in Conjecture~\ref{conj:main},
as an extended nonnegative real number:
\[
 \mathfrak C
 =\sup_{\substack{n\geq3,\ A=A^T,\ s_3(A)\geq1\\
                  X\ {\rm independent},\ 0\leq f_i\leq1\\
                  b\in\mathbb R^n,\ c\in\mathbb R}}
       \|p_{X^TAX+b^TX+c}\|_\infty .
\]
Every density in this definition exists by
Lemma~\ref{geom:ac}.  The next identities are valid whether
$\mathfrak C$ is finite or infinite.

\begin{lemma}[Approximate-identity recovery]\label{fg:kernel-recovery}
Let $K\geq0$ be any probability density in $L^1(\mathbb R)$,
and set $K_r(x)=r^{-1}K(x/r)$.  For every probability
density $p$ and every $r_0>0$,
\begin{equation}\label{fg:kernel-recovery-eq}
 \sup_{0<r\leq r_0}\|K_r*p\|_\infty=\|p\|_\infty .
\end{equation}
\end{lemma}
\begin{proof}
Young's inequality gives the upper bound.
Translation continuity in $L^1$ and dominated convergence
give
\[
 \|K_r*p-p\|_1
 \leq\int_{\mathbb R}K(z)\,
           \|p(\,\cdot-rz)-p\|_1\,dz\longrightarrow0 .
\]
If the left side of \eqref{fg:kernel-recovery-eq} were a
finite number $M$ smaller than $\|p\|_\infty$, take
$r_j\downarrow0$ and an almost-everywhere convergent
subsequence of the $L^1$ approximations.
The limit would satisfy $p\leq M$ almost everywhere,
a contradiction.  The same argument treats
$\|p\|_\infty=\infty$.
\end{proof}

\begin{theorem}[An exact fixed-kernel formulation]
\label{fg:fixed-kernel}
For every fixed probability density $K\in L^1(\mathbb R)$,
\begin{equation}\label{fg:fixed-kernel-eq}
 \mathfrak C
 =\sup_{\substack{A,X,b,c\ {\rm as\ above}\\s_3(A)\geq1}}
       s_3(A)\,\|K*p_Q\|_\infty .
\end{equation}
In particular,
\begin{equation}\label{fg:fixed-box-eq}
 \mathfrak C
 =\sup_{\substack{A,X,b,c\ {\rm as\ above}\\s_3(A)\geq1,\ y\in\mathbb R}}
       \frac{s_3(A)}2\,\mathbb P(|Q-y|\leq1).
\end{equation}
\end{theorem}

\begin{proof}
First observe the equivalent weighted form
\[
 \mathfrak C=\sup s_3(A)\|p_Q\|_\infty .
\]
The weighted supremum is at least the original one because
$s_3(A)\geq1$.  For the reverse inequality, normalize each
instance by $Q\mapsto Q/s_3(A)$, including its affine term.
The new matrix has third singular value one and its density
height is $s_3(A)\|p_Q\|_\infty$.

Let the right side of \eqref{fg:fixed-kernel-eq} be
$\mathfrak C_K$.  Young's inequality gives
$\mathfrak C_K\leq\mathfrak C$.
Fix an admissible instance and write $s=s_3(A)$.
For $0<r\leq s$, the form $Q/r$ is still admissible,
with matrix $A/r$ and density $p_r(u)=r p_Q(ru)$.
Direct substitution gives
\[
 (K*p_r)(y/r)=r(K_r*p_Q)(y),
 \qquad
 \frac{s}{r}\|K*p_r\|_\infty
     =s\|K_r*p_Q\|_\infty .
\]
Taking the supremum over $0<r\leq s$ and applying
Lemma~\ref{fg:kernel-recovery} shows
$\mathfrak C_K\geq s\|p_Q\|_\infty$.
The supremum over the original instance proves equality.

Take $K=\frac12{\bf1}_{[-1,1]}$ for
\eqref{fg:fixed-box-eq}.  For an absolutely continuous
law the endpoints carry no mass, and the box convolution
is the continuous function
$y\mapsto\frac12\mathbb P(|Q-y|\leq1)$.
Its essential supremum is therefore its pointwise
supremum, as used in that identity.
\end{proof}

\begin{lemma}[Stability at a fixed interval width]
\label{fg:fixed-box-stability}
Suppose probability laws $\mu_j$ converge weakly to an
atomless probability law $\mu$.  For every fixed $a>0$,
\begin{equation}\label{fg:box-stability}
 \sup_{y\in\mathbb R}
 \bigl|\mu_j([y-a,y+a])-\mu([y-a,y+a])\bigr|
 \longrightarrow0.
\end{equation}
\end{lemma}
\begin{proof}
Let $F_j,F$ be the distribution functions.
Because $\mu$ is atomless, $F$ is continuous.
Weak convergence gives $F_j(t)\to F(t)$ at every $t$.
A finite grid whose successive $F$ values differ by at
most a prescribed $\epsilon$, together with two tail
points of total $\mu$ mass at most $2\epsilon$, proves by
monotonicity that $\sup_t|F_j(t)-F(t)|\to0$.
For clarity, one may choose a bounded interval containing
all but $\epsilon$ of each tail of $\mu$, subdivide it
finely using uniform continuity of $F$, and bracket each
$t$ between adjacent grid points.
The same uniform bound controls $F_j(t-)$, since
$F(t-)=F(t)$.
Now
$\mu_j([y-a,y+a])=F_j(y+a)-F_j((y-a)-)$.
Subtract the corresponding expression for $\mu$ to obtain
\eqref{fg:box-stability}.
\end{proof}

\subsection{A countable heat-kernel formulation}

Let $\mathcal U$ be the countable class of finite unions of
bounded half-open intervals with rational endpoints and
total Lebesgue measure exactly one.
For $n\geq3$, let $\mathcal A_n$ consist of real symmetric
matrices with rational entries such that every entry and
every principal minor is nonzero, the spectrum is simple,
and $s_3(A)>1$.
Put
\[
 \gamma_\epsilon(t)=
    \frac{1}{\sqrt{2\pi}\epsilon}
                 \exp\left(-\frac{t^2}{2\epsilon^2}\right).
\]

\begin{theorem}[A countable exact formulation]\label{fg:countable-heat}
The full optimal constant satisfies
\begin{equation}\label{fg:countable-heat-eq}
 \mathfrak C=
 \sup_{\substack{
       n\geq3,\ A\in\mathcal A_n,\ b\in\mathbb Q^n\\
       E_1,\ldots,E_n\in\mathcal U\\
       \epsilon\in\mathbb Q_{>0},\ y\in\mathbb Q}}
 \int_{E_1\times\cdots\times E_n}
       \gamma_\epsilon(x^TAx+b^Tx-y)\,dx .
\end{equation}
There is no bound here on $n$, the number of intervals,
their locations, or their rational denominators.
\end{theorem}

\begin{proof}
Lemma~\ref{fg:kernel-recovery}, with a Gaussian kernel, shows
that for each output law
\[
 \|p_Q\|_\infty
   =\sup_{\epsilon>0,\ y\in\mathbb R}
                   \mathbb E\gamma_\epsilon(Q-y).
\]
The convolution at every positive $\epsilon$ is continuous,
so a pointwise supremum in $y$ is appropriate.
Constants in $Q$ can be absorbed into $y$.
Every value on the right side of
\eqref{fg:countable-heat-eq} comes from an admissible product
law, and is at most $\mathfrak C$.
We prove the reverse inequality by approximating an
arbitrary fixed smoothed value.

First, for any bounded nonnegative measurable payoff $h$,
\begin{equation}\label{fg:bathtub-linear}
 \sup_{\substack{0\leq f\leq1\\\int f=1}}\int hf
       =\sup_{|E|=1}\int_Eh .
\end{equation}
Here the suprema need not be attained.
One elementary proof chooses
$\tau=\inf\{t\geq0:|\{h>t\}|\leq1\}$.
Monotone convergence from levels above $\tau$ gives
$|\{h>\tau\}|\leq1$.
For $\delta>0$, choose a set $E$ of measure one that contains
$\{h>\tau\}$ and whose remaining points have
$h\geq\tau-\delta$; if $\tau=0$, fill with arbitrary points.
Such a set exists by the definition of $\tau$ and the
nonatomic nature of Lebesgue measure.
In the identity
$\int h(f-{\bf1}_E)=\int(h-\tau)(f-{\bf1}_E)$,
the contribution from $\{h>\tau\}$ is nonpositive.
The contribution outside $E$ is also nonpositive, and
the remaining contribution is at most $\delta$.
This proves \eqref{fg:bathtub-linear} by letting
$\delta\downarrow0$.

Apply \eqref{fg:bathtub-linear} successively to each original
coordinate of the bounded payoff
$\gamma_\epsilon(x^TAx+b^Tx-y)$, averaging over the others.
At each step the smoothed value can be preserved up to an
arbitrarily small prescribed error.
After finitely many steps the input laws are uniform on
measurable sets $E_i$ of measure one.

Each such set can be approximated in symmetric difference
by a member of $\mathcal U$.
Indeed, Lebesgue regularity gives a finite bounded union of
rational intervals $F$ with $|E_i\mathbin\triangle F|$
arbitrarily small.  Its total length is rational.
If that length exceeds one, remove or shorten interval
pieces, retaining rational endpoints, until the length is
one.  If it is less than one, add a disjoint rational
interval of the missing rational length.
The additional change has measure $||F|-1|$, which is at
most $|E_i\mathbin\triangle F|$.
Thus the final symmetric difference can also be made
arbitrarily small.
The telescoping inequality
\[
 \left\|\prod_i f_i-\prod_i g_i\right\|_1
       \leq\sum_i\|f_i-g_i\|_1
\]
shows that these replacements alter the fixed smoothed
value by at most the last sum times
$\|\gamma_\epsilon\|_\infty$.

The resulting product set is bounded.
Next replace $A$ by $(1+\delta)A$ and then by a nearby
rational symmetric matrix $A_\delta$ whose entries and
principal minors are all nonzero and whose spectrum is
simple.  Such matrices are dense: each forbidden condition
is the zero set of a nonzero polynomial in the independent
matrix entries, including the spectral discriminant.
Their finite union has empty interior, and its complement
is open and contains rational points in every sufficiently
small ball.
Choose
$\|A_\delta-(1+\delta)A\|<\delta/2$.
Singular-value stability gives
$s_3(A_\delta)>1$, while $A_\delta\to A$.
Also approximate $b$ by rational vectors.
On the fixed bounded product set, the corresponding
quadratic polynomials converge uniformly.
Continuity of the Gaussian payoff therefore preserves
the smoothed integral in the limit.

Finally approximate $\epsilon>0$ and $y$ by positive
rational and rational parameters, respectively.
The kernels are uniformly bounded in a neighborhood of
the fixed positive $\epsilon$, and depend continuously
on all their parameters.  Dominated convergence preserves
the integral.
All approximation errors can be chosen arbitrarily small.
Every smoothed value in the full problem is therefore
bounded by the right side of
\eqref{fg:countable-heat-eq}, proving equality.
\end{proof}

\begin{remark}
The fixed-kernel and countable formulations are exact
identities for the optimal density bound.  Neither identity
by itself gives a finite bound for the supremum.
The fixed-width stability in
Lemma~\ref{fg:fixed-box-stability} also does not assert
uniform control as the interval width tends to zero.
\end{remark}

\section{Further coordinate block estimates and spectral stability}
\label{sec:additional-results}

The estimates in this section apply to blocks of original independent
coordinates with densities bounded by one.  Their conclusions are
uniform in affine terms.  The deterministic spectral inequalities
control perturbations of coordinate submatrices.

\subsection{A linear bank with two crossing directions}

\begin{lemma}[Negative moments from a density bound]
\label{supp:negative-moment}
If a probability law on $\mathbb R^d$ has density at most $K$, then for
$0<p<d$ and every $z\in\mathbb R^d$,
\[
 \mathbb E|Y-z|^{-p}\le\frac d{d-p}(K v_d)^{p/d},
\]
where $v_d$ is the volume of the unit ball.
\end{lemma}
\begin{proof}
For every $r>0$, the probability of the ball of radius $r$ is at most
$\min\{1,Kv_dr^d\}$.  Integrate this bound in
\[
 \mathbb E|Y-z|^{-p}
 =\int_0^\infty\mathbb P\{|Y-z|<t^{-1/p}\}\,dt
\]
and split at $t=(Kv_d)^{p/d}$.  The two contributions are
$(Kv_d)^{p/d}$ and $p(Kv_d)^{p/d}/(d-p)$.
\end{proof}

\begin{theorem}[An arbitrary-rank crossing block with a linear bank]
\label{supp:null-bank}
Let $X_I$ and $X_J$ be disjoint original coordinate banks and let
\[
 Q=F(X_I)+2X_I^{\mathsf T}M X_J+b_J^{\mathsf T}X_J,
\]
where $F$ is any real measurable function.  If $s_2(M)>0$, then
\begin{equation}\label{eq:supp-null-bank}
 \|p_Q\|_\infty\le
 \frac{2\sqrt\pi}{\sqrt{s_1(M)s_2(M)}}.
\end{equation}
The rank and all smaller singular values of $M$ are unrestricted.
\end{theorem}
\begin{proof}
Condition on $X_I=x$.  The remaining polynomial is linear in $X_J$,
with coefficient vector $c(x)=2M^{\mathsf T}x+b_J$.  The
one-dimensional product-projection estimate of \citet{LPP2016} gives
a conditional density bound $\sqrt2/|c(x)|$ whenever $c(x)\ne0$.
Let $u_1,u_2$ and $v_1,v_2$ be corresponding left and right singular
vectors for the two largest singular values $s_1,s_2$ of $M$.
The vector $U=(u_1^{\mathsf T}X_I,u_2^{\mathsf T}X_I)$ has density
bounded by two.  Therefore
\[
 Z=(2s_1U_1+v_1^{\mathsf T}b_J,
       2s_2U_2+v_2^{\mathsf T}b_J)
\]
has density at most $1/(2s_1s_2)$, and $|c(X_I)|\ge|Z|$.
Lemma~\ref{supp:negative-moment} with $d=2,p=1$ gives
\[
 \mathbb E|c(X_I)|^{-1}
 \le2\sqrt{\frac\pi{2s_1s_2}}.
\]
In particular, the exceptional set $c(X_I)=0$ has probability zero.
Averaging the conditional density inequality proves
\eqref{eq:supp-null-bank}.  This argument retains the full matrix $M$;
the projection onto two singular directions is used only for a lower
bound on the length of $c(X_I)$.
\end{proof}

The condition that the $J$ bank enters linearly is essential to this
proof.  Theorem~\ref{supp:null-bank} includes a bilinear form with
an arbitrary quadratic or measurable addition on one bank.  Its
conditioning argument is closely related to the independent-block
density estimates in \citet{BGKS2024}.

\subsection{A singular physical pair and a random slope}

\begin{theorem}[Rank-one physical pair]
\label{supp:rank-one-pair}
Let $X,Y,Z_1,\ldots,Z_m$ be independent original coordinates.  Suppose
\begin{align*}
 Q={}&aX^2+2bXY+cY^2+2X\langle u,Z\rangle
       +2Y\langle v,Z\rangle+F(Z)\\
    &+2\alpha X+2\beta Y,
\end{align*}
where $a\ne0$, $ac=b^2$, $F$ is real measurable, and
$w=v-(b/a)u\ne0$.  Then
\begin{equation}\label{eq:supp-rank-one-pair}
 \|p_Q\|_\infty\le
 \frac{4\,2^{1/4}}{\sqrt{|a|\,|w|}}.
\end{equation}
The dependence of $F$ on all coordinates of $Z$ is unrestricted.
\end{theorem}
\begin{proof}
Complete the square only in the original coordinate $X$.  With
$\eta=\beta-b\alpha/a$ this writes
\[
 Q=a\left(X+\frac{bY+\langle u,Z\rangle+\alpha}{a}\right)^2
       +2Y(\langle w,Z\rangle+\eta)+\widetilde F(Z).
\]
Conditioning on $(Y,Z)$ and using the scalar square density bound gives
\[
 p_Q(t)\le |a|^{-1/2}
  \mathbb E k_a\bigl(t-2Y(\langle w,Z\rangle+\eta)
                     -\widetilde F(Z)\bigr),
 \quad k_a(u)=|u|^{-1/2}\mathbf1_{\{au>0\}},
\]
for almost every $t$.  For any probability density $f\le M$,
\[
 \sup_t\int_0^\infty y^{-1/2}f(t\pm y)\,dy\le2\sqrt M,
\]
by filling the interval of length $1/M$ nearest zero.
If $\ell\ne0$, the variable $2\ell Y$ has density at most
$1/(2|\ell|)$.  Hence the expectation of the one-sided kernel
over $Y$ is at most $\sqrt2\,|\ell|^{-1/2}$, uniformly in its
translate.
The scalar projection $\langle w,Z\rangle+\eta$ has density at
most $\sqrt2/|w|$ by \citet{LPP2016}.
Lemma~\ref{supp:negative-moment} with $d=1,p=1/2$ gives
\[
 \mathbb E|\langle w,Z\rangle+\eta|^{-1/2}
 \le2\sqrt{2\sqrt2/|w|}.
\]
Multiplication by the preceding bound for the slope gives
$\sqrt2\cdot2\sqrt{2\sqrt2}=4\,2^{1/4}$, as asserted.
All changes of variables can be integrated against nonnegative output
test functions, so the almost-everywhere conditional formulas introduce
no exceptional-level assumption.
\end{proof}

\begin{corollary}[An aggregate bound for singular pairs]
\label{supp:singular-pair-volume}
Under Theorem~\ref{supp:rank-one-pair}, put
\[
 W=a^2\sum_{i\ne j}w_i^2w_j^2.
\]
If $W>0$, then
\[
 \|p_Q\|_\infty\le
     4\,2^{1/4}|a|^{-1/4}W^{-1/8}.
\]
\end{corollary}
\begin{proof}
The inequality $|w|^4\ge\sum_{i\ne j}w_i^2w_j^2=W/a^2$
inserted into \eqref{eq:supp-rank-one-pair} proves the statement.
\end{proof}

\subsection{Spectral energy in a coordinate set}

For a coordinate set $H$, let $P_H:\mathbb R^{|H|}\to\mathbb R^n$
be the coordinate inclusion.  Compression is written
$T_{HH}=P_H^{\mathsf T}TP_H$.  Eigenvalues of a positive
semidefinite matrix are listed in decreasing order.

\begin{lemma}[Spectral-square spill]
\label{supp:spill}
Let $A$ be symmetric, $C=A_{HH}$, $B=A_{H,H^c}$, and
$0\le G\le A^2$.  Suppose $\lambda_k(G_{HH})\ge g^2$.
For $0\le r<k$ put $c=s_{r+1}(C)$.  If $c<g$, then
\[
 s_{k-r}(B)\ge\sqrt{g^2-c^2}.
\]
In particular, one may use $G=\Lambda^2
\mathbf1_{\{|A|\ge\Lambda\}}$ for any $\Lambda>0$.
\end{lemma}
\begin{proof}
The exact block identity is $(A^2)_{HH}=C^2+BB^{\mathsf T}$.
Write $C=C_r+E$, where $C_r$ consists of its $r$ largest
absolute eigenmodes.  Then $C_rE=0$, $\operatorname{rank}C_r\le r$,
and $E^2\le c^2I$.  Thus
\[
 BB^{\mathsf T}+C_r^2=(A^2)_{HH}-E^2\ge G_{HH}-c^2I.
\]
The left side has at least $k$ eigenvalues at least $g^2-c^2$.
Adding a positive semidefinite matrix of rank at most $r$ can
increase that count by at most $r$.  Consequently
$BB^{\mathsf T}$ has at least $k-r$ eigenvalues of that size,
which is the assertion.  The stated spectral projection gives an
admissible $G$ by the spectral theorem.
\end{proof}

\begin{lemma}[A pair in an approximate rank-two self block]
\label{supp:pair-extraction}
Let $C=C^{\mathsf T}$ be $m$ by $m$, $m\ge3$, and write
$N=\binom m2$.  If $s_2(C)\ge b$ and $s_3(C)<a$, then some
principal two-coordinate block $K$ satisfies
\[
 s_{\min}(K)>b/N-a.
\]
\end{lemma}
\begin{proof}
Let $C_2=\lambda_1uu^{\mathsf T}+\lambda_2vv^{\mathsf T}$ be the
two largest absolute eigenmodes.  The orthonormality of $u,v$ gives
$\sum_{i<j}(u_iv_j-u_jv_i)^2=1$.  Choose a pair $I$ whose
summand is at least $1/N$.  Then
\[
 |\det(C_2)_{II}|\ge s_1(C)s_2(C)/N,
 \qquad \|(C_2)_{II}\|\le s_1(C).
\]
It follows that $s_{\min}((C_2)_{II})\ge b/N$.
The remaining matrix has norm $s_3(C)<a$, so the singular-value
perturbation inequality proves the claim.
\end{proof}

\begin{theorem}[Six spectral-square directions on a coordinate set]
\label{supp:six-carrier}
Let $s_3(A)\ge s>0$ and let $H$ have size $m\ge6$.  Suppose
$0\le G\le A^2$ and $\lambda_6(G_{HH})\ge g^2$.
For $N=\binom m2$, choose $0<a<b/N$ and $0<b<g$, and set
$\kappa=b/N-a$.  Then
\begin{equation}\label{eq:supp-six-carrier}
 \|p_Q\|_\infty\le\max\left\{
 \frac{2\sqrt{6\pi m}}a,
 \frac{9.2+\log_+(10000\kappa/s)}\kappa,
 \frac{5\sqrt{2e}}{\sqrt{g^2-b^2}}\right\}.
\end{equation}
All entries of the actual matrix $A$ outside $H$ are retained.
\end{theorem}
\begin{proof}
Put $C=A_{HH}$ and $B=A_{H,H^c}$.  If $s_3(C)\ge a$,
condition on the complementary original coordinates and apply
Corollary~\ref{cor:finite-core}.  This gives the first term.
If $s_3(C)<a$ and $s_2(C)\ge b$,
Lemma~\ref{supp:pair-extraction} gives a physical pair $K$ with
$s_{\min}(K)>\kappa$, hence $|\det K|>\kappa^2$.
A definite pair gives $\pi/\kappa$ by Lemma~\ref{geom:definite}.
An indefinite pair gives the estimate in Theorem~\ref{geom:fullpair}.
The function
\[
 x\longmapsto\frac{9.2+\log_+(10000x/s)}x
\]
is decreasing on $(0,\infty)$: on its logarithmic branch its
derivative has numerator $1-9.2-\log(10000x/s)<0$, and
elsewhere this is immediate.  Substitute $x=\sqrt{|\det K|}
>\kappa$; the second term in \eqref{eq:supp-six-carrier}
also exceeds $\pi/\kappa$.

The remaining case has $s_2(C)<b$.
Lemma~\ref{supp:spill} with $k=6,r=1$ gives
$s_5(B)>\sqrt{g^2-b^2}$.
The five-volume case of Corollary~\ref{four:d-cases} gives
\[
 \|p_Q\|_\infty\le5\sqrt{2e}\,
       e_5(BB^{\mathsf T})^{-1/10}
 \le5\sqrt{2e}/s_5(B),
\]
which proves the third term.  These three alternatives are exhaustive.
\end{proof}

The carrier hypothesis itself implies $s_6(A)\ge g$, so one may
always take $s=g$ if no stronger global spectral lower bound is
available.  In particular, this theorem does not extract six directions
from the three-direction premise of Conjecture~\ref{conj:main}.

\subsection{Stability under additional spectrum}

\begin{proposition}[Localized anticommutator bound]
\label{supp:anticommutator}
Let $A=A_0+R$ be symmetric and define
\[
 \nu_H=\max\{0,-\lambda_{\min}((A_0R+RA_0)_{HH})\}.
\]
For every $j\le|H|$,
\[
 \lambda_j((A^2)_{HH})\ge
       \lambda_j((A_0^2)_{HH})-\nu_H.
\]
If $A_0R+RA_0=0$, all these inequalities hold with zero loss.
This includes disjoint spectral packets, for which $A_0R=0$.
\end{proposition}
\begin{proof}
The identity $A^2=A_0^2+A_0R+RA_0+R^2$, together with
$R^2\ge0$, gives $(A^2)_{HH}\ge(A_0^2)_{HH}-\nu_HI$.
Apply monotonicity of ordered eigenvalues under this order.
The zero-loss assertions follow from the same identity.
\end{proof}

\begin{theorem}[Low-rank additions with an operator-norm remainder]
\label{supp:rank-reserve}
Let $A=A_0+S+R$ be symmetric, $\|R\|\le\varepsilon$, and
$r_H=\operatorname{rank}(SP_H)$.  If $k+r_H\le|H|$, then
\begin{equation}\label{eq:supp-rank-reserve}
 \sqrt{\lambda_k((A^2)_{HH})}\ge
 \left(\sqrt{\lambda_{k+r_H}((A_0^2)_{HH})}
                                  -\varepsilon\right)_+.
\end{equation}
No restriction is imposed on the signs or magnitudes of the
nonzero eigenvalues of $S$.
\end{theorem}
\begin{proof}
For any rectangular matrices $T,E$ with $\operatorname{rank}E\le r$,
the singular-value variational principle gives
\begin{equation}\label{eq:supp-rank-interlace}
 s_k(T+E)\ge s_{k+r}(T).
\end{equation}
Indeed, take a $(k+r)$-dimensional subspace on which
$|Tx|\ge s_{k+r}(T)|x|$ and intersect it with $\ker E$.
The intersection has dimension at least $k$, and $T+E=T$ there.
Apply \eqref{eq:supp-rank-interlace} to $T=A_0P_H$, $E=SP_H$,
then use $\|RP_H\|\le\varepsilon$.  Finally,
\[
 s_j(MP_H)^2=\lambda_j((M^2)_{HH})
\]
for symmetric $M$, which proves \eqref{eq:supp-rank-reserve}.
\end{proof}

\begin{remark}[The amount of spectral reserve]
A different argument notes that
\[
 \operatorname{range}((A_0+S)^2-A_0^2)
 \subset \operatorname{range}(S)+A_0\operatorname{range}(S).
\]
Rank interlacing of the squares then gives a valid but weaker loss
of at most $2\operatorname{rank}S$.  The fact that this difference
can have rank $2\operatorname{rank}S$ does not make that spectral
reserve loss necessary.  Theorem~\ref{supp:rank-reserve} uses the
rectangular operators before squaring and loses at most
$\operatorname{rank}S$, and often only $r_H$.
\end{remark}

\begin{corollary}[An adaptive split of a signed addition]
\label{supp:adaptive-rank}
Let $E=E^{\mathsf T}$, let $S_h$ contain its eigenmodes with
absolute eigenvalue greater than $h\ge0$, and put
$r_H(h)=\operatorname{rank}(S_hP_H)$.  If $\|T\|\le\delta$
and $A=A_0+E+T$, then, whenever $k+r_H(h)\le|H|$,
\[
 \sqrt{\lambda_k((A^2)_{HH})}\ge
 \left(\sqrt{\lambda_{k+r_H(h)}((A_0^2)_{HH})}
                                      -h-\delta\right)_+.
\]
The threshold $h$ may be chosen to maximize the bound.
\end{corollary}
\begin{proof}
The remaining part $E-S_h$ has operator norm at most $h$.
Apply Theorem~\ref{supp:rank-reserve} with remainder
$E-S_h+T$.
\end{proof}

\begin{proposition}[Separate positive and negative reserves]
\label{supp:signed-reserve}
Let $C=C_0+P-N+T$ be symmetric $m$ by $m$, where $P,N\ge0$
and $\|T\|\le\delta$.  Write $p_j=\lambda_j(P)$ and
$n_j=\lambda_j(N)$ in decreasing order, and set
$p_{m+1}=n_{m+1}=0$.  For $1\le u,v\le m$,
\begin{align*}
 \lambda_u(C)&\ge
 \max_{0\le r\le m-u}
    \{\lambda_{u+r}(C_0)-n_{r+1}\}-\delta,\\
 \lambda_v(-C)&\ge
 \max_{0\le r\le m-v}
    \{\lambda_{v+r}(-C_0)-p_{r+1}\}-\delta.
\end{align*}
If the active bounds are positive and supply a total of at least
three positive and negative directions, their minimum is a lower
bound for $s_3(C)$.
\end{proposition}
\begin{proof}
Write $N=N_r+N_{\mathrm{tail}}$, where $N_r$ has rank at most $r$
and $\|N_{\mathrm{tail}}\|=n_{r+1}$.  The favorable term $P\ge0$
can only increase ordered eigenvalues, so
\[
 \lambda_{u+r}(C_0+P+T-N_{\mathrm{tail}})
 \ge\lambda_{u+r}(C_0)-\delta-n_{r+1}.
\]
Subtracting $N_r$ loses at most $r$ indices by rank interlacing,
which proves the first inequality after maximizing in $r$.
Apply the same argument to $-C$ for the second.  Strictly positive
positive and negative spectral directions are disjoint, proving
the final assertion.
\end{proof}

\subsection{A complementary Schur block protects three directions}

\begin{theorem}[Complementary Schur protection]
\label{supp:schur-protection}
On an original coordinate block $H$ of size $m$, suppose
\[
 A_{HH}=B+U\Theta U^{\mathsf T}+T,
 \qquad B=B^{\mathsf T},\quad\Theta=\Theta^{\mathsf T},
 \quad\|T\|\le\delta.
\]
Let the columns of $W$ form an orthonormal basis of
$\operatorname{range}U$, and let the columns of $Z$ form an
orthonormal basis of its orthogonal complement.
Suppose $D=Z^{\mathsf T}BZ$ is invertible and $s_3(D)\ge\mu>0$.
Put
\[
 R=(W^{\mathsf T}BZ)D^{-1},\quad\rho=\|R\|,\quad
 \ell=\frac{\sqrt{\rho^2+4}-\rho}{2}.
\]
If $\delta<\ell^2\mu$, then for every $\Theta$,
\[
 \|p_Q\|_\infty\le
 \frac{2\sqrt{6\pi m}}{\ell^2\mu-\delta}.
\]
\end{theorem}
\begin{proof}
The orthogonal basis $[W\ Z]$ is used for matrix algebra only.
In this basis the block without $T$ equals
\[
 L\begin{pmatrix}S&0\\0&D\end{pmatrix}L^{\mathsf T},
 \quad L=\begin{pmatrix}I&R\\0&I\end{pmatrix},
 \quad
 S=W^{\mathsf T}(B+U\Theta U^{\mathsf T})W-RDR^{\mathsf T}.
\]
No invertibility assumption on $S$ is needed.  A singular-value
decomposition of $R$ decomposes $L$ into two-dimensional shears
$\bigl(\begin{smallmatrix}1&r\\0&1\end{smallmatrix}\bigr)$
and identity blocks.  The smaller singular value of such a shear
is $(\sqrt{r^2+4}-r)/2$, decreasing in $r\ge0$, so
$s_{\min}(L)=\ell$.  Multiplication on each side by an invertible
matrix reduces every singular value by at most its smallest singular
value.  Consequently
\[
 s_3(B+U\Theta U^{\mathsf T})
 \ge\ell^2s_3(\operatorname{diag}(S,D))\ge\ell^2\mu.
\]
The perturbation $T$ costs at most $\delta$.  Condition outside $H$
and apply Corollary~\ref{cor:finite-core} to finish.
\end{proof}

\begin{corollary}[Many large eigenvalues of one sign]
\label{supp:codimension-one}
If $A$ has at least $n-1$ eigenvalues at least $s>0$, or at least
$n-1$ eigenvalues at most $-s$, and $n\ge4$, then
\[
 \|p_Q\|_\infty\le\frac{2\sqrt{24\pi}}s.
\]
The remaining eigenvalue may have any sign and magnitude.
\end{corollary}
\begin{proof}
Every principal four-coordinate block has at least three eigenvalues
of the indicated sign and magnitude by interlacing.  Apply
Corollary~\ref{cor:finite-core} with $m=4$.
\end{proof}

These results give bounds under additional spectral hypotheses.
They do not establish that $s_3(A)\ge1$ alone guarantees a bounded
coordinate block, linear bank, or complementary Schur block satisfying
their respective hypotheses.

\subsection{Limitations of density interpolation and joint caps}

\begin{proposition}[An $L^3$ bound and a square-root interval bound]
\label{supp:interpolation-obstruction}
There are probability densities $p_N$, $N\ge2$, such that
\[
 \|p_N\|_3\le2^{1/3},\qquad
 \sup_z\int_{z-\varepsilon}^{z+\varepsilon}p_N(x)\,dx
        \le3\sqrt\varepsilon\quad(\varepsilon>0),
\]
while $\|p_N\|_\infty=N$.
\end{proposition}
\begin{proof}
Set
\[
 p_N=N\mathbf1_{[0,N^{-3}]}
                +(1-N^{-2})\mathbf1_{[2,3]}.
\]
Its mass is one, its essential supremum is $N$, and
$\int p_N^3=1+(1-N^{-2})^3\le2$.
An interval of length $2\varepsilon$ receives at most
$\min\{N^{-2},2N\varepsilon\}$ from the spike.  Dividing by
$\sqrt\varepsilon$, the maximum over $\varepsilon>0$ occurs
at $\varepsilon=(2N^3)^{-1}$ and equals
$\sqrt2N^{-1/2}\le1$.  If $0<\varepsilon\le1$, the other
part has mass at most $\min\{1,2\varepsilon\}\le2\sqrt\varepsilon$.
For $\varepsilon\ge1$, the total mass is already at most
$1\le3\sqrt\varepsilon$.  This proves all assertions.
\end{proof}

Thus a uniform $L^3$ bound and a uniform square-root interval bound
do not imply a uniform density bound.  An additional hypothesis is
needed to combine these estimates for quadratic forms.

\begin{proposition}[A moving quadratic ridge under a joint density cap]
\label{supp:moving-ridge}
For every $0<\varepsilon<1$ there are a uniform $X$ on $[0,1]$
and a triple $(Y,Z,W)$ independent of $X$, with joint probability density
bounded by one such that
\[
 T=4XY+Z^2-YW,\qquad \|p_T\|_\infty\ge\frac1{9\varepsilon}.
\]
The quadratic matrix has spectrum
$\{1,0,\sqrt{17}/2,-\sqrt{17}/2\}$.  The triple is not required
to have independent coordinates.
\end{proposition}
\begin{proof}
For $L>0$, define
\[
 S_{\varepsilon,L}=
 \{(y,z,w):\varepsilon\le y\le2\varepsilon,
       \ L\le w\le2L,
       \ \sqrt{yw}\le z\le\sqrt{yw+\varepsilon}\}.
\]
Its volume
\[
 V_\varepsilon(L)=\int_\varepsilon^{2\varepsilon}
       \int_L^{2L}
       (\sqrt{yw+\varepsilon}-\sqrt{yw})\,dw\,dy
\]
is positive and continuous.  It tends to zero as $L\downarrow0$,
because it is at most $\varepsilon^{3/2}L$.
For $L\ge1$, rationalization and $yw\le4\varepsilon L$ give
\[
 V_\varepsilon(L)\ge
    \frac{\varepsilon^{3/2}\sqrt L}{2\sqrt5},
\]
so it tends to infinity as $L\to\infty$.  Choose $L$ for which
$V_\varepsilon(L)=1$ and let the triple have density
$\mathbf1_{S_{\varepsilon,L}}$.
On this set, $0\le Z^2-YW\le\varepsilon$, while
$0\le4XY\le8\varepsilon$.  Thus $T\in[0,9\varepsilon]$
almost surely.  Given the triple, $4XY$ is uniform on an interval
of positive length; consequently $T$ has a density.  A probability
density supported on an interval of length $9\varepsilon$ has
essential supremum at least its reciprocal.

In coordinates $(X,Y,Z,W)$ the matrix is
\[
 A_*=\begin{pmatrix}
 0&2&0&0\\2&0&0&-1/2\\0&0&1&0\\0&-1/2&0&0
 \end{pmatrix}.
\]
The $(X,Y,W)$ block has eigenvalues
$0,\pm\sqrt{2^2+(1/2)^2}$, and $Z$ contributes eigenvalue one.
This proves the stated spectrum.
\end{proof}

Proposition~\ref{supp:moving-ridge} concerns a weaker independence
assumption than Conjecture~\ref{conj:main}.  It shows that keeping
one independent hub and a bounded joint density for the remaining
three statistics does not preserve enough information for the
desired conclusion.  It is not a counterexample to that conjecture.

\section{Gaussian limits and rank-three examples}
\label{sec:gaussian-benchmarks}

This section gives a sharp Gaussian density bound and explicit examples
for specified input classes or output levels.  For related Gaussian
quadratic-form estimates, see \citet{GNSU2019,BNU2020,CK2026}.

\subsection{The sharp Gaussian ceiling, including affine terms}

\begin{lemma}[Extracting a central square]\label{bench:central-square}
For every \(\lambda\geq1\) and \(m\in\mathbb R\), the law of
\(\lambda(Z+m)^2\), where \(Z\) is standard normal, is a convolution of
the law of a central \(\chi_1^2\) variable and a probability measure
on \([0,\infty)\).
\end{lemma}
\begin{proof}
The elementary Gaussian integral gives, for \(t\geq0\),
\[
 \mathbb E e^{-t\lambda(Z+m)^2}
 =(1+2\lambda t)^{-1/2}
       \exp\left(-\frac{\lambda m^2t}{1+2\lambda t}\right).
\]
First consider the centered factor.  For \(\lambda>1\), the finite
nonnegative measure
\[
 \nu_\lambda(dx)=\frac{e^{-x/(2\lambda)}-e^{-x/2}}{2x}
                 \mathbf1_{\{x>0\}}\,dx
\]
has mass \(\tfrac12\log\lambda\).  A compound Poisson sum with this
intensity measure has Laplace transform
\[
 \exp\left(-\int_0^\infty(1-e^{-tx})\nu_\lambda(dx)\right)
 =\left(\frac{1+2t}{1+2\lambda t}\right)^{1/2}.
\]
The identity follows by integrating
\(\int_0^\infty(e^{-ax}-e^{-bx})\,dx/x=\log(b/a)\);
it may also be checked by differentiation in \(t\) and evaluation at
zero.  At \(\lambda=1\), use the constant zero variable.  Multiplication
by \((1+2t)^{-1/2}\) extracts the central square.

The remaining noncentral factor is the Laplace transform of an
independent compound Poisson sum: take Poisson intensity \(m^2/2\)
and jumps distributed as \(\lambda\chi_2^2\).  Indeed its transform is
\[
 \exp\left[\frac{m^2}{2}\left((1+2\lambda t)^{-1}-1\right)\right]
 =\exp\left(-\frac{\lambda m^2t}{1+2\lambda t}\right).
\]
Uniqueness of Laplace transforms proves the stated convolution identity.
All variables in this distributional construction can be chosen
independently.
\end{proof}

\begin{theorem}[Sharp Gaussian density ceiling]\label{bench:gaussian}
Let \(G\sim N(\mu,\Sigma)\), where \(\Sigma\succeq\sigma^2I\) and
\(\sigma>0\).  If \(A\) is real symmetric and \(s_3(A)\geq a>0\),
then, for arbitrary \(b,c\), the quadratic polynomial
\(G^{\mathsf T}AG+b^{\mathsf T}G+c\) has a continuous density vanishing
at infinity and
\begin{equation}\label{bench:gaussian-bound}
 \|p_{G^{\mathsf T}AG+b^{\mathsf T}G+c}\|_\infty
 \leq\frac{1}{2\sqrt2\,\sigma^2a}.
\end{equation}
The constant is sharp.  In particular, \(\Sigma\succeq I/12\) and
\(s_3(A)\geq1\) give the bound \(3\sqrt2\).
\end{theorem}
\begin{proof}
Write \(G=\mu+\Sigma^{1/2}Z\).  The standard Gaussian vector \(Z\)
may be orthogonally rotated to diagonalize
\(B=\Sigma^{1/2}A\Sigma^{1/2}\).  Its coordinates remain independent.
The singular-value product inequality gives
\[
 s_3(B)\geq s_{min}(\Sigma^{1/2})^2s_3(A)
          \geq\sigma^2a.
\]
For completeness, this inequality follows by writing
\(A=\Sigma^{-1/2}B\Sigma^{-1/2}\) and applying
\(s_j(UDV)\leq\|U\|\|V\|s_j(D)\), which itself follows from the
best rank-\((j-1)\) approximation characterization of \(s_j\).

Put \(d=\sigma^2a\), and divide the output by \(d\).
Select three eigenvalues of \(B/d\) with absolute values at least
one.  Completing the square in each corresponding coordinate and
using Lemma~\ref{bench:central-square} shows that the entire output
is distributed as
\[
 \epsilon_1Z_1^2+\epsilon_2Z_2^2+\epsilon_3Z_3^2+W,
 \qquad \epsilon_j\in\{-1,1\},
\]
where the three displayed standard normals are independent of \(W\).
The constant terms, the extracted Poisson terms with their signs, all
other eigenvalue coordinates, and linear Gaussian terms from the kernel
are absorbed into \(W\).  This argument uses independence after a
Gaussian rotation only.

Convolution cannot increase a density supremum.  If all three signs
coincide, the density is a reflection of the \(\chi_3^2\) density,
whose maximum is
\[
 \max_{x>0}\frac{x^{1/2}e^{-x/2}}{\sqrt{2\pi}}
 =\frac1{\sqrt{2\pi e}}<\frac1{2\sqrt2}.
\]
For the mixed signature, reflection reduces to
\(D=Z_1^2+Z_2^2-Z_3^2\).  The first two squares sum to an exponential
variable of rate \(1/2\).  Therefore, for \(t\geq0\),
\[
 p_D(t)=\frac{e^{-t/2}}{2\sqrt{2\pi}}
                \int_0^\infty v^{-1/2}e^{-v}\,dv
        =\frac{e^{-t/2}}{2\sqrt2}.
\]
For \(u>0\), the same convolution gives
\[
 p_D(-u)=\frac{e^{-u/2}}{2\sqrt{2\pi}}
       \int_0^\infty(u+v)^{-1/2}e^{-v}\,dv
       \leq\frac{e^{-u/2}}{2\sqrt2}.
\]
These formulas give a continuous density, with maximum
\(p_D(0)=1/(2\sqrt2)\), vanishing at infinity.  Convolution with a
probability measure preserves these properties.  Rescaling by \(d\)
proves \eqref{bench:gaussian-bound}.  Equality holds for a centered
isotropic three-dimensional Gaussian and
\(A=a\,\diag(1,1,-1)\).
\end{proof}

\subsection{A global ceiling for diffuse Lindeberg arrays}

\begin{theorem}[Uniform convergence of diffuse quadratic densities]
\label{bench:lindeberg}
For each \(m\), let \(R_m\) be an \(n_m\)-by-\(3\) matrix with
orthonormal columns and rows \(r_{mi}^{\mathsf T}\), such that
\[
 \tau_m:=\max_i\|r_{mi}\|^2\longrightarrow0.
\]
Fix \(\Lambda=\diag(\lambda_1,\lambda_2,\lambda_3)\) with
\(|\lambda_j|\geq1\).  Let the physical coordinates \(X_{mi}\) be
independent, with densities bounded by one, finite means \(\mu_{mi}\),
and finite variances \(v_{mi}\).  Suppose
\begin{equation}\label{bench:lindeberg-limits}
 h_m:=\sum_i r_{mi}\mu_{mi}\longrightarrow h,\qquad
 \Sigma_m:=\sum_i v_{mi}r_{mi}r_{mi}^{\mathsf T}
                        \longrightarrow\Sigma,
\end{equation}
and, for every \(\varepsilon>0\),
\begin{equation}\label{bench:lindeberg-condition}
 \sum_i\mathbb E\left[
   \|r_{mi}(X_{mi}-\mu_{mi})\|^2
   \mathbf1_{\{\|r_{mi}(X_{mi}-\mu_{mi})\|>\varepsilon\}}
                  \right]\longrightarrow0.
\end{equation}
Then the forms
\[
 Q_m=X_m^{\mathsf T}R_m\Lambda R_m^{\mathsf T}X_m
\]
have continuous bounded densities for all sufficiently large \(m\).
If \(G\sim N(h,\Sigma)\) and \(Q_\infty=G^{\mathsf T}\Lambda G\),
then
\begin{equation}\label{bench:uniform-limit}
 \|p_{Q_m}-p_{Q_\infty}\|_\infty\longrightarrow0,
 \qquad
 \limsup_m\|p_{Q_m}\|_\infty\leq3\sqrt2.
\end{equation}
The last constant is sharp within this class.
\end{theorem}
\begin{proof}
Every cap-one density with finite variance satisfies
\(\operatorname{Var}(X)\geq1/12\).  One direct proof, writing
\(\mu=\mathbb EX\), is
\[
 \mathbb P(|X-\mu|>t)\geq1-2t\quad(0\leq t\leq1/2),
 \qquad
 \mathbb E(X-\mu)^2
 \geq\int_0^{1/2}2t(1-2t)\,dt=\frac1{12}.
\]
Consequently \(\Sigma_m\succeq I_3/12\) and
\(\Sigma\succeq I_3/12\).

Here are details of the required vector central limit argument.
Fix \(u\in\mathbb R^3\), and put
\(\xi_{mi}=u^{\mathsf T}r_{mi}(X_{mi}-\mu_{mi})\).
The scalar Lindeberg condition follows from
\eqref{bench:lindeberg-condition}, and
\(\sum_i\mathbb E\xi_{mi}^2\to u^{\mathsf T}\Sigma u\).
The expansion of \(e^{ix}\) to order two, splitting its remainder
at \(|x|=\delta\), gives
\[
 \sum_i\left|
  \mathbb Ee^{i\xi_{mi}}-1+\tfrac12\mathbb E\xi_{mi}^2
            \right|\longrightarrow0.
\]
Indeed, the small part is at most a constant times
\(\delta\sum_i\mathbb E\xi_{mi}^2\), and the large part tends to
zero by Lindeberg; then let \(\delta\downarrow0\).
The same splitting shows
\(\max_i\mathbb E\xi_{mi}^2\to0\).
Taking logarithms of the product of the scalar characteristic
functions therefore proves
\[
 R_m^{\mathsf T}X_m\ \Longrightarrow\ N(h,\Sigma).
\]
The continuity theorem for characteristic functions and the continuous
mapping theorem now give pointwise convergence
\(\varphi_{Q_m}(t)\to\varphi_{Q_\infty}(t)\).

For large \(m\), all row leverages are at most \(1/2\), so
Theorem~\ref{four:diffuse} applies to the actual finite product laws.
Since \(|\det\Lambda|\geq1\), it gives the common integrable bound
\[
 |\varphi_{Q_m}(t)|\leq
 g(t):=\min\{1,(\pi e)^{3/2}|t|^{-3/2}\}.
\]
The pointwise limit satisfies the same bound.  Dominated convergence
and Fourier inversion give
\[
 \|p_{Q_m}-p_{Q_\infty}\|_\infty
 \leq\frac1{2\pi}\int_{\mathbb R}
       |\varphi_{Q_m}(t)-\varphi_{Q_\infty}(t)|\,dt
 \longrightarrow0.
\]
Theorem~\ref{bench:gaussian}, with the covariance floor already proved,
gives the final ceiling.

For sharpness, take \(n_m=3m\), with \(m\) rows equal to
\(e_j/\sqrt m\) for each \(j=1,2,3\), and use independent centered
unit uniforms.  The hypotheses hold with \(h=0\),
\(\Sigma=I_3/12\), and \(\Lambda=\diag(1,1,-1)\).
The limiting density has maximum \(3\sqrt2\), so uniform convergence
gives \(\|p_{Q_m}\|_\infty\to3\sqrt2\).
\end{proof}

\begin{remark}[Hypotheses of the limiting assertion]
The nonzero spectrum of every finite matrix
\(R_m\Lambda R_m^{\mathsf T}\) is exactly the spectrum of \(\Lambda\).
The three projected coordinates need not be independent before the
Gaussian limit.  The proof also allows
\(Q_m=Y_m^{\mathsf T}\Lambda Y_m+\beta_m^{\mathsf T}Y_m+c_m\),
where \(Y_m=R_m^{\mathsf T}X_m\), \(\beta_m\to\beta\), and
\(c_m\to c\), because the physical linear terms are separable phases
in Theorem~\ref{four:diffuse}.  Finite variance, convergence of the
covariances and means, and the Lindeberg condition are hypotheses of
this limiting assertion.
Weak convergence alone would not supply
\eqref{bench:uniform-limit}; the integrable product-law Fourier bound
is the step that excludes narrow surviving density peaks.
\end{remark}

\subsection{One heavy coordinate and two Gaussianized banks}

\begin{theorem}[A uniform gap at the cone tip]\label{bench:one-heavy}
Let \(0<a\leq1\), \(d=\sqrt{1-a^2}\), and
\[
 Q_a=-aX^2+2dXB+aB^2+C^2,
\]
where \(B,C\) are independent \(N(0,1/12)\), independent of \(X\).
Over all probability densities \(0\leq f_X\leq1\), the unique
maximizer of \(p_{Q_a}(0)\), up to null sets, is the centered unit
interval density.  Write the maximum as \(C(a)\).  Then
\begin{equation}\label{bench:one-heavy-bound}
 C(a)\leq
 \frac{\sqrt{6\pi}\,\operatorname{erf}(\sqrt{3/2})+3\sqrt2}{2}
 <4.111374<3\sqrt2.
\end{equation}
This is a bound at output zero.
\end{theorem}
\begin{proof}
Conditional on \(X=x\), complete the square:
\[
 Q_a=a(B+dx/a)^2+C^2-x^2/a.
\]
The density of \((\sqrt a(B+dx/a),C)\) is
\[
 \frac6{\pi\sqrt a}
       \exp\{-6[(u/\sqrt a-dx/a)^2+v^2]\}.
\]
Polar coordinates in \((u,v)\) therefore give the conditional
zero-density kernel, for almost every \(x\),
\begin{equation}\label{bench:heavy-kernel}
 K_a(x)=\frac3{\pi\sqrt a}\int_0^{2\pi}
                      e^{-6F_a(\theta)x^2}\,d\theta,
 \qquad
 F_a(\theta)=\frac{(\cos\theta-d)^2}{a^2}
                         +\frac{\sin^2\theta}{a}.
\end{equation}
The conditional density at any level is at most \(6/\sqrt a\).
For each \(x\ne0\) it is continuous in the level near zero, so
dominated convergence against \(f_X\) proves continuity of the
unconditional density there and identifies its value with
\(\int f_X(x)K_a(x)\,dx\).

The function \(F_a\) is strictly positive.  Thus \(K_a\) is even
and strictly decreasing in \(|x|\).  Filling the interval closest to
zero is consequently the unique cap-one maximizer: explicitly, with
\(I=[-1/2,1/2]\),
\[
 \int (\mathbf1_I-f_X)(K_a-K_a(1/2))\,dx\geq0,
\]
and equality forces \(f_X=\mathbf1_I\) almost everywhere.  Put
\[
 L(s)=\int_{-1/2}^{1/2}e^{-6sx^2}\,dx,
 \qquad R(s)=\sqrt{1+s}\,L(s).
\]
Then \(C(a)=\frac3{\pi\sqrt a}\int L(F_a(\theta))\,d\theta\).
We claim that \(R(0)=1\) and \(R\) is strictly decreasing on
\((0,\infty)\).  If
\(M(s)=\int_{-1/2}^{1/2}x^2e^{-6sx^2}\,dx\), integration by parts
gives \(L(s)-12sM(s)=e^{-3s/2}\), while \(e^y>1+y\) gives
\[
 L(s)=e^{-3s/2}\int_{-1/2}^{1/2}e^{6s(1/4-x^2)}\,dx
       >(1+s)e^{-3s/2}\qquad(s>0).
\]
Since \(L'=-6M\), these identities imply
\[
 \frac{d}{ds}\log R(s)
 =\frac1{2(1+s)}-
        \frac{L(s)-e^{-3s/2}}{2sL(s)}<0.
\]

Write \(\alpha=\sqrt{1+a}\), \(\beta=\sqrt{1-a}\).
Direct expansion gives the perfect-square identity
\[
 a^2(1+F_a(\theta))=(\alpha-\beta\cos\theta)^2.
\]
The tangent-half-angle substitution, or an elementary rational
integral, gives
\[
 \int_0^{2\pi}\frac{d\theta}{\alpha-\beta\cos\theta}
       =\frac{2\pi}{\sqrt{\alpha^2-\beta^2}}.
\]
Consequently
\[
 d\nu_a(\theta)=
 \frac{(1+F_a(\theta))^{-1/2}}{\pi\sqrt{2a}}\,d\theta
\]
is a probability measure and
\begin{equation}\label{bench:heavy-average}
 \frac{C(a)}{3\sqrt2}=\int R(F_a)\,d\nu_a.
\end{equation}

We next show \(\nu_a\{F_a\geq1\}\geq1/2\).
At \(a=1\), \(F_a\equiv1\).  For \(0<a<1\), put
\[
 c_a=\frac{\alpha-a\sqrt2}{\beta},\qquad
 \theta_0=\arccos c_a,\qquad
 \kappa=\sqrt{\frac{\alpha+\beta}{\alpha-\beta}}.
\]
Here \(0<c_a<1\): the upper inequality follows from
\(\alpha+\beta>\sqrt2\), and the lower one also follows from
the comparison below.  The event \(F_a<1\) is exactly
\(\cos\theta>c_a\).  The same substitution gives
\[
 \nu_a\{F_a<1\}
 =\frac2\pi\arctan\left(\kappa\tan\frac{\theta_0}{2}\right).
\]
But
\[
 c_a-\frac\beta\alpha
 =\frac{a(2-\sqrt2\,\alpha)}{\alpha\beta}\geq0.
\]
Equivalently,
\(\kappa\tan(\theta_0/2)\leq1\), so the displayed probability
is at most \(1/2\).  As \(R\leq1\) everywhere and
\(R(s)\leq R(1)\) for \(s\geq1\),
\eqref{bench:heavy-average} yields
\[
 C(a)\leq\frac{3\sqrt2}{2}(1+R(1)),\qquad
 R(1)=\sqrt{\frac\pi3}\,
              \operatorname{erf}(\sqrt{3/2})<1.
\]
This proves the exact bound and its strict gap below \(3\sqrt2\).

The decimal upper bound can also be certified without numerical
quadrature.  Expanding the exponential and integrating gives the
alternating series
\[
 L(1)=\sum_{j=0}^{\infty}
       \frac{(-3/2)^j}{j!(2j+1)}.
\]
Its term magnitudes decrease, so the even partial sum through
\(j=12\) is an upper bound.  Since
\(\sqrt2<14142135623731/10^{13}\), exact rational arithmetic gives
\[
 \frac{3\sqrt2}{2}+3L(1)
 <\frac{3}{2}\frac{14142135623731}{10^{13}}
       +3\sum_{j=0}^{12}\frac{(-3/2)^j}{j!(2j+1)}
 <4.111374.
\]
The expression on the left equals the exact ceiling in
\eqref{bench:one-heavy-bound}.
\end{proof}

\begin{remark}[Exact finite product realization]
The matrix on the three displayed statistics is
\[
 M_a=\begin{pmatrix}-a&d&0\\d&a&0\\0&0&1\end{pmatrix},
 \qquad M_a^2=I_3,\qquad \operatorname{spec}(M_a)=\{1,1,-1\}.
\]
Replace \(B,C\) by two independent normalized sums of \(m\)
centered unit uniforms, retaining the independent heavy coordinate
\(X\).  The corresponding matrix in the \(1+2m\) original
coordinates is \(T_mM_aT_m^{\mathsf T}\), where \(T_m\) has one
row \(e_1^{\mathsf T}\), \(m\) rows \(e_2^{\mathsf T}/\sqrt m\),
and \(m\) rows \(e_3^{\mathsf T}/\sqrt m\).  Its columns are
orthonormal, so its nonzero spectrum is exactly \(\{1,1,-1\}\).
For each fixed \(a>0\) and fixed heavy density, its cone-tip density
converges to that of \(Q_a\).

The bank characteristic function is
\(\psi_m(t)=[\sin(t/(2\sqrt m))/(t/(2\sqrt m))]^m\).
For \(|t|\leq\sqrt m\), the elementary bound
\(|\sin u/u|\leq e^{-c u^2}\) on \(|u|\leq1/2\), for some
\(c>0\), gives a common Gaussian majorant.  For
\(\sqrt m\leq|t|\leq4\sqrt m\), compactness gives the bound
\(\rho^m\), with a fixed \(\rho<1\).  For
\(|t|\geq4\sqrt m\), the bound
\((2\sqrt m/|t|)^m\) has integral tending to zero.  The central
limit expansion on compact intervals and these three bounds prove
\(L^1\) convergence of \(\psi_m\) to the Gaussian characteristic
function.  Thus the bank densities converge uniformly, and are
uniformly bounded, by Fourier inversion.  In the conditional polar
formula used for \eqref{bench:heavy-kernel}, replace the product of
Gaussian bank densities by the product of these two bank densities.
Uniform convergence of the products, the finite angular range, and
\(\int f_X=1\) prove the asserted cone-tip limit.
Neither this argument nor \eqref{bench:one-heavy-bound} gives a
global bound over all output levels, or a convergence rate uniform
as \(a\downarrow0\).
\end{remark}

\subsection{An exact dense Householder example}

\begin{proposition}[The Householder cube cone tip]\label{bench:householder}
Let \(X_1,X_2,X_3\) be independent centered unit uniforms, and put
\[
 e=\frac1{\sqrt3}(1,1,1)^{\mathsf T},\qquad
 A=I-2ee^{\mathsf T}
 =\frac13\begin{pmatrix}1&-2&-2\\-2&1&-2\\-2&-2&1\end{pmatrix}.
\]
Then \(\operatorname{spec}(A)=\{1,1,-1\}\), and the density of
\(Q=X^{\mathsf T}AX\) is continuous at zero, with
\begin{equation}\label{bench:householder-tip}
 p_Q(0)=6\sqrt3\,\operatorname{artanh}
              \left(\frac{\sqrt2-1}{\sqrt3}\right)
       =2.5343552034\ldots.
\end{equation}
In particular, this is also the limit of
\(\mathbb P(|Q|\leq\varepsilon)/(2\varepsilon)\).
\end{proposition}
\begin{proof}
The matrix acts by \(-1\) on \(e\) and by \(+1\) on
\(e^\perp\), proving the spectral statement.  Choose an orthonormal
basis \(u,v\) of \(e^\perp\), and let
\(u_\theta=u\cos\theta+v\sin\theta\).
For \(\mathcal C=[-1/2,1/2]^3\), polar coordinates in the positive
two-plane give the following density version:
\[
 p_Q(s)=\frac12\int_{\mathbb R}\int_0^{2\pi}
  \mathbf1_{\{s+z^2>0\}}
  \mathbf1_{\mathcal C}\bigl(\sqrt{s+z^2}\,u_\theta+ze\bigr)
                                                   \,d\theta\,dz.
\]
Indeed the radial measure is \(r\,dr\), and changing from
\(r\) to \(s=r^2-z^2\) contributes the factor \(1/2\).
Whenever the cube indicator is nonzero, \(|z|\leq\sqrt3/2\).
As \(s\to0\), its integrand converges except at \(z=0\) or
where the limiting ray meets the cube boundary.  For each \(\theta\)
there are only finitely many such nonzero radii.  Dominated convergence
therefore proves continuity at zero.

Put \(d_\pm(\theta)=u_\theta\pm e\) and
\(M_\pm(\theta)=\max_i|d_{\pm,i}(\theta)|\).
At zero, each ray remains in the cube up to
\(r=1/(2M_\pm(\theta))\), so
\[
 p_Q(0)=\frac14\int_0^{2\pi}
        \left(\frac1{M_+(\theta)}+\frac1{M_-(\theta)}\right)d\theta
       =\frac12\int_0^{2\pi}\frac{d\theta}{M_+(\theta)}.
\]
The last equality uses
\(M_-(\theta)=M_+(\theta+\pi)\).
With a suitable origin of the angular coordinate, the three
coordinates of \(d_+\) are phase shifts by \(2\pi/3\) of
\[
 a_0+b_0\cos\theta,\qquad
 a_0=1/\sqrt3,\quad b_0=\sqrt{2/3}.
\]
Every negative coordinate has absolute value at most
\(b_0-a_0<a_0\), whereas the largest positive coordinate is at
least \(a_0\), because their sum is \(\sqrt3\).
Thus the largest absolute coordinate is the largest positive one.
Each coordinate is maximal on a sector of width \(2\pi/3\), giving
\[
 p_Q(0)=\frac32\int_{-\pi/3}^{\pi/3}
                    \frac{d\theta}{a_0+b_0\cos\theta}.
\]
Finally \(t=\tan(\theta/2)\) yields
\[
 \int_{-\pi/3}^{\pi/3}\frac{d\theta}{a_0+b_0\cos\theta}
 =\frac4{\sqrt{b_0^2-a_0^2}}
   \operatorname{artanh}\left(
      \frac1{\sqrt3}\sqrt{\frac{b_0-a_0}{b_0+a_0}}\right).
\]
Since \(\sqrt{b_0^2-a_0^2}=1/\sqrt3\) and
\(\sqrt{(b_0-a_0)/(b_0+a_0)}=\sqrt2-1\), this proves
\eqref{bench:householder-tip}.  Continuity proves the stated
small-ball limit.
\end{proof}

The Householder calculation concerns the cone tip and does not locate
the global maximum of that example.  The one-heavy calculation has the
same output-level restriction.  The diffuse Lindeberg theorem is global
in the output level, under its stated array hypotheses.  These distinctions
leave the unrestricted indefinite product-law endpoint open.

\section{Counterexamples to stronger assertions and open cases}\label{sec:limitations}

The following counterexamples delimit several stronger intermediate estimates. None is a counterexample to Conjecture~\ref{conj:main}.

\subsection{A bounded joint density cannot replace a product law}

\begin{proposition}[A thin quadratic shell]\label{lim:shell}
There is no bound for the density of a nondegenerate three-variable quadratic that depends only on an upper bound for the joint input density and a lower bound for the three eigenvalue moduli.
\end{proposition}
\begin{proof}
Fix $0<\delta<1$ and put $R=e^{1/\delta}$. In the change of variables
\[
 x_1=r,\quad x_3=z,\quad x_2=(t-z^2)/r,
 \qquad 1<r<R,\quad0<z<1,\quad0<t<\delta,
\]
the Jacobian is $1/r$. The resulting region $E_\delta\subset\mathbb R^3$ has volume $\delta\log R=1$. Give $X$ the density $\mathbf1_{E_\delta}\le1$. Then $x_1x_2+x_3^2=t$ is uniform on $(0,\delta)$. The quadratic
\[
 Q=2X_1X_2+2X_3^2
\]
has matrix eigenvalues $1,-1,2$, yet its density equals $1/(2\delta)$ on $(0,2\delta)$. This tends to infinity. The coordinates of $X$ are dependent.
\end{proof}

\begin{example}[Conditional densities can be arbitrarily high]\label{lim:conditional}
Let $X_1,X_2$ be independent uniforms on $(0,1)$, and set $L=(X_1+X_2)/\sqrt2$, $V=(X_1-X_2)/\sqrt2$. For $0<\ell<1/\sqrt2$, the conditional law of $V$ given $L=\ell$ is uniform on $(-\ell,\ell)$ and has density $1/(2\ell)$. The marginal density of $L$ there is $2\ell$. Thus the conditional density norm is unbounded, and its $p$th moment is infinite for every $p\ge2$. Both coordinates before rotation have density one. An unconditional projection theorem therefore cannot be substituted for a conditional-density estimate.
\end{example}

\subsection{Uniform pointwise Fourier decay is too strong}

\begin{proposition}\label{lim:pointwise}
Even for $Q=X_1^2+X_2^2+X_3^2$, there is no function $b(t)\to0$ as $t\to\infty$ such that $|\mathbb E e^{itQ}|\le b(t)$ uniformly over all independent coordinate densities bounded by one.
\end{proposition}
\begin{proof}
Fix $t>0$ and $0<\eta<\pi/2$. The set
\[
 F=\{x>0:\operatorname{dist}(tx^2,2\pi\mathbb Z)<\eta\}
\]
has infinite measure. Indeed, its intervals near $\sqrt{2\pi j/t}$ have lengths comparable, for large $j$, to $\eta/(\sqrt t\sqrt j)$, whose sum diverges. Choose a measurable subset $E\subset F$ of measure one, and take all three coordinate densities equal to $\mathbf1_E$. Each expectation $\mathbb E e^{itX_i^2}$ has real part at least $\cos\eta$, hence modulus at least $\cos\eta$. Independence gives $|\mathbb E e^{itQ}|\ge(\cos\eta)^3$. Since $\eta$ is arbitrary, the supremum over admissible laws equals one for every fixed $t>0$. The laws may depend on $t$. This is consistent with the Riemann--Lebesgue lemma for each fixed law and with the bounded density in Theorem~\ref{geom:diagonal}.
\end{proof}

\subsection{An arbitrarily small aligned quadratic phase can destroy an integral bound}

\begin{theorem}[Aligned-twist obstruction]\label{lim:twist}
Let $p=(1,1)/\sqrt2$ and fix $\kappa\ne0$. Among independent coordinate densities $f_1,f_2\le1$, the quantity
\[
 J_\kappa=\int_{\mathbb R}\left|\mathbb E
 e^{iu\langle p,X\rangle+i\kappa\langle p,X\rangle^2}\right|\,du
\]
is unbounded. In contrast, $J_0\le2\pi\sqrt2$ for every such pair. Thus small operator norm of a quadratic phase is insufficient for a law-uniform bound on this Fourier integral, even when both frame leverages equal $1/2$.
\end{theorem}
\begin{proof}
For an integer $N\ge2$, set $\sigma=1/N$. Choose a spacing $D>0$ sufficiently large below, and take
\[
 f_1(x)=\frac1N\sum_{j=0}^{N-1}\gamma_\sigma(x-\sqrt2Dj),
 \qquad
 f_2(x)=\frac1N\sum_{k=0}^{N-1}\gamma_\sigma(x-\sqrt2DNk),
\]
where $\gamma_\sigma$ is the centered Gaussian density of variance $\sigma^2$. Each isolated bump has height $1/\sqrt{2\pi}<1$. The tails of other bumps converge uniformly to zero as $D$ grows. More explicitly, choose the nearest center: all other centers are at least $(m-1/2)\sqrt2D$ away for integers $m\ge1$, so their contribution is bounded by a convergent Gaussian tail series tending to zero. Thus $f_1,f_2\le1$ for sufficiently large $D$.

The linear statistic $L=\langle p,X\rangle$ is an equally weighted mixture of $N^2$ Gaussians with variance $\sigma^2$ and distinct means $m_{jk}=D(j+Nk)$. For one component of mean $m$, completing the square in its Gaussian integral gives
\[
 |T_m(u)|=(1+4\kappa^2\sigma^4)^{-1/4}
 \exp\left[-\frac{\sigma^2(u+2\kappa m)^2}
 {2(1+4\kappa^2\sigma^4)}\right].
\]
Its $L^1$ norm is
\[
 A_N=\frac{\sqrt{2\pi}}\sigma(1+4\kappa^2\sigma^4)^{1/4}
 \ge\sqrt{2\pi}N.
\]
The frequency centers are separated by $2|\kappa|D$. Choose $D$ so large that the disjoint windows
$U_m=\{|u+2\kappa m|<|\kappa|D\}$ each contain at least $(1-\eta)A_N$ of their own component, with a fixed $\eta<1/4$. The reverse triangle inequality on these windows yields
\[
 J_\kappa\ge\frac1{N^2}\sum_m
 \left(\int_{U_m}|T_m|-\int_{U_m^c}|T_m|\right)
 \ge(1-2\eta)A_N.
\]
To see the first inequality, sum the cross-component errors over disjoint windows: for each component, their union is contained in the complement of its own window. The lower bound grows linearly with $N$.

For $\kappa=0$, independence and Cauchy--Schwarz give
\[
 J_0=\sqrt2\int|\widehat f_1(v)\widehat f_2(v)|\,dv
 \le\sqrt2\,\|\widehat f_1\|_2\|\widehat f_2\|_2
 \le2\pi\sqrt2,
\]
by Plancherel and $\int f_i^2\le\int f_i=1$.
\end{proof}

The input laws in this construction consist of many narrow separated components. Their final density caps arise from separation of component locations. They need not obey the average component-peak condition in the Gaussian-mixture theorem of Section~\ref{sec:hubs}. That theorem and the present obstruction have different hypotheses.

Two further failures are proved in the supplement: Proposition~\ref{supp:interpolation-obstruction} shows that a bounded $L^3$ density norm and a square-root interval profile do not imply a bounded density, and Proposition~\ref{supp:moving-ridge} shows that an independent hub together with a bounded joint density for a dependent triple can still leave an unbounded output density. Neither provides an admissible product-law counterexample.

\subsection{An independent-block estimate for separated phases}

\begin{lemma}[Phase-weighted projection]\label{lim:phase-L2}
Let $Z=T X$ be a rank-$k$ linear image of a product vector, and let $\psi(X)$ be any real measurable phase. Then
\[
 \int_{\mathbb R^k}|\mathbb E e^{i\langle u,Z\rangle+i\psi(X)}|^2\,du
 \le(2\pi)^k\|p_Z\|_\infty
 \le(2\pi)^k2^{k/2}\det(TT^T)^{-1/2}.
\]
\end{lemma}
\begin{proof}
The complex pushforward has density $m(z)=\mathbb E[e^{i\psi(X)}\mid Z=z]p_Z(z)$, so $|m|\le p_Z$ and $\int|m|^2\le\|p_Z\|_\infty$. Apply Plancherel and the projection theorem of \citet{LPP2016}. This uses a bound on the complex pushforward, not on a conditional input density.
\end{proof}

\begin{theorem}[Two independent blocks]\label{lim:block-phase}
Let the coordinates split as $K\sqcup I\sqcup J$, and let $P$ be an $n\times k$ matrix. Put $G_K=P_K^TP_K$ and $G_I=P_I^TP_I$, and assume both are positive definite. Let $E$ be any symmetric matrix with $E_{KI}=0$. Then, allowing also an arbitrary extra phase depending on $X_J$,
\[
 \int_{\mathbb R^k}|\mathbb E e^{i\langle Pu,X\rangle+iX^TEX+i\psi(X_J)}|\,du
 \le(2\pi)^k2^{k/2}(\det G_K\det G_I)^{-1/4}.
\]
\end{theorem}
\begin{proof}
Condition on $X_J$. The phase splits into one function of $X_K$ and one function of $X_I$, together with a constant. Independence factors the conditional expectation. Apply Cauchy--Schwarz in $u$ and Lemma~\ref{lim:phase-L2} to each factor. The bound is independent of the frozen coordinates, so the triangle inequality and Fubini finish the proof.
\end{proof}

\begin{corollary}\label{lim:block-quadratic}
Let $U$ have three orthonormal columns, let $\Lambda=\operatorname{diag}(\lambda_1,\lambda_2,\lambda_3)$ be invertible, and let $A=U\Lambda U^T+B$. If a split $K\sqcup I\sqcup J$ has $B_{KI}=0$ and positive definite $G_K=U_K^TU_K$, $G_I=U_I^TU_I$, then every affine perturbation satisfies
\[
 \|p_Q\|_\infty\le6|\lambda_1\lambda_2\lambda_3|^{-1/3}
 (\det G_K\det G_I)^{-1/6}.
\]
\end{corollary}
\begin{proof}
Fresnel-linearize the three displayed squares, using the regularization justified in Theorem~\ref{four:fresnel}. The preceding theorem, with the phase $tX^TBX$ and the affine phase absorbed into its coordinate blocks, gives
\[
 |\varphi_Q(t)|\le(2\pi)^{3/2}|\det\Lambda|^{-1/2}
 (\det G_K\det G_I)^{-1/4}|t|^{-3/2}.
\]
Integrating its minimum with one gives $(3/\pi)$ times the $2/3$ power of the displayed coefficient, which is the stated bound. Arbitrary within-block quadratic terms are permitted, but the cross-block zero condition is exact.
\end{proof}

\subsection{The remaining uniform estimate}

A normalized family illustrating the remaining difficulty is
\begin{equation}\label{eq:remaining-core}
 Q_\varepsilon=2X_0L+V^2+2\varepsilon LP,
\end{equation}
where $X_0$ is an independent physical coordinate and $L,V,P$ are orthonormal linear statistics of the other coordinates with a spread frame. In the orthonormal basis formed by the hub coordinate and these three vectors, its matrix is
\[
 \begin{pmatrix}
 0&1&0&0\\1&0&0&\varepsilon\\0&0&1&0\\0&\varepsilon&0&0
 \end{pmatrix}.
\]
The characteristic polynomial is $z(z-1)(z^2-1-\varepsilon^2)$, so its three nonzero eigenvalues are $1,\sqrt{1+\varepsilon^2},-\sqrt{1+\varepsilon^2}$. The third eigenvalue modulus does not deteriorate when $\varepsilon\downarrow0$.

For $\varepsilon=0$, the independent-hub/transverse-square argument gives a uniform bound. For every fixed nonzero $\varepsilon$, Fourier decay gives a bounded density. The bounds for arbitrary cap-one coordinate laws deteriorate as $\varepsilon\downarrow0$, even after symmetric-window recovery. The Gaussian-mixture results give uniform bounds under the additional hypotheses in Section~\ref{sec:hubs}.

An estimate for \eqref{eq:remaining-core} alone does not cover the affine, hub-diagonal, and light-block terms that arise in more general matrices. In addition, the sufficient conditions above do not imply that every no-pivot matrix has a coordinate decomposition satisfying one of them.

To deduce the unrestricted conjecture from the results above, one would need a bound for the uncovered analytic cases and a quantitative decomposition covering all admissible matrices. A uniform estimate for the adapted-hub functional, extended to the additional quadratic terms, could address the analytic cases. The constants must be independent of dimension and of a vanishing leakage parameter. No counterexample to the unrestricted conjecture follows from the results above.


\bibliographystyle{plainnat}
\bibliography{bibliography}
\end{document}
